\chapter{The field of norms}

\section{APF extensions}

\begin{definition}[The APF openness condition]
\label{fw:apf-openness}
\dcref{Wintenberger1983, Definition 1.2.1}
\source{wintenberger:page-0005}
\uses{fw:absolute-upper-projections}
\lean{FontaineWintenberger.apfSubgroup, FontaineWintenberger.IsAPFSubextension}
\leanok
For every intermediate field $L\subset K^{\mathrm{sep}}$, without a normality
assumption on $L/K$, put
\[
 H_L(u)=G_L\vee G_K^u\subset G_K.
\]
The extension $L/K$ is APF if $H_L(u)$ is open for every $u\geq-1$.
\end{definition}

\begin{lemma}[Compact APF product subgroups]
\label{fw:apf-products}
\dcref{Wintenberger1983, Definition 1.2.1; compact subgroup formulation}
\source{wintenberger:page-0005}
\uses{fw:apf-openness,fw:profinite-galois-foundations}
\lean{FontaineWintenberger.coe_apfSubgroup, FontaineWintenberger.coe_apfSubgroup_upper_mul, FontaineWintenberger.isCompact_apfSubgroup, FontaineWintenberger.isClosed_apfSubgroup, FontaineWintenberger.apfSubgroup_antitoneOn, FontaineWintenberger.apfSubgroup_neg_one}
\leanok
For every $L\subset K^{\mathrm{sep}}$ and every $u$, the subgroup $H_L(u)$
has underlying set $G_LG_K^u=G_K^uG_L$, and is compact and closed.
The groups $H_L(u)$ decrease for $u\geq-1$ and $H_L(-1)=G_K$.
\end{lemma}
\begin{proof}
The upper subgroup $G_K^u$ is normal, so multiplying $G_L$ and $G_K^u$
in either order gives the subgroup they generate. The fixing subgroup
$G_L$ is closed in the Krull topology: it is the intersection, over all
$x\in L$, of the closed subgroups fixing $K(x)$, which are open of
finite index. Thus both factors are compact. Their product is compact
as the continuous multiplication image of their compact product, and
is closed in the Hausdorff group $G_K$. Decreasingness follows from
that of the upper groups, and $G_K^{-1}=G_K$ gives the endpoint.
\end{proof}

\begin{proposition}[Elementary APF properties]
\label{fw:apf-elementary}
\dcref{Wintenberger1983, Definition 1.2.1 and Proposition 1.2.3}
\source{wintenberger:page-0005}
\uses{fw:apf-products}
\lean{FontaineWintenberger.isOpen_apfSubgroup_iff_finiteIndex, FontaineWintenberger.isAPFSubextension_iff_finiteIndex, FontaineWintenberger.IsAPFSubextension.index_pos, FontaineWintenberger.isAPFSubextension_of_finite, FontaineWintenberger.IsAPFSubextension.of_le}
\leanok
Every finite separable $L/K$ is APF, including $L=K$.
If $M\subset L$ and $L/K$ is APF, then $M/K$ is APF.
For arbitrary separable $L/K$, APF is equivalent to finite index of every
$H_L(u)$, $u\geq-1$. In particular the corresponding natural-number
indices are positive.
\end{proposition}
\begin{proof}
For finite $L/K$, the subgroup $G_L$ is open. A subgroup containing an
open subgroup is open, since it is the union of its cosets. Hence every
$H_L(u)$ is open. If $M\subset L$, then $G_L\subset G_M$, so
$H_L(u)\subset H_M(u)$, and the same observation proves inheritance.
For the last assertion, an open subgroup of a compact group has finite
index: its discrete coset space is a continuous image of a compact space
and therefore finite. Conversely, each $H_L(u)$ is closed. If it has
finite index, its complement is a finite union of closed cosets, hence
closed, so the subgroup is open. A finite coset space is nonempty,
which gives positivity of its cardinality.
\end{proof}

\begin{lemma}[Conjugacy of the APF openness condition]
\label{fw:apf-conjugacy}
\dcref{Wintenberger1983, Definition 1.2.1; independence of the chosen separable embedding}
\source{wintenberger:page-0005}
\uses{fw:apf-elementary,fw:absolute-upper-projections}
\lean{FontaineWintenberger.apfSubgroup_map, FontaineWintenberger.apfSubgroup_map_index, FontaineWintenberger.isAPFSubextension_map_iff}
\leanok
For $L\subset K^{\mathrm{sep}}$ and $\sigma\in G_K$, the fixed-closure
APF condition holds for $\sigma(L)/K$ if and only if it holds for $L/K$.
Indeed $H_{\sigma(L)}(u)=\sigma H_L(u)\sigma^{-1}$, so their actual
subgroup indices agree.
\end{lemma}
\begin{proof}
An automorphism fixes $\sigma(L)$ pointwise exactly when its conjugate by
$\sigma^{-1}$ fixes $L$. Thus $G_{\sigma(L)}=\sigma G_L\sigma^{-1}$.
Normality gives $\sigma G_K^u\sigma^{-1}=G_K^u$, so conjugation carries
$H_L(u)$ onto $H_{\sigma(L)}(u)$. It induces a bijection on the actual
coset spaces, preserving indices. The proved equivalence between
finite index and openness of these closed subgroups proves the result.
\end{proof}

\begin{proposition}[APF for an abstract separable extension]
\label{fw:apf-abstract}
\dcref{Wintenberger1983, Definition 1.2.1; intrinsic extension formulation}
\source{wintenberger:page-0005}
\uses{fw:apf-conjugacy}
\lean{FontaineWintenberger.nonempty_separableClosure_embeddings, FontaineWintenberger.exists_aut_comp_separableClosure_embeddings, FontaineWintenberger.isAPFSubextension_fieldRange_iff, FontaineWintenberger.IsAPF, FontaineWintenberger.isAPF_iff_embedding, FontaineWintenberger.isAPF_iff_exists_embedding, FontaineWintenberger.isAPF_iff_subextension, FontaineWintenberger.isAPF_of_finite, FontaineWintenberger.IsAPF.of_algHom, FontaineWintenberger.isAPF_algEquiv}
\leanok
Let $L/K$ be an abstract separable algebraic field extension. There exists
a $K$-embedding $i:L\hookrightarrow K^{\mathrm{sep}}$. Define $L/K$ to be
APF if the fixed-closure condition holds for the image of some such $i$.
For every chosen embedding $i$, this is equivalent to the fixed-closure
condition for $i(L)$, and equivalently it holds for all embeddings.
Every finite separable extension is APF. The APF property descends along
every $K$-algebra embedding $M\hookrightarrow L$, and is invariant under
every $K$-algebra isomorphism $L\simeq M$.
\end{proposition}
\begin{proof}
The separable closure admits an embedding of each separable algebraic
extension. To see this directly, order embeddings of intermediate
subfields of $L/K$ by extension. Unions of chains are embeddings. A
maximal one must contain every element of $L$: otherwise the separable
minimal polynomial of a missing element has a root in the separably
closed target, extending the embedding. This is a contradiction.

For two embeddings $i,j$, the isomorphism $i(L)\simeq j(L)$ extends to
an embedding of $K^{\mathrm{sep}}$ into itself by the same maximal-extension
argument. It is surjective by normality: on the finite set of roots of
the minimal polynomial over $K$ of any element, it is an injective
self-map and hence a permutation. Thus it is some $\sigma\in G_K$ with
$\sigma\circ i=j$. The preceding conjugacy lemma identifies the two
fixed-closure APF conditions. This proves both embedding-independence
and the asserted equivalences, with existence already established.
For a finite extension, every image is finite, so the earlier finite
APF theorem applies. For an embedding $f:M\hookrightarrow L$, choose an
embedding $i$ of $L$ into the separable closure. Then $i(f(M))\subset i(L)$;
the earlier subextension theorem and embedding-independence show that
APF descends to $M$. Applying this to both directions of a $K$-algebra
isomorphism proves isomorphism invariance.
\end{proof}

\begin{lemma}[Finite positive APF inertia indices]
\label{fw:apf-inertia-indices}
\dcref{Wintenberger1983, Definition 1.2.1}
\source{wintenberger:page-0005}
\uses{fw:apf-elementary,fw:apf-products}
\lean{FontaineWintenberger.apfInertiaSubgroup, FontaineWintenberger.apfInertiaIndex, FontaineWintenberger.apfInertiaSubgroup_eq_inf, FontaineWintenberger.apfInertiaSubgroup_le_zero, FontaineWintenberger.apfInertiaSubgroup_antitoneOn, FontaineWintenberger.apfInertiaIndex_zero, FontaineWintenberger.IsAPFSubextension.inertiaIndex_pos, FontaineWintenberger.IsAPFSubextension.inertia_isFiniteRelIndex, FontaineWintenberger.IsAPFSubextension.inertiaIndex_monotoneOn}
\leanok
For $u\geq0$, set $J_L(u)=(G_L\cap G_K^0)\vee G_K^u$. Then
\[
 J_L(u)=H_L(u)\cap G_K^0\subset G_K^0.
\]
If $L/K$ is APF, the actual index
\[
 a_L(u)=[G_K^0:J_L(u)]
\]
is finite and positive. Moreover $a_L(0)=1$ and $a_L$ is nondecreasing
on $[0,+\infty)$.
\end{lemma}
\begin{proof}
Write $U=G_K^0$ and $N=G_K^u\subset U$. Normality of $N$ identifies both
suprema with subgroup products. If $hn\in H_L(u)\cap U$, where
$h\in G_L$ and $n\in N$, then $h=(hn)n^{-1}\in U$. Thus
$hn\in(G_L\cap U)N$, giving the displayed equality; the reverse
inclusion follows directly from $N\subset U$.
The APF condition makes $H_L(u)$ have finite index in $G_K$. Intersecting
with any subgroup preserves finite relative index, so its intersection
with $U$ has finite index in $U$. Its coset space is nonempty, proving
positivity. At $u=0$, one has $J_L(0)=U$, so the index is one.
For $0\leq u\leq v$, decreasingness gives $J_L(v)\subset J_L(u)\subset U$.
Multiplicativity of finite subgroup indices gives
\[
 a_L(v)=[J_L(u):J_L(v)]a_L(u).
\]
The first factor is a positive integer, proving monotonicity.
\end{proof}

\begin{lemma}[Intrinsic APF inertia index]
\label{fw:apf-index-embedding}
\dcref{Wintenberger1983, Definition 1.2.1; independence of separable embedding}
\source{wintenberger:page-0005}
\uses{fw:apf-abstract,fw:apf-inertia-indices}
\lean{FontaineWintenberger.apfInertiaSubgroup_map, FontaineWintenberger.apfInertiaIndex_map, FontaineWintenberger.apfInertiaIndex_fieldRange_eq, FontaineWintenberger.apfIndex, FontaineWintenberger.apfIndex_eq_embedding, FontaineWintenberger.apfIndex_zero, FontaineWintenberger.IsAPF.apfIndex_pos, FontaineWintenberger.IsAPF.apfIndex_monotoneOn, FontaineWintenberger.apfIndex_algEquiv}
\leanok
For an abstract separable algebraic extension $L/K$, choose an embedding
$i:L\hookrightarrow K^{\mathrm{sep}}$ and put
$a_{L/K}(u)=a_{i(L)}(u)$. This number is independent of the embedding
and is invariant under $K$-algebra isomorphisms of extensions.
If $L/K$ is APF, it is positive and nondecreasing for $u\geq0$, and
$a_{L/K}(0)=1$.
\end{lemma}
\begin{proof}
Any two embeddings differ by an absolute automorphism $\sigma$. Both
$G_K^0$ and $G_K^u$ are normal, so conjugation by $\sigma$ carries
$(G_{i(L)}\cap G_K^0)\vee G_K^u$ onto
$(G_{\sigma(i(L))}\cap G_K^0)\vee G_K^u$, and preserves $G_K^0$.
It therefore induces a bijection of the actual relative coset spaces.
Their indices agree. Positivity, monotonicity and the value at zero
follow from the preceding fixed-closure index theorem applied to $i(L)$;
the abstract APF equivalence supplies its hypothesis. Composing an
embedding with a $K$-algebra isomorphism leaves its image unchanged,
which proves isomorphism invariance.
\end{proof}

\begin{definition}[APF extensions]
\label{fw:apf}
\dcref{Wintenberger1983, Definition 1.2.1}
\source{wintenberger:page-0005}
\uses{fw:apf-abstract,fw:apf-openness,fw:apf-elementary,fw:apf-index-embedding,fw:apf-extended-bound,fw:absolute-upper-projections,fw:upper-inverse-integral}
\lean{FontaineWintenberger.IsAPF, FontaineWintenberger.apfPsi, FontaineWintenberger.apfPsi_of_nonpos, FontaineWintenberger.apfPsi_eq_integral_index, FontaineWintenberger.apfPhi, FontaineWintenberger.apfPsi_apfPhi, FontaineWintenberger.apfPhi_apfPsi, FontaineWintenberger.apfRamificationBound}
\leanok
An extension $L/K$ is \emph{arithmetically profinite}, or \emph{APF},
if $G_K^uG_L$ is open in $G_K$ for every real $u\geq-1$.
For such an extension put $G_L^0=G_L\cap G_K^0$ and
\[
 \psi_{L/K}(u)=
 \begin{cases}
  u,&-1\leq u\leq0,\\
  \displaystyle\int_0^u [G_K^0:G_L^0G_K^v]\,dv,&u\geq0.
 \end{cases}
\]
Write $\varphi_{L/K}=\psi_{L/K}^{-1}$ and
\[
 i(L/K)=\sup\{u\geq-1:G_K^uG_L=G_K\}.
\]
The supremum defining $i(L/K)$ is taken in $[-1,+\infty]$;
in particular $i(K/K)=+\infty$. The functions $\psi_{L/K}$ and $\varphi_{L/K}$ use the same normalization
as the finite Herbrand functions.
\end{definition}

\begin{lemma}[Extended ramification bound]
\label{fw:apf-extended-bound}
\dcref{Wintenberger1983, Definition 1.2.1; extended-real endpoint convention}
\source{wintenberger:page-0005}
\uses{fw:apf-abstract,fw:absolute-upper-projections}
\lean{FontaineWintenberger.apfSubgroup_fieldRange_index_eq, FontaineWintenberger.apfSubgroup_fieldRange_eq_top_iff, FontaineWintenberger.apfRamificationBound, FontaineWintenberger.apfRamificationBound_eq_embedding, FontaineWintenberger.le_apfRamificationBound, FontaineWintenberger.neg_one_le_apfRamificationBound, FontaineWintenberger.apfRamificationBound_algEquiv, FontaineWintenberger.apfRamificationBound_self}
\leanok
The bound $i(L/K)$ is independent of the chosen separable embedding,
is invariant under $K$-algebra isomorphisms, and satisfies
$-1\leq i(L/K)\leq+\infty$. Every $u\geq-1$ with
$G_K^uG_L=G_K$ satisfies $u\leq i(L/K)$. For the trivial extension,
$i(K/K)=+\infty$.
\end{lemma}
\begin{proof}
The set whose extended-real supremum defines $i(L/K)$ contains $-1$,
because $G_K^{-1}=G_K$. Thus its supremum is at least $-1$ and bounds
every member. Changing the embedding conjugates $G_LG_K^u$ and
preserves the condition that it equals $G_K$; hence it leaves this set
unchanged. Composing an embedding with an isomorphism also leaves the
image unchanged. For $L=K$ its fixing subgroup is $G_K$, so every
real $u\geq-1$ belongs to the set. The set is unbounded above in
$\mathbb R$, and its supremum in the extended reals is $+\infty$.
\end{proof}

\begin{proposition}[The APF Herbrand homeomorphism]
\label{fw:apf-herbrand-homeomorphism}
\dcref{Wintenberger1983, Definition 1.2.1; elementary integral consequences}
\source{wintenberger:page-0005}
\uses{fw:apf-index-embedding,fw:apf}
\lean{FontaineWintenberger.apfHerbrandSlope, FontaineWintenberger.apfHerbrandSlope_of_nonpos, FontaineWintenberger.IsAPF.one_le_apfHerbrandSlope, FontaineWintenberger.IsAPF.monotone_apfHerbrandSlope, FontaineWintenberger.IsAPF.intervalIntegrable_apfHerbrandSlope, FontaineWintenberger.apfPsi, FontaineWintenberger.apfPsi_of_nonpos, FontaineWintenberger.apfPsi_zero, FontaineWintenberger.apfPsi_eq_integral_index, FontaineWintenberger.IsAPF.continuous_apfPsi, FontaineWintenberger.IsAPF.apfPsi_sub_ge, FontaineWintenberger.IsAPF.strictMono_apfPsi, FontaineWintenberger.IsAPF.le_apfPsi, FontaineWintenberger.IsAPF.surjective_apfPsi, FontaineWintenberger.apfHerbrandOrderIso, FontaineWintenberger.apfHerbrandOrderIso_apply, FontaineWintenberger.apfPhi, FontaineWintenberger.apfPsi_apfPhi, FontaineWintenberger.apfPhi_apfPsi, FontaineWintenberger.apfPhi_of_nonpos, FontaineWintenberger.continuous_apfPhi, FontaineWintenberger.strictMono_apfPhi, FontaineWintenberger.IsAPF.apfPsi_mem_Ici, FontaineWintenberger.apfPhi_mem_Ici, FontaineWintenberger.apfPsi_algEquiv, FontaineWintenberger.apfPhi_algEquiv}
\leanok
For every APF extension $L/K$, the function $\psi_{L/K}$ is continuous,
strictly increasing and maps $[-1,+\infty)$ onto itself. Its inverse
$\varphi_{L/K}$ is continuous and strictly increasing, and both functions
are the identity on $[-1,0]$. Both are invariant under $K$-algebra
isomorphisms of the extension. Moreover
\[
 \psi_{L/K}(t)-\psi_{L/K}(s)\geq t-s
 \quad(-1\leq s\leq t).
\]
\end{proposition}
\begin{proof}
Extend $a_{L/K}$ to negative arguments by the constant value one. The
resulting function $b(v)=a_{L/K}(\max(v,0))$ is nondecreasing and at
least one. It is bounded and integrable on each compact interval.
Its primitive $P(t)=\int_0^t b(v)\,dv$ is continuous. For $t\leq0$,
the integrand is one on $[t,0]$, so $P(t)=t$; for $t\geq0$ this is
the defining integral of $\psi_{L/K}$. For $s\leq t$, additivity of
interval integrals and $b\geq1$ give
$P(t)-P(s)=\int_s^t b(v)\,dv\geq t-s$. Thus $P$ is strictly
increasing. It tends to $+\infty$ to the right and equals the identity
to the left, so the intermediate value theorem makes it surjective
on $\mathbb R$. A continuous increasing bijection of $\mathbb R$ is an
order isomorphism and has continuous increasing inverse. Restricting
to $[-1,+\infty)$, which it preserves in both directions, proves all
claims. Isomorphism invariance of the intrinsic indices identifies
the integrals, and uniqueness of the inverse identifies their inverses.
\end{proof}

\begin{proposition}[APF indices on finite Galois layers]
\label{fw:finite-apf-indices}
\dcref{Wintenberger1983, 1.1.2 and Definition 1.2.1; finite Galois specialization}
\source{wintenberger:page-0005,serre-1979-local-fields:page-0081}
\uses{fw:apf,fw:absolute-upper-projections,fw:upper-inverse-integral}
\lean{FontaineWintenberger.isAPF_of_finite, FontaineWintenberger.apfSubgroup_eq_comap_upper, FontaineWintenberger.apfInertiaIndex_eq_upperRelIndex, FontaineWintenberger.apfIndex_eq_upperRelIndex, FontaineWintenberger.apfPsi_eq_inverseHerbrandFunction, FontaineWintenberger.apfPhi_eq_herbrandFunction}
\leanok
Every finite Galois extension $E/K$ is APF. For $G=\Gal(E/K)$ and
$u\geq0$,
\[
 [G_K^0:(G_E\cap G_K^0)G_K^u]=[G^0:G^u].
\]
Consequently the function $\psi_{E/K}$ in the APF definition is the
inverse of the finite Herbrand function, and $\varphi_{E/K}$ is the
finite Herbrand function itself on $[-1,+\infty)$.
\end{proposition}
\begin{proof}
The subgroup $G_E$ is open in $G_K$. Thus $G_EG_K^u$, which is a
subgroup since $G_K^u$ is normal, contains an open subgroup and is open.
This proves APF. Surjective absolute upper projections at $0$ and $u$
make the restriction map $G_K^0\to G^0$ surjective with kernel
$G_E\cap G_K^0$ and with image of $G_K^u$ equal to $G^u$.
The inverse image of $G^u$ is precisely $(G_E\cap G_K^0)G_K^u$:
if an element has the same restriction as some element of $G_K^u$,
their quotient lies in the kernel. The subgroup index identity follows.
Integrating it and using the inverse Herbrand integral gives the last
claim for $u\geq0$; both definitions are the identity on $[-1,0]$.
For an abstract finite Galois extension choose a separable embedding.
The isomorphism to its image identifies the intrinsic APF index and
carries the actual finite upper groups onto those of the image, proving
the same identity without an ambient-field choice. Uniqueness of inverse
functions then identifies $\varphi_{E/K}$ with the finite Herbrand function.
\end{proof}

\begin{proposition}[Finite non-Galois APF change of parameter]
\label{fw:finite-nonnormal-apf}
\dcref{Wintenberger1983, Definition 1.2.1 and Proposition 1.2.3; finite overfield formulation}
\source{wintenberger:page-0005}
\uses{fw:finite-apf-indices,fw:herbrand-subgroup-integral,fw:apf-inertia-indices}
\lean{FontaineWintenberger.apfSubgroup_map_eq_comap_sup_upper, FontaineWintenberger.fixingSubgroup_sup_upper_inf_zero, FontaineWintenberger.apfInertiaIndex_map_eq_subgroupRelIndex, FontaineWintenberger.apfIndex_eq_subgroupRelIndex, FontaineWintenberger.apfPsi_eq_subgroup_herbrand_comp_inverse, FontaineWintenberger.apfPhi_eq_herbrand_comp_subgroup_inverse}
\leanok
Let $E/K$ be a finite Galois extension inside $K^{\mathrm{sep}}$, let
$M$ be any intermediate field of $E/K$, and put $G=\Gal(E/K)$ and
$H=\Gal(E/M)$. With the actual actions on $\mathcal O_E$, for $u\geq0$,
\[
 a_{M/K}(u)=[G_0:(H\cap G_0)G^u],\qquad
 \psi_{M/K}(u)=\varphi_H(\psi_G(u)).
\]
The second equality holds also for $-1\leq u\leq0$. In particular it
requires no normality of $M/K$. Equivalently, for $v\geq-1$,
$\varphi_{M/K}(v)=\varphi_G(\psi_H(v))$.
\end{proposition}
\begin{proof}
Let $\rho:G_K\to G$ be ordinary restriction. Its kernel is $G_E$,
and the fixing subgroup of $M$ in $G_K$ is $\rho^{-1}(H)$.
The finite upper projection theorem and the kernel identity imply
\[
 G_MG_K^u=\rho^{-1}(HG^u).
\]
Indeed, restriction takes the product into $HG^u$. Conversely, given
an element of the preimage, lift its $G^u$ factor from $G_K^u$;
the remaining factor restricts into $H$, so fixes $M$. At $u=0$,
restriction takes $G_K^0$ onto $G_0$. For $u\geq0$, intersection with
$G_K^0$ therefore takes the defining APF subgroup onto
\[
 (HG^u)\cap G_0=(H\cap G_0)G^u.
\]
For this last equality use $G^u\subset G_0$: a product in $G_0$
has its $H$ factor in $G_0$. Passing to the actual coset spaces under
the surjective restriction gives the index identity. Integrating it
and applying Lemma~\ref{fw:herbrand-subgroup-integral} proves the
change of parameter for $u\geq0$. Both sides equal $u$ on $[-1,0]$.
Independence of the chosen embedding identifies the APF functions of
the abstract field $M$ with those of its image in $K^{\mathrm{sep}}$.
Composing the change of parameter with the proved inverse functions
gives the final inverse identity, since both maps preserve $[-1,+\infty)$.
\end{proof}

\begin{lemma}[Intrinsic finite indices in an arbitrary Galois overfield]
\label{fw:finite-galois-overfield-apf}
\uses{fw:finite-nonnormal-apf,fw:ramification-field-equiv,fw:apf-index-embedding,fw:apf-extended-bound}
\lean{FontaineWintenberger.fixingSubgroup_map_autCongr, FontaineWintenberger.apfIndex_eq_subgroupRelIndex_of_finiteGalois, FontaineWintenberger.apfPsi_eq_subgroup_herbrand_comp_inverse_of_finiteGalois, FontaineWintenberger.apfPhi_eq_herbrand_comp_subgroup_inverse_of_finiteGalois, FontaineWintenberger.apfSubgroup_fieldRange_index_eq_fixingSubgroup_sup_upper_index, FontaineWintenberger.apfSubgroup_fieldRange_eq_top_iff_fixingSubgroup_sup_upper_eq_top, FontaineWintenberger.apfRamificationBound_eq_sSup_fixingSubgroup_sup_upper}
\leanok
Let $\Omega/K$ be any finite Galois extension and $M$ an actual
intermediate field. Put $G=\operatorname{Gal}(\Omega/K)$,
$H=\operatorname{Gal}(\Omega/M)$ and let $G^u$ be the actual
upper groups on its spectral integer ring. For $u\geq0$,
\[
 a_{M/K}(u)=[G_0:(H\cap G_0)G^u],
\]
and for $u\geq-1$,
\[
 \psi_{M/K}(u)=\varphi_H(\psi_G(u)).
\]
For $v\geq-1$, the inverse is
\[
 \varphi_{M/K}(v)=\varphi_G(\psi_H(v)).
\]
For every $K$-embedding $j:\Omega\hookrightarrow K^{\rm sep}$,
ordinary absolute restriction also gives
\[
 [G_K:G_{j(M)}G_K^u]=[G:HG^u]\qquad(u\geq-1).
\]
In particular the absolute product is the whole group exactly
when $HG^u=G$. The actual extended first bound is
\[
 i(M/K)=\sup\{u\in[-1,\infty):HG^u=G\},
\]
with supremum taken in the extended reals. No embedding of
$\Omega$ in a preassigned closure, or normality of $M/K$, is
part of the assertion.
\end{lemma}
\begin{proof}
Choose a $K$-embedding $j:\Omega\hookrightarrow K^{\rm sep}$,
which exists because the actual finite Galois extension is
separable. Its image $E$ is a finite Galois intermediate field.
The induced isomorphism $\Omega\simeq_K E$ identifies $M$
with its actual image $M'\subseteq E$. Conjugation by this
isomorphism identifies the full automorphism groups. An
automorphism fixes every element of $M$ if and only if its
conjugate fixes every element of $M'$, so the actual fixing
subgroups correspond. Uniqueness of the extended norm identifies
the spectral integer rings and all displacement orders; therefore
the lower and upper groups and their Herbrand functions correspond.
The finite nonnormal index and Herbrand formulas for $M'\subseteq E$
now transport to the displayed formulas, by intrinsic invariance
of the APF index under $M\simeq_K M'$. Taking inverses of the
composition gives the displayed formula for $\varphi_{M/K}$.
Finally ordinary restriction to $E$ is surjective, and identifies
the absolute APF product for $M'$ with the inverse image of
$\operatorname{Gal}(E/M')\operatorname{Gal}(E/K)^u$.
Surjective restriction identifies the left cosets of that inverse
image with the cosets of the finite product, proving the index
identity. That inverse image is the whole group exactly when the
product is the whole finite group. Substitution in the intrinsic
first-bound supremum and the same fixing-group isomorphism give
the last formula, including both infinite endpoints.
\end{proof}

\begin{proposition}[Upper groups after a finite base change]
\label{fw:finite-base-upper}
\dcref{Wintenberger1983, Proposition 1.2.3; finite upper-group change of parameter}
\source{wintenberger:page-0005}
\uses{fw:finite-galois-overfield-apf,fw:ramification-subgroups,fw:ramification-dvr-equiv,fw:spectral-base-integers}
\lean{FontaineWintenberger.finiteBase_upper_map_fixingSubgroupEquiv, FontaineWintenberger.finiteBase_upperRamificationGroup_of_finiteGalois, FontaineWintenberger.finiteBase_upperRamificationGroup}
\leanok
Let $E/K$ be any finite Galois extension, let $M$ be any
actual intermediate field, and write $G=\Gal(E/K)$, $H=\Gal(E/M)$, identified
with the subgroup of $G$ fixing $M$. Give the fields their compatible
extended norms. For every $u\geq-1$,
\[
 \Gal(E/M)^{\psi_{M/K}(u)}=H\cap\Gal(E/K)^u.
\]
The group on the left uses the integer ring constructed over $M$;
the group on the right uses that constructed over $K$.
\end{proposition}
\begin{proof}
The extended norm of $E$ is the unique multiplicative norm extending
that of the complete field $K$. The inherited norm makes the finite
intermediate field $M$ complete, and the same norm of $E$ extends it.
Uniqueness of the spectral extension over $M$ therefore identifies the
two norms and their unit balls. Thus their integer rings are equal as
subrings of $E$. The identification of the subgroup fixing $M$ with
$\Gal(E/M)$ acts on this very ring by the same field automorphisms.
It consequently preserves every displacement valuation, lower group,
and Herbrand function.

Put $t=\psi_G(u)$ and $v=\psi_{M/K}(u)$. The preceding proposition gives
$v=\varphi_H(t)$ and $t\geq-1$, so $\psi_H(v)=t$. Lower subgroup
compatibility and the definition of upper numbering now give
\[
 H^v=H_{\psi_H(v)}=H_t=H\cap G_t=H\cap G^u.
\]
The integer-ring and automorphism identifications in the first paragraph
make the left side exactly the actual upper group over $M$, proving
the assertion including the endpoints.
\end{proof}

\begin{lemma}[Finite base change through an abstract embedded base]
\label{fw:finite-tower-upper}
\dcref{Wintenberger1983, Proposition 1.2.3; intrinsic finite-base formulation}
\source{wintenberger:page-0005}
\uses{fw:finite-base-upper,fw:apf-herbrand-homeomorphism}
\lean{FontaineWintenberger.finiteTower_autEquivImage, FontaineWintenberger.finiteTower_spectralIntegerRing_eq, FontaineWintenberger.finiteTower_spectralIntegerRingEquiv, FontaineWintenberger.finiteTower_spectralIntegerRingEquiv_smul, FontaineWintenberger.finiteTower_upper_map_autEquivImage, FontaineWintenberger.finiteTower_fixingSubtype_comp, FontaineWintenberger.finiteTower_range_restrictScalarsHom, FontaineWintenberger.finiteTower_upper_map_restrictScalars, FontaineWintenberger.finiteTower_upper_map_restrictScalars_range, FontaineWintenberger.finiteTower_mem_upper_iff_restrictScalars}
\leanok
Let $E/K$ be finite Galois inside $K^{\mathrm{sep}}$, and let $M/K$
be a finite separable extension with a $K$-embedding $i:M\hookrightarrow E$.
Give $M$ its compatible extended norm and regard $E$ as an $M$-algebra
through $i$. Write $\iota:\Gal(E/M)\hookrightarrow\Gal(E/K)$ for
restriction of scalars. Then, for every $u\geq-1$,
\[
 \iota\bigl(\Gal(E/M)^{\psi_{M/K}(u)}\bigr)
   =\iota(\Gal(E/M))\cap\Gal(E/K)^u.
\]
The upper group on the left is constructed using its own actual
spectral integer ring over $M$. Equivalently, an automorphism of $E/M$
belongs to the left upper group exactly when its restriction of scalars
belongs to $\Gal(E/K)^u$.
\end{lemma}
\begin{proof}
Put $S=i(M)\subset E$. The actual field isomorphism $M\simeq S$
identifies their intrinsic APF functions over $K$, so
$\psi_{M/K}=\psi_{S/K}$. An automorphism of $E$ fixes $M$ through
$i$ exactly when it fixes every element of $S$. Thus changing scalars
from $M$ to $S$ identifies their actual automorphism groups while
preserving the underlying maps on $E$.

The compatible norm of $E$ extends the norms of both $M$ and $S$.
Spectral uniqueness identifies both constructed integer rings with the
same unit ball of $E$. The scalar-change isomorphism therefore preserves
the actual integer action, every displacement, and upper ramification
groups. Apply the preceding finite-base proposition to $S\subset E$
and then restrict scalars to $K$. The image fixing subgroup is exactly
$\iota(\Gal(E/M))$, which gives the displayed equality.
\end{proof}

\begin{lemma}[A common separable closure after base change]
\label{fw:apf-base-separable-closure}
\dcref{Wintenberger1983, Proposition 1.2.3; identification of ambient separable closures}
\source{wintenberger:page-0005}
\lean{FontaineWintenberger.isSepClosure_over_subextension, FontaineWintenberger.separableClosureBaseChangeEquiv, FontaineWintenberger.absoluteGaloisBaseChangeEquiv, FontaineWintenberger.coe_absoluteGaloisBaseChangeEquiv, FontaineWintenberger.absoluteGaloisBaseChangeEquiv_apply}
\leanok
For every intermediate field $M\subset K^{\mathrm{sep}}$, the field
$K^{\mathrm{sep}}$ is a separable closure of $M$. There is an
$M$-isomorphism
$e:M^{\mathrm{sep}}\simeq K^{\mathrm{sep}}$. Conjugation by $e$,
followed by restriction of scalars, identifies the actual absolute
Galois group of $M$ with the subgroup of $G_K$ fixing $M$.
For every $g\in G_M$ and $x\in M^{\mathrm{sep}}$, its image satisfies
$\widetilde g(e(x))=e(g(x))$.
\end{lemma}
\begin{proof}
The separably closed field $K^{\mathrm{sep}}$ is separable algebraic over
$M$: the minimal polynomial over $M$ of any element divides its
separable minimal polynomial over $K$. It is therefore a separable
closure of $M$. Uniqueness of separable closures gives the actual
$M$-isomorphism $e$. Conjugation sends an $M$-automorphism $g$ to
$ege^{-1}$, an $M$-automorphism of $K^{\mathrm{sep}}$. Restricting its
scalars to $K$ identifies it with an element fixing $M$, and every such
element is obtained, by reversing conjugation. These maps preserve
multiplication and are inverse. The displayed equality is evaluation
of $ege^{-1}$ at $e(x)$. No finiteness or normality of $M/K$ is needed.
\end{proof}

\begin{lemma}[Residue fields after finite normed base extension]
\label{fw:finite-base-residue}
\dcref{Wintenberger1983, Proposition 1.2.3; residue-field prerequisite}
\source{wintenberger:page-0005}
\uses{fw:normed-intermediate-tower}
\lean{FontaineWintenberger.normedAlgebra_residueField_finiteDimensional, FontaineWintenberger.normedAlgebra_residueField_perfectField}
\leanok
Let $M/K$ be finite separable, with its compatible extended norm.
The residue field of its actual norm unit ball is finite dimensional
over $k_K$. If $k_K$ is perfect, this residue field is perfect.
\end{lemma}
\begin{proof}
The actual unit ball $\mathcal O_M$ is the integral closure of
$\mathcal O_K$ and is finite as an $\mathcal O_K$-module. The compatible
valuation map is local, so it induces the canonical map on residue
fields. Quotienting the finite module and then factoring its scalars
through $\mathcal O_K/\mathfrak m_K$ proves finite dimensionality of
$k_M/k_K$. This extension is algebraic, and every algebraic extension
of a perfect field is perfect.
\end{proof}

\begin{lemma}[Testing absolute upper groups on cofinal finite layers]
\label{fw:cofinal-upper-conditions}
\dcref{Serre1979, Chapter IV, Section 3, Remark 1; cofinal consequence of the infinite-upper definition}
\source{serre-1979-local-fields:page-0082}
\uses{fw:finite-tower-upper-restriction,fw:infinite-upper-groups}
\lean{FontaineWintenberger.mem_upperRamificationGroupInfinite_iff_cofinal}
\leanok
Let $\Omega/K$ be a field extension and let $(E_j)_j$ be a cofinal
family of finite Galois $K$-subextensions of $\Omega$. For $u\geq-1$,
an automorphism $g$ belongs to the actual infinite upper group $G^u$
if and only if $g|_{E_j}\in\Gal(E_j/K)^u$ for every $j$.
\end{lemma}
\begin{proof}
Membership in the infinite group implies every finite restriction
condition by definition. Conversely let $D/K$ be any finite Galois
subextension and choose $E_j$ containing $D$. The restriction of
$g|_{E_j}$ to $D$ is $g|_D$. The finite upper quotient theorem maps
$\Gal(E_j/K)^u$ onto $\Gal(D/K)^u$, so the assumed condition at
$E_j$ gives the required condition at $D$. Testing every $D$ is the
definition of membership in $G^u$.
\end{proof}

\begin{lemma}[Cofinal Galois overfields for a finite base extension]
\label{fw:cofinal-base-galois-layers}
\dcref{Wintenberger1983, Proposition 1.2.3; finite normal-closure reduction}
\source{wintenberger:page-0005}
\lean{FontaineWintenberger.finiteGaloisNormalClosure, FontaineWintenberger.le_finiteGaloisNormalClosure, FontaineWintenberger.exists_finiteGalois_over_base_layer, FontaineWintenberger.exists_finiteGalois_containing_base, FontaineWintenberger.finiteGaloisLayerOverBase, FontaineWintenberger.finiteGaloisLayerOverBase_restrictScalars}
\leanok
Let $M/K$ be a finite subextension of $\Omega=K^{\mathrm{sep}}$.
The finite extensions $E/K$ in $\Omega$ that are Galois over $K$ and
contain $M$ are cofinal among the finite Galois extensions over $M$
in $\Omega$. They are also cofinal among finite Galois $K$-subextensions
of $\Omega$. Each such $E/M$ is finite Galois.
\end{lemma}
\begin{proof}
If $D/M$ is finite Galois in $\Omega$, then $D/K$ is finite separable:
finiteness and separability hold in the finite tower $D/M/K$.
Choose finitely many $K$-generators of $D$ and adjoin to $K$ all their
conjugates in $\Omega$. The resulting splitting field $E$ is finite
Galois over $K$ and contains $D$, hence contains $M$.
For a finite Galois $D/K$, instead choose generators of the finite
compositum $DM/K$ and take the same finite normal closure.
Finally every minimal polynomial over $M$ of an element of $E$ divides
a separable polynomial over $K$ that splits in $E$. It is separable
and splits in $E$, proving that $E/M$ is finite Galois.
\end{proof}

\begin{proposition}[Absolute upper groups after finite base change]
\label{fw:absolute-finite-base-upper}
\dcref{Wintenberger1983, Proposition 1.2.3; absolute upper-group change of parameter}
\source{wintenberger:page-0005}
\uses{fw:finite-tower-upper,fw:apf-base-separable-closure,fw:cofinal-base-galois-layers,fw:cofinal-upper-conditions,fw:finite-base-residue,fw:infinite-upper-isomorphism}
\lean{FontaineWintenberger.mem_upperRamificationGroupInfinite_base_iff, FontaineWintenberger.mem_upperRamificationGroupInfinite_containing_base_iff, FontaineWintenberger.finiteBase_restriction_mem_upper_iff, FontaineWintenberger.mem_upperRamificationGroupInfinite_finiteBase_iff, FontaineWintenberger.absoluteGaloisBaseChange_mem_upper_iff, FontaineWintenberger.absoluteGaloisBaseChange_map_upper}
\leanok
Let $M/K$ be a finite subextension of $\Omega=K^{\mathrm{sep}}$, with its
compatible extended norm. Let $H\subset G_K$ be the subgroup fixing $M$,
and let $\theta:G_M\simeq H$ be the actual separable-closure identification.
For every $u\geq-1$, inclusion of $H$ into $G_K$ gives
\[
 \theta\bigl(G_M^{\psi_{M/K}(u)}\bigr)=H\cap G_K^u.
\]
No normality of $M/K$ or finiteness of the residue field is required.
\end{proposition}
\begin{proof}
The finite field $M$ is complete for its inherited norm, its unit ball
is a discrete valuation ring, and that ring is finite integral over
$\mathcal O_K$. Its residue field is therefore a finite algebraic
extension of $k_K$, hence perfect. Thus the absolute upper-group
construction applies over $M$ with these actual valued structures.
Identify $M^{\mathrm{sep}}$ with $\Omega$ by the preceding closure lemma.
Conjugation carries its finite restriction conditions to the same
conditions on actual finite Galois $M$-subextensions of $\Omega$.

Put $v=\psi_{M/K}(u)$. For $g\in\operatorname{Aut}(\Omega/M)$, membership
in the upper group at $v$ can be tested on the finite Galois overfields
$E/K$ containing $M$. Indeed they are cofinal by the preceding lemma;
restriction of an upper group to a smaller finite Galois $M$-layer is
an upper group by the finite quotient theorem. For each such $E$, the
intrinsic finite-tower theorem gives
\[
 g|_E\in\Gal(E/M)^v
 \quad\Longleftrightarrow\quad
 g|_E\in\Gal(E/K)^u,
\]
because $g$ already fixes $M$. Testing the right condition on these
$E$ is equivalent to testing it on every finite Galois $K$-subextension:
they are cofinal on the $K$ side as well, and finite upper restriction
preserves the parameter. This last condition is precisely $g\in G_K^u$.
Transporting by the actual group identification proves the identity.
\end{proof}

\begin{lemma}[Isomorphisms preserve infinite upper groups]
\label{fw:infinite-upper-isomorphism}
\dcref{Serre1979, Chapter IV, Section 3, Remark 1; elementary isomorphism consequence}
\source{serre-1979-local-fields:page-0082}
\uses{fw:ramification-field-equiv,fw:absolute-upper-projections}
\lean{FontaineWintenberger.finiteGaloisMap, FontaineWintenberger.finiteGaloisMapEquiv, FontaineWintenberger.coe_finiteGaloisMapEquiv, FontaineWintenberger.restrictNormal_finiteGaloisMap, FontaineWintenberger.mem_upperRamificationGroupInfinite_of_autCongr, FontaineWintenberger.mem_upperRamificationGroupInfinite_autCongr_iff, FontaineWintenberger.upperRamificationGroupInfinite_map_autCongr}
\leanok
Let $f:L\simeq M$ be a $K$-algebra isomorphism of algebraic Galois
extensions. Conjugation $s\mapsto fsf^{-1}$ carries the actual infinite
upper group $\Gal(L/K)^u$ onto $\Gal(M/K)^u$ for every $u\geq-1$.
\end{lemma}
\begin{proof}
For each finite Galois subextension $D/K$ of $L$, its image $f(D)$ is
finite Galois over $K$, and $f$ restricts to $e:D\simeq f(D)$.
Ordinary restrictions commute with conjugation:
\[
 (fsf^{-1})|_{f(D)}=e(s|_D)e^{-1}.
\]
The finite isomorphism theorem identifies the two actual upper groups
under this conjugation, using their own integer rings. If $fsf^{-1}$
satisfies every finite upper condition in $M$, its condition on $f(D)$
therefore implies the condition on $D$ for $s$. Since $D$ was arbitrary,
$s$ belongs to the infinite upper group in $L$. Apply the same argument
to $f^{-1}$ for the converse. These implications identify the image
under the actual group isomorphism, proving the assertion.
\end{proof}

\begin{proposition}[A finite-level criterion for Galois APF extensions]
\label{fw:galois-apf-criterion}
\dcref{Wintenberger1983, Definition 1.2.1; Galois reformulation}
\source{wintenberger:page-0005}
\uses{fw:apf,fw:absolute-upper-projections,fw:infinite-upper-isomorphism}
\lean{FontaineWintenberger.isOpen_upperRamificationGroupInfinite_iff_finiteIndex, FontaineWintenberger.apfSubgroup_eq_comap_upperInfinite, FontaineWintenberger.apfInertiaIndex_eq_upperInfiniteRelIndex, FontaineWintenberger.isAPFSubextension_iff_upperInfinite_finiteIndex, FontaineWintenberger.isAPF_iff_upperInfinite_finiteIndex, FontaineWintenberger.isAPF_iff_upperInfinite_isOpen, FontaineWintenberger.apfIndex_eq_upperInfiniteRelIndex}
\leanok
For an algebraic Galois extension $L/K$, put $\Gamma=\Gal(L/K)$ and
let $\Gamma^u$ be the inverse limit of the finite upper groups.
Then $L/K$ is APF if and only if every $\Gamma^u$ is open, equivalently
of finite index, in $\Gamma$. For $u\geq0$,
\[
 [G_K^0:(G_L\cap G_K^0)G_K^u]=[\Gamma^0:\Gamma^u].
\]
Here indices take values in $\mathbb N\cup\{+\infty\}$; the infinite
case means simultaneous infinitude of the two coset spaces.
\end{proposition}
\begin{proof}
The absolute projection theorem identifies $\Gamma^u$ with the image
of $G_K^u$. Thus the inverse image of $\Gamma^u$ under
$G_K\to\Gamma$ is $G_LG_K^u$. A continuous surjection of compact
Hausdorff groups is a quotient map, so this inverse image is open if
and only if $\Gamma^u$ is open. The subgroup $\Gamma^u$ is closed;
its quotient is compact Hausdorff, and is finite if and only if it is
discrete. Hence finite index and openness are equivalent.
Restricting the quotient map to $G_K^0$ gives kernel $G_L\cap G_K^0$,
image $\Gamma^0$, and image of $G_K^u$ equal to $\Gamma^u$.
The same inverse-image argument as for a finite Galois layer gives a
bijection of coset spaces. It proves equality of finite indices and
simultaneous infinitude otherwise.
For an abstract algebraic Galois extension choose an embedding into
$K^{\mathrm{sep}}$. The isomorphism to its image identifies the APF
condition and the intrinsic inertia index. The infinite isomorphism
lemma identifies its actual upper groups with those of the image;
the induced bijections of coset spaces preserve both indices. Thus all
assertions hold for the abstract extension independently of this choice.
\end{proof}

\begin{definition}[Strictly APF extensions]
\label{fw:strict-apf}
\dcref{Wintenberger1983, Definition 1.2.1}
\source{wintenberger:page-0005}
\uses{fw:apf}
\lean{FontaineWintenberger.apfStrictnessRatio, FontaineWintenberger.apfStrictnessConstant, FontaineWintenberger.IsStrictlyAPF, FontaineWintenberger.apfStrictnessConstant_le, FontaineWintenberger.apfStrictnessRatio_algEquiv, FontaineWintenberger.apfStrictnessConstant_algEquiv, FontaineWintenberger.isStrictlyAPF_algEquiv}
\leanok
An APF extension $L/K$ is \emph{strictly APF} if
\[
 \liminf_{u\to+\infty}
   \frac{\psi_{L/K}(u)}{[G_K^0:G_L^0G_K^u]}>0.
\]
For an APF extension with $i(L/K)>0$, define
\[
 c(L/K)=\inf_{u\geq i(L/K)}
    \frac{\psi_{L/K}(u)}{[G_K^0:G_L^0G_K^u]},
\]
with the infimum in $[0,+\infty]$ and the convention $\inf\varnothing=+\infty$.
The lower limit is interpreted in $[0,+\infty]$; its arguments are
nonnegative for all sufficiently large $u$. These data are invariant under
$K$-algebra isomorphisms of extensions: the intrinsic index, integral and
bound are invariant, hence so are their ratio, lower limit and infimum.
\end{definition}

\begin{lemma}[An unbounded first ramification parameter]
\label{fw:apf-unbounded-first-parameter}
\dcref{Wintenberger1983, Definition 1.2.1; extended endpoint}
\source{wintenberger:page-0005}
\uses{fw:strict-apf,fw:apf-extended-bound,fw:apf-herbrand-homeomorphism,fw:apf-index-embedding}
\lean{FontaineWintenberger.apfSubgroup_eq_top_of_bound_eq_top, FontaineWintenberger.isAPF_of_bound_eq_top, FontaineWintenberger.apfIndex_eq_one_of_bound_eq_top, FontaineWintenberger.apfPsi_eq_self_of_bound_eq_top, FontaineWintenberger.apfPhi_eq_self_of_bound_eq_top, FontaineWintenberger.apfStrictnessConstant_of_bound_eq_top, FontaineWintenberger.isStrictlyAPF_of_bound_eq_top, FontaineWintenberger.isStrictlyAPF_self, FontaineWintenberger.apfStrictnessConstant_self}
\leanok
If $i(L/K)=+\infty$, then $H_L(u)=G_K$ for every $u\geq-1$, the
extension is APF, $a_{L/K}(u)=1$ for $u\geq0$, both Herbrand functions
are the identity, $c(L/K)=+\infty$, and the extension is strictly APF.
In particular the trivial extension is strictly APF with constant infinity.
\end{lemma}
\begin{proof}
The set defining $i$ is unbounded. Given $u\geq-1$, choose $v>u$ in
that set. Decreasingness gives $H_L(u)\supset H_L(v)=G_K$. Thus all
these subgroups are the whole group and are open, proving APF.
Intersecting with inertia gives $a(u)=1$ for $u\geq0$. The extended
slope is identically one, so its integral and inverse are the identity.
No real parameter satisfies $u\geq i$, so $c=\inf\varnothing=+\infty$.
The ratio equals $u$ for $u\geq0$, and in particular is at least one
for all $u\geq1$. Its lower limit is positive. Apply the already proved
$i(K/K)=+\infty$ for the final assertion.
\end{proof}

\begin{lemma}[Bounded index growth implies strict APF]
\label{fw:apf-bounded-index-growth}
\uses{fw:strict-apf,fw:apf-herbrand-homeomorphism,fw:apf-index-embedding}
\lean{FontaineWintenberger.IsAPF.window_mul_index_le_apfPsi, FontaineWintenberger.IsAPF.apfPsi_div_index_ge_of_boundedGrowth, FontaineWintenberger.IsAPF.eventually_le_apfStrictnessRatio_of_boundedGrowth, FontaineWintenberger.isStrictlyAPF_of_apfIndex_boundedGrowth}
\leanok
Let $L/K$ be an abstract separable APF extension and let
$a(u)=a_{L/K}(u)$ be its intrinsic inertia index. Suppose there are
real constants $\ell>0$, $C>0$, and $U$ such that
\[
 a(u+\ell)\leq C a(u)\qquad(u\geq U).
\]
Then $L/K$ is strictly APF. More precisely, eventually
$\psi_{L/K}(v)/a(v)\geq\ell/C$.
\end{lemma}
\begin{proof}
For $v\geq\max\{U+\ell,\ell\}$ put $w=v-\ell$.
Then $w\geq0$, $w\geq U$, and $v=w+\ell$.
The actual index is positive and nondecreasing on $[0,\infty)$,
and its Herbrand integral is $\psi(v)=\int_0^v a(t)\,dt$.
The integrand is nonnegative, so
\[
 \psi(v)\geq\int_w^v a(t)\,dt
       \geq\ell a(w).
\]
The growth hypothesis gives $a(v)\leq C a(w)$.
Since $a(v)>0$ and $C>0$, division proves
$\psi(v)/a(v)\geq\ell/C>0$. Thus the nonnegative extended-real
strictness ratio has a positive constant eventual lower bound.
Its lower limit is at least that constant, proving strict APF
with its intrinsic definition, including finite extensions.
\end{proof}

\begin{proposition}[Positive-bound strictness criterion]
\label{fw:strict-apf-positive-bound}
\dcref{Wintenberger1983, Definition 1.2.1; positivity assertion}
\source{wintenberger:page-0005}
\uses{fw:strict-apf,fw:apf-unbounded-first-parameter,fw:apf-extended-bound,fw:apf-herbrand-homeomorphism,fw:apf-index-embedding}
\lean{FontaineWintenberger.isStrictlyAPF_of_constant_pos_of_bound_eq_coe, FontaineWintenberger.IsStrictlyAPF.constant_pos_of_bound_eq_coe, FontaineWintenberger.isStrictlyAPF_iff_constant_pos}
\leanok
Suppose $L/K$ is APF and $i(L/K)>0$. Then it is strictly APF if and
only if $c(L/K)>0$. If $i(L/K)=+\infty$, then $a_{L/K}(u)=1$ for
$u\geq0$, both Herbrand functions are the identity, $c(L/K)=+\infty$,
and the extension is strictly APF. In particular these conclusions
hold for $K/K$.
\end{proposition}
\begin{proof}
Write $a=a_{L/K}$, $P=\psi_{L/K}$ and $i=i(L/K)$.
If $i=+\infty$, apply Lemma~\ref{fw:apf-unbounded-first-parameter}.

Otherwise $i$ is a positive real number. If $c>0$, the defining
infimum gives $P(u)/a(u)\geq c$ for every $u\geq i$, so its lower
limit is positive. Conversely, strictness supplies $\delta>0$ and
$T\geq i$ such that $P(u)/a(u)\geq\delta$ for all $u\geq T$.
For $i\leq u\leq T$, positivity and monotonicity of the actual finite
indices, and $P(u)\geq u$, give
\[
 \frac{P(u)}{a(u)}\geq\frac{i}{a(T)}>0.
\]
Hence every term in the infimum for $c$ is at least
$\min(\delta,i/a(T))>0$. This proves $c>0$.
\end{proof}

\begin{proposition}[Wild ramification and the first APF bound]
\label{fw:apf-wild-ramification}
\dcref{Wintenberger1983, Definition 1.2.1; totally ramified $p$-extension clause}
\source{wintenberger:page-0005}
\uses{fw:apf}
\notready
An APF extension is a totally ramified $p$-extension if and only if
$i(L/K)>0$.
\end{proposition}
\begin{proof}
Let $P_K$ be wild inertia. In a finite Galois quotient, the positive
upper groups for sufficiently small positive parameters equal wild
inertia, a $p$-group, and the tame inertia quotient has order prime to
$p$. A totally ramified $p$-power-degree intermediate field has wild
inertia transitive on its embeddings: inertia is transitive by total
ramification, and normality of wild inertia makes all its orbits have
the same $p$-power size. Their number is therefore a $p$-power and also
divides the prime-to-$p$ tame quotient. Thus there is one orbit. Conversely, transitivity of wild
inertia implies total ramification and $p$-power degree.

Apply this criterion to all finite subextensions of $L$. It says that
$L/K$ is a totally ramified $p$-extension exactly when $G_LP_K=G_K$.
For the converse implication here, use the orbit criterion in a finite
Galois normal closure of each subextension; for the forward implication,
compactness and separation by finite Galois quotients give the equality.
If $G_LG_K^u=G_K$ for some $u>0$, containment $G_K^u\subset P_K$ gives
$G_LP_K=G_K$. Conversely, choose any $\varepsilon>0$; APF makes
$H_L(\varepsilon)$ open. Pass to the finite Galois quotient by its
normal core. The image of $G_LP_K$ is the whole quotient, and positive
upper groups with sufficiently small positive parameter have the same
image as $P_K$. Choose such a parameter with $0<u\leq\varepsilon$.
The normal core lies in $H_L(\varepsilon)\subset H_L(u)$, so surjectivity
of its quotient image gives $H_L(u)=G_K$. The admissible set
is decreasing in its parameter, so this is equivalent to its supremum
being positive.
\end{proof}

\begin{lemma}[Finite-index restriction and normal products]
\label{fw:finite-index-base-restriction}
\dcref{Wintenberger1983, Proposition 1.2.3; group-index calculation}
\source{wintenberger:page-0005}
\lean{Subgroup.finiteIndex_iff_comap_of_range, Subgroup.comap_sup_eq_of_normal_right_of_le_range}
\leanok
Let $f:G'\to G$ be a group homomorphism. If its image $H$ has finite
index, a subgroup $J\subset G$ has finite index if and only if
$f^{-1}(J)$ has finite index in $G'$. If $A\subset H$ and $U$ is a
normal subgroup of $G$, then
\[
 f^{-1}(AU)=f^{-1}(A)f^{-1}(U).
\]
\end{lemma}
\begin{proof}
The cosets of $f^{-1}(J)$ in $G'$ correspond to the cosets of
$J\cap H$ in $H$. Finite index of $J$ implies finiteness of this
coset set. Conversely, if it is finite, multiplication of indices
with the finite index of $H$ makes $J\cap H$ have finite index in
$G$, and its containing subgroup $J$ has finite index as well.
For the product identity, one inclusion follows from the homomorphism
law. If $f(g)=au$ with $a\in A$, choose $a'\in G'$ with $f(a')=a$,
using $A\subset H$. Then $f(a'^{-1}g)=u$, so
$g=a'(a'^{-1}g)$ belongs to the right product. Normality of $U$ and
its preimage ensures that the products are subgroups.
\end{proof}

\begin{proposition}[APF under a finite change of the base field]
\label{fw:apf-finite-base}
\dcref{Wintenberger1983, Proposition 1.2.3(i); APF assertion}
\source{wintenberger:page-0005}
\uses{fw:absolute-finite-base-upper,fw:finite-index-base-restriction,fw:apf,fw:apf-herbrand-homeomorphism}
\lean{FontaineWintenberger.baseChangedSeparableEmbedding, FontaineWintenberger.le_baseChangedSeparableEmbedding_fieldRange, FontaineWintenberger.comap_fixingSubgroup_baseChangedEmbedding, FontaineWintenberger.absoluteGaloisBaseChange_range, FontaineWintenberger.absoluteGaloisBaseChange_range_finiteIndex, FontaineWintenberger.comap_absoluteUpper_finiteBase, FontaineWintenberger.comap_apfSubgroup_finiteBase, FontaineWintenberger.isAPF_finiteBase_iff}
\leanok
Let $M/K$ be a finite subextension of $K^{\mathrm{sep}}$ with its
compatible norm. For an arbitrary abstract separable extension $N/K$
and a scalar tower $K\to M\to N$, one has
\[
 N/K\text{ is APF}\quad\Longleftrightarrow\quad N/M\text{ is APF}.
\]
Neither $N/K$ nor $M/K$ is required to be normal.
\end{proposition}
\begin{proof}
Choose an $M$-embedding $i:N\to M^{\mathrm{sep}}$ and let $j:N\to
K^{\mathrm{sep}}$ be its composite with the actual closure
identification. Its image contains $M$. Write $\rho:G_M\to G_K$ for
the resulting injective homomorphism, with image $H$ fixing $M$.
An automorphism fixes every $i(x)$ if and only if its image under
$\rho$ fixes every $j(x)$, by the actual equivariance of the closure
identification. Thus the fixing subgroup over $M$ is the preimage of
the fixing subgroup over $K$.

At $v=\psi_{M/K}(u)$, absolute base change also identifies the new upper
group with $\rho^{-1}(G_K^u)$. The old fixing subgroup lies in $H$.
The normal-product preimage identity therefore identifies the new APF
product with the preimage of the old one. Since $M/K$ is finite, $H$
has finite index. The finite-index preimage criterion makes finiteness
of the two APF-product indices equivalent. Finally $\psi_{M/K}$ maps
$[-1,+\infty)$ bijectively onto itself, and APF can be tested using any
specified separable embedding. Quantifying over the parameters proves
the equivalence for the public abstract extension predicates.
\end{proof}

\begin{lemma}[The actual finite-extension fixing index]
\label{fw:finite-extension-fixing-index}
\dcref{Wintenberger1983, Proposition 1.2.3(ii); Galois correspondence}
\source{wintenberger:page-0005}
\uses{fw:profinite-galois-foundations}
\lean{FontaineWintenberger.fixingSubgroup_relIndex_eq_finrank, FontaineWintenberger.fixingSubgroup_relIndex_pos}
\leanok
Let $\Omega/K$ be an algebraic Galois extension. For arbitrary
$M\subset N\subset\Omega$, the relative index
of $G_N$ in $G_M$ equals $[N:M]$, with both infinite cardinalities
represented by zero when taking their natural cardinal. In particular,
a finite extension $N/M$ gives a positive finite relative index.
\end{lemma}
\begin{proof}
The same $\Omega$ is algebraic Galois over $M$.
Restriction of scalars identifies its actual automorphism group over
$M$ with $G_M$. Under this group isomorphism, the subgroup fixing the
actual intermediate field $N/M$ is exactly $G_N$, since both fixing
conditions require each element of the same carrier $N$ to be fixed.
The Galois correspondence identifies the cosets of this fixing subgroup
with the $M$-embeddings of $N$ into $\Omega$. If $N/M$ is finite,
separability makes their number exactly $[N:M]$, a positive integer.
If $N/M$ is infinite, it contains finite subextensions of arbitrarily
large degree: otherwise the dimensions of all spans of finite subsets
would be bounded. Restriction onto the embeddings of each such finite
subextension is surjective, since its embeddings extend to automorphisms of the common
algebraic normal extension $\Omega/M$. Thus the coset set is infinite as well.
Its natural cardinal and the natural-valued field dimension are then
both zero, which proves the stated identity in this case too.
\end{proof}

\begin{lemma}[Relative indices of normal products]
\label{fw:normal-product-relative-index}
\dcref{Wintenberger1983, Proposition 1.2.3; relative-index calculation}
\source{wintenberger:page-0005}
\lean{Subgroup.relIndex_sup_sup_normal, Subgroup.sup_inf_eq_sup_inf_of_le_of_normal, Subgroup.relIndex_sup_sup_eq_inf}
\leanok
Let $A,B,U$ be subgroups of a group $G$, with $U$ normal. Then
\[
 [BU:AU\cap BU]=[B:AU\cap B].
\]
If $A\subset B$, then
\[
 (AU)\cap B=A(U\cap B),\qquad
 [BU:AU]=[B:A(U\cap B)].
\]
These statements include infinite indices.
\end{lemma}
\begin{proof}
Pass to the actual quotient $q:G\to G/U$. The preimage of $q(A)$ is
$AU$, and that of $q(B)$ is $BU$. The coset sets in the first displayed
equality are both identified with the cosets of $q(A)\cap q(B)$ in
$q(B)$, by the canonical quotient maps. This proves the first claim
without a finiteness assumption. If $A\subset B$, an element $au\in B$
has $u=a^{-1}(au)\in U\cap B$. This proves the nontrivial inclusion
in the intersection identity; the other is immediate. Substitute this
identity into the first equality to obtain the last formula.
\end{proof}

\begin{lemma}[The nested APF index factor]
\label{fw:apf-nested-index-factor}
\dcref{Wintenberger1983, Proposition 1.2.3; inertia-index factorization}
\source{wintenberger:page-0005}
\uses{fw:normal-product-relative-index,fw:apf-inertia-indices,fw:apf-index-embedding}
\lean{FontaineWintenberger.apfTowerIndexFactor, FontaineWintenberger.apfTowerIndexFactor_zero, FontaineWintenberger.apfInertiaSubgroup_antitone_extension, FontaineWintenberger.apfInertiaIndex_eq_mul_towerIndexFactor, FontaineWintenberger.apfTowerIndexFactor_eq_inertiaRelIndex, FontaineWintenberger.IsAPFSubextension.towerIndexFactor_pos, FontaineWintenberger.IsAPFSubextension.towerIndexFactor_monotoneOn, FontaineWintenberger.apfIndex_eq_mul_towerIndexFactor}
\leanok
Let $M\subset N\subset K^{\mathrm{sep}}$, and suppose $N/K$ is APF.
Put $I=G_K^0$, $A_0=G_N\cap I$, $B_0=G_M\cap I$, and $U_u=G_K^u$.
For $u\geq0$ define
\[
 d(u)=[B_0:A_0(U_u\cap B_0)].
\]
This is a positive finite integer, $d(0)=1$, it is nondecreasing, and
\[
 a_{N/K}(u)=a_{M/K}(u)d(u).
\]
Equivalently, $d(u)$ is the relative index of $A_0U_u$ in $B_0U_u$.
\end{lemma}
\begin{proof}
The preceding normal-product index lemma gives the equivalent formula.
The chain $A_0U_u\subset B_0U_u\subset I$ and multiplication of
indices give the displayed factorization. The APF property makes the
index of $A_0U_u$ in $I$ finite and positive. Hence both factors are
finite and positive. At $u=0$, both product groups are $I$, so $d(0)=1$.
Finally $U_u$ decreases with $u$. Therefore $A_0(U_u\cap B_0)$ also
decreases as a subgroup of $B_0$. Its finite positive index is
nondecreasing, which proves the asserted monotonicity.
\end{proof}

\begin{proposition}[Strict APF descent and its quantitative bounds]
\label{fw:strict-apf-descent}
\dcref{Wintenberger1983, Proposition 1.2.3(iii)--(iv); subextension assertions}
\source{wintenberger:page-0005}
\uses{fw:apf-nested-index-factor,fw:strict-apf,fw:apf-herbrand-homeomorphism,fw:apf-extended-bound}
\lean{FontaineWintenberger.apfPsi_le_towerIndexFactor_mul, FontaineWintenberger.apfStrictnessRatio_le_of_le, FontaineWintenberger.IsStrictlyAPF.of_le, FontaineWintenberger.apfRamificationBound_le_of_le, FontaineWintenberger.apfStrictnessConstant_le_of_le, FontaineWintenberger.IsStrictlyAPF.of_algHom, FontaineWintenberger.apfRamificationBound_le_of_algHom, FontaineWintenberger.apfStrictnessConstant_le_of_algHom}
\leanok
For $M\subset N\subset K^{\mathrm{sep}}$ with $N/K$ APF, and $u\geq0$,
\[
 \psi_{N/K}(u)\leq d(u)\psi_{M/K}(u),\qquad
 \frac{\psi_{N/K}(u)}{a_{N/K}(u)}
 \leq\frac{\psi_{M/K}(u)}{a_{M/K}(u)}.
\]
Consequently strict APF descends along every embedding of abstract
separable extensions $M\hookrightarrow N$ over $K$.
Moreover $i(N/K)\leq i(M/K)$, and if $N/K$ is APF and $i(N/K)>0$,
then $c(N/K)\leq c(M/K)$.
\end{proposition}
\begin{proof}
For $0\leq v\leq u$, the preceding factorization and monotonicity give
$a_N(v)=a_M(v)d(v)\leq a_M(v)d(u)$. Integrate over $[0,u]$ to get
the first inequality. All indices are positive, and
$a_N(u)=a_M(u)d(u)$, so division gives the second. Taking nonnegative
extended-real lower limits preserves this inequality on the tail.
Strict positivity of the larger-extension lower limit therefore gives
strict positivity for the subextension, whose APF property already
descends.

Since $G_N\subset G_M$, every parameter for which $G_NG_K^u=G_K$
is also one for which $G_MG_K^u=G_K$. Inclusion of the defining
admissible sets proves the bound on $i$, including infinite suprema.
If $i(N/K)>0$, every real $u\geq i(M/K)$ is nonnegative and lies
also above $i(N/K)$. At each such $u$ the second inequality bounds the
ratio for $M/K$ below by $c(N/K)$. Taking the extended infimum proves
the bound on $c$; an empty domain has value $+\infty$.
For abstract fields choose a separable-closure embedding of $N$ and
compose it with $M\hookrightarrow N$. Their images are nested actual
subfields. The proved invariance of the functions, bounds and strict
predicate under the two image isomorphisms transfers all conclusions.
\end{proof}

\begin{proposition}[APF under finite extension above an arbitrary subfield]
\label{fw:apf-finite-extension}
\dcref{Wintenberger1983, Proposition 1.2.3(ii)}
\source{wintenberger:page-0005}
\uses{fw:finite-extension-fixing-index,fw:normal-product-relative-index,fw:apf-nested-index-factor,fw:strict-apf-descent}
\lean{FontaineWintenberger.finiteDimensional_fieldRange_extendScalars, FontaineWintenberger.isAPFSubextension_finiteExtension_iff, FontaineWintenberger.inertiaFixingRelIndex_pos, FontaineWintenberger.apfTowerIndexFactor_le_inertiaFixingRelIndex, FontaineWintenberger.apfPsi_mono_extension, FontaineWintenberger.apfStrictnessRatio_le_finiteExtension_mul, FontaineWintenberger.isStrictlyAPF_finiteExtension_iff, FontaineWintenberger.isAPF_finite_extension_iff, FontaineWintenberger.isStrictlyAPF_finite_extension_iff}
\leanok
Let $M\subset N\subset K^{\mathrm{sep}}$ and $[N:M]<\infty$.
Then $N/K$ is APF, respectively strictly APF, if and only if $M/K$ is.
This also holds for arbitrary abstract separable fields in a finite
extension tower $K\subset M\subset N$.
More precisely, put $I=G_K^0$ and
$C=[G_M\cap I:G_N\cap I]$. This is a positive finite integer,
$d(u)\leq C$ for $u\geq0$, and whenever $N/K$ is APF,
\[
 \psi_M(u)\leq\psi_N(u),\qquad R_M(u)\leq C R_N(u).
\]
\end{proposition}
\begin{proof}
Write $A=G_N$, $B=G_M$, $U=G_K^u$. The actual fixing-degree identity
gives $[B:A]=[N:M]<\infty$. The normal-product coset identity gives
$[BU:AU]\leq[B:A]$. Hence $AU$ has finite index in $G_K$ if and only
if $BU$ does. The actual products are closed, so finite index is
equivalent to openness; this proves the APF equivalence.

Intersection with $I$ injects the coset set of $A\cap I$ in
$B\cap I$ into the cosets of $A$ in $B$. Thus $C$ is positive and
finite. The intrinsic expression
$d(u)=[B\cap I:(A\cap I)(U\cap B\cap I)]$ gives $d(u)\leq C$.
Since also $d(u)\geq1$, the actual factorization
$a_N(u)=a_M(u)d(u)$ gives $a_M\leq a_N\leq C a_M$.
Integrating the first inequality gives $\psi_M\leq\psi_N$.
All indices are positive, so division yields
$R_M\leq C R_N$. If the lower limit of $R_M$ is positive, the lower
limit of $R_N$ is positive: otherwise the last inequality and
commutation of the lower limit with multiplication by the finite
constant $C$ force the former lower limit to be zero. The reverse
strict implication is the preceding descent theorem.

For an abstract finite tower choose a separable embedding of $N$ into
$K^{\mathrm{sep}}$ and restrict it to $M$. Their image fields form an
actual nested pair. The two image isomorphisms give a semilinear
bijection from $N$ to its image, whose scalar map is the actual
isomorphism from $M$ onto its image. It carries a finite spanning set
to a finite spanning set, so the image extension is finite.
Apply the proved assertions to those image fields and transport APF
and strict APF through the actual image isomorphisms.
\end{proof}

\begin{lemma}[The actual finite-base APF index identity]
\label{fw:apf-finite-base-index}
\dcref{Wintenberger1983, Proposition 1.2.3(i); inertia-index comparison}
\source{wintenberger:page-0005}
\uses{fw:apf-finite-base,fw:absolute-finite-base-upper,fw:apf-nested-index-factor}
\lean{FontaineWintenberger.comap_apfInertiaSubgroup_finiteBase, FontaineWintenberger.apfIndex_finiteBase_eq_towerIndexFactor, FontaineWintenberger.apfIndex_eq_mul_finiteBase}
\leanok
Let $M/K$ be finite and let $N$ be any abstract separable extension in
a tower $K\subset M\subset N$. For $u\geq0$,
\[
 a_{N/K}(u)=a_{M/K}(u)a_{N/M}(\psi_{M/K}(u)).
\]
No APF or normality assumption on $N$ is required for this identity of
natural-valued indices; an infinite index has value zero.
\end{lemma}
\begin{proof}
Choose an actual embedding of $N$ into the separable closure of $M$
and transport it through the actual isomorphism of separable closures.
Its image $D$ in $K^{\mathrm{sep}}$ contains $M$.
Let $\rho:G_M\to G_K$ be the actual restriction homomorphism, whose
range is the subgroup $B$ fixing $M$. Put $I=G_K^0$ and
$A=G_D\subset B$. The proved upper-group base change at $u=0$
identifies the new inertia group with $B\cap I$, and at $u$ identifies
the new upper group with $B\cap G_K^u$ at parameter
$\psi_{M/K}(u)$. The actual fixing-subgroup preimage is the subgroup
fixing the chosen image of $N$. Therefore the new inertia product is
the preimage of $(A\cap I)G_K^u$. Relative indices under $\rho$
identify its index in new inertia with
\[
 [(B\cap I):((A\cap I)G_K^u)\cap B\cap I]=d(u).
\]
The normal-product index identity gives the equality with the intrinsic
tower factor. The previously proved factorization of the two old-base
indices now gives the displayed result. All steps identify actual
coset sets and use natural cardinal identities, so no finiteness
assumption on these indices is needed. Finally the public index is
independent of the chosen actual separable embedding, giving the
claimed abstract-field statement.
\end{proof}

\begin{proposition}[The actual finite-base Herbrand composition]
\label{fw:apf-finite-base-herbrand}
\dcref{Wintenberger1983, Proposition 1.2.3(i); integration of the inertia identity}
\source{wintenberger:page-0005}
\uses{fw:apf-finite-base-index,fw:apf-finite-base,fw:apf-herbrand-homeomorphism}
\lean{FontaineWintenberger.IsAPF.hasDerivAt_apfPsi, FontaineWintenberger.IsAPF.exists_apfPsi_derivative_countable, FontaineWintenberger.IsAPF.exists_finiteBase_comp_derivative_countable, FontaineWintenberger.IsAPF.apfPsi_finiteBase_comp}
\leanok
Let $M/K$ be finite and let $N/K$ be any abstract APF separable
extension in a tower $K\subset M\subset N$. Then, on the whole real
line with the identity continuation below zero,
\[
 \psi_{N/K}=\psi_{N/M}\circ\psi_{M/K}.
\]
\end{proposition}
\begin{proof}
All three extensions are APF, by finite-base APF invariance and finite
APF. Their slope functions, continued by one below zero, are
nondecreasing and hence locally integrable. Such a monotone real
function is continuous except on a countable set. At each point where
a slope is continuous, the fundamental theorem of calculus gives the
derivative of its actual integral equal to that slope.

Put $f=\psi_{M/K}$ and $g=\psi_{N/M}$. Let $S_M$ and $S_{N/M}$ be
the respective countable sets of discontinuities of their slopes.
Since $f$ is injective, $S=S_M\cup f^{-1}(S_{N/M})$ is countable.
For every $u>0$ outside $S$, both required derivatives exist and the
chain rule, followed by the actual finite-base index identity, gives
\[
 (g\circ f)'(u)=a_{N/M}(f(u))a_{M/K}(u)=a_{N/K}(u).
\]
For each $t\geq0$, the composite is continuous on $[0,t]$ and the
right-hand slope is integrable there. The fundamental theorem of
calculus with a countable exceptional set therefore identifies its
increment with $\int_0^t a_{N/K}(u)\,du$. The composite vanishes at
zero, so this is exactly $\psi_{N/K}(t)$. For $t\leq0$, all three
functions are the identity, which proves the equality everywhere.
\end{proof}

\begin{lemma}[Finite-base indices stabilize and their integral is affine]
\label{fw:apf-finite-slope}
\dcref{Wintenberger1983, Proposition 1.2.3(i); finite-index bound}
\source{wintenberger:page-0005}
\uses{fw:apf,fw:apf-inertia-indices,fw:apf-index-embedding,fw:apf-herbrand-homeomorphism}
\lean{FontaineWintenberger.fixingSubgroup_inertiaRelIndex_pos, FontaineWintenberger.apfIndex_le_fixingSubgroup_inertiaRelIndex, FontaineWintenberger.exists_apfIndex_finite_constant, FontaineWintenberger.exists_apfPsi_finite_affine}
\leanok
Let $M/K$ be finite inside $K^{\mathrm{sep}}$, and put
$I=G_K^0$ and $C_0=[I:G_M\cap I]$.
Then $C_0$ is a positive finite integer, and for $u\geq0$,
$1\leq a_{M/K}(u)\leq C_0$. There are a positive integer $C$ and a
real $b\geq0$ such that for every $u\geq b$,
\[
 a_{M/K}(u)=C,\qquad
 \psi_{M/K}(u)=\psi_{M/K}(b)+C(u-b).
\]
\end{lemma}
\begin{proof}
The actual fixing subgroup $G_M$ is open and of finite index. The
inertia cosets inject into its cosets in $G_K$, so $C_0$ is finite and
positive. The product $(G_M\cap I)G_K^u$ contains $G_M\cap I$ inside
$I$, giving the upper bound on its index; APF gives the lower bound.
The function $a_{M/K}$ is nondecreasing. Its positive integer values
lie in the finite set $\{1,\ldots,C_0\}$, so its range on $[0,\infty)$
has a largest value $C$, attained at some $b\geq0$. Monotonicity
forces its value to remain $C$ at every $u\geq b$.
Finally split the defining integral at $b$. On $[b,u]$ its actual
integrand is the constant $C$, whose integral is $C(u-b)$.
This proves the displayed affine formula.
\end{proof}

\begin{proposition}[Strict APF is invariant under finite base change]
\label{fw:strict-apf-finite-base}
\dcref{Wintenberger1983, Proposition 1.2.3(i)}
\source{wintenberger:page-0005}
\uses{fw:apf-finite-base,fw:apf-finite-base-index,fw:apf-finite-base-herbrand,fw:apf-finite-slope,fw:strict-apf}
\lean{FontaineWintenberger.IsAPF.liminf_comp_apfPsi, FontaineWintenberger.IsAPF.apfStrictnessRatio_finiteBase, FontaineWintenberger.IsAPF.exists_liminf_finiteBase_mul, FontaineWintenberger.isStrictlyAPF_finiteBase_iff}
\leanok
Let $M/K$ be finite and $N$ any abstract separable extension in a tower
$K\subset M\subset N$. Then $N/K$ is strictly APF if and only if $N/M$
is strictly APF. When $N/K$ is APF, for $u\geq0$,
\[
 R_{N/M}(\psi_{M/K}(u))=a_{M/K}(u)R_{N/K}(u).
\]
There is a positive finite integer $C$ such that
\[
 \liminf_{v\to\infty}R_{N/M}(v)
   =C\liminf_{u\to\infty}R_{N/K}(u)
\]
as an identity of nonnegative extended real numbers.
\end{proposition}
\begin{proof}
APF is already invariant under the finite base change. Assume that
property holds. Herbrand composition identifies the two numerator
integrals, and the actual index identity identifies their positive
denominators. Division gives the displayed pointwise ratio identity.

The finite-base index has a positive eventual constant value $C$.
The ratio identity is therefore eventually
$R_{N/M}\circ\psi_{M/K}=C R_{N/K}$.
The function $\psi_{M/K}$ is an increasing order isomorphism of the
real line, so it carries the filter at positive infinity onto itself.
Consequently composition with it leaves the lower limit unchanged.
Taking lower limits of the eventual identity and commuting the lower
limit with multiplication by the positive finite constant $C$ gives
the claimed extended-real identity. Multiplication by such a constant
preserves strict positivity, including the value $+\infty$. Thus the
two strictness conditions are equivalent. If APF fails, it fails on
both sides and both strict APF predicates are false.
\end{proof}

\begin{proposition}[Finite base change with the original compatible norm]
\label{fw:abstract-normed-finite-base}
\source{wintenberger:page-0005}
\uses{fw:apf-finite-base,fw:absolute-finite-base-upper,fw:apf-finite-base-index,fw:apf-finite-base-herbrand,fw:strict-apf-finite-base,fw:spectral-base-integers}
\lean{FontaineWintenberger.isAPF_finiteBase_iff_of_normedAlgebra, FontaineWintenberger.apfIndex_eq_mul_finiteBase_of_normedAlgebra, FontaineWintenberger.IsAPF.apfPsi_finiteBase_comp_of_normedAlgebra, FontaineWintenberger.IsAPF.apfStrictnessRatio_finiteBase_of_normedAlgebra, FontaineWintenberger.isStrictlyAPF_finiteBase_iff_of_normedAlgebra}
\leanok
Let $M/K$ be a finite separable extension equipped with its original
compatible complete norm, and let $K\subseteq M\subseteq N$ be any
abstract separable field tower. The actual integer and residue fields
of $M$ carry their original scalar maps. Then $N/K$ is APF if and
only if $N/M$ is APF, and the same equivalence holds for strict APF.
For $u\geq0$,
\[
 a_{N/K}(u)=a_{M/K}(u)a_{N/M}(\psi_{M/K}(u)).
\]
If $N/K$ is APF, then on all real parameters
$\psi_{N/M}\circ\psi_{M/K}=\psi_{N/K}$, and for $u\geq0$,
\[
 R_{N/M}(\psi_{M/K}(u))=a_{M/K}(u)R_{N/K}(u).
\]
\end{proposition}
\begin{proof}
Embed the original $M$ into $K^{\rm sep}$ and use this embedding
as its scalar map on that closure. Since $M/K$ is finite separable,
$K^{\rm sep}$ is a separable closure of $M$. Its unique extension
norm restricts to the original compatible norm on $M$. The actual
isomorphism $M^{\rm sep}\simeq K^{\rm sep}$ over $M$ therefore
preserves all finite-layer integer rings and their ramification groups.
Every finite Galois $M$-subfield of the common closure lies in a
finite Galois $K$-subfield containing the actual image of $M$:
restrict its finite extension to $K$ and take its finite normal closure.
Thus these old-base normal fields test all new-base upper conditions.
The finite upper-group change of base gives, for $u\geq-1$,
\[
 \rho^{-1}(G_K^u)=G_M^{\psi_{M/K}(u)},
\]
where $\rho:G_M\hookrightarrow G_K$ has finite-index image equal
to the actual subgroup fixing the embedded $M$.

Embed $N$ in $M^{\rm sep}$ and transport it through this common
closure. Equivariance identifies the preimage of its old fixing
subgroup with its new fixing subgroup. Hence the new APF product
is the preimage of the old APF product at the displayed parameter.
Finite-index range of $\rho$ makes one product have finite index
exactly when the other does. Since $\psi_{M/K}$ carries
$[-1,\infty)$ bijectively onto itself, this proves APF equivalence.
At zero the same preimage identifies the actual inertia groups.
The inertia coset-index identity then gives the displayed factorization.
Integrating it by the chain rule outside the countable discontinuity
sets of the monotone slopes, and using continuity and value zero at
zero, gives the Herbrand composition. Below zero all functions are
the identity. Dividing the composition by the positive index factors
gives the ratio identity. Finally $a_{M/K}$ is eventually a positive
finite constant; composition by its Herbrand order isomorphism
preserves lower limits. Taking lower limits in the ratio identity
therefore multiplies the strictness lower limit by that positive
constant, proving strict APF equivalence.
\end{proof}

\begin{lemma}[The actual eventual ramification slope of a finite base]
\label{fw:finite-base-eventual-ramification-slope}
\uses{fw:finite-galois-overfield-apf,fw:inertia-relative-index,fw:finite-extension-e-d-invariance,fw:spectral-base-integers,fw:apf-finite-slope}
\lean{FontaineWintenberger.intermediate_spectral_ramificationIdx_eq_relIndex, FontaineWintenberger.exists_intermediate_apfIndex_eq_ramificationIdx, FontaineWintenberger.exists_apfIndex_eq_spectral_ramificationIdx, FontaineWintenberger.exists_apfPsi_spectral_ramificationIdx_affine, FontaineWintenberger.exists_apfPsi_spectral_ramificationIdx_shift, FontaineWintenberger.normed_ramificationIdx_eq_spectral, FontaineWintenberger.exists_apfIndex_eq_normed_ramificationIdx, FontaineWintenberger.exists_apfPsi_normed_ramificationIdx_affine, FontaineWintenberger.exists_apfPsi_normed_ramificationIdx_shift}
\leanok
Let $M/K$ be any finite separable extension with its original
compatible norm and actual integer rings. Write $e=e(M/K)$.
There exists $U\geq0$ such that for all real $u\geq U$,
\[
 a_{M/K}(u)=e,\qquad
 \psi_{M/K}(u)=\psi_{M/K}(U)+e(u-U).
\]
No normality or total ramification of $M/K$ is assumed.
\end{lemma}
\begin{proof}
Embed $M$ into a separable closure and place its actual image
$M'$ in the finite normal closure $\Omega/K$. Write
$G=\operatorname{Gal}(\Omega/K)$, $H=\operatorname{Gal}(\Omega/M')$
and $I=G_0$. Choose $U\geq0$ after the last upper break of
$\Omega/K$. For every $u\geq U$ the actual upper group is trivial,
so the finite-overfield inertia-relative formula gives
$a_{M'/K}(u)=H\mathbin{\mathrm{relIndex}} I$.
The actual fixed-field inertia formula identifies this relative
index with $e(M'/K)$ on the restricted spectral integer ring.
Uniqueness of the extension integer ring identifies that index
with its own spectral ramification index. The original field
isomorphism transports it to $e(M/K)$. A compatible original
norm has the same spectral unit ball, so this is precisely the
original ramification index in the statement. Invariance of the
intrinsic APF index gives $a_{M/K}(u)=e$ for every $u\geq U$.
Integrate this actual constant slope from $U$ to $u$ and use
additivity of the defining integral to obtain the displayed
formula for $\psi$. This includes the trivial extension, with
$e=1$, and requires no assertion about the residue degree.
\end{proof}

\begin{lemma}[The extended Herbrand order isomorphism]
\label{fw:apf-herbrand-extended}
\dcref{Wintenberger1983, Proposition 1.2.3(iv); infinite endpoints}
\source{wintenberger:page-0005}
\uses{fw:apf-herbrand-homeomorphism}
\lean{FontaineWintenberger.apfHerbrandERealOrderIso, FontaineWintenberger.apfHerbrandERealOrderIso_coe, FontaineWintenberger.apfHerbrandERealOrderIso_top, FontaineWintenberger.apfHerbrandERealOrderIso_bot, FontaineWintenberger.apfHerbrandERealOrderIso_symm_coe, FontaineWintenberger.IsAPF.le_apfHerbrandERealOrderIso}
\leanok
For an APF abstract separable extension $L/K$, the actual function
$\psi_{L/K}$ extends to an order isomorphism $\overline\psi_{L/K}$ of
the extended real line by fixing $-\infty$ and $+\infty$.
Its inverse agrees with $\varphi_{L/K}$ at every finite real input,
and $x\leq\overline\psi_{L/K}(x)$ for every extended real $x$.
\end{lemma}
\begin{proof}
Extend the proved real order isomorphism first by a greatest point and
then by a least point. Extending its inverse in the same manner gives
the inverse, since the original two maps are inverse on the real
line and both added endpoints are fixed. The inequality on real inputs
is the proved bound $u\leq\psi_{L/K}(u)$; at the two endpoints it is
equality. These cases exhaust the extended real line.
\end{proof}

\begin{proposition}[The extended finite-base ramification and strictness bounds]
\label{fw:apf-finite-base-bounds}
\dcref{Wintenberger1983, Proposition 1.2.3(iv); finite-base assertions}
\source{wintenberger:page-0005}
\uses{fw:apf-herbrand-extended,fw:apf-finite-base,fw:absolute-finite-base-upper,fw:strict-apf-finite-base,fw:apf-extended-bound}
\lean{FontaineWintenberger.apfSubgroup_finiteBase_eq_top, FontaineWintenberger.apfHerbrandERealOrderIso_bound_le, FontaineWintenberger.apfRamificationBound_le_finiteBase, FontaineWintenberger.apfStrictnessConstant_le_finiteBase}
\leanok
For finite $M/K$ and an abstract separable tower $K\subset M\subset N$,
\[
 i(N/M)\geq\overline\psi_{M/K}(i(N/K))\geq i(N/K).
\]
If $N/K$ is APF and $i(N/K)>0$, then $c(N/M)\geq c(N/K)$.
The bounds include infinite endpoints and empty infimum domains.
\end{proposition}
\begin{proof}
Choose an actual separable embedding of $N$ over $M$ and transport it
to the common closure over $K$. The proved APF-product preimage
identity shows that whenever $G_NG_K^u=G_K$, the new-base APF product
at $\psi_{M/K}(u)$ is the whole new absolute group. This parameter is
at least $-1$, so it is admissible for $i(N/M)$. The extended order
isomorphism preserves suprema. Apply it to the defining supremum of
the old admissible set and use the just-proved inclusion of its image
in the new admissible set. This gives the first inequality, while the
extended lower bound on $\overline\psi$ gives the second.

For the strictness bound, consider any real $v\geq i(N/M)$ and put
$u=\varphi_{M/K}(v)$. Applying the inverse extended order isomorphism
to the first bound gives $u\geq i(N/K)>0$. The pointwise ratio identity
and the positive integer index then give
\[
 R_{N/M}(v)=a_{M/K}(u)R_{N/K}(u)
            \geq R_{N/K}(u)\geq c(N/K).
\]
Take the defining infimum over all such $v$. If no such real $v$
exists, this is the empty extended infimum $+\infty$, and the same
infimum argument proves the assertion without a separate convention.
\end{proof}

\begin{proposition}[Inheritance of APF and strict APF]
\label{fw:apf-inheritance}
\dcref{Wintenberger1983, Proposition 1.2.3}
\source{wintenberger:page-0005}
\uses{fw:apf-finite-base,fw:absolute-finite-base-upper,fw:apf,fw:strict-apf,fw:apf-inertia-indices,fw:apf-herbrand-homeomorphism,fw:finite-nonnormal-apf,fw:finite-index-base-restriction,fw:normal-product-relative-index,fw:apf-nested-index-factor,fw:strict-apf-descent,fw:apf-finite-extension,fw:apf-finite-base-index,fw:apf-finite-base-herbrand,fw:apf-finite-slope,fw:strict-apf-finite-base,fw:apf-finite-base-bounds}
\lean{FontaineWintenberger.isAPF_finiteBase_iff, FontaineWintenberger.isStrictlyAPF_finiteBase_iff, FontaineWintenberger.isAPF_finite_extension_iff, FontaineWintenberger.isStrictlyAPF_finite_extension_iff, FontaineWintenberger.IsAPF.of_algHom, FontaineWintenberger.IsStrictlyAPF.of_algHom, FontaineWintenberger.apfRamificationBound_le_of_algHom, FontaineWintenberger.apfHerbrandERealOrderIso_bound_le, FontaineWintenberger.apfRamificationBound_le_finiteBase, FontaineWintenberger.IsAPF.le_apfHerbrandERealOrderIso, FontaineWintenberger.apfStrictnessConstant_le_of_algHom, FontaineWintenberger.apfStrictnessConstant_le_finiteBase}
\leanok
Let $K\subset M\subset N\subset K^{\mathrm{sep}}$. In expressions
with an infinite argument, extend a finite-base function by
$\psi_{M/K}(+\infty)=+\infty$. Infima over an empty set have value
$+\infty$.
\begin{enumerate}
\item If $M/K$ is finite, then $N/K$ is APF, respectively strictly APF,
      if and only if $N/M$ has the same property.
\item If $N/M$ is finite, then $N/K$ is APF, respectively strictly APF,
      if and only if $M/K$ has the same property.
\item If $N/K$ is APF, respectively strictly APF, then $M/K$ has the same
      property.
\item If $N/K$ is APF, then $i(M/K)\geq i(N/K)$; if also $M/K$ is finite,
      then
      \[
       i(N/M)\geq\psi_{M/K}(i(N/K))\geq i(N/K).
      \]
      If in addition $i(N/K)>0$, then $c(M/K)\geq c(N/K)$, and when
      $M/K$ is finite also $c(N/M)\geq c(N/K)$.
\end{enumerate}
\end{proposition}
\begin{proof}
Write $G=G_K$, $I=G_K^0$, $U_u=G_K^u$, $A=G_N$ and $B=G_M$;
then $A\subset B$. All products with $U_u$ below are subgroups because
$U_u$ is normal. Closedness of these products makes openness equivalent
to finite index.

First suppose $M/K$ is finite. The actual absolute change-of-base
isomorphism identifies $G_M^{\psi_{M/K}(u)}$ with $B\cap U_u$ and
identifies the subgroup fixing $N$ with $A$. Since $A\subset B$,
\[
 (AU_u)\cap B=A(B\cap U_u).
\]
Indeed, in a product $au\in B$ the factor $u=a^{-1}(au)$ lies in $B$.
Because $B$ has finite index in $G$, $AU_u$ has finite index in $G$
exactly when this intersection has finite index in $B$. The increasing
bijection $\psi_{M/K}$ covers every parameter at least $-1$. This proves
the APF equivalence in (1).

For (2), if $N/M$ is finite put $d=[B:A]$. The natural map from the
left cosets of $A$ in $B$ onto the cosets of $AU_u$ in $BU_u$ shows
$[BU_u:AU_u]\leq d$. Thus finiteness of either $[G:AU_u]$ or
$[G:BU_u]$ implies finiteness of the other. For arbitrary $M\subset N$,
$AU_u\subset BU_u$, so finite index of the former implies finite index
of the latter. This proves all the APF assertions.

For the strict assertions all relevant APF indices are now positive
finite integers. Put
\[
 A_0=A\cap I,\quad B_0=B\cap I,\quad
 a_N(u)=[I:A_0U_u],\quad a_M(u)=[I:B_0U_u]\qquad(u\geq0).
\]
The ordinary index formula and the product-coset bijection give
\[
 a_N(u)=a_M(u)d(u),\qquad
 d(u)=[B_0:A_0(B_0\cap U_u)].
\]
The groups defining $d(u)$ decrease with $u$, so $d$ is nondecreasing.
For $0\leq v\leq u$ we have $d(v)\leq d(u)$, and integration gives
\[
 \psi_{N/K}(u)\leq d(u)\psi_{M/K}(u),\qquad
 R_N(u):=\frac{\psi_{N/K}(u)}{a_N(u)}
 \leq R_M(u):=\frac{\psi_{M/K}(u)}{a_M(u)}.
\]
Taking lower limits proves strict descent in (3). If $[B:A]=d$ is
finite, then $1\leq d(u)\leq d$. Therefore
\[
 a_M\leq a_N\leq da_M,\qquad
 \psi_{M/K}\leq\psi_{N/K}\leq d\psi_{M/K},\qquad
 d^{-1}R_M\leq R_N\leq dR_M.
\]
Positivity of one lower limit is equivalent to positivity of the other,
which proves the strict assertion in (2).

Return to finite $M/K$, and write $f=\psi_{M/K}$. The inertia group
over $M$ identifies with $B_0$, and its upper group at $f(u)$ is
$B_0\cap U_u$ for $u\geq0$. The same coset calculation gives
\[
 a_{N/K}(u)=a_{M/K}(u)a_{N/M}(f(u)).
\]
The preceding Herbrand composition theorem integrates this actual index
identity, using the chain rule off the countable sets of slope
discontinuities and the fundamental theorem of calculus with a
countable exceptional set. Both primitives vanish at zero, so
\[
 \psi_{N/K}=\psi_{N/M}\circ f.
\]
They are both the identity on $[-1,0]$ as well.
The finite-slope lemma gives a positive constant $C$ equal to
$a_{M/K}(u)$ for all sufficiently large $u$, and $f$ is eventually
affine with slope $C$. Consequently
\[
 R_{N/K}(u)=C^{-1}R_{N/M}(f(u))
\]
there. Since $f$ is an increasing homeomorphism tending to infinity,
positivity of the two lower limits is equivalent. This proves the
strict assertion in (1).

Finally, $AU_u=G$ implies $BU_u=G$, so the admissible sets defining
$i$ give $i(M/K)\geq i(N/K)$. In the finite-base case, intersecting
$AU_u=G$ with $B$ gives $A(B\cap U_u)=B$. Absolute base change makes
$f(u)$ admissible for $i(N/M)$. Continuity and increasing bijectivity
of $f$, including its stated extension at infinity, give
\[
 i(N/M)\geq f(i(N/K))\geq i(N/K).
\]
For the first inequality take suprema over the admissible $u$; for the
second use $f(u)\geq u$ and continuity. This also covers $i=-1$ and
$i=+\infty$.

Assume $i(N/K)>0$. For every real $u\geq i(M/K)$, also
$u\geq i(N/K)$ and $R_M(u)\geq R_N(u)\geq c(N/K)$.
Taking the defining infimum proves $c(M/K)\geq c(N/K)$.
For finite $M/K$, the two index and function identities above give
\[
 R_{N/M}(f(u))=a_{M/K}(u)R_{N/K}(u)\geq R_{N/K}(u)
 \quad(u\geq0).
\]
If a real $v\geq i(N/M)$ exists, the bound just proved implies
$u=f^{-1}(v)\geq i(N/K)>0$. The last inequality bounds its ratio below
by $c(N/K)$; take the infimum over such $v$. If an index bound is
infinite, its admissible infimum domain is empty and its value is
$+\infty$; the preceding bound on $i$ gives the same conclusion.
This proves the remaining assertions in (4).
\end{proof}


\begin{theorem}[The mixed characteristic Lie criterion]
\label{fw:lie-apf}
\dcref{Wintenberger1983, Examples 1.2.2; Sen1972}
\source{wintenberger:page-0005}
\uses{fw:strict-apf}
\notready
Let $K$ be a finite extension of $\Qp$ and let $L/K$ be Galois.
If $k_L/k_K$ is finite and $\Gal(L/K)$ is a $p$-adic Lie group,
then $L/K$ is strictly APF.
This is the mixed characteristic specialization of Wintenberger's
Examples 1.2.2, using Sen's ramification theorem \cite{Sen1972}.
\end{theorem}

\section{Construction from norm-compatible families}

\begin{lemma}[The ideal norm of the maximal ideal]
\label{fw:ideal-norm-maximal}
\lean{Ideal.relNorm_eq_pow_of_isMaximal, Ideal.inertiaDeg_eq_of_isMaximal}
\mathlibok
Let $A\subset B$ be a finite injective local extension of discrete
valuation rings whose fraction-field extension is separable.
Write $\mathfrak m_A,\mathfrak m_B$ for their maximal ideals and
$f=[k_B:k_A]$. For the relative ideal norm,
\[
 N_{B/A}(\mathfrak m_B)=\mathfrak m_A^f.
\]
\end{lemma}

\begin{lemma}[Normalized valuations of algebraic norms]
\label{fw:norm-valuation}
\uses{fw:ideal-norm-maximal,fw:normalized-field-valuation}
\lean{FontaineWintenberger.addVal_intNorm_uniformizer, FontaineWintenberger.addVal_intNorm, FontaineWintenberger.intValuation_intNorm, FontaineWintenberger.normalizedValuation_norm, FontaineWintenberger.extensionValuationSubring_normalizedValuation_norm, FontaineWintenberger.extensionValuationSubring_normalizedValuation_norm_of_residueDegree_one}
\leanok
Let $E'/E$ be a finite separable extension of complete discretely
valued fields, with their actual valuation rings $B/A$, and put
$f=[k_{E'}:k_E]$. The normalized additive valuations satisfy
\[
 v_E(N_{E'/E}(x))=f\,v_{E'}(x)\qquad(x\in E').
\]
Equivalently, their normalized multiplicative valuations satisfy
$w_E(N_{E'/E}(x))=w_{E'}(x)^f$, including $x=0$.
In particular the normalized valuation is preserved by norms in a
totally ramified extension. No normality or finite-residue-field
hypothesis is needed.
\end{lemma}
\begin{proof}
First let $x\in B$ be nonzero and put $n=v_{E'}(x)$. The principal
ideal $(x)$ is $\mathfrak m_B^n$. The norm of a principal ideal is
the principal ideal of the integral norm, and the integral norm is
the restriction of the field norm. Multiplicativity of the relative
ideal norm and Lemma~\ref{fw:ideal-norm-maximal} give
\[
 (N_{E'/E}(x))
 =N_{B/A}((x))
 =N_{B/A}(\mathfrak m_B)^n
 =\mathfrak m_A^{fn}.
\]
In a discrete valuation ring the exponent of a nonzero principal ideal
is the normalized valuation of its generator, proving the identity
for integral nonzero $x$. For arbitrary nonzero $x\in E'$, write
$x=b/c$ with nonzero $b,c\in B$. The norm is multiplicative, so
subtracting the identities for $b$ and $c$ proves the assertion.
The residue degree $f$ is a positive integer, and both valuations of
zero are infinite; hence the zero case also follows. Exponentiation
translates the additive identity into the multiplicative identity.
\end{proof}

\begin{definition}[The multiplicative norm limit over all finite layers]
\label{fw:norm-limit-multiplicative}
\dcref{Wintenberger1983, 2.1.1}
\source{wintenberger:page-0008}
\lean{FontaineWintenberger.FiniteLayer, FontaineWintenberger.FiniteLayer.finiteDimensional, FontaineWintenberger.FiniteLayer.bottom, FontaineWintenberger.FiniteLayer.bottom_le, FontaineWintenberger.FiniteLayer.sup, FontaineWintenberger.FiniteLayer.le_sup_left, FontaineWintenberger.FiniteLayer.le_sup_right, FontaineWintenberger.FiniteLayer.norm, FontaineWintenberger.FiniteLayer.norm_eq_zero_iff, FontaineWintenberger.FiniteLayer.norm_self, FontaineWintenberger.FiniteLayer.norm_norm, FontaineWintenberger.FiniteLayer.normUnits, FontaineWintenberger.FiniteLayer.normUnits_val, FontaineWintenberger.FiniteLayer.normUnits_self, FontaineWintenberger.FiniteLayer.normUnits_normUnits, FontaineWintenberger.normLimitGroup, FontaineWintenberger.NormLimit, FontaineWintenberger.normLimitEval, FontaineWintenberger.normLimitEval_coe, FontaineWintenberger.normLimitEval_eq_zero_iff, FontaineWintenberger.NormLimit.ext, FontaineWintenberger.normLimitEval_norm, FontaineWintenberger.existsUnique_normLimit_of_compatible}
\leanok
For any field extension $L/K$, let $\mathcal E_{L/K}$ be the finite
intermediate extensions, ordered by inclusion. This is a nonempty
directed set: it contains $K$ and is closed under finite composita.
For $E\subset E'$ in this set, the ordinary field norm
$N_{E'/E}:E'^\times\to E^\times$ is a group homomorphism. These maps
are identity on equal layers and compose in towers.

Let $G$ be the subgroup of $\prod_{E\in\mathcal E_{L/K}}E^\times$
of families satisfying $N_{E'/E}(\alpha_{E'})=\alpha_E$ for every
inclusion. Define $X_K(L)=G\sqcup\{0\}$ with componentwise
multiplication and absorbing zero. Every coordinate defines a
multiplicative evaluation $X_K(L)\to E$ preserving zero and one.
The evaluations jointly distinguish elements, and evaluation at any
one layer is zero exactly on the zero element.

This identifies $X_K(L)$ with the families
$(\alpha_E)\in\prod_E E$ satisfying the same norm equations: each
such family is either identically zero or has every coordinate nonzero.
\end{definition}
\begin{proof}
Finite composita of finite extensions are finite. If $E\subset E'$,
a finite spanning set for $E'/K$ also spans $E'/E$, so the norm is
the norm of a genuinely finite extension. It preserves multiplication,
one and nonzeroness; its identity and tower identities give the
transition laws. Applying multiplicativity and the inverse law to
the compatibility equations shows that they define a subgroup of
the displayed product.

Extend the coordinate homomorphisms on units by sending the adjoined
zero to zero. Their multiplicative laws and joint injectivity follow
coordinatewise; a coordinate of a unit is nonzero, proving the stated
zero test. Finally consider a compatible family of field elements.
For every layer $E$, compatibility with $K$ gives
$N_{E/K}(\alpha_E)=\alpha_K$. A norm in a finite field extension
vanishes exactly on zero. Thus if $\alpha_K=0$, all coordinates
are zero; otherwise all are nonzero and determine unique units.
The norm equations for these units are the original equations,
giving the asserted identification in both directions.
\end{proof}

\begin{lemma}[Invariance of the actual finite residue degree]
\label{fw:finite-extension-residue-invariance}
\uses{fw:ramification-field-equiv,fw:chosen-extension-integers,fw:spectral-integer-residue-maps}
\lean{FontaineWintenberger.spectral_residueDegree_algEquiv, FontaineWintenberger.extensionValuationSubring_residueDegree_eq_spectral}
\leanok
For a finite separable extension $E/K$, its actual residue-field
degree is preserved by every $K$-algebra isomorphism and by
replacement of the spectral valuation with any compatible
extension valuation on the same field.
\end{lemma}
\begin{proof}
The original field isomorphism restricts to an isomorphism of
its actual spectral integer rings over $\mathcal O_K$.
For a change of valuation, uniqueness of the integral closure
gives the same kind of actual integer-ring isomorphism.
Each is local, carries the maximal ideal onto the maximal ideal,
and therefore induces an isomorphism of the actual residue
fields over the base residue field. The associated linear
isomorphism preserves their finite dimensions, proving both
claims.
\end{proof}

\begin{lemma}[Actual finite APF stages]
\label{fw:apf-finite-stages}
\uses{fw:apf-abstract,fw:absolute-upper-projections,fw:finite-upper-groups}
\lean{FontaineWintenberger.exists_absoluteUpper_le_fixingSubgroup_of_finite, FontaineWintenberger.apfFiniteStage, FontaineWintenberger.fixingSubgroup_apfFiniteStage, FontaineWintenberger.apfFiniteStage_le, FontaineWintenberger.finiteDimensional_apfFiniteStage, FontaineWintenberger.apfFiniteStage_mono, FontaineWintenberger.exists_le_apfFiniteStage, FontaineWintenberger.exists_finiteLayer_map_eq_apfFiniteStage, FontaineWintenberger.IsAPF.exists_finiteLayer_ge_map_eq_apfFiniteStage}
\leanok
Embed any APF separable algebraic extension $L/K$ into $K^{\rm sep}$,
write $H$ for its actual fixing subgroup and put
\[
 E_u=(K^{\rm sep})^{H G_K^u}\qquad(u\geq-1).
\]
Then $E_u\subseteq L$ is finite over $K$, its actual fixing
subgroup is $H G_K^u$, and $E_u\subseteq E_v$ for $u\leq v$.
These fields are cofinal among the actual finite intermediate
fields of $L/K$. Pulling them back along the original embedding
gives actual finite subfields of the original abstract $L$.
\end{lemma}
\begin{proof}
The product $H G_K^u$ is a subgroup because $G_K^u$ is normal.
Both factors are compact, hence their product is closed. The
infinite Galois correspondence therefore identifies it with
the fixing subgroup of its actual fixed field. It contains $H$,
so its fixed field lies in $L$. APF makes the product open, and
the open-subgroup correspondence makes this fixed field finite.
Decreasingness of the upper groups and reversal of inclusions
under fixed fields give $E_u\subseteq E_v$ for $u\leq v$.

Given a finite $F\subseteq L$, choose an actual finite Galois
overfield $\Omega/K$ containing $F$ in $K^{\rm sep}$. Its upper
groups are eventually trivial. Ordinary projection of absolute
upper groups onto these finite upper groups therefore gives
$G_K^u\subseteq\operatorname{Gal}(K^{\rm sep}/\Omega)$ for all
sufficiently large $u$. Both $H$ and $G_K^u$ fix $F$, so their
product fixes $F$. Thus $F\subseteq E_u$, proving cofinality.
Finally, the original embedding is an isomorphism from $L$
onto its actual image. Restrict its inverse to each $E_u$ and
take its field image in the original $L$. This is an actual
finite intermediate field, isomorphic to $E_u$; mapping it
back gives precisely $E_u$. These inverse restrictions preserve
all the stated inclusions.
\end{proof}

\begin{lemma}[Cofinal growth of the actual relative first bounds]
\label{fw:apf-deep-finite-tail}
\source{wintenberger:page-0005,wintenberger:page-0008,wintenberger:page-0013}
\uses{fw:apf-finite-stages,fw:abstract-normed-finite-base,fw:apf-extended-bound,fw:apf-herbrand-homeomorphism}
\lean{FontaineWintenberger.AbstractFiniteBase.apfPsi_le_apfRamificationBound_of_range_le, FontaineWintenberger.apfPsi_le_relativeFirstBound_of_apfFiniteStage_le, FontaineWintenberger.IsAPF.exists_finiteLayer_relativeFirstBound_ge}
\leanok
Let $L/K$ be an abstract separable APF extension, let
$j:L\hookrightarrow K^{\rm sep}$, and let $u\geq0$.
For actual finite intermediate fields $E\subseteq F\subseteq L$,
if $j(E)$ contains the actual stage
$(K^{\rm sep})^{G_LG_K^u}$, then
\[
 u\leq\Psi_{E/K}(u)\leq i(F/E).
\]
The right side is the intrinsic extended first bound, so the
statement also includes $F=E$, when that bound is $+\infty$.
Consequently, for every real $c$, there is an actual finite
$T\subseteq L$ such that $c\leq i(F/E)$ for every actual
finite pair $T\subseteq E\subseteq F\subseteq L$.
\end{lemma}
\begin{proof}
This is a local deduction from the actual finite-base upper comparison;
the cited later pages describe its norm-limit consumer.
Write $H_L$, $H_F$, and $H_E$ for the actual fixing subgroups
of their images in $K^{\rm sep}$. The stage containment gives
\[
 H_E\subseteq H_LG_K^u\subseteq H_FG_K^u.
\]
Choose a separable closure of the original normed field $E$.
The closure isomorphism extending the original embedding
$j|_E$ induces restriction
$\rho:G_E^{\rm abs}\longrightarrow G_K$ with image $H_E$.
If $i:F\hookrightarrow E^{\rm sep}$ is obtained by composing
$j|_F$ with the inverse closure isomorphism, its resulting
embedding into $K^{\rm sep}$ is precisely $j|_F$.
Thus the preimage of $H_FG_K^u$ under $\rho$ is the whole
group. The actual finite-base upper comparison identifies
that preimage with
$G_F^{\rm abs}G_E^{\Psi_{E/K}(u)}$.
This relative product is therefore the whole group, so its
parameter $\Psi_{E/K}(u)$ belongs to the defining set of the
intrinsic extended supremum $i(F/E)$.
The finite-extension inverse Herbrand function satisfies
$\Psi_{E/K}(u)\geq u$, proving the two inequalities.
For the final assertion, choose $u\geq\max(c,0)$ and pull its
actual finite stage back along $j$ to the original $L$.
Every $E$ above that pullback satisfies the displayed bound.
This proves cofinal growth without any normality or residue-field
finiteness assumption. The argument supplies explicitly the
cofinal precision used in the norm-limit construction.
\end{proof}

\begin{lemma}[Countable cofinal finite layers of an APF extension]
\label{fw:countable-apf-finite-layers}
\uses{fw:apf-finite-stages}
\lean{FontaineWintenberger.IsAPF.exists_monotone_cofinal_finiteLayers_map, FontaineWintenberger.IsAPF.exists_monotone_cofinal_finiteLayers}
\leanok
Every abstract separable APF extension $L/K$ has an increasing
sequence $(E_n)_{n\in\mathbb N}$ of actual finite intermediate
fields which is cofinal among all finite intermediate fields.
The sequence may be chosen so that, for any specified embedding
$j:L\hookrightarrow K^{\rm sep}$, its $n$-th image is the actual
fixed field of $G_LG_K^n$.
\end{lemma}
\begin{proof}
For each $n\geq0$ take the actual finite stage at real parameter
$n$ and restrict the inverse of $j$ to that field.
The exact image identities and injectivity of $j$ transport
monotonicity of the stages to monotonicity of $(E_n)$.
Given an actual finite $D\subseteq L$, the finite-stage
cofinality proof supplies a real $U\geq0$ such that every
stage at parameter $u\geq U$ contains $j(D)$.
Choose an integer $n\geq U$. Then $j(D)\subseteq j(E_n)$,
and injectivity of $j$ gives $D\subseteq E_n$.
\end{proof}

\begin{lemma}[The finite zero index and residue stabilization]
\label{fw:apf-residue-stabilization}
\source{wintenberger:page-0008}
\uses{fw:finite-extension-residue-invariance,fw:finite-galois-overfield-apf,fw:inertia-residue-index,fw:apf-finite-stages}
\lean{FontaineWintenberger.apfSubgroup_index_zero_eq_residueDegree_of_finiteGalois, FontaineWintenberger.apfSubgroup_index_zero_eq_residueDegree, FontaineWintenberger.apfSubgroup_zero_eq_of_apfFiniteStage_le, FontaineWintenberger.residueDegree_eq_zeroStage_of_apfFiniteStage_le, FontaineWintenberger.spectralResidueInclusion_zeroStage_surjective, FontaineWintenberger.spectralResidueInclusion_surjective_of_zeroStage_le, FontaineWintenberger.IsAPF.exists_finiteLayer_residueDegree_stable, FontaineWintenberger.IsAPF.exists_finiteLayer_residueInclusion_surjective}
\leanok
For every finite separable extension $M/K$ and every embedding
$j:M\hookrightarrow K^{\rm sep}$,
\[
 [G_K:G_MG_K^0]=[k_M:k_K].
\]
If $L/K$ is APF, put $E_0=(K^{\rm sep})^{G_LG_K^0}$.
Every finite $E$ with $E_0\subseteq E\subseteq L$ satisfies
$[k_E:k_K]=[k_{E_0}:k_K]$. Consequently the actual residue
extension $k_E/k_{E_0}$ has degree one.
\end{lemma}
\begin{proof}
Choose a finite Galois normal closure $\Omega/K$ of $M/K$
in $K^{\rm sep}$ and write $H=\operatorname{Gal}(\Omega/M)$.
Projection onto $G=\operatorname{Gal}(\Omega/K)$ identifies
the displayed full-group index with $[G:HG^0]$.
The finite zero-orbit-index identity identifies this with
the residue degree of the restricted valuation on $M$.
The actual integer-ring isomorphisms of
Lemma~\ref{fw:finite-extension-residue-invariance} identify
that residue degree with the original spectral residue degree
of $M/K$. Conjugacy of separable-closure embeddings gives
the formula for the specified $j$.

By Lemma~\ref{fw:apf-finite-stages}, $E_0/K$ is finite and
$G_{E_0}=G_LG_K^0$. If $E_0\subseteq E\subseteq L$, then
$G_L\subseteq G_E\subseteq G_{E_0}$ and $G_K^0\subseteq G_{E_0}$.
Thus $G_EG_K^0=G_LG_K^0=G_{E_0}$. Applying the finite
formula to $E$ and $E_0$ proves equality of their residue
degrees over $k_K$. In the actual residue-field tower,
$[k_E:k_K]=[k_{E_0}:k_K][k_E:k_{E_0}]$.
The first factor is a positive finite dimension; cancellation
gives $[k_E:k_{E_0}]=1$.
\end{proof}

\begin{lemma}[An actual wild tail of finite APF layers]
\label{fw:apf-wild-finite-tail}
\source{wintenberger:page-0008}
\uses{fw:apf-finite-stages,fw:apf-residue-stabilization,fw:absolute-upper-projections,fw:wild-ramification-p-group,fw:local-dvr-degree-identity}
\lean{FontaineWintenberger.relIndex_sup_eq_relIndex_of_normal_right, FontaineWintenberger.relIndex_eq_prime_pow_of_le_sup, FontaineWintenberger.finrank_eq_prime_pow_of_apfFiniteStage_one_le, FontaineWintenberger.FiniteLayer.ramificationIdx_eq_finrank_of_residueInclusion_surjective, FontaineWintenberger.IsAPF.exists_finiteLayer_residue_surjective_and_degree_prime_pow, FontaineWintenberger.IsAPF.exists_finiteLayer_wild_tail}
\leanok
Let the perfect residue field of $K$ have characteristic $p>0$,
and let $L/K$ be an abstract separable APF extension.
There is an actual finite intermediate field $T\subseteq L$
such that every finite pair $T\subseteq E\subseteq F\subseteq L$
has surjective actual residue inclusion and $[F:E]=p^n$ for
some $n\geq0$. In particular every such actual finite step is
totally ramified.
\end{lemma}
\begin{proof}
Embed $L$ into $K^{\rm sep}$, put $H=G_L$, and take the actual
finite stage $T=(K^{\rm sep})^{HG_K^1}$, pulled back to the
original $L$. The zero stage is contained in $T$, so residue
stabilization proves surjectivity of every indicated residue map.
Its actual residue degree is one. The finite local degree identity
therefore gives total ramification of every step.

For such a pair, choose a finite Galois normal closure $\Omega/K$
containing the actual image of $F$, and let $\rho:G_K\to G$ be
the ordinary restriction, where $G=\operatorname{Gal}(\Omega/K)$.
Put $A=\rho(H)$, $B=\operatorname{Gal}(\Omega/F)$,
$C=\operatorname{Gal}(\Omega/E)$ and $U=G^1$.
Surjectivity of restriction from each actual fixing subgroup gives
$A\subseteq B\subseteq C$. The containment $T\subseteq E$
gives $G_E\subseteq HG_K^1$; projection and the upper-group
projection theorem give $C\subseteq AU$.
The lower inverse parameter at $1$ is at least $1$, so
$U\subseteq G_1$. Wild inertia $G_1$ is a finite $p$-group,
and hence so is $U$. Normality of $U$ makes $AU$ a subgroup,
and its actual $A$-cosets are the $U/(A\cap U)$ cosets.
Thus $[AU:A]$ is a power of $p$. The finite index chain
\[
 [AU:A]=[AU:C][C:B][B:A]
\]
shows that $[C:B]$ is a power of $p$. Finite Galois
correspondence identifies this actual index with $[F:E]$.
The original field embeddings are isomorphisms onto their images,
so this is precisely the degree of the original finite step.
\end{proof}

\begin{lemma}[The stable normalized valuation on norm-compatible families]
\label{fw:norm-limit-stable-valuation}
\source{wintenberger:page-0008}
\uses{fw:apf-residue-stabilization,fw:norm-valuation,fw:norm-limit-multiplicative}
\lean{FontaineWintenberger.FiniteLayer.normalizedValuation_norm_of_residueInclusion_surjective, FontaineWintenberger.normLimitCoordinateValuation, FontaineWintenberger.IsAPF.exists_normLimitCoordinateValuation_stable, FontaineWintenberger.normLimitValuation, FontaineWintenberger.IsAPF.exists_normLimitValuation_eq_coordinate, FontaineWintenberger.normLimitValuation_eq_zero_iff}
\leanok
Let $L/K$ be an abstract separable APF extension and let $X$ be its
multiplicative norm limit, with zero adjoined. There is an actual finite
layer $E_0$ such that for all finite $E_0\subseteq E\subseteq F\subseteq L$,
\[
 w_E(N_{F/E}(z))=w_F(z)\qquad(z\in F),
\]
where $w_E$ and $w_F$ are the normalized multiplicative valuations of
the actual spectral integer rings. Thus $w_E(\alpha_E)$ is independent
of $E\supseteq E_0$ for every $\alpha\in X$. This defines a
multiplicative map $w:X\to\mathbb Z^{\rm mult}\cup\{0\}$ with
$w(0)=0$, $w(1)=1$, and $w(\alpha)=0$ exactly when $\alpha=0$.
\end{lemma}
\begin{proof}
Choose the actual finite layer supplied by residue stabilization.
Every finite inclusion above it induces a surjection of the actual
residue fields. As a field embedding this map is bijective, so the
actual residue degree is one. The original integer-ring maps are
local and finite; the norm-valuation identity of
Lemma~\ref{fw:norm-valuation} therefore proves the displayed equality
on the original fields. Apply it to the compatibility equation
$N_{F/E}(\alpha_F)=\alpha_E$.
For incomparable $E$ and $F$, their actual finite compositum is
above both, so two such equalities prove independence. Evaluate at
any one finite layer above $E_0$ to define $w$. Evaluation and its
normalized field valuation are multiplicative and preserve zero
and one, so their composite has those properties. The normalized
field valuation vanishes exactly on zero, and any coordinate of
the zero-adjoined unit limit vanishes exactly when the element is
zero. These two facts prove the asserted zero test.
\end{proof}

\begin{definition}[The norm valuation]
\label{fw:norm-families}
\dcref{Wintenberger1983, 2.1.1--2.1.2}
\source{wintenberger:page-0008}
\uses{fw:apf,fw:apf-finite-stages,fw:norm-valuation,fw:norm-limit-multiplicative,fw:apf-finite-unramified-tame-parts,fw:norm-limit-stable-valuation}
\lean{FontaineWintenberger.apfMaximalUnramified, FontaineWintenberger.apfMaximalTame, FontaineWintenberger.normLimitValuation, FontaineWintenberger.normLimitValuation_eq_coordinate_of_maximalUnramified_le}
\leanok
For infinite APF $L/K$, let $K_0$ and $K_1$ be the maximal unramified and
maximal tamely ramified subextensions of $L/K$; both are finite over $K$.
For $E\supset K_0$ define the norm valuation by
\[
 v((\alpha_E)_E)=v_E(\alpha_E),\qquad v(0)=+\infty.
\]
\end{definition}
\begin{proof}
The maximal unramified subextension is fixed by
$G_LG_K^0$, which is open by APF. If $P_K$ is the wild inertia group,
the maximal tamely ramified subextension is fixed by $G_LP_K$.
Since $G_K^1\subset P_K$, this subgroup contains the open subgroup
$G_LG_K^1$ and is also open. Hence $K_0/K$ and $K_1/K$ are finite.
To check independence of $E$ in the valuation formula,
first take $K_0\subset E\subset E'\subset L$ finite over $K$.
Since $k_E=k_{E'}=k_L$, the residue degree of $E'/E$ is one.
Lemma~\ref{fw:norm-valuation} gives
$v_E(N_{E'/E}(x))=v_{E'}(x)$.
Applying this to $x=\alpha_{E'}$ proves the required equality.
For two incomparable fields use their compositum, which is another finite
subextension of $L/K$ containing both.
\end{proof}

\begin{definition}[Residue constants]
\label{fw:residue-constants}
\dcref{Wintenberger1983, 2.1.2}
\source{wintenberger:page-0008}
\uses{fw:norm-families,fw:perfect-multiplicative-residue-lifts,fw:apf-finite-unramified-tame-parts}
\lean{FontaineWintenberger.sourceResidueDegree, FontaineWintenberger.sourceResidueRoot, FontaineWintenberger.sourceResidueIntegralCoordinate, FontaineWintenberger.sourceResidueCoordinate, FontaineWintenberger.sourceResidueCoordinate_norm, FontaineWintenberger.sourceResidueConstants, FontaineWintenberger.normLimitEval_sourceResidueConstants}
\leanok
Let $L/K$ be infinite APF. For $x\in k_L=k_{K_0}$ write $[x]$ for its
multiplicative representative in $K_0$.
For $E\supset K_1$, the integer $[E:K_1]=p^r$ is a power of $p$;
put
\[
 x_E=[x^{p^{-r}}].
\]
The compatible family $(x_E)_{E\supset K_1}$ defines an element
$f_{L/K}(x)\in X_K(L)$.
\end{definition}
\begin{proof}
The inverse Frobenius $x\mapsto x^{p^{-r}}$ is defined because $k_L$ is
perfect. If $E\subset E'$ both contain $K_1$, then $x_{E'}\in K_0\subset E$
and therefore
\[
 N_{E'/E}(x_{E'})=x_{E'}^{[E':E]}=x_E.
\]
The subextensions containing $K_1$ are cofinal: the compositum of any
finite $F/K$ in $L$ with $K_1$ is such an extension.
For a finite $F/K$ choose $E\supset F K_1$ and define
$x_F=N_{E/F}(x_E)$; compatibility makes this independent of $E$ by passing
to a compositum of any two choices. This extends the family uniquely to
all finite subextensions and proves the assertion.
\end{proof}

\begin{lemma}[The finite attained bound of a totally ramified subextension]
\label{fw:finite-totally-ramified-first-bound}
\uses{fw:finite-first-upper-break,fw:finite-galois-overfield-apf,fw:inertia-relative-index}
\lean{FontaineWintenberger.exists_apfRamificationBound_eq_coe_of_fixingSubgroup, FontaineWintenberger.exists_apfRamificationBound_eq_coe_of_totallyRamified}
\leanok
Let $\Omega/K$ be finite Galois and let $M$ be a nontrivial
actual intermediate field totally ramified over $K$. There is
$b\in\mathbb R$, $b\geq0$, with $i(M/K)=b$ and
\[
 \operatorname{Gal}(\Omega/M)\operatorname{Gal}(\Omega/K)^u
   =\operatorname{Gal}(\Omega/K)\quad\Longleftrightarrow\quad u\leq b
 \qquad(u\geq-1).
\]
More generally the same conclusion holds whenever the actual
fixing subgroup $H$ is proper and $HG_0=G$.
\end{lemma}
\begin{proof}
Total ramification is equivalent to $HG_0=G$, by the actual
inertia-relative index formula and field degree. Nontriviality
makes $H$ proper by Galois correspondence. The finite first-break
lemma supplies $b\geq0$ and the asserted equivalence for $u\geq0$.
If $-1\leq u\leq0$, then $G^u\supseteq G^0=G_0$; hence
$HG^u=G$ and also $u\leq b$. Thus the same equivalence holds
throughout the required half-line. The arbitrary-overfield
first-bound formula is the supremum of these actual parameters.
All of them are at most $b$, and $b$ itself is admissible.
The extended-real supremum is therefore exactly the finite value $b$.
\end{proof}

\begin{lemma}[The finite Galois interpretation of the ramification bound]
\label{fw:finite-galois-ramification-bound}
\uses{fw:apf-extended-bound,fw:finite-apf-indices,fw:ramification-field-equiv}
\lean{FontaineWintenberger.apfSubgroup_eq_top_iff_finite_upper_eq_top, FontaineWintenberger.apfSubgroup_fieldRange_eq_top_iff_finite_upper_eq_top, FontaineWintenberger.apfRamificationBound_eq_sSup_finite_upper}
\leanok
Let $E'/E$ be finite Galois with group $G$. For $u\geq-1$, the
condition $G_E^uG_{E'}=G_E$ is equivalent to $G^u=G$. Consequently
\[
 i(E'/E)=\sup\{u\geq-1:G^u=G\}.
\]
\end{lemma}
\begin{proof}
Choose a separable embedding of $E'$ into a separable closure of
$E$. Absolute restriction onto $G$ is surjective and has kernel
$G_{E'}$. Its image on the absolute upper group $G_E^u$ is the
actual finite upper group $G^u$. The preimage of this image is
therefore $G_E^uG_{E'}$. Since restriction is surjective, that
preimage equals the whole absolute group exactly when its image
equals $G$. Taking the same extended-real supremum on the two
equal admissible sets proves the formula. Independence of the
chosen embedding gives the assertion for the abstract extension.
\end{proof}

\begin{lemma}[The first break and trace bound in prime degree]
\label{fw:cyclic-prime-trace-bound}
\source{wintenberger:page-0010,serre-1979-local-fields:page-0090}
\dcref{Wintenberger, proof of Proposition 2.2.1; Serre, Chapter V, Section 3, Lemma 4}
\uses{fw:finite-galois-ramification-bound,fw:integral-order-function,fw:herbrand-functions,fw:hilbert-different,fw:trace-ideal}
\lean{FontaineWintenberger.ramificationOrder_eq_of_prime_card, FontaineWintenberger.exists_common_ramificationOrder_of_prime_card, FontaineWintenberger.exists_common_ramificationOrder_succ_of_prime_card, FontaineWintenberger.lowerRamificationGroup_nat_of_common_order, FontaineWintenberger.lowerRamificationGroupReal_of_common_order, FontaineWintenberger.herbrandFunction_common_break, FontaineWintenberger.upperRamificationGroup_of_common_order, FontaineWintenberger.sum_ramificationOrder_of_common_order, FontaineWintenberger.spectralIntegerRing_lower_zero_eq_top_of_totallyRamified, FontaineWintenberger.spectralIntegerRing_exists_common_break_of_prime_degree, FontaineWintenberger.spectralIntegerRing_different_order_of_common_order, FontaineWintenberger.apfRamificationBound_eq_common_break, FontaineWintenberger.spectralIntegerRing_intTraceIdeal_of_common_break, FontaineWintenberger.spectralIntegerRing_intTrace_order_bound, FontaineWintenberger.spectralIntegerRing_prime_degree_ramification}
\leanok
Let $E'/E$ be totally ramified Galois of prime degree $p$.
There is an integer $t\geq0$ such that every nonidentity element
of $G=\Gal(E'/E)$ has intrinsic order $t+1$.
One has
\[
 i(E'/E)=t,\qquad d(E'/E)=(p-1)(t+1),\qquad
 \operatorname{Tr}_{E'/E}(\mathcal O_{E'})
     =\mathfrak m_E^{\lfloor(p-1)(t+1)/p\rfloor}.
\]
In particular every $x\in\mathcal O_{E'}$ satisfies
$v_E(\operatorname{Tr}(x))\geq (p-1)i(E'/E)/p$.
\end{lemma}
\begin{proof}
The group has order $p$, so every proper subgroup is trivial:
its order divides $p$ and hence is one. Thus each lower group
is either $G$ or the identity subgroup. Total ramification gives
$G_0=G$, and the lower groups are eventually trivial. Their
monotonicity and the intrinsic integral-order membership test
give a single integer $t\geq0$ with $G_n=G$ for $0\leq n\leq t$
and $G_n=1$ for $n>t$. Each nonidentity order is therefore $t+1$.
With the real lower convention $G_s=G_{\lceil s\rceil}$, the
group is full through $s=t$ and trivial immediately after $t$.
The Herbrand integral is $\varphi(s)=s$ on $[-1,t]$ and
$\varphi(s)=t+(s-t)/p$ for $s\geq t$. Its inverse sends
$u\leq t$ into $s\leq t$ and $u>t$ into $s>t$.
Thus $G^u=G$ exactly when $u\leq t$, for $u\geq-1$.
The preceding lemma identifies the public bound with $t$.
Hilbert's formula now sums $p-1$ identical nonidentity orders,
giving $d=(p-1)(t+1)$. The trace formula with $n=0$ and $e=p$
gives the asserted ideal. Finally
\[
 \left\lfloor\frac{(p-1)(t+1)}p\right\rfloor
       =t-\left\lfloor\frac tp\right\rfloor
       \geq\frac{p-1}{p}t.
\]
This proves the bound even when the trace of $1$ is zero.
\end{proof}

\begin{lemma}[A norm defect with a controlled trace representative]
\label{fw:norm-defect-controlled-trace}
\source{wintenberger:page-0010,serre-1979-local-fields:page-0091}
\uses{fw:cyclic-prime-trace-bound,fw:cyclic-prime-norm-congruence}
\lean{FontaineWintenberger.exists_prod_add_sub_eq_sum_smul_mem, FontaineWintenberger.spectralIntegerRing_exists_intNorm_add_sub_eq_intTrace_mem_pow, FontaineWintenberger.spectralIntegerRing_intNorm_sub_sub_mem_intTraceIdeal_pow}
\leanok
Let a finite group $G$ of prime order act by automorphisms on a
commutative ring $B$, and let $J\subseteq B$ be a $G$-stable ideal.
For $a\in J$ and $b\in B$, there is $z\in J$ with
\[
 \prod_{g\in G}g(a+b)-\prod_{g\in G}g(a)-\prod_{g\in G}g(b)
       =\sum_{g\in G}g(z).
\]
Consequently, in an actual finite Galois extension of prime degree,
if $a\in\mathfrak m_{E'}^r$ and $b\in\mathcal O_{E'}$, the
actual integral norm defect belongs to
$\operatorname{Tr}_{E'/E}(\mathfrak m_{E'}^r)$.
\end{lemma}
\begin{proof}
Expand the left product over all subsets of $G$. Remove the empty
and full subsets, which give the two subtracted products. Each
remaining monomial contains at least one factor $g(a)\in J$,
so lies in $J$. The action on nonempty proper subsets is free:
a stabilizer is trivial or all of the prime-order group, and a
nonempty subset fixed by all translations is the whole group.
Choose one monomial from each orbit and sum these choices to
obtain $z\in J$. Equivariance of the subset monomials makes the
sum on each free orbit the sum of all conjugates of its chosen
monomial. Adding orbitwise proves the displayed identity.
For the last assertion take $B=\mathcal O_{E'}$ and
$J=\mathfrak m_{E'}^r$. Every automorphism preserves the
actual maximal ideal and hence its powers. The actual integral
norm and trace are the indicated product and sum, so the same
identity gives the asserted trace-ideal membership.
\end{proof}

\begin{lemma}[Prime-step norm continuity at the Herbrand precision]
\label{fw:prime-norm-herbrand-continuity}
\source{wintenberger:page-0012,wintenberger:page-0013}
\uses{fw:norm-defect-controlled-trace,fw:cyclic-prime-trace-bound,fw:trace-ideal,fw:norm-valuation,fw:intrinsic-prime-degree-apf-profile,fw:integral-finite-layer-norms}
\lean{FontaineWintenberger.ceil_prime_herbrand_piece_le_precision, FontaineWintenberger.exists_spectralIntegerRing_prime_norm_sub_precision, FontaineWintenberger.spectralIntegerRing_prime_norm_sub_congruence, FontaineWintenberger.FiniteLayer.integerNorm_prime_sub_congruence}
\leanok
Let $E'/E$ be totally ramified Galois of prime degree $p$,
with its actual common lower break $t\in\mathbb N$.
For every $r\in\mathbb N$ and $x,y\in\mathcal O_{E'}$ with
$x-y\in\mathfrak m_{E'}^r$,
\[
 N(x)-N(y)\in\mathfrak m_E^{\lceil\varphi_{E'/E}(r)\rceil_+}.
\]
\end{lemma}
\begin{proof}
Put $a=x-y$ and $b=y$. The controlled representative lemma
gives $N(x)-N(y)=N(a)+\operatorname{Tr}(z)$ with
$z\in\mathfrak m_{E'}^r$. The norm valuation identity and
residue degree one give $N(a)\in\mathfrak m_E^r$.
The different exponent is $d=(p-1)(t+1)$, and the exact trace
ideal formula gives
$\operatorname{Tr}(z)\in\mathfrak m_E^{\lfloor(r+d)/p\rfloor}$.
Hence their sum lies in the power with exponent
$\min(r,\lfloor(r+d)/p\rfloor)$.
The intrinsic prime-degree profile gives a real break $b$ with
$d=(p-1)(b+1)$. Comparing with $d=(p-1)(t+1)$ and cancelling
$p-1>0$ gives $b=t$. Its actual intrinsic Herbrand function is
therefore $r$ for $r\leq t$ and $t+(r-t)/p$ for $r\geq t$. If $r\leq t$, then
$r+d-pr=(p-1)(t+1-r)\geq0$, so $\lfloor(r+d)/p\rfloor\geq r$.
If $r\geq t$, put $q=r+(p-1)t$. Then
$\lfloor(r+d)/p\rfloor=\lfloor(q+p-1)/p\rfloor
=\lceil q/p\rceil=\lceil\varphi_{E'/E}(r)\rceil$,
and $\varphi_{E'/E}(r)\leq r$.
These two cases prove the required exponent inequality, including
$r=t$ and $r=0$. The argument concerns $1-N(x/y)$ when expressed
using units; it uses the norm difference directly and needs no
identity with $N(1-x/y)$.
\end{proof}

\begin{lemma}[Herbrand continuity along actual wild prime towers]
\label{fw:wild-tower-norm-continuity}
\source{wintenberger:page-0012,wintenberger:page-0013}
\uses{fw:prime-norm-herbrand-continuity,fw:relative-prime-field-tower,fw:wild-relative-total-ramification,fw:abstract-normed-finite-base,fw:integral-finite-layer-norms}
\lean{FontaineWintenberger.FiniteLayer.integerNorm_sub_congruence_of_wild_field}
\leanok
Let $\Omega/K$ be finite Galois, let $P$ be its actual first
lower ramification group, and let $\Omega^P\subseteq E\subseteq F$
be actual finite intermediate fields. Then for $r\geq0$ integral
and $x,y\in\mathcal O_F$,
\[
 x-y\in\mathfrak m_F^r\quad\Longrightarrow\quad
 N_{F/E}(x)-N_{F/E}(y)
 \in\mathfrak m_E^{\lceil\varphi_{F/E}(r)\rceil_+}.
\]
\end{lemma}
\begin{proof}
The actual fixing group of $E$ is a subgroup of the finite
$p$-group $P$. The actual prime-degree field-chain lemma
therefore gives a finite increasing chain from $E$ to $F$
whose successive steps are Galois of degree $p$.
Every such step is totally ramified, since both its fields
contain $\Omega^P$. Apply prime-step norm continuity on
the original integer rings at each step.
For a tower $A\subseteq B\subseteq C$, the actual inverse
Herbrand comparison gives
$\Psi_{C/A}=\Psi_{C/B}\circ\Psi_{B/A}$; taking inverses gives
$\varphi_{C/A}=\varphi_{B/A}\circ\varphi_{C/B}$.
Suppose the norm from $C$ to $B$ gives precision
$s=\lceil\varphi_{C/B}(r)\rceil_+$.
The norm from $B$ to $A$ then gives precision
$\lceil\varphi_{B/A}(s)\rceil_+$.
Since $s\geq\varphi_{C/B}(r)$ and the actual Herbrand
functions are increasing, this integer is at least
$\lceil\varphi_{C/A}(r)\rceil_+$.
The actual integral norms compose by norm transitivity,
so induction along the finite chain proves the assertion.
For a chain of length zero, the norm and Herbrand function
are both the identity, so the same assertion holds.
\end{proof}

\begin{lemma}[Herbrand scaling for the actual tame compositum]
\label{fw:wild-compositum-herbrand-scaling}
\source{wintenberger:page-0005,wintenberger:page-0010,wintenberger:page-0011,wintenberger:page-0013}
\dcref{Local deduction from finite-base transitivity and tame-compositum reduction}
\uses{fw:wild-tame-compositum,fw:wild-fixed-tame-field,fw:abstract-normed-finite-base,fw:ramification-subgroups}
\lean{FontaineWintenberger.wild_compositum_extendScalars_eq_fixedField, FontaineWintenberger.wild_fixedField_apfPsi, FontaineWintenberger.wild_compositum_apfPsi_over_right, FontaineWintenberger.wild_compositum_ramificationIdx_eq, FontaineWintenberger.wild_compositum_apfPsi_over_right_of_sup_eq_top, FontaineWintenberger.wild_compositum_apfPsi_scaling, FontaineWintenberger.wild_compositum_apfPhi_scaling, FontaineWintenberger.wild_compositum_herbrand_scaling_of_totallyRamified}
\leanok
Let $\Omega/E$ be finite Galois over a complete discrete valuation
field with arbitrary perfect residue field of characteristic $p$.
Let $P$ be its actual first lower group and put $A=\Omega^P$.
Let $F\subseteq\Omega$ be an actual totally ramified intermediate
field of $p$-power degree over $E$, put $M=AF$, and set
$e=e(A/E)>0$. Then $e(M/F)=e$ and, for every $u\geq0$,
\[
 \Psi_{M/F}(u)=eu,\qquad
 \Psi_{M/A}(eu)=e\Psi_{F/E}(u),\qquad
 \varphi_{M/A}(eu)=e\varphi_{F/E}(u).
\]
\end{lemma}
\begin{proof}
Write $H=\operatorname{Gal}(\Omega/F)$.
Write $D$ for the actual inertia group of $\Omega/E$.
The actual tame-compositum result gives $HP=G$.
Since $P\subseteq D$, each $d\in D$ can be written $d=hp$
with $p\in P$ and $h=dp^{-1}\in H\cap D$.
Thus $D=(H\cap D)P$, and the original coset map induces a
bijection from $(H\cap D)/(H\cap P)$ to $D/P$.
Consequently
$[D:P]=[H\cap D:H\cap P]$.
Actual lower-group compatibility, including the own integer-ring
identification, carries the inertia and first lower groups
for $\Omega/F$ onto $H\cap D$ and $H\cap P$.
The tame fixed-field ramification-index formula therefore gives
$e(M/F)=[H\cap D:H\cap P]=[D:P]=e(A/E)=e$.
The fixing group of $M=A F$ is $P\cap H$.
Actual lower-group compatibility identifies $P\cap H$ with the
first lower group of $\operatorname{Gal}(\Omega/F)$; the own
integer-ring identification preserves that equality.
Thus $M/F$ is exactly the actual wild fixed extension for
$\Omega/F$. The proved tame inverse Herbrand formula gives
$\Psi_{M/F}(u)=e(M/F)u=eu$.
The same formula gives $\Psi_{A/E}(u)=eu$.
Composition through the two actual finite bases therefore yields
\[
 \Psi_{M/A}(eu)=\Psi_{M/A}(\Psi_{A/E}(u))
 =\Psi_{M/E}(u)
 =\Psi_{M/F}(\Psi_{F/E}(u))=e\Psi_{F/E}(u).
\]
For the last equality use $\Psi_{F/E}(u)\geq u\geq0$.
To invert, let $v=\varphi_{F/E}(u)$.
Monotonicity and the value at zero give $v\geq0$.
The preceding identity gives $\Psi_{M/A}(ev)=eu$.
Applying its actual inverse gives
$\varphi_{M/A}(eu)=ev=e\varphi_{F/E}(u)$.
These arguments use the original finite fields throughout.
\end{proof}

\begin{lemma}[Herbrand continuity for the actual nonnormal wild norm]
\label{fw:nonnormal-wild-norm-continuity}
\source{wintenberger:page-0012,wintenberger:page-0013}
\uses{fw:wild-tower-norm-continuity,fw:wild-compositum-herbrand-scaling,fw:wild-tame-compositum,fw:ceil-precision-restriction,fw:norm-valuation}
\lean{FontaineWintenberger.intNorm_sub_congruence_of_totallyRamified_fixedField, FontaineWintenberger.intNorm_sub_congruence_of_totallyRamified_intermediateField, FontaineWintenberger.intNorm_sub_congruence_of_totallyRamified}
\leanok
Let $F/E$ be a finite separable totally ramified extension of
$p$-power degree, over a complete discrete valuation field
with arbitrary perfect residue field of characteristic $p$.
For every integer $r\geq0$ and $x,y\in\mathcal O_F$,
\[
 x-y\in\mathfrak m_F^r\quad\Longrightarrow\quad
 N_{F/E}(x)-N_{F/E}(y)
 \in\mathfrak m_E^{\lceil\varphi_{F/E}(r)\rceil_+}.
\]
No normality of $F/E$ is required.
\end{lemma}
\begin{proof}
Take an actual finite Galois normal closure $\Omega/E$ and
put $A=\Omega^P$, where $P$ is its actual wild inertia group.
Put $M=AF$ and $e=e(A/E)$.
The actual tame-compositum lemma makes $A$ and $F$ linearly
disjoint over $E$, identifies $[M:A]=[F:E]$, and identifies
the base change of the original norm with $N_{M/A}$.
The extension $M/A$ is totally ramified. Comparing ramification
indices in the two actual towers from $E$ to $M$ gives
$e(M/F)[F:E]=e[M:A]$, hence $e(M/F)=e$.
Moreover $M/F$ is the actual wild fixed field of $\Omega/F$:
its fixing group is $P\cap\operatorname{Gal}(\Omega/F)$,
which is precisely the first lower group for that subgroup.
The tame inverse Herbrand formulas therefore give
$\Psi_{A/E}(u)=eu$ and $\Psi_{M/F}(u)=eu$ for $u\geq0$.
Composition in the two actual towers yields
$\Psi_{M/A}(eu)=e\Psi_{F/E}(u)$, and inversion yields
\[
 \varphi_{M/A}(er)=e\varphi_{F/E}(r)\qquad(r\geq0).
\]
The inclusion $F\hookrightarrow M$ multiplies normalized
valuations by $e$, so the images of $x,y$ differ in
$\mathfrak m_M^{er}$. The preceding wild-tower lemma gives
$N_{M/A}(x)-N_{M/A}(y)\in
\mathfrak m_A^{\lceil e\varphi_{F/E}(r)\rceil_+}$.
Norm base change identifies this difference with the image
of the original $N_{F/E}(x)-N_{F/E}(y)$.
The original inclusion $E\hookrightarrow A$ also multiplies
normalized valuations by $e>0$.
The actual ceiling descent lemma therefore places the original
difference in
$\mathfrak m_E^{\lceil\varphi_{F/E}(r)\rceil_+}$, as required.
\end{proof}

\begin{lemma}[Norm precision on the original finite intermediate fields]
\label{fw:finite-layer-norm-precision}
\source{wintenberger:page-0009,wintenberger:page-0012,wintenberger:page-0013}
\dcref{Original-field transport of Proposition 2.2.1 and Lemma 2.3.2.1}
\uses{fw:nonnormal-wild-norm-continuity,fw:norm-estimates,fw:integral-finite-layer-norms,fw:finite-extension-e-d-invariance,fw:ideal-power-valuation}
\lean{FontaineWintenberger.FiniteLayer.integerNorm_sub_congruence_of_totallyRamified, FontaineWintenberger.FiniteLayer.exists_integerNorm_add_congruence_of_totallyRamified, FontaineWintenberger.FiniteLayer.integerNorm_add_sub_eq_zero_of_finrank_eq_one, FontaineWintenberger.FiniteLayer.integerNorm_add_congruence_of_le_apfRamificationBound}
\leanok
Let $E\subseteq F\subseteq L$ be actual finite intermediate fields
in an abstract separable extension of $K$. Equip $E$ and $F$ with
the canonical spectral norms over $K$, and write $R_E,R_F$ for
their original integer rings. Suppose the actual index of
$R_E\longrightarrow R_F$ is $e(F/E)=[F:E]=p^n$.
Then its actual integral norm satisfies, for $r\in\mathbb N$,
\[
 x-y\in\mathfrak m_F^r\quad\Longrightarrow\quad
 N_{F/E}(x)-N_{F/E}(y)\in
 \mathfrak m_E^{\lceil\varphi_{F/E}(r)\rceil_+}.
\]
For every real $b\leq i(F/E)$ it also satisfies
\[
 N_{F/E}(x+y)-N_{F/E}(x)-N_{F/E}(y)
 \in\mathfrak m_E^{\lceil(p-1)b/p\rceil_+}
 \qquad(x,y\in R_F).
\]
The first bound $i(F/E)$ has its actual extended-real value:
if $[F:E]=1$, the fields are equal, their norm is the identity,
and the addition defect is zero at every finite precision.
\end{lemma}
\begin{proof}
The canonical norm of $F$ is compatible with that of $E$.
Consequently the valuation ring obtained by extending the norm of
$E$ to $F$ is the same unit ball as $R_F$; the norm unit ball of
$E$ is $R_E$. These identifications preserve the original field
inclusions. A normalized valuation is determined by its DVR,
so the identifications preserve all normalized values and ideal powers.
They also preserve the actual ramification index: evaluate either
integer-ring inclusion on the same base uniformizer and use the
normalized valuation formula for that index.
Finally the actual integral norms agree, because under inclusion
in the original field they are both the determinant field norm.
Thus the preceding nonnormal continuity theorem applies on exactly
$R_E,R_F$ and gives the first assertion with the intrinsic relative
Herbrand function.

If $[F:E]>1$, the norm-estimate theorem supplies the actual finite
first bound $c=i(F/E)$ and addition precision
$\lceil(p-1)c/p\rceil_+$. Since $p\geq2$, multiplication by
$(p-1)/p$, taking the nonnegative ceiling, and then decreasing the
power of the maximal ideal carry this to the stated precision for
every $b\leq c$. If $[F:E]=1$, the inclusion and equality of finite
$K$-dimensions give $E=F$ as actual intermediate fields.
Its norm is the identity, so its addition defect is exactly zero.
This proves the bound for every finite $b$, without replacing the
infinite first bound of the identity step by a finite value.
\end{proof}

\begin{lemma}[The cyclic prime-degree norm congruence]
\label{fw:cyclic-prime-norm-congruence}
\source{wintenberger:page-0010,serre-1979-local-fields:page-0091}
\dcref{Wintenberger, Lemma 2.2.2.3; Serre, Chapter V, Section 3, Lemma 5}
\uses{fw:cyclic-prime-trace-bound}
\lean{FontaineWintenberger.exists_prod_add_sub_eq_sum_smul, FontaineWintenberger.spectralIntegerRing_algebraMap_intNorm_eq_prod, FontaineWintenberger.spectralIntegerRing_algebraMap_intTrace_eq_sum, FontaineWintenberger.spectralIntegerRing_exists_intNorm_add_sub_eq_intTrace, FontaineWintenberger.spectralIntegerRing_intNorm_add_sub_mem_intTraceIdeal, FontaineWintenberger.spectralIntegerRing_prime_degree_norm_estimate, FontaineWintenberger.spectralIntegerRing_prime_degree_norm_add_congruence}
\leanok
Let $E'/E$ be totally ramified Galois of prime degree $p$.
For all $\alpha,\beta\in\mathcal O_{E'}$, one has
\[
 N(\alpha+\beta)-N(\alpha)-N(\beta)
   \in\operatorname{Tr}_{E'/E}(\mathcal O_{E'}).
\]
Consequently its normalized valuation is at least
$(p-1)i(E'/E)/p$.
\end{lemma}
\begin{proof}
Write $G=\Gal(E'/E)$. Expanding the product formula for the
actual field norm gives
\[
 N(\alpha+\beta)=
 \sum_{S\subseteq G}
  \left(\prod_{g\in S}g(\alpha)\right)
  \left(\prod_{g\notin S}g(\beta)\right).
\]
The empty and full subsets contribute $N(\beta)$ and $N(\alpha)$.
The group $G$ acts on the remaining nonempty proper subsets by
left translation. The stabilizer of any such subset is a subgroup
of the prime-order group, hence is trivial or full. A nonempty
subset fixed by every translation contains every group element:
translate one of its elements to any prescribed element.
Since our subset is proper, its stabilizer is trivial.
Thus each middle orbit is indexed bijectively by $G$.
For an orbit representative $S$, put
\[
 x_S=\left(\prod_{g\in S}g(\alpha)\right)
     \left(\prod_{g\notin S}g(\beta)\right)\in\mathcal O_{E'}.
\]
Translation by $h$ replaces this monomial by $h(x_S)$.
The sum over its orbit is therefore
$\sum_{h\in G}h(x_S)=\operatorname{Tr}_{E'/E}(x_S)$.
Adding over the finitely many middle orbits proves the asserted
membership in the actual trace ideal. The preceding trace bound
applies to their sum. This expands both summands symmetrically;
it proves the source congruence without first dividing one
summand by the other.
\end{proof}

\begin{lemma}[Finite-precision norm congruences in a tower]
\label{fw:norm-congruence-tower}
\source{wintenberger:page-0009,wintenberger:page-0010}
\dcref{Wintenberger, Lemma 2.2.2.1 and its proof}
\uses{fw:norm-valuation}
\lean{FontaineWintenberger.intNorm_mem_maximalIdeal_pow, FontaineWintenberger.intNorm_sub_mem_maximalIdeal_pow_of_add_congruence, FontaineWintenberger.intNorm_add_congruence_tower}
\leanok
Let $R\hookrightarrow S\hookrightarrow T$ be finite injective
local maps of DVRs whose fraction-field extensions are separable,
and let $r\geq0$ be an integer. Integral norms send
$\mathfrak m_S^r$ into $\mathfrak m_R^r$.
If
\[
 N_{S/R}(x+y)-N_{S/R}(x)-N_{S/R}(y)\in\mathfrak m_R^r
 \quad(x,y\in S),
\]
then $x-y\in\mathfrak m_S^r$ implies
$N_{S/R}(x)-N_{S/R}(y)\in\mathfrak m_R^r$.
If the analogous addition congruence also holds for $T/S$ at
the same precision $r$, it holds for $T/R$ at precision $r$.
\end{lemma}
\begin{proof}
For $x\in\mathfrak m_S^r$, the normalized norm identity gives
$v_R(N(x))=f(S/R)v_S(x)\geq v_S(x)\geq r$, since the
finite residue degree $f(S/R)$ is a positive integer. This
also includes $x=0$, with infinite valuation.
For the second assertion put $z=x-y$. The assumed addition
congruence applied to $y$ and $z$, together with
$N(z)\in\mathfrak m_R^r$, gives
$N(x)-N(y)\in\mathfrak m_R^r$.

For $a,b\in T$ put $u=N_{T/S}(a)$, $v=N_{T/S}(b)$ and
$w=N_{T/S}(a+b)$. The top addition congruence gives
$w-(u+v)\in\mathfrak m_S^r$. The second assertion therefore
gives $N_{S/R}(w)-N_{S/R}(u+v)\in\mathfrak m_R^r$.
The bottom addition congruence gives
$N_{S/R}(u+v)-N_{S/R}(u)-N_{S/R}(v)\in\mathfrak m_R^r$.
Add these two containments and use the actual transitivity
identity $N_{S/R}N_{T/S}=N_{T/R}$ to conclude.
This finite-precision form applies successively to a tower once
the genuine addition congruence is proved for each step.
\end{proof}

\begin{lemma}[The field fixed by wild inertia]
\label{fw:wild-fixed-tame-field}
\source{wintenberger:page-0010}
\dcref{Wintenberger, paragraph 2.2.2.5}
\uses{fw:wild-ramification-p-group,fw:tame-inertia-quotient,fw:local-inertia-size,fw:finite-tower-upper-restriction,fw:upper-inverse-integral,fw:herbrand-functions}
\lean{FontaineWintenberger.upperRamificationGroup_le_lower_one, FontaineWintenberger.lowerRamificationGroup_subgroup_zero_eq_top, FontaineWintenberger.inverseHerbrandFunction_eq_card_inertia_mul, FontaineWintenberger.wild_fixedField_top_ramificationIdx_eq_finrank, FontaineWintenberger.wild_fixedField_upper_eq_bot, FontaineWintenberger.wild_fixedField_ramificationIdx_eq_relIndex, FontaineWintenberger.wild_fixedField_inverseHerbrandFunction, FontaineWintenberger.wild_fixedField_ramificationIdx_coprime, FontaineWintenberger.wild_fixedField_isPGroup}
\leanok
Let $\Omega/E$ be finite Galois, where $E$ is complete discretely
valued with perfect residue field of characteristic $p$. Put
$G=\Gal(\Omega/E)$, $D=G_0$, $P=G_1$, and $E_1=\Omega^P$.
Then $\Omega/E_1$ is totally ramified with $p$-group Galois group,
$e_1=e(E_1/E)=|D/P|$ is prime to $p$, and for $u\geq0$,
\[
 \psi_{E_1/E}(u)=e_1u.
\]
For $u>0$ the actual upper group $G^u$ is contained in $P$.
\end{lemma}
\begin{proof}
The subgroup $P$ is normal and is a $p$-group. Its inertia group
for $\Omega/E_1$ is $P\cap D=P$: the action on the actual
valuation ring is the restricted original action. The cardinality
of inertia is its ramification index; since the Galois group is
$P$, this proves total ramification. In the quotient $G/P$, inertia
is the image of $D$, hence has order $|D/P|$. The tame inertia
homomorphism identifies $D/P$ with a finite subgroup of the
multiplicative group of the residue field. Its order is prime to
$p$, giving the assertion about $e_1$.

For $u>0$, positivity of the inverse Herbrand function gives
$\psi_{\Omega/E}(u)>0$. Its ceiling is at least one, so lower
monotonicity gives $G^u=G_{\psi(u)}\subseteq G_1=P$.
Upper quotient compatibility therefore makes every positive upper
group of $\Gal(E_1/E)$ trivial. Its upper group at zero is its
inertia group, of order $e_1$. The derivative of its inverse
Herbrand function on the positive half-line is consequently the
constant $e_1$. Integrating from zero, where the function is zero,
proves $\psi_{E_1/E}(u)=e_1u$, also at $u=0$.
\end{proof}

\begin{lemma}[Actual total ramification above the wild fixed field]
\label{fw:wild-relative-total-ramification}
\uses{fw:ramification-subgroups,fw:local-inertia-size,fw:spectral-base-integers,fw:spectral-integer-residue-maps,fw:total-ramification-tower}
\lean{FontaineWintenberger.finiteBase_lower_zero_eq_top, FontaineWintenberger.wildField_lower_zero_eq_top, FontaineWintenberger.wildField_top_ramificationIdx_eq_finrank, FontaineWintenberger.wildField_step_ramificationIdx_eq_finrank}
\leanok
Let $\Omega/K$ be finite Galois, put $P=G_1$ and
$F_1=\Omega^P$. For every actual intermediate field
$F_1\subseteq A\subseteq\Omega$, the actual lower zero group
$\operatorname{Gal}(\Omega/A)_0$ is the whole group and
$e(\Omega/A)=[\Omega:A]$. The integer ring on the left is the
spectral integer ring constructed over $A$. Consequently every
actual finite tower between $F_1$ and $\Omega$, in particular each
step of the cyclic prime tower, is totally ramified.
\end{lemma}
\begin{proof}
Galois correspondence gives
$\operatorname{Gal}(\Omega/A)\subseteq P\subseteq G_0$.
The subgroup lower-numbering theorem, through the identity of the
actual spectral integer rings, therefore makes its own lower zero
group the whole group. Its inertia cardinality is its actual
ramification index, and the Galois degree theorem identifies the
full group cardinality with $[\Omega:A]$.
In particular $\Omega/F_1$ is totally ramified. The local tower
criterion applied to its actual integer-ring towers makes both
factors totally ramified for each intermediate field, and hence
makes every step of any such finite tower totally ramified.
\end{proof}

\begin{lemma}[The actual tame compositum and its norm]
\label{fw:wild-tame-compositum}
\source{wintenberger:page-0010}
\dcref{Wintenberger, paragraph 2.2.2.5}
\uses{fw:inertia-relative-index,fw:wild-fixed-tame-field,fw:finite-extension-fixing-index}
\lean{FontaineWintenberger.fixedField_sup_wild_eq_top, FontaineWintenberger.fixedField_wild_linearDisjoint, FontaineWintenberger.fixedField_wild_compositum_finrank, FontaineWintenberger.fixedField_wild_norm_baseChange}
\leanok
Let $E'/E$ be finite separable totally ramified of degree $q=p^n$,
let $\Omega/E$ be a finite Galois extension containing $E'$, and
put $G=\Gal(\Omega/E)$, $H=\Gal(\Omega/E')$, $P=G_1$,
$E_1=\Omega^P$, and $M=E'E_1$.
Then $HP=G$, the actual fields $E'$ and $E_1$ are linearly
disjoint over $E$, and
\[
 [M:E_1]=[E':E],\qquad
 N_{M/E_1}(\alpha)=N_{E'/E}(\alpha)\quad(\alpha\in E'),
\]
with both sides of the norm formula compared under the actual
original inclusions. No normality of $E'/E$ is assumed.
\end{lemma}
\begin{proof}
Write $D=G_0$. The actual inertia-relative index formula gives
$[D:H\cap D]=e(E'/E)=q=[G:H]$, hence $HD=G$.
The subgroup $(H\cap D)P$ has index in $D$ dividing both
$q$ and $[D:P]$. The latter is prime to $p$, so this index
is one. Thus $D=(H\cap D)P$ and $HP=G$.
Since $P$ is normal, any element of $E'\cap E_1$ is fixed
by $H$ and $P$, hence by $HP=G$, and therefore lies in $E$.
The finite Galois extension $E_1/E$ is consequently linearly
disjoint from $E'/E$. Equivalently the subgroup index formula
gives $[M:E_1]=[P:H\cap P]=[G:H]=q$.

An $E$-basis of $E'$ spans $M$ over $E_1$ and has $q$ elements;
the displayed degree makes it an $E_1$-basis. Multiplication
by $\alpha\in E'$ has the same matrix after this scalar
extension. Equality of the determinants gives the actual norm
formula. This uses coprimality of $e(E_1/E)$ and does not assert
coprimality of the entire degree $[E_1:E]$, whose unramified
part may have $p$-power degree.
\end{proof}

\begin{lemma}[Finite first-bound base change in an actual overfield]
\label{fw:finite-overfield-first-bound-base-change}
\uses{fw:finite-base-upper,fw:finite-galois-overfield-apf,fw:apf-herbrand-extended,fw:normal-product-relative-index}
\lean{FontaineWintenberger.finiteOverfield_fixing_sup_upper_eq_top, FontaineWintenberger.apfHerbrandERealOrderIso_bound_le_of_finiteGalois, FontaineWintenberger.apfRamificationBound_le_finiteBase_of_finiteGalois}
\leanok
Let $\Omega/K$ be any finite Galois extension and let
$M\subseteq N$ be actual intermediate fields. Then
\[
 \overline\psi_{M/K}(i(N/K))\leq i(N/M),\qquad
 i(N/K)\leq i(N/M).
\]
The first bounds use the actual fields with their compatible
norms and include both extended-real endpoints.
\end{lemma}
\begin{proof}
Put $G=\Gal(\Omega/K)$, $J=\Gal(\Omega/N)$ and
$H=\Gal(\Omega/M)$, so $J\subseteq H$. If $u\geq-1$ and
$JG^u=G$, normality of $G^u$ gives
\[
 (JG^u)\cap H=J(G^u\cap H)=H.
\]
Indeed a product $jg$ belongs to $H$ exactly when $g$ does,
since $j\in J\subseteq H$. The canonical fixing-group
isomorphism and finite upper base change identify the last
product with the actual product over $M$ at parameter
$\psi_{M/K}(u)$. This parameter is at least $-1$ and is thus
admissible for $i(N/M)$. The actual extended Herbrand order
isomorphism preserves arbitrary suprema; applying it to the
actual-overfield admissible-set formula for $i(N/K)$ gives
the first inequality. Its pointwise lower bound by the identity
gives the second inequality, including infinite endpoints.
\end{proof}

\begin{corollary}[A common bound for actual intermediate steps]
\label{fw:finite-overfield-step-first-bound}
\uses{fw:finite-overfield-first-bound-base-change,fw:apf-inheritance}
\lean{FontaineWintenberger.apfRamificationBound_le_step_of_finiteGalois}
\leanok
In any actual finite Galois overfield $\Omega/K$, let
$A\subseteq B\subseteq N$ be actual intermediate fields. Then
\[
 i(N/K)\leq i(B/A).
\]
\end{corollary}
\begin{proof}
Finite base change gives $i(N/K)\leq i(N/A)$.
The actual inclusion $B\hookrightarrow N$ is an $A$-embedding,
so descent of the intrinsic first bound gives
$i(N/A)\leq i(B/A)$. Composing proves the assertion, including
trivial steps and extended-real endpoints.
\end{proof}

\begin{lemma}[First-bound scaling by the actual tame field]
\label{fw:wild-compositum-first-bound}
\uses{fw:finite-totally-ramified-first-bound,fw:finite-base-upper,fw:wild-fixed-tame-field,fw:ramification-subgroups,fw:wild-tame-compositum}
\lean{FontaineWintenberger.exists_apfRamificationBound_wild_compositum, FontaineWintenberger.exists_apfRamificationBound_wild_compositum_of_totallyRamified}
\leanok
Let $\Omega/K$ be finite Galois, $G=\Gal(\Omega/K)$,
$P=G_1$, $A=\Omega^P$, and let a proper subgroup $H$ satisfy
$HP=G$. Put $F=\Omega^H$ and $M=AF$, regarded as an actual
extension of $A$. For $e=e(A/K)$ one has
\[
 i(M/A)=e\,i(F/K).
\]
Both first bounds are finite nonnegative real numbers. In particular
this applies to every nontrivial totally ramified $p$-power-degree
$F/K$, by transitivity of actual wild inertia.
\end{lemma}
\begin{proof}
Since $P\subseteq G_0$ and $HP=G$, the finite attained first-bound
lemma gives $b\geq0$ for $F/K$. The fixing subgroup of $M$ is
$H\cap P$. Inside $P$ it is proper: otherwise $P\subseteq H$
and $G=HP=H$, contradicting properness. The own-ring inertia of
$\Omega/A$ is all its Galois group. Applying the same finite
first-bound lemma over $A$ gives $c\geq0$ for $M/A$.

For $u>0$ put $U=G^u\subseteq P$. The group $U$ is normal in
$G$. If $HU=G$, write each $g\in P$ as $g=hu$ with $h\in H$
and $u\in U$. Then $h=gu^{-1}\in H\cap P$, so
$(H\cap P)U=P$. Conversely this last equality and $HP=G$
imply $HU=G$. The canonical fixing-group isomorphism identifies
this product in $P$ with the actual product over $A$. Finite base
change and the actual tame inverse Herbrand formula identify its
upper parameter with $\psi_{A/K}(u)=eu$. Hence, for every $u>0$,
\[
 u\leq b\quad\Longleftrightarrow\quad eu\leq c.
\]
The positive integer $e$ is nonzero by the inertia-cardinality
formula. If $b>0$, substitution of $u=b$ yields $eb\leq c$;
if $b=0$ this follows from $c\geq0$. If $c>0$, substitution of
$u=c/e$ yields $c\leq eb$; if $c=0$ this follows from $b\geq0$.
Thus $c=eb$, including the zero-break case. The final assertion
uses the proved wild-inertia transitivity of the actual totally
ramified $p$-power subextension.
\end{proof}

\begin{lemma}[Tame reduction of a nonnormal wild layer]
\label{fw:nonnormal-tame-norm-reduction}
\source{wintenberger:page-0010,wintenberger:page-0011}
\dcref{Wintenberger, paragraph 2.2.2.5}
\uses{fw:wild-compositum-first-bound,fw:wild-tame-compositum,fw:wild-fixed-tame-field,fw:relative-prime-field-tower,fw:wild-relative-total-ramification,fw:local-inertia-size,fw:finite-base-upper,fw:normalized-valuation-scaling,fw:apf-finite-base-bounds,fw:finite-overfield-first-bound-base-change,fw:finite-overfield-step-first-bound}
\notready
Let $E'/E$ be finite separable totally ramified of degree $q=p^n$.
Choose a finite Galois closure $\Omega/E$ containing $E'$, put
$P=\Gal(\Omega/E)_1$, $E_1=\Omega^P$, and $M=E'E_1$.
Then $E'$ and $E_1$ are linearly disjoint over $E$,
$[M:E_1]=q$, and $M/E_1$ is totally ramified and admits a finite
tower of cyclic Galois degree-$p$ extensions. Writing
$e_1=e(E_1/E)$, one has
\[
 i(M/E_1)=e_1 i(E'/E),\qquad
 N_{M/E_1}(\alpha)=N_{E'/E}(\alpha)\quad(\alpha\in E'),
\]
where the second equality uses the actual inclusions into $M$ and
$E_1$. The first ramification bound of every step in the displayed
prime-degree tower is at least $i(M/E_1)$.
\end{lemma}
\begin{proof}
The tame-compositum lemma supplies linear disjointness, the actual
relative degree and the actual norm identity. Retain its notation
$G,H,D,P$ and $HP=G$.

The wild-fixed-field lemma makes $\Omega/E_1$ totally ramified with
$p$-group Galois group. All its intermediate extensions are
totally ramified, by multiplicativity of ramification indices
and residue degrees and the identity $[B:A]=e(B/A)f(B/A)$.
The actual relative prime-field tower gives the claimed cyclic
degree-$p$ chain from $E_1$ to $M$.

For $u>0$ put $N=G^u\subseteq P$. Since $HP=G$, one has
\[
 HN=G\quad\Longleftrightarrow\quad (H\cap P)N=P.
\]
Indeed intersecting $HN$ with $P$ gives $(H\cap P)N$; the
converse follows by multiplying by $H$. Finite base change
identifies the upper group of $\Omega/E_1$ at
$\psi_{E_1/E}(u)=e_1u$ with $N$. Thus the transitivity condition
defining the first ramification bound is equivalent at parameters
$u$ and $e_1u$. At zero both actions are transitive because both
extensions are totally ramified. Taking suprema proves the bound
identity. Finally, for an intermediate step $A\subset B$ in the
chain, descent gives $i(B/A)\geq i(M/A)$, and finite base change
gives
$i(M/A)\geq\psi_{A/E_1}(i(M/E_1))\geq i(M/E_1)$, since this inverse
Herbrand function has slopes at least one. This proves the common
lower bound for all steps.
\end{proof}

\begin{lemma}[Integral norms on all actual finite layers]
\label{fw:integral-finite-layer-norms}
\uses{fw:norm-limit-multiplicative,fw:spectral-integer-residue-maps,fw:norm-valuation,fw:norm-congruence-tower}
\lean{FontaineWintenberger.FiniteLayer.integerNorm, FontaineWintenberger.FiniteLayer.coe_integerNorm, FontaineWintenberger.FiniteLayer.integerNorm_self, FontaineWintenberger.FiniteLayer.integerNorm_integerNorm, FontaineWintenberger.FiniteLayer.integerNorm_mem_maximalIdeal_pow}
\leanok
Let $L/K$ be separable algebraic and index all its actual finite
intermediate fields, including nonnormal ones. For $E\subseteq F$,
the actual field norm restricts to a multiplicative map
\[
 N^{\mathcal O}_{F/E}:\mathcal O_F\longrightarrow\mathcal O_E.
\]
Its inclusion into $E$ is precisely the original field norm.
These maps are the identity when $E=F$ and compose in every
actual finite field tower. Moreover, for every $r\geq0$,
\[
 N^{\mathcal O}_{F/E}(\mathfrak m_F^r)\subseteq\mathfrak m_E^r.
\]
\end{lemma}
\begin{proof}
The actual field inclusion gives the local spectral integer-ring
inclusion over the base integers. Both integer rings are finite
over the base integers, so the top ring is also finite over the
bottom one. The integral norm is therefore defined on these
actual rings; its image under the fraction-field inclusions is
the field norm. This gives the asserted restriction, uniquely
because the inclusion into $E$ is injective. Identity and
composition follow from the actual field-norm identities and
injectivity of that same inclusion. For $x\in\mathfrak m_F^r$,
the normalized valuation formula gives
$v_E(N(x))=f(F/E)v_F(x)\geq r$, since the actual residue degree
$f(F/E)$ is a positive integer. This is exactly membership in
$\mathfrak m_E^r$, including $x=0$.
\end{proof}

\begin{lemma}[Norm congruences along an actual intermediate-field chain]
\label{fw:actual-field-norm-congruence-chain}
\uses{fw:integral-finite-layer-norms,fw:norm-congruence-tower,fw:spectral-integer-residue-maps,fw:spectral-base-integers}
\lean{FontaineWintenberger.FiniteLayer.integerNorm_add_congruence_tower, FontaineWintenberger.FiniteLayer.integerNorm_add_congruence_of_chain}
\leanok
Let $\Omega/K$ be separable algebraic, let $M\subseteq N$ be
actual finite intermediate fields, and fix $r\geq0$. Suppose there is a
finite increasing chain from $M$ to $N$ such that the actual
integral norm of each step is additive modulo the $r$-th power
of the bottom maximal ideal of that step. Then the actual
integral norm from $N$ to $M$ is additive modulo
$\mathfrak m_M^r$.
\end{lemma}
\begin{proof}
Use the actual spectral integers of each intermediate field over
$K$. Uniqueness of extended norms identifies them with their own
unit balls. The actual field inclusions give local maps of these
rings, and maps in a triple compose by the actual field inclusions.
All rings are DVRs finite over the original base integers; hence
each ring is also finite over the preceding one. Their fraction
fields are the original intermediate fields, so the extensions
are separable. In a chain of length zero the norm is the identity
and its addition defect is zero. For an additional step, the
actual three rings form a scalar tower. Apply the proved integral
norm-congruence tower lemma at precision $r$ to the preceding
norm and this step norm. Its norm identity is transitivity of
the actual integral norms on the same fields. Induction over the
finite chain proves the assertion. Identity of the spectral and
own integer rings transports the conclusion to the stated norm.
\end{proof}

\begin{lemma}[The common integer precision of prime-degree norms]
\label{fw:prime-norm-common-precision}
\uses{fw:cyclic-prime-trace-bound,fw:cyclic-prime-norm-congruence}
\lean{FontaineWintenberger.mul_div_le_prime_norm_precision, FontaineWintenberger.ceil_mul_div_le_prime_norm_precision, FontaineWintenberger.spectralIntegerRing_prime_norm_add_congruence_of_bound}
\leanok
Let $L/K$ be an actual totally ramified Galois extension of prime
degree $p$. If a real number $b$ satisfies $b\leq i(L/K)$, then
for every integral $x,y\in\mathcal O_L$,
\[
 N(x+y)-N(x)-N(y)\in
 \mathfrak m_K^{\lceil (p-1)b/p\rceil_+},
\]
where $\lceil\cdot\rceil_+$ is the nonnegative integer ceiling.
\end{lemma}
\begin{proof}
The actual cyclic prime-degree producer gives an integer $t\geq0$
with $i(L/K)=t$ and norm addition congruence modulo
$\mathfrak m_K^q$, where $q=\lfloor(p-1)(t+1)/p\rfloor$.
Write $(p-1)(t+1)=pq+a$ with $0\leq a\leq p-1$.
Subtracting $p-1$ gives $(p-1)t\leq pq$, so
$(p-1)t/p\leq q$. Since $b\leq t$ and $(p-1)/p\geq0$,
$\lceil(p-1)b/p\rceil_+\leq q$. Powers of the maximal ideal
decrease with their exponent, so the actual congruence at $q$
implies the claimed congruence. This also includes negative $b$,
whose nonnegative ceiling is zero.
\end{proof}

\begin{lemma}[Prime-step congruences on the original finite-layer integers]
\label{fw:finite-layer-prime-norm-congruence}
\uses{fw:integral-finite-layer-norms,fw:prime-norm-common-precision,fw:spectral-base-integers,fw:finite-base-residue}
\lean{FontaineWintenberger.FiniteLayer.integerNorm_prime_add_congruence_of_bound}
\leanok
For actual finite layers $E\subseteq F$ of a separable extension
of $K$, suppose $F/E$ is Galois, totally ramified and of prime
degree $p$. If $b\leq i(F/E)$, their original integral norm satisfies
\[
 N^{\mathcal O}_{F/E}(x+y)-N^{\mathcal O}_{F/E}(x)
    -N^{\mathcal O}_{F/E}(y)
 \in\mathfrak m_E^{\lceil(p-1)b/p\rceil_+}
 \qquad(x,y\in\mathcal O_F).
\]
\end{lemma}
\begin{proof}
Give the algebraic ambient extension its spectral norm over $K$
and give $E$ and $F$ their inherited norms. These are compatible
with the spectral norms by uniqueness and the proved restriction
identity. The complete field $E$ is discretely valued
and its residue field is perfect, as a finite extension of the
perfect residue field of $K$. Its own integer ring is the
original spectral integer ring of $E$ over $K$. The two spectral
constructions of the integer ring of $F$ agree as actual subrings
of $F$, since both are the unit ball for this same norm.
The identity-ring equivalences identify their maximal ideals and
all their powers. The integral norms agree under these equivalences:
both include into $E$ as the same actual field norm, and that
inclusion is injective. Apply the actual prime-degree congruence
at common precision $b$ over $E$, then transport along these
identity equivalences. This gives the claimed statement on the
original finite-layer integer rings.
\end{proof}

\begin{lemma}[The actual wild-field norm congruence]
\label{fw:actual-wild-field-norm-congruence}
\source{wintenberger:page-0010,wintenberger:page-0011}
\dcref{Wintenberger, paragraphs 2.2.2.5--2.2.2.6}
\uses{fw:relative-prime-field-tower,fw:wild-relative-total-ramification,fw:finite-overfield-step-first-bound,fw:finite-layer-prime-norm-congruence,fw:actual-field-norm-congruence-chain}
\lean{FontaineWintenberger.FiniteLayer.integerNorm_add_congruence_of_wild_field}
\leanok
Let $\Omega/K$ be finite Galois with perfect residue field of
characteristic $p$. Put $P=\operatorname{Gal}(\Omega/K)_1$.
For actual intermediate fields $\Omega^P\subseteq E\subseteq F$
and any real $b\leq i(F/E)$, the original integral norm satisfies
\[
 N_{F/E}(x+y)-N_{F/E}(x)-N_{F/E}(y)
 \in\mathfrak m_E^{\lceil(p-1)b/p\rceil_+}
 \qquad(x,y\in\mathcal O_F).
\]
\end{lemma}
\begin{proof}
Galois correspondence embeds $\operatorname{Gal}(\Omega/E)$
in the finite $p$-group $P$. The actual relative prime-field tower
therefore gives a finite increasing chain from $E$ to $F$ whose
steps are cyclic Galois of degree $p$. Every intermediate field
in this increasing chain contains $\Omega^P$, so every step is
totally ramified on its actual integer rings. For a step
$A\subset B$, both fields are contained in $F$; the actual
intermediate-step bound gives
$b\leq i(F/E)\leq i(B/A)$. The prime-step congruence for the
original finite-layer integers consequently holds at the common
integer precision $\lceil(p-1)b/p\rceil_+$ for each step.
The proved actual norm-chain congruence composes these same
integral norm maps. For a chain of length zero, the norm is the
identity and its addition defect is zero. This proves the
assertion in all cases.
\end{proof}

\begin{lemma}[Finite perfect-residue power digits]
\label{fw:finite-perfect-residue-digits}
\uses{fw:ideal-power-valuation}
\lean{FontaineWintenberger.exists_pow_digit_uniformizer, FontaineWintenberger.exists_pow_digits_uniformizer, FontaineWintenberger.exists_pow_digits_sub_mem_maximalIdeal_pow}
\leanok
Let $R$ be a DVR with uniformizer $\pi$ and perfect residue field
of characteristic $p$. For every $n,r\geq0$ and $a\in R$ there
exist $b_0,\ldots,b_{r-1},c\in R$ such that
\[
 a=\sum_{j<r}b_j^{p^n}\pi^j+\pi^r c.
\]
In particular the difference from the displayed finite sum belongs
to $\mathfrak m_R^r$. No completeness or finiteness of the residue
field is required.
\end{lemma}
\begin{proof}
The $n$-fold Frobenius of the actual perfect residue field is a
bijection. Given $a$, lift a $p^n$-th root of its residue to $b$.
Then $a-b^{p^n}\in\mathfrak m_R=(\pi)$, so
$a=b^{p^n}+\pi a_1$ for some $a_1\in R$. Repeating this
construction $r$ times and substituting the successive identities
gives the displayed formula, with $c=a_r$. For $r=0$ it is the
empty sum and $c=a$. The remainder is divisible by $\pi^r$,
which generates $\mathfrak m_R^r$.
\end{proof}

\begin{lemma}[Finite residue digits adapted to a wild norm]
\label{fw:finite-norm-residue-digits}
\source{wintenberger:page-0011}
\dcref{Wintenberger, paragraph 2.2.2.6, with finite residue lifts}
\uses{fw:norm-valuation,fw:finite-perfect-residue-digits,fw:local-dvr-degree-identity,fw:chosen-extension-integers}
\lean{FontaineWintenberger.irreducible_intNorm_of_inertiaDeg_eq_one, FontaineWintenberger.intNorm_sum_sub_sum_intNorm_mem_maximalIdeal_pow, FontaineWintenberger.exists_intNorm_sub_mem_maximalIdeal_pow, FontaineWintenberger.extensionValuationSubring_exists_intNorm_sub_mem_maximalIdeal_pow}
\leanok
Let $E'/E$ be finite separable totally ramified of degree $q=p^n$,
with perfect residue field of characteristic $p$. Let $r\geq0$
be an integer and suppose its integral norm is additive modulo
$\mathfrak m_E^r$. For every $a\in\mathcal O_E$ there is
$\alpha\in\mathcal O_{E'}$ with
$N_{E'/E}(\alpha)-a\in\mathfrak m_E^r$.
\end{lemma}
\begin{proof}
Choose a uniformizer $\pi'$ of $E'$ and put $\pi=N(\pi')$.
The normalized norm identity and residue degree one give
$v_E(\pi)=1$, so $\pi$ is a uniformizer of $E$.
If $r=0$, choose $\alpha=0$. For positive $r$, construct
$b_0,\ldots,b_{r-1}\in\mathcal O_E$ recursively. Start with
$a_0=a$. Since the residue field is perfect, choose the residue
of $b_j$ to be a $q$-th root of the residue of $a_j$.
Then $a_j-b_j^q\in(\pi)$, so there is $a_{j+1}\in\mathcal O_E$
with $a_j=b_j^q+\pi a_{j+1}$. Substitution in these $r$ identities
gives
\[
 a=\sum_{j<r}b_j^q\pi^j+\pi^r a_r.
\]
Set $\alpha=\sum_{j<r}b_j(\pi')^j$. Repeated use of the assumed
norm addition congruence, actual norm multiplicativity and the
scalar norm identity gives
\[
 N(\alpha)\equiv\sum_{j<r}N(b_j)N(\pi')^j
   =\sum_{j<r}b_j^q\pi^j\equiv a\pmod{\mathfrak m_E^r}.
\]
Only finitely many residue lifts have been used.
\end{proof}

\begin{lemma}[Ceiling precision under restriction of valuations]
\label{fw:ceil-precision-restriction}
\uses{fw:normalized-valuation-scaling,fw:ideal-power-valuation}
\lean{FontaineWintenberger.mem_maximalIdeal_pow_of_algebraMap_mem_ceil}
\leanok
Let $R\longrightarrow S$ be an injective local map of DVRs with
actual ramification index $e>0$. For $a\in R$ and $c\in\mathbb R$,
\[
 a\in\mathfrak m_S^{\lceil ec\rceil_+}
 \quad\Longrightarrow\quad
 a\in\mathfrak m_R^{\lceil c\rceil_+},
\]
where $a$ on the left denotes its actual image in $S$.
\end{lemma}
\begin{proof}
Normalized valuation restriction gives $v_S(a)=e v_R(a)$.
If $a=0$, both memberships hold. Otherwise put $v_R(a)=n\geq0$.
The left membership implies $\lceil ec\rceil_+\leq en$, hence
$ec\leq en$. Positivity of $e$ gives $c\leq n$, and the
nonnegative integer $n$ consequently satisfies
$\lceil c\rceil_+\leq n$. This is precisely the right membership.
\end{proof}

\begin{lemma}[Norm addition after actual tame descent]
\label{fw:tame-norm-addition-descent}
\source{wintenberger:page-0010,wintenberger:page-0011}
\dcref{Wintenberger, paragraph 2.2.2.5}
\uses{fw:wild-compositum-first-bound,fw:wild-tame-compositum,fw:actual-wild-field-norm-congruence,fw:ceil-precision-restriction,fw:integral-finite-layer-norms}
\lean{FontaineWintenberger.exists_intNorm_add_congruence_of_totallyRamified_fixedField}
\leanok
Let $\Omega/K$ be any finite Galois extension, and let
$F=\Omega^H$ be a nontrivial totally ramified extension of degree
$p^n$, without requiring $H$ normal. Its actual first bound is a
finite real $b\geq0$, and its original integral norm satisfies
\[
 N_{F/K}(x+y)-N_{F/K}(x)-N_{F/K}(y)
 \in\mathfrak m_K^{\lceil(p-1)b/p\rceil_+}
 \qquad(x,y\in\mathcal O_F).
\]
\end{lemma}
\begin{proof}
Put $P=\operatorname{Gal}(\Omega/K)_1$, $A=\Omega^P$,
$M=AF$ and $e=e(A/K)$. The actual tame-compositum producers
give $HP=G$, linear disjointness and
$N_{M/A}(z)=N_{F/K}(z)$ under the original field inclusions.
First-bound scaling gives $i(F/K)=b\geq0$ and $i(M/A)=eb$.
Include $x,y$ in the actual integers of $M$. Apply the actual
wild-field congruence from $A$ to $M$ at its own bound $eb$.
It puts the resulting defect in
$\mathfrak m_A^{\lceil(p-1)eb/p\rceil_+}$.
Each integral norm includes into its original field as the actual
field norm. Thus the field-norm base-change identity and
injectivity of the integer-ring inclusion identify this defect
with the actual image of the original norm defect in
$\mathcal O_K$. Since $(p-1)eb/p=e((p-1)b/p)$, ceiling precision
restriction along $\mathcal O_K\longrightarrow\mathcal O_A$
gives exactly the asserted original precision. This uses the
actual ramification index $e$ and allows an unramified part of
$A/K$; it makes no assertion that $[A:K]$ is prime to $p$.
\end{proof}

\begin{proposition}[Norm estimates in a wild finite layer]
\label{fw:norm-estimates}
\dcref{Wintenberger1983, Proposition 2.2.1}
\source{wintenberger:page-0009}
\uses{fw:apf,fw:norm-valuation,fw:tame-norm-addition-descent,fw:finite-norm-residue-digits,fw:finite-extension-e-d-invariance,fw:apf-extended-bound,fw:ramification-field-equiv}
\lean{FontaineWintenberger.exists_intNorm_add_congruence_of_totallyRamified_intermediateField, FontaineWintenberger.exists_intNorm_add_congruence_of_totallyRamified, FontaineWintenberger.exists_intNorm_estimates_of_totallyRamified}
\leanok
Let $E'/E$ be a nontrivial finite separable totally ramified extension of
$p$-power degree, and put $i=i(E'/E)$.
For $\alpha,\beta\in\mathcal O_{E'}$ one has
\[
 v_E\bigl(N_{E'/E}(\alpha+\beta)
            -N_{E'/E}(\alpha)-N_{E'/E}(\beta)\bigr)
       \geq\frac{p-1}{p}i.
\]
For every $a\in\mathcal O_E$, there is $\alpha\in\mathcal O_{E'}$ with
\[
 v_E\bigl(N_{E'/E}(\alpha)-a\bigr)\geq\frac{p-1}{p}i.
\]
\end{proposition}
\begin{proof}
Embed the original finite separable extension in a separable closure
and take its actual finite Galois closure $\Omega/E$. Its image $F$
is isomorphic to the original extension over $E$. The canonical
isomorphism of their actual integer rings preserves the
ramification index, ordinary field degree and intrinsic first
bound. It also preserves integral norms: both include as the
same actual field norms under the original field isomorphism.
Thus $F/E$ is a nontrivial totally ramified extension of the same
$p$-power degree. The proved tame norm-addition descent gives a
finite $b\geq0$ equal to $i(F/E)=i(E'/E)$, and gives addition
congruence at $r=\lceil(p-1)b/p\rceil_+$. Transport the original
integers through this isomorphism. Equality of integral norms
makes the same congruence hold for the original extension.
Membership in $\mathfrak m_E^r$ implies the first displayed
normalized valuation bound.

Apply the finite residue-digit lemma to this actual original
extension at the same precision $r$, using its proved norm
addition congruence. For every $a\in\mathcal O_E$ it produces
$\alpha\in\mathcal O_{E'}$ with
$N(\alpha)-a\in\mathfrak m_E^r$. Its normalized valuation is at
least $r\geq(p-1)b/p$, which is the second assertion. Only
finitely many residue lifts are used; perfectness is sufficient
and no finiteness of the residue field is required.
\end{proof}

\begin{lemma}[Continuity of the original finite norms]
\label{fw:finite-norm-continuity}
\source{wintenberger:page-0012,wintenberger:page-0013}
\dcref{Local determinant argument for the norms in Section 2.3}
\uses{fw:spectral-norm-extension,fw:norm-limit-multiplicative}
\lean{FontaineWintenberger.continuous_algebraNorm_of_finiteDimensional, FontaineWintenberger.FiniteLayer.continuous_norm}
\leanok
Let $K\subseteq E\subseteq F$ be finite extensions of a complete
nontrivially normed field, with compatible norms. The original
field norm $N_{F/E}:F\longrightarrow E$ is continuous.
In particular, every actual finite-layer norm used in the
multiplicative norm limit is continuous for the canonical spectral norms.
\end{lemma}
\begin{proof}
Choose a finite $E$-basis of $F$. Left multiplication by $x\in F$
has a matrix whose entries are $E$-linear functions of $x$, hence
$K$-linear functions. Every linear map from a finite-dimensional
normed space over the complete field $K$ is continuous.
Thus the matrix depends continuously on $x$. Its determinant is a
finite sum of finite products of its entries and is continuous.
The determinant is the original field norm by its defining formula.
Canonical spectral norms are compatible norms on the actual finite
intermediate fields, so the same argument applies to each such norm.
\end{proof}

\begin{lemma}[Ideal powers in the original discrete valuation topology]
\label{fw:ideal-powers-topology}
\source{wintenberger:page-0012,wintenberger:page-0013}
\dcref{Local topology used in the norm-limit construction}
\uses{fw:normed-intermediate-tower,fw:normalized-field-valuation}
\lean{FontaineWintenberger.norm_uniformizer_pos_lt_one, FontaineWintenberger.mem_maximalIdeal_pow_iff_norm_le, FontaineWintenberger.image_maximalIdeal_pow_eq_closedBall, FontaineWintenberger.isClosed_image_maximalIdeal_pow, FontaineWintenberger.nhds_zero_hasBasis_maximalIdeal_pow, FontaineWintenberger.isComplete_image_maximalIdeal_pow, FontaineWintenberger.FiniteLayer.isClosed_image_maximalIdeal_pow, FontaineWintenberger.FiniteLayer.nhds_zero_hasBasis_maximalIdeal_pow}
\leanok
Let $E$ be a complete nonarchimedean normed field whose unit ball
$\mathcal O_E$ is a discrete valuation ring, and let $\pi$ be a
uniformizer. For every $n\in\mathbb N$, the image of
$\mathfrak m_E^n$ in $E$ is exactly
\[
 \{x\in E:|x|\leq|\pi|^n\}.
\]
These sets are closed neighborhoods of zero and form a neighborhood
basis at zero. The integer ring and each of these ideal powers
are complete for their induced uniform structures.
\end{lemma}
\begin{proof}
The maximal ideal is generated by $\pi$, so its $n$th power is
generated by $\pi^n$. Every element $x=\pi^n y$ with
$y\in\mathcal O_E$ satisfies $|x|\leq|\pi|^n$.
Conversely, if $|x|\leq|\pi|^n$, then
$|x/\pi^n|\leq1$. This quotient lies in the actual unit ball,
so $x$ is the image of an element of $(\pi^n)$.
This also proves the equality at $n=0$.
The uniformizer is nonzero and belongs to the maximal ideal,
so $0<|\pi|<1$. Each displayed closed ball is closed by continuity
of the norm and contains the open ball of the same positive radius.
Since $|\pi|^n\longrightarrow0$, for every positive radius
$\varepsilon$ some such closed ball lies in the open
$\varepsilon$-ball. This proves the asserted neighborhood basis.
Finally each set is closed in the complete field $E$, hence complete;
the case $n=0$ gives completeness of the integer ring.
\end{proof}

\begin{lemma}[The ideal-power Cauchy criterion]
\label{fw:ideal-power-cauchy-criterion}
\source{wintenberger:page-0012,wintenberger:page-0013}
\dcref{Cauchy criterion for the discrete topology in Section 2.3}
\uses{fw:ideal-powers-topology}
\lean{FontaineWintenberger.cauchy_map_of_eventually_sub_mem_maximalIdeal_pow, FontaineWintenberger.exists_tendsto_of_eventually_sub_mem_maximalIdeal_pow}
\leanok
Let $E$ be a nonarchimedean normed field whose unit ball is a
DVR, and let $(z_i)$ be a net indexed by a nonempty directed
ordered set. Suppose that, for each $n\in\mathbb N$, there is
$i_n$ such that $z_i-z_j$ belongs to the image of
$\mathfrak m_E^n$ whenever $i,j\geq i_n$.
Then the net is Cauchy. If $E$ is complete, it converges in $E$.
\end{lemma}
\begin{proof}
The difference map identifies entourages of the additive uniform
structure with neighborhoods of zero. The actual ideal powers
are a neighborhood basis there. For each such power, the assumed
common tail therefore sends all pairs into the corresponding
entourage, which is the Cauchy condition.
Nonemptiness and directedness make the tail filter nontrivial.
Completeness supplies a limit of its Cauchy image filter, hence
convergence of the original net.
\end{proof}

\begin{lemma}[Uniform Herbrand control below an APF extension]
\label{fw:apf-herbrand-continuity-control}
\source{wintenberger:page-0005,wintenberger:page-0013}
\dcref{Local consequence of the intrinsic APF indices and their inverses}
\uses{fw:apf-finite-extension,fw:abstract-normed-finite-base,fw:apf-herbrand-homeomorphism}
\lean{FontaineWintenberger.apfPsi_le_of_algHom, FontaineWintenberger.apfPhi_le_of_algHom, FontaineWintenberger.le_ceil_apfPhi_of_apfPsi_le, FontaineWintenberger.FiniteLayer.le_ceil_apfPhi_of_apfPsi_le}
\leanok
Let $L/K$ be an abstract separable APF extension and let
$E\subseteq F\subseteq L$ be actual finite intermediate fields.
For every $t\geq0$,
\[
 \Psi_{F/E}(t)\leq\Psi_{L/E}(t),\qquad
 \varphi_{F/E}(t)\geq\varphi_{L/E}(t).
\]
In particular, if $N,s\in\mathbb N$ satisfy
$s\geq\Psi_{L/E}(N)$, then
$\lceil\varphi_{F/E}(s)\rceil_+\geq N$.
\end{lemma}
\begin{proof}
Finite-base inheritance makes $L/E$ APF.
For any actual common separable embedding over $E$, the fixing
group of $L$ is contained in that of $F$. Intersecting the
corresponding products with inertia and taking their finite
relative indices gives $A_{F/E}(u)\leq A_{L/E}(u)$ for $u\geq0$.
Integrating from zero to $t$ gives the first inequality.
Write $v=\varphi_{L/E}(t)\geq0$. Then
$\Psi_{F/E}(v)\leq\Psi_{L/E}(v)=t$.
Applying the increasing inverse of $\Psi_{F/E}$ gives
$v\leq\varphi_{F/E}(t)$, proving the second inequality.
For the last assertion,
$\varphi_{F/E}(s)\geq\varphi_{L/E}(s)
\geq\varphi_{L/E}(\Psi_{L/E}(N))=N$;
taking the nonnegative ceiling preserves the bound.
\end{proof}

\begin{lemma}[Integral coordinates of the original norm families]
\label{fw:norm-limit-integrality}
\source{wintenberger:page-0008}
\dcref{Integral part of the stable valuation construction}
\uses{fw:norm-limit-stable-valuation,fw:integral-finite-layer-norms}
\lean{FontaineWintenberger.IsAPF.exists_finiteLayer_normLimitEval_mem_of_le_one, FontaineWintenberger.normLimitEval_mem_of_valuation_le_one, FontaineWintenberger.normLimitValuation_le_one_iff}
\leanok
For an abstract APF extension $L/K$ with perfect residue field,
a norm-compatible family has nonnegative stable additive valuation
if and only if every original finite coordinate is integral.
\end{lemma}
\begin{proof}
On a finite stable tail, the stable valuation is the actual normalized
coordinate valuation. The valuation ring of that coordinate valuation
is its original integer ring, giving integrality on the stable tail.
For an arbitrary finite $E$, choose a finite compositum $F$ containing
$E$ and a stable-tail witness. Its coordinate is integral, so its
actual integral norm belongs to $\mathcal O_E$.
Original-field inclusion identifies that norm with the original field
norm, which is the coordinate at $E$ by compatibility.
Conversely, evaluate an everywhere-integral family at any finite
stable layer. Its actual normalized valuation is integral there,
and equals the stable valuation, proving the assertion.
\end{proof}

\begin{lemma}[Integral norm-sum limits on the original families]
\label{fw:integral-addition-limit}
\source{wintenberger:page-0012,wintenberger:page-0013}
\dcref{Direct Cauchy form of the norm-limit argument in Section 2.3}
\uses{fw:norm-limit-multiplicative,fw:norm-limit-integrality,fw:apf-wild-finite-tail,fw:apf-deep-finite-tail,fw:finite-layer-norm-precision,fw:apf-herbrand-continuity-control,fw:finite-norm-continuity,fw:ideal-powers-topology,fw:ideal-power-cauchy-criterion}
\lean{FontaineWintenberger.IsAPF.exists_finiteLayer_integerNorm_add_congruence, FontaineWintenberger.IsAPF.exists_finiteLayer_normSumNet_sub_mem, FontaineWintenberger.IsAPF.exists_finiteLayer_normSumNet_cauchy, FontaineWintenberger.exists_tendsto_normSumNet_of_valuation_le_one, FontaineWintenberger.existsUnique_normLimit_integral_add}
\leanok
Let $L/K$ be an abstract separable APF extension with perfect
residue field of characteristic $p$. Let $\alpha,\beta$ be actual
norm-compatible families whose stable normalized additive valuations
are nonnegative. For every actual finite $E\subseteq L$, the net
\[
 z_F=N_{F/E}(\alpha_F+\beta_F)\qquad(F\supseteq E)
\]
converges in the original field $E$. Its limits form an actual
norm-compatible integral family $\gamma$ on the same multiplicative
norm limit. For every integer $s\geq0$, there is an actual finite
$T\subseteq L$ such that, for all finite $E\supseteq T$,
\[
 \gamma_E-\alpha_E-\beta_E\in\mathfrak m_E^s.
\]
\end{lemma}
\begin{proof}
The preceding lemma makes every original coordinate integral.
Choose an actual finite layer $T_0$ above the wild-tail witness.
Every finite step above $T_0$ is totally ramified of $p$-power degree. First fix $E\supseteq T_0$.
To prove the Cauchy assertion at any integer output precision $N$,
choose an integer $s\geq\Psi_{L/E}(N)$.
Cofinal growth of relative first bounds supplies an actual finite
$T\supseteq E$ such that, for all $T\subseteq B\subseteq C$,
\[
 i(C/B)\geq\frac p{p-1}s.
\]
If $C=B$, the norm defect is zero. Otherwise the actual finite
norm estimate gives
\[
 N_{C/B}(\alpha_C+\beta_C)-\alpha_B-\beta_B
 \in\mathfrak m_B^s.
\]
Indeed the exponent supplied there is
$\lceil(p-1)i(C/B)/p\rceil_+\geq s$.
Apply actual nonnormal Herbrand norm continuity to this difference
in $B$, using the two integral inputs
$N_{C/B}(\alpha_C+\beta_C)$ and $\alpha_B+\beta_B$.
The preceding uniform control lemma gives
\[
 z_C-z_B\in\mathfrak m_E^N.
\]
For any $B_1,B_2\supseteq T$, take their actual finite compositum
$C\supseteq B_1,B_2$. The two displayed estimates and closure of
$\mathfrak m_E^N$ under subtraction give
$z_{B_1}-z_{B_2}\in\mathfrak m_E^N$.
These ideal powers are a neighborhood basis at zero, so the net
is Cauchy. Completeness of the original finite field $E$ gives its
limit $\gamma_E$. The integer ring is closed, so the limit is integral.

For a finite $E$ below or incomparable with $T_0$, choose
$E^*=ET_0$. On the actual tail $F\supseteq E^*$, norm transitivity
writes $z_F=N_{E^*/E}(N_{F/E^*}(\alpha_F+\beta_F))$.
The finite field norm is continuous: in any finite basis it is a
polynomial determinant in the coordinates. Thus the same net
converges in $E$ to $N_{E^*/E}(\gamma_{E^*})$.
Its uniqueness makes this definition independent of the auxiliary
common upper field. Norm transitivity and the same continuity show
$N_{F/E}(\gamma_F)=\gamma_E$ for every actual finite inclusion.
The proved multiplicative compatible-family bridge therefore places
$\gamma$ on the original norm-limit carrier, including the case
that its coordinates are all zero.

Finally fix an integer $s\geq0$ and enlarge $T_0$ using the deep-tail
bound with $ps/(p-1)$.
For every $E$ above that layer and every $F\supseteq E$, the actual
norm estimate gives
$N_{F/E}(\alpha_F+\beta_F)-\alpha_E-\beta_E\in\mathfrak m_E^s$.
This ideal power is closed in the complete discrete valuation field.
Pass to the proved limit to obtain the final assertion.
This is the Cauchy argument used for the norm-limit inverse map,
applied directly to the original compatible families; it uses no
compactness of the residue field and introduces no second carrier.
\end{proof}

\begin{theorem}[Addition as a limit of norms]
\label{fw:addition-limit}
\dcref{Wintenberger1983, Theorem 2.1.3(i)}
\source{wintenberger:page-0008}
\uses{fw:norm-limit-multiplicative,fw:norm-limit-stable-valuation,fw:integral-addition-limit}
\lean{FontaineWintenberger.exists_tendsto_normSumNet, FontaineWintenberger.existsUnique_normLimit_add, FontaineWintenberger.normLimitAdd, FontaineWintenberger.tendsto_normSumNet_normLimitAdd}
\leanok
Let $L/K$ be infinite APF and $\alpha,\beta\in X_K(L)$.
For every finite $E/K$ in $L$, the net
\[
 N_{E'/E}(\alpha_{E'}+\beta_{E'})
 \qquad(E\subset E'\subset L,\ [E':K]<\infty)
\]
converges in $E$. Its limits $\gamma_E$ form a norm-compatible family.
Define $\alpha+\beta=(\gamma_E)_E$.
\end{theorem}
\begin{proof}
If one family is zero, its norm-sum net is exactly the other
family coordinate by coordinate. Otherwise choose $\delta$ to be
the one with smaller stable additive valuation. Then
$\alpha/\delta$ and $\beta/\delta$ are integral families, so their
norm-sum nets converge by the preceding lemma.
On every actual finite level, norm multiplicativity and compatibility
give the exact equality
\[
 N_{F/E}(\alpha_F+\beta_F)
 =\delta_E N_{F/E}(\alpha_F/\delta_F+\beta_F/\delta_F).
\]
Multiplication by the fixed $\delta_E$ is continuous, so these
nets converge. Continuity of each actual finite norm and its
transitivity prove compatibility of the limits, as in the integral
case. The existing compatible-family bridge supplies the resulting
family on the same norm-limit carrier. Thus no choice of a norm-family
uniformizer is required for this construction.
\end{proof}


\begin{remark}
The following arguments derive the algebra and valuation laws directly
on the multiplicative norm limit. They adapt the norm-continuity and
asymptotic-coordinate arguments of \cite[Sections 2.2--2.3]{Wintenberger1983}:
comparison in changing finite fields is always transported to one fixed
finite field before separatedness is used. No surjectivity of the stable
valuation or description of the residue field is needed here.
Throughout these arguments $L/K$ is an abstract infinite separable APF
extension, the residue field is arbitrary perfect of characteristic
$p>0$, and $v$ denotes the stable normalized additive valuation, with
$v(0)=+\infty$.
\end{remark}

\begin{lemma}[The exact symmetry and zero identities for addition]
\label{fw:norm-limit-addition-identities}
\source{wintenberger:page-0008}
\dcref{Direct consequences of the norm-limit formula in Theorem 2.1.3(i)}
\uses{fw:addition-limit}
\lean{FontaineWintenberger.normLimitAdd_comm, FontaineWintenberger.normLimitAdd_zero_left, FontaineWintenberger.normLimitAdd_zero_right}
\leanok
For all $\alpha,\beta\in X_K(L)$,
\[
 \alpha+\beta=\beta+\alpha,\qquad \alpha+0=0+\alpha=\alpha.
\]
\end{lemma}
\begin{proof}
At each fixed finite base $E$, the two nets defining
$\alpha+\beta$ and $\beta+\alpha$ are identical by commutativity
in every upper finite field. Their limits are equal.
The net defining $\alpha+0$ is constantly $\alpha_E$, since
$N_{F/E}(\alpha_F)=\alpha_E$; its limit is therefore $\alpha_E$.
Joint injectivity of the original coordinate evaluations proves both
identities, and commutativity gives the left zero identity.
\end{proof}

\begin{lemma}[Distributivity of the original norm-limit multiplication]
\label{fw:norm-limit-distributivity}
\source{wintenberger:page-0008,wintenberger:page-0012}
\dcref{Direct multiplicative form of the norm-limit construction}
\uses{fw:addition-limit,fw:norm-limit-multiplicative}
\lean{FontaineWintenberger.mul_normLimitAdd}
\leanok
For all $\delta,\alpha,\beta\in X_K(L)$,
\[
 \delta(\alpha+\beta)=\delta\alpha+\delta\beta.
\]
\end{lemma}
\begin{proof}
Fix an actual finite base $E$. For every finite $F\supseteq E$,
ordinary norm multiplicativity gives
\[
 N_{F/E}(\delta_F\alpha_F+\delta_F\beta_F)
 =N_{F/E}(\delta_F)N_{F/E}(\alpha_F+\beta_F)
 =\delta_E N_{F/E}(\alpha_F+\beta_F).
\]
This also holds when one of the coordinates is zero.
Multiplication by the fixed element $\delta_E$ is continuous in $E$.
Taking the already established limits identifies the $E$-coordinate
of the right side of the assertion with that of the left side.
Joint injectivity proves the assertion. Commutativity of multiplication
also gives right distributivity.
\end{proof}

\begin{lemma}[Arbitrary-pair asymptotic coordinate approximation]
\label{fw:norm-limit-coordinate-approximation}
\source{wintenberger:page-0013}
\dcref{Multiplicatively rescaled form of Remark 2.3.3.1}
\uses{fw:integral-addition-limit,fw:norm-limit-stable-valuation,fw:norm-limit-distributivity,fw:norm-limit-addition-identities}
\lean{FontaineWintenberger.normLimitAdd_sub_mem}
\leanok
For any $\alpha,\beta\in X_K(L)$ and any integer $n\geq0$, there is an
actual finite $T\subseteq L$ such that every finite $E\supseteq T$
satisfies
\[
 v_E\bigl((\alpha+\beta)_E-\alpha_E-\beta_E\bigr)\geq n.
\]
The subtractions in this formula are in the original finite field $E$.
The nonnegative precisions are cofinal; the same statement for every integer
follows by using precision zero when $n<0$.
\end{lemma}
\begin{proof}
Choose a member $\delta$ of $\{1,\alpha,\beta\}$ with least
stable additive valuation. It is nonzero, has valuation $c\leq0$,
and the original families $a=\alpha\delta^{-1}$ and
$b=\beta\delta^{-1}$ are integral. Put $q=-c\geq0$.
The integral approximation with precision $n+q$ gives, on an
actual finite tail,
\[
 v_E((a+b)_E-a_E-b_E)\geq n+q.
\]
Enlarge this tail so that $v_E(\delta_E)=c\geq-q$.
Distributivity and the original coordinatewise multiplication give
\[
 (\alpha+\beta)_E-\alpha_E-\beta_E
 =\delta_E\bigl((a+b)_E-a_E-b_E\bigr).
\]
Its valuation is at least $-q+(n+q)=n$. The computation includes a
zero error, whose valuation is infinite, and zero input families.
No norm-family uniformizer is chosen.
\end{proof}

\begin{lemma}[Separation of compatible families by asymptotic coordinates]
\label{fw:norm-limit-asymptotic-separation}
\source{wintenberger:page-0012,wintenberger:page-0013}
\dcref{Fixed-base norm-continuity form of the uniqueness in Section 2.3}
\uses{fw:norm-limit-multiplicative,fw:norm-limit-stable-valuation,fw:norm-limit-integrality,fw:apf-wild-finite-tail,fw:nonnormal-wild-norm-continuity,fw:apf-herbrand-continuity-control,fw:ideal-powers-topology,fw:finite-layer-norm-precision}
\lean{FontaineWintenberger.normLimit_ext_of_integral_eventually_sub_mem, FontaineWintenberger.normLimit_ext_of_eventually_sub_mem}
\leanok
Let $\alpha,\beta\in X_K(L)$. Suppose that for every $n\geq0$
there is an actual finite $T_n\subseteq L$ such that
\[
 v_E(\alpha_E-\beta_E)\geq n
 \qquad\text{for every finite }E\supseteq T_n.
\]
Then $\alpha=\beta$.
\end{lemma}
\begin{proof}
Choose a nonzero family $\delta$ with least stable valuation among
the nonzero members of $\{1,\alpha,\beta\}$. Since $1$ is one
candidate, $v(\delta)\leq0$. The families
$a=\alpha\delta^{-1}$, $b=\beta\delta^{-1}$ and $\delta^{-1}$
are integral, including when $\alpha$ or $\beta$ is zero.
The integral-coordinate lemma therefore puts every original
coordinate of $\delta^{-1}$ in its actual integer ring.
For each $s\geq0$, apply the hypothesis at $n=s$.
Multiplication of its ideal-power error by this integral coordinate
gives
\[
 a_F-b_F=(\delta^{-1})_F(\alpha_F-\beta_F)
 \in\mathfrak m_F^s
\]
on the same actual finite tail. This also covers a zero difference
and needs no choice of a family uniformizer.

Fix an arbitrary original finite layer $D$. Choose a single actual
finite $B\supseteq D$ containing a wild-tail witness. This field $B$
is fixed for the remainder of the argument. Given an integer $N\geq0$,
choose an integer $s\geq\Psi_{L/B}(N)$, and choose a finite
$F\supseteq B$ on the tail just obtained for this $s$.
Every original coordinate of $a$ and $b$ is integral.
The actual extension $F/B$ is totally ramified of $p$-power degree.
Nonnormal norm continuity and the intrinsic APF comparison therefore give
\[
 \begin{split}
 v_B(a_B-b_B)
 &=v_B\bigl(N_{F/B}(a_F)-N_{F/B}(b_F)\bigr)\\
 &\geq\lceil\varphi_{F/B}(s)\rceil_+\geq N.
 \end{split}
\]
Thus the single fixed element $a_B-b_B$ lies in every power of
$\mathfrak m_B$. The separated discrete valuation on $B$ forces
$a_B=b_B$, hence $\alpha_B=\beta_B$ after multiplication by $\delta_B$.
Applying the actual norm $N_{B/D}$ and compatibility gives
$\alpha_D=\beta_D$. Since $D$ was arbitrary, joint injectivity of
all the original coordinates proves the result.
\end{proof}

\begin{proposition}[The commutative additive monoid on the norm limit]
\label{fw:norm-limit-additive-monoid}
\source{wintenberger:page-0012,wintenberger:page-0013}
\dcref{Direct-coordinate adaptation of the ring laws in Proposition 2.2.4}
\uses{fw:norm-limit-addition-identities,fw:norm-limit-distributivity,fw:norm-limit-coordinate-approximation,fw:norm-limit-asymptotic-separation,fw:integral-addition-limit,fw:norm-limit-stable-valuation}
\lean{FontaineWintenberger.normLimitAdd_assoc, FontaineWintenberger.normLimitAddCommMonoid, FontaineWintenberger.normLimitValuation_normLimitAdd_le_one}
\leanok
The addition on the original multiplicative norm limit is associative.
Together with its original zero it makes $X_K(L)$ a commutative monoid.
The families of nonnegative stable valuation are closed under addition.
\end{proposition}
\begin{proof}
Closure follows from the integral addition limit and uniqueness of the
coordinate limits defining addition.
For arbitrary norm-compatible families $a,b,c$ and a precision $n\geq0$,
apply coordinate approximation to the four pairs
$(a,b)$, $(b,c)$, $(a+b,c)$ and $(a,b+c)$.
An actual finite compositum of their four witnesses supplies one tail
on which every error belongs to $\mathfrak m_E^n$. Write
\[
\begin{aligned}
 e_1&=((a+b)+c)_E-(a+b)_E-c_E,&
 e_2&=(a+b)_E-a_E-b_E,\\
 e_3&=(a+(b+c))_E-a_E-(b+c)_E,&
 e_4&=(b+c)_E-b_E-c_E.
\end{aligned}
\]
Ordinary associativity in the original finite field $E$ gives
\[
 ((a+b)+c)_E-(a+(b+c))_E=e_1+e_2-e_3-e_4
 \in\mathfrak m_E^n.
\]
Asymptotic separation proves $(a+b)+c=a+(b+c)$.
The symmetry and zero identities complete the monoid axioms.
\end{proof}

\begin{lemma}[Finite sums of integral families approximate coordinate sums]
\label{fw:norm-limit-finite-sum-approximation}
\source{wintenberger:page-0013}
\dcref{Finite-sum consequence of the additive approximation in Remark 2.3.3.1}
\uses{fw:norm-limit-additive-monoid,fw:norm-limit-coordinate-approximation}
\lean{FontaineWintenberger.normLimitValuation_sum_le_one, FontaineWintenberger.normLimitEval_sum_sub_mem, FontaineWintenberger.normLimitEval_nsmul_sub_mem}
\leanok
Let $a_1,\ldots,a_m$ be integral norm-compatible families and put
$a=\sum_{i=1}^m a_i$, with the empty sum equal to zero.
Then $a$ is integral. For every integer $n\geq0$, on an actual
finite tail,
\[
 a_E-\sum_{i=1}^m(a_i)_E\in\mathfrak m_E^n.
\]
The coordinate approximation also holds for arbitrary families $a_i$,
without the integrality assertion.
\end{lemma}
\begin{proof}
Integrality follows by induction from closure under addition.
For coordinate approximation, allow arbitrary families and induct on $m$.
The empty sum has zero coordinate and zero error.
If the assertion holds for $a=\sum_{i=1}^m a_i$, then for a fixed
$n$ take a common upper bound
of the induction witness and the approximation witness for $(a,a_{m+1})$.
On this tail,
\[
 (a+a_{m+1})_E-\sum_{i=1}^{m+1}(a_i)_E
 =\bigl((a+a_{m+1})_E-a_E-(a_{m+1})_E\bigr)
  +\bigl(a_E-\sum_{i=1}^m(a_i)_E\bigr).
\]
Both summands lie in the same ideal power, proving the induction step.
\end{proof}

\begin{lemma}[Unbounded actual finite degrees and ramification]
\label{fw:apf-unbounded-finite-ramification}
\source{wintenberger:page-0012}
\dcref{Finite-degree explanation of the divergence used in Proposition 2.2.4}
\uses{fw:norm-limit-multiplicative,fw:apf-wild-finite-tail}
\lean{FontaineWintenberger.FiniteLayer.exists_tail_finrank_ge, FontaineWintenberger.IsAPF.exists_finiteLayer_ramificationIdx_unbounded}
\leanok
For an infinite algebraic extension $L/K$ and any actual finite
$D\subseteq L$, the degrees $[E:D]$, over finite $E\supseteq D$,
tend to infinity along the finite-layer tail filter.
If $L/K$ is APF and $D$ contains a wild-tail witness, every such
step is totally ramified, and its actual ramification index
$e(E/D)=[E:D]$ tends to infinity as well.
\end{lemma}
\begin{proof}
No finite intermediate field $F$ equals $L$, since that would make
$L/K$ finite. Choose $x\in L\setminus F$. Algebraicity makes $F(x)/K$
finite, and $[F(x):F]\geq2$. Starting with $D$, repeat this construction
$M$ times. Degree multiplicativity shows that the resulting finite
$T\supseteq D$ has $[T:D]\geq M+1$; indeed each step increases the
positive integer degree by at least one. For any finite $E\supseteq T$,
\[
 [E:D]=[E:T][T:D]\geq[T:D]\geq M.
\]
This proves the asserted eventual bound for every $M$.
On the APF wild tail the actual residue degree of $E/D$ is one,
and the proved local degree identity gives $e(E/D)=[E:D]$.
Consequently the same eventual bounds hold for its actual
ramification index.
\end{proof}

\begin{lemma}[The residue prime becomes arbitrarily small]
\label{fw:apf-prime-valuation-diverges}
\source{wintenberger:page-0012}
\dcref{Direct valuation form of the characteristic argument in Proposition 2.2.4}
\uses{fw:apf-unbounded-finite-ramification,fw:normalized-valuation-scaling,fw:spectral-integer-residue-maps}
\lean{FontaineWintenberger.IsAPF.exists_finiteLayer_natCast_mem_maximalIdeal_pow, FontaineWintenberger.IsAPF.exists_finiteLayer_natCast_mem_image_maximalIdeal_pow}
\leanok
For every integer $n\geq0$, on an actual finite tail,
\[
 v_E(p)\geq n.
\]
Here $p$ denotes its image in the original finite field $E$.
If $K$ has characteristic $p$, that image is zero and its valuation
is infinite at every finite layer.
\end{lemma}
\begin{proof}
If $p=0$ in $K$, its image in every extension is zero, proving the
assertion. Otherwise fix a finite $D$ containing a wild-tail witness.
The integer $p$ belongs to $\mathfrak m_D$: its image in the actual
residue field is zero because that field has characteristic $p$.
It is nonzero in $D$, so $v_D(p)$ is a positive integer.
Normalized valuation scaling on the actual inclusions gives
\[
 v_E(p)=e(E/D)v_D(p)\geq e(E/D)\qquad(E\supseteq D).
\]
The preceding divergence of these actual ramification indices gives
the required tail for each $n$. This covers mixed characteristic
without imposing finiteness of the residue field.
\end{proof}

\begin{lemma}[The sum of $p$ copies of one vanishes]
\label{fw:norm-limit-prime-sum-zero}
\source{wintenberger:page-0012,wintenberger:page-0013}
\dcref{Direct-coordinate adaptation of the characteristic argument in Proposition 2.2.4}
\uses{fw:norm-limit-finite-sum-approximation,fw:apf-prime-valuation-diverges,fw:norm-limit-asymptotic-separation,fw:norm-limit-stable-valuation}
\lean{FontaineWintenberger.normLimit_natCast_prime_eq_zero}
\leanok
In the commutative additive monoid $X_K(L)$,
\[
 \underbrace{1+\cdots+1}_{p\ \mathrm{terms}}=0.
\]
\end{lemma}
\begin{proof}
Write $s$ for this sum. The family $1$ is integral since its stable
valuation is zero. For each $n\geq0$, finite-sum approximation gives
$s_E-p\in\mathfrak m_E^n$ on an actual finite tail.
On another such tail $p\in\mathfrak m_E^n$ by the preceding lemma.
Their common finite upper bound therefore gives
$s_E\in\mathfrak m_E^n$ throughout its tail.
Asymptotic separation, applied to $s$ and the original zero family,
proves $s=0$.
\end{proof}

\begin{theorem}[The field structure on the original multiplicative norm limit]
\label{fw:norm-limit-field}
\source{wintenberger:page-0008,wintenberger:page-0012,wintenberger:page-0013}
\dcref{Direct-coordinate deduction from the multiplicative limit and the ring laws}
\uses{fw:norm-limit-multiplicative,fw:norm-limit-additive-monoid,fw:norm-limit-distributivity,fw:norm-limit-prime-sum-zero}
\lean{FontaineWintenberger.normLimitField, FontaineWintenberger.normLimitField_charP, FontaineWintenberger.normLimit_pred_nsmul_add_eq_zero}
\leanok
The preceding addition and the original multiplication make the same
set $X_K(L)$ a field of characteristic $p$. The additive inverse of
$\alpha$ is the sum of $p-1$ copies of $\alpha$.
\end{theorem}
\begin{proof}
The proved additive laws and distributivity, together with the original
commutative multiplication, give a commutative semiring.
For each $m\geq0$, induction using distributivity shows that the sum
of $m$ copies of $\alpha$ is $\alpha$ times the sum of $m$ copies of $1$.
In particular the sum of $p$ copies of $\alpha$ is zero.
Since $p\geq2$, associativity expresses that sum as $\alpha$ plus the
sum of $p-1$ copies of $\alpha$; commutativity gives both inverse
identities. Thus addition is an abelian group and the semiring is a
commutative ring.

Every nonzero element is already a unit in the original multiplicative
norm limit, with its original coordinatewise inverse. Its zero and one
are distinct, since they evaluate to distinct elements of any finite
layer. This commutative ring is therefore a field. Its characteristic
is positive and divides $p$, because the sum of $p$ copies of one is
zero. A nontrivial field has prime characteristic, so that
characteristic is $p$. In particular the inverse formula also holds
when $p=2$, when it reads $-\alpha=\alpha$.
\end{proof}

\begin{theorem}[The stable norm valuation is an ultrametric field valuation]
\label{fw:norm-limit-field-valuation}
\source{wintenberger:page-0012,wintenberger:page-0013}
\dcref{Direct-coordinate form of the valuation assertion in Proposition 2.2.4}
\uses{fw:norm-limit-field,fw:norm-limit-stable-valuation,fw:integral-addition-limit,fw:norm-limit-distributivity}
\lean{FontaineWintenberger.normLimitValuation_add_le_max, FontaineWintenberger.normLimitFieldValuation, FontaineWintenberger.normLimitFieldValuation_apply, FontaineWintenberger.normLimitValued}
\leanok
On this field the existing stable normalized additive valuation satisfies
\[
 v(\alpha+\beta)\geq\min(v(\alpha),v(\beta)),\qquad
 v(\alpha\beta)=v(\alpha)+v(\beta),\qquad
 v(\alpha)=+\infty\Longleftrightarrow\alpha=0.
\]
Equivalently its existing stable normalized multiplicative map $w$
is a field valuation, with
$w(\alpha+\beta)\leq\max(w(\alpha),w(\beta))$.
\end{theorem}
\begin{proof}
Multiplicativity and the zero test were proved on the original
multiplicative carrier, so hold for exactly the same map on this field.
If either summand is zero, the additive inequality follows from the zero
identity. Otherwise choose $\delta$ to be the summand with smaller
stable valuation and put $a=\alpha\delta^{-1}$,
$b=\beta\delta^{-1}$. Both are integral families, and the integral
addition theorem makes $a+b$ integral. Distributivity gives
$\alpha+\beta=\delta(a+b)$. Therefore
\[
 v(\alpha+\beta)=v(\delta)+v(a+b)
 \geq v(\delta)=\min(v(\alpha),v(\beta)).
\]
If $a+b=0$, the left side is infinite and the inequality still holds.
Together with the already proved values at zero and one this verifies
all the valuation axioms. Reversing the order of additive values gives
the stated multiplicative ultrametric inequality.
\end{proof}

\subsection{Truncated norms and completion of the original field}

\begin{lemma}[Surjective truncated norms at an actual high wild break]
\label{fw:truncated-wild-norm-map}
\source{wintenberger:page-0009,wintenberger:page-0011}
\dcref{Proposition 2.2.1 and paragraph 2.2.3.3, with explicit source precision}
\uses{fw:norm-estimates,fw:finite-norm-residue-digits,fw:norm-valuation,fw:ideal-power-valuation,fw:finite-layer-norm-precision}
\lean{FontaineWintenberger.FiniteLayer.truncatedNorm, FontaineWintenberger.FiniteLayer.truncatedNorm_mk, FontaineWintenberger.FiniteLayer.truncatedNorm_surjective, FontaineWintenberger.truncatedAddVal, FontaineWintenberger.FiniteLayer.truncatedNorm_truncatedAddVal, FontaineWintenberger.FiniteLayer.truncatedNorm_comp, FontaineWintenberger.FiniteLayer.truncatedNorm_self, FontaineWintenberger.FiniteLayer.exists_surjective_truncatedNorm_self}
\leanok
Let $F/E$ be a finite separable totally ramified extension of degree
$p^d$ over a complete discrete valuation field with arbitrary perfect
residue field of characteristic $p$. Let $0\leq r\leq s$ be integers.
Suppose either $F=E$, or
\[
 r\leq\left\lceil\frac{p-1}{p}i(F/E)\right\rceil_+.
\]
The actual field norm induces a surjective unital ring homomorphism
\[
 \overline N_{F/E}^{s,r}:
 \mathcal O_F/\mathfrak m_F^s\longrightarrow
 \mathcal O_E/\mathfrak m_E^r,\qquad
 [x]\longmapsto[N_{F/E}(x)].
\]
For an integral representative $x$, its image has truncated normalized
valuation $\min(v_F(x),r)$. For a tower $E\subseteq F\subseteq M$
and integers $r\leq s\leq t$ for which the three displayed maps exist,
\[
 \overline N_{F/E}^{s,r}\circ\overline N_{M/F}^{t,s}
 =\overline N_{M/E}^{t,r}.
\]
The equal-field map is the ordinary reduction of precision.
\end{lemma}
\begin{proof}
In the equal-field case every assertion is the corresponding statement
for the identity norm followed by quotient reduction. Otherwise the
actual wild norm estimate makes the norm additive modulo
$\mathfrak m_E^r$, since the supplied precision is at least $r$.
It preserves multiplication, zero and one exactly. If $x-y$ belongs
to $\mathfrak m_F^s$, norm additivity applied to $y$ and $x-y$ gives
\[
 N(x)-N(y)\equiv N(x-y)\pmod{\mathfrak m_E^r}.
\]
Total ramification gives $v_E(N(x-y))=v_F(x-y)\geq s\geq r$.
Thus the map is well defined on the asserted source quotient and is
a unital ring homomorphism, also when $r=0$.

To verify surjectivity, take an arbitrary representative
$a\in\mathcal O_E$ of the desired target class. Choose a uniformizer
$\pi'$ of $F$ and put $\pi=N(\pi')$, an actual uniformizer of $E$.
Perfectness gives finite digits $b_0,\ldots,b_{r-1}\in\mathcal O_E$
with $a\equiv\sum_{j<r}b_j^{p^d}\pi^j$ modulo $\mathfrak m_E^r$.
For $x=\sum_{j<r}b_j(\pi')^j$, repeated norm additivity at that same
precision, multiplicativity and the scalar norm identity give
$N(x)\equiv a$ modulo $\mathfrak m_E^r$.
This supplies an actual preimage of each target class, including the
empty digit sum when $r=0$. Finally normalized norm valuation gives
$\min(v_E(N(x)),r)=\min(v_F(x),r)$. Applying ordinary norm
transitivity to each integral representative proves the composition
identity on the quotients. No finiteness of the residue field is used.
\end{proof}

\begin{lemma}[A cofinal tower with explicit growing norm precisions]
\label{fw:deep-truncated-finite-tower}
\source{wintenberger:page-0011}
\dcref{Countable cofinal form of paragraphs 2.2.3.1--2.2.3.3}
\uses{fw:countable-apf-finite-layers,fw:apf-wild-finite-tail,fw:apf-deep-finite-tail,fw:apf-prime-valuation-diverges,fw:norm-limit-stable-valuation,fw:truncated-wild-norm-map}
\lean{FontaineWintenberger.IsAPF.exists_deep_finiteTower, FontaineWintenberger.deepFiniteTower, FontaineWintenberger.deepFiniteTower_monotone, FontaineWintenberger.deepFiniteTower_cofinal, FontaineWintenberger.le_deepFiniteTower, FontaineWintenberger.normLimitValuation_eq_deepFiniteTower, FontaineWintenberger.natCast_mem_deepFiniteTower_maximalIdeal_pow, FontaineWintenberger.deepFiniteTower_wild_tail, FontaineWintenberger.deepFiniteTowerQuotient_charP, FontaineWintenberger.deepFiniteTowerQuotient_maximalIdeal_pow, FontaineWintenberger.deepFiniteTowerNorm, FontaineWintenberger.deepFiniteTowerNorm_surjective, FontaineWintenberger.deepFiniteTowerTransition, FontaineWintenberger.deepFiniteTowerTransition_surjective, FontaineWintenberger.deepFiniteTowerTransition_comp, FontaineWintenberger.deepFiniteTowerTransition_self, FontaineWintenberger.deepFiniteTowerTransition_truncatedAddVal}
\leanok
Let $L/K$ be infinite separable APF and let $D\subseteq L$ be any
specified actual finite layer. There is an increasing cofinal sequence
of actual finite layers $E_n\supseteq D$ such that, with $h_n=n+2$:
all finite inclusions above $E_0$ are totally ramified of $p$-power
degree; $v(\alpha)=v_{E_n}(\alpha_{E_n})$ for every original family
$\alpha$; $v_{E_n}(p)\geq h_n$; and every finite $F\supseteq E_n$
has norm additive and surjective modulo $\mathfrak m_{E_n}^{h_n}$.
Consequently, for $m\geq n$, actual norms induce surjective ring maps
\[
 q_{n,m}:B_m\longrightarrow B_n,\qquad
 B_n=\mathcal O_{E_n}/\mathfrak m_{E_n}^{h_n},
\]
with $q_{n,n}=1$ and $q_{i,n}q_{n,m}=q_{i,m}$. Every $B_n$ has
characteristic $p$, and its maximal ideal has $h_n$th power zero.
\end{lemma}
\begin{proof}
Take an increasing cofinal sequence $C_n$ of actual finite layers.
Fix finite witnesses for the wild tail and the stable valuation.
For each $n$, the deep-tail theorem supplies a finite $T_n$ above which
all relative first bounds are at least $ph_n/(p-1)$; enlarge it by a
witness making the original valuation of $p$ at least $h_n$.
Recursively let $E_n$ contain $D$, $C_n$, $T_n$, both fixed witnesses,
and $E_{n-1}$ when $n>0$, by taking their actual finite compositum.
This sequence is increasing and cofinal because it contains $C_n$.
The wild and stable-valuation assertions hold by construction.
For $F\supseteq E_n$, the deep bound and actual finite norm estimates
give the claimed precision $h_n$; the identity step has zero addition
defect and is surjective at that precision as well.
Apply the preceding truncated-norm lemma with source precision $h_m$
and target precision $h_n$. Its tower identity gives every asserted
transition identity. Finally $p\in\mathfrak m_{E_n}^{h_n}$, and $h_n>0$ makes
$B_n$ nonzero. It has characteristic $p$: this nonzero unital ring has
positive characteristic dividing the prime $p$. Its maximal ideal
is $\mathfrak m_{E_n}/\mathfrak m_{E_n}^{h_n}$, whose $h_n$th power
is zero. All steps use actual finite fields in the given abstract $L$.
\end{proof}

\begin{lemma}[The truncated inverse limit and its surjective projections]
\label{fw:truncated-norm-inverse-limit}
\source{wintenberger:page-0011,wintenberger:page-0012}
\dcref{The inverse limit in paragraph 2.2.3.3 and its projections}
\uses{fw:deep-truncated-finite-tower}
\lean{FontaineWintenberger.SequentialInverseLimit.subring, FontaineWintenberger.SequentialInverseLimit.projection, FontaineWintenberger.SequentialInverseLimit.projection_surjective, FontaineWintenberger.SequentialInverseLimit.isClosed_subring, FontaineWintenberger.SequentialInverseLimit.completeSpace, FontaineWintenberger.SequentialInverseLimit.t2Space, FontaineWintenberger.SequentialInverseLimit.nhds_zero_hasBasis_ker_projection}
\leanok
For such a tower put
\[
 A=\{(a_n)\in\prod_{n\geq0}B_n:
                    q_{n,n+1}(a_{n+1})=a_n\text{ for all }n\}.
\]
This is a unital ring with coordinatewise operations. Every projection
$\operatorname{pr}_n:A\to B_n$ is surjective. With the uniformity
induced from the product of the discrete additive groups $B_n$, the
ring $A$ is Hausdorff and complete. Its neighborhoods of zero are
cofinally the kernels of these individual projections.
\end{lemma}
\begin{proof}
The transition maps are unital ring homomorphisms, so the displayed
compatibility equations are preserved by all ring operations.
Fix $n$ and any $b\in B_n$. Set $a_n=b$. For each $m\geq n$,
choose $a_{m+1}$ in the nonempty fiber of $a_m$ under the proved
surjective map $q_{m,m+1}$. This recursive choice constructs all
upper coordinates. For $m<n$ set $a_m=q_{m,n}(b)$.
Transition composition proves every adjacent compatibility equation,
including the one at $n-1$. Thus $(a_m)\in A$ projects to the prescribed
$b$, proving actual surjectivity without any cardinality assumption
on the quotient rings.

Every discrete uniform additive group is complete: a nontrivial Cauchy filter
contains a set all of whose pairwise differences are zero, hence
converges to its unique element. Products of complete Hausdorff
uniform groups are complete and Hausdorff. Each adjacent compatibility
equation defines a closed subset of this product, since its two maps
have discrete target. Their intersection $A$ is therefore closed and
complete. The basic product neighborhoods prescribe finitely many
coordinates; if $N$ is their greatest index, compatibility makes
$\ker(\operatorname{pr}_N)$ contained in all those coordinate kernels.
This proves the last assertion. This is a closed-product argument,
not a compactness argument for finite residue fields.
\end{proof}

\begin{lemma}[Reconstruction by limits of the actual original norms]
\label{fw:truncated-norm-reconstruction}
\source{wintenberger:page-0012,wintenberger:page-0013}
\dcref{Countable original-field form of Proposition 2.3.1}
\uses{fw:truncated-norm-inverse-limit,fw:nonnormal-wild-norm-continuity,fw:apf-herbrand-continuity-control,fw:finite-norm-continuity,fw:ideal-powers-topology,fw:ideal-power-cauchy-criterion,fw:norm-limit-multiplicative,fw:norm-limit-integrality}
\lean{FontaineWintenberger.deepNormReconstruction, FontaineWintenberger.deepNormReconstruction_valuation_le_one, FontaineWintenberger.normLimitEval_deepNormReconstruction, FontaineWintenberger.deepNormReconstruction_sub_mem, FontaineWintenberger.deepNormReconstruction_tendsto, FontaineWintenberger.existsUnique_normLimit_of_truncatedNorms}
\leanok
Write $\mathcal O_X=\{\alpha\in X_K(L):v(\alpha)\geq0\}$.
For $a=(a_n)\in A$, choose arbitrary integral representatives
$x_n\in\mathcal O_{E_n}$ of $a_n$. For every fixed $j$ the sequence
\[
 N_{E_n/E_j}(x_n)\qquad(n\geq j)
\]
converges in the original complete field $E_j$. Its limits extend
uniquely to an actual family $\eta(a)\in\mathcal O_X$ on the original
norm-limit carrier. This family is independent of all representative
choices and satisfies
\[
 \eta(a)_{E_j}\equiv x_j\pmod{\mathfrak m_{E_j}^{h_j}}.
\]
\end{lemma}
\begin{proof}
Compatibility in $A$ gives, for $m\geq n$,
\[
 N_{E_m/E_n}(x_m)-x_n\in\mathfrak m_{E_n}^{h_n}.
\]
Fix $j$ and an integer output precision $N\geq0$.
Choose $n_0\geq j$ with $h_{n_0}\geq\Psi_{L/E_j}(N)$.
For $m\geq n\geq n_0$, apply actual nonnormal norm continuity to
the integral inputs $N_{E_m/E_n}(x_m)$ and $x_n$, then the uniform
intrinsic APF comparison. It gives
\[
 N_{E_m/E_j}(x_m)-N_{E_n/E_j}(x_n)
       \in\mathfrak m_{E_j}^{N}.
\]
The integer ideal powers are a neighborhood basis, so this sequence is
Cauchy and converges in $E_j$. Its terms are integral, and the original
integer ring is closed, so its limit $\gamma_j$ is integral.
For $i\leq j$, continuity of the original norm and norm transitivity
give $N_{E_j/E_i}(\gamma_j)=\gamma_i$.
For any actual finite $D\subseteq L$, choose $j$ with $D\subseteq E_j$
and define $\gamma_D=N_{E_j/D}(\gamma_j)$. Passing to the larger of two
possible indices and using these identities proves independence of $j$.
It also proves all original norm compatibilities. The compatible-family
bridge gives a unique element $\eta(a)$ of the same original carrier;
its integral coordinates give membership in $\mathcal O_X$.

If $x'_n$ is another set of representatives, then
$x_n-x'_n\in\mathfrak m_{E_n}^{h_n}$. The identical fixed-base norm
continuity argument makes the two norm sequences differ by a sequence
tending to zero in each fixed $E_j$. Hence their limits coincide,
proving independence of representatives. Finally, for every $n\geq j$,
compatibility already gives
$N_{E_n/E_j}(x_n)-x_j\in\mathfrak m_{E_j}^{h_j}$.
This fixed ideal power is closed. Passing to the proved limit gives
the stated congruence. No separatedness inference is made directly
from estimates in changing finite fields.
\end{proof}

\begin{theorem}[Truncation identifies the same integral norm families]
\label{fw:norm-limit-truncation-equivalence}
\source{wintenberger:page-0012,wintenberger:page-0013}
\dcref{Proposition 2.3.1 with the original norm-limit carrier}
\uses{fw:truncated-norm-reconstruction,fw:norm-limit-field-valuation,fw:integral-addition-limit,fw:deep-truncated-finite-tower,fw:ideal-powers-topology,fw:addition-limit}
\lean{FontaineWintenberger.normLimitTruncationEquiv, FontaineWintenberger.normLimitTruncation_eq_zero_iff, FontaineWintenberger.normLimitTruncationHom, FontaineWintenberger.deepIntegralReconstruction, FontaineWintenberger.normLimitTruncationHom_deepIntegralReconstruction, FontaineWintenberger.deepIntegralReconstruction_normLimitTruncationHom}
\leanok
The actual coordinate truncations define a unital ring isomorphism
\[
 \xi:\mathcal O_X\xrightarrow{\ \sim\ }A,\qquad
 \xi(\alpha)_n=[\alpha_{E_n}],
\]
whose inverse is the original-field reconstruction $\eta$.
For every $n$ its projection kernel is exactly
\[
 \ker(\operatorname{pr}_n\circ\xi)
        =\{\alpha\in\mathcal O_X:v(\alpha)\geq h_n\}.
\]
\end{theorem}
\begin{proof}
Every original coordinate of an integral family is integral, and its
norm compatibility gives an element of $A$. Multiplication, zero and
one are preserved exactly. To check addition at a fixed index $n$,
let $\alpha,\beta\in\mathcal O_X$. Every later actual norm over $E_n$
is additive modulo $\mathfrak m_{E_n}^{h_n}$, so for every finite
$F\supseteq E_n$,
\[
 N_{F/E_n}(\alpha_F+\beta_F)-\alpha_{E_n}-\beta_{E_n}
              \in\mathfrak m_{E_n}^{h_n}.
\]
The original norm-sum limit defines $\alpha+\beta$.
Closedness of this fixed ideal power proves the same congruence for
its limit, and therefore additivity of the $n$th truncation.
Thus $\xi$ is a ring homomorphism.

The reconstruction congruence gives $\xi\eta(a)=a$.
Conversely, for $\xi(\alpha)$ one may choose precisely the original
coordinates $\alpha_{E_n}$ as representatives. Their norms down to
any $E_j$ are constantly $\alpha_{E_j}$. Independence of representatives
then gives $\eta\xi(\alpha)=\alpha$. This proves the isomorphism on
the original families, including zero.
Finally, $\xi(\alpha)_n=0$ means
$v_{E_n}(\alpha_{E_n})\geq h_n$, which is equivalent to
$v(\alpha)\geq h_n$ by the chosen stable-coordinate identity.
This proves the exact kernel formula.
\end{proof}

\begin{theorem}[An actual norm-family uniformizer and the full normalized value group]
\label{fw:norm-limit-uniformizer}
\source{wintenberger:page-0009,wintenberger:page-0012}
\dcref{The value-group assertion of Theorem 2.1.3 and Proposition 2.2.4}
\uses{fw:truncated-norm-reconstruction,fw:deep-truncated-finite-tower,fw:truncated-norm-inverse-limit,fw:norm-limit-field-valuation,fw:normalized-field-valuation}
\lean{FontaineWintenberger.exists_normLimitValuation_eq_exp_neg_one, FontaineWintenberger.normLimitUniformizer_zpow_valuation, FontaineWintenberger.normLimitValuation_surjective, FontaineWintenberger.normLimitFieldValuation_valueGroup_eq_top, FontaineWintenberger.normLimitFieldValuation_isRankOneDiscrete, FontaineWintenberger.normLimitFieldValuation_isDiscreteValuationRing, FontaineWintenberger.normLimitUniformizer_isUniformizer, FontaineWintenberger.normLimitUniformizer_maximalIdeal, FontaineWintenberger.normLimitFieldValuation_eq_normalized}
\leanok
There is an actual original norm-compatible family $\Pi$ with $v(\Pi)=1$.
Thus $v(X_K(L)^\times)=\mathbb Z$, its actual valuation ring
$\mathcal O_X$ is a DVR with uniformizer $\Pi$, and its canonical
normalized valuation agrees with the existing stable valuation on $X_K(L)$.
\end{theorem}
\begin{proof}
Choose a uniformizer $\pi_0$ in the original finite DVR $\mathcal O_{E_0}$.
Surjectivity of $\operatorname{pr}_0:A\to B_0$, with $h_0=2$, supplies
a compatible truncated family whose coordinate is $[\pi_0]$.
Apply the proved reconstruction $\eta$ to obtain an actual $\Pi\in\mathcal O_X$.
Its $E_0$-coordinate differs from $\pi_0$ by an element of
$\mathfrak m_{E_0}^2$. Since $v_{E_0}(\pi_0)=1$, the ordinary finite-field
valuation inequality gives $v_{E_0}(\Pi_{E_0})=1$.
The stable-coordinate identity therefore gives $v(\Pi)=1$.
In particular $\Pi\ne0$. For every $m\in\mathbb Z$, its original
field power $\Pi^m$ has stable valuation $m$, proving surjectivity.

An integral nonzero $\alpha$ has an integer valuation $n\geq0$, and
$u=\alpha\Pi^{-n}$ has value zero, so both $u$ and $u^{-1}$ are
integral. Thus $\alpha=\Pi^n u$ with $u\in\mathcal O_X^\times$.
Every nonzero ideal has a member of least finite nonnegative valuation,
by well ordering. Such a member divides every other member, using the
valuation criterion for integrality of the quotient, so it generates
the ideal. The maximal ideal consists exactly of the positive-valued
elements and is $(\Pi)$. The valuation ring is consequently a local
principal ideal domain with nonzero maximal ideal, hence a DVR.
Finally every nonzero field element is $\Pi^m u$ for an integer $m$
and an integral unit $u$. Its canonical normalized additive valuation
and the original stable valuation both give $m$. They also agree at
zero, proving equality of the actual valuations on the same field.
\end{proof}

\begin{lemma}[The residue field of the original norm-limit valuation]
\label{fw:norm-limit-residue-field}
\source{wintenberger:page-0009,wintenberger:page-0011}
\dcref{The residue-field assertion in Theorem 2.1.3}
\uses{fw:norm-limit-truncation-equivalence,fw:abstract-extension-residue-field,fw:apf-residue-stabilization,fw:norm-limit-uniformizer}
\lean{FontaineWintenberger.normLimitResidueLayerEquiv, FontaineWintenberger.normLimitResidueLayerEquiv_residue, FontaineWintenberger.normLimitResidueEquiv, FontaineWintenberger.normLimitResidueEquiv_residue, FontaineWintenberger.normLimitResidue_perfectField}
\leanok
Let $\mathcal O_X$ be the valuation ring of the original norm-limit
field and let $E_0$ be the first field of a deep truncated tower.
Reduction of the original coordinate at $E_0$ induces an isomorphism
\[
 \mathcal O_X/\mathfrak m_X \simeq k_{E_0}.
\]
Choosing the tower above an actual residue-stability stage gives
$k_{E_0}\simeq k_L$ and hence an isomorphism of the original residue
field with the intrinsic residue field of $L$. No assertion of
independence of the chosen stage is required for these isomorphisms.
\end{lemma}
\begin{proof}
The actual truncation isomorphism identifies $\mathcal O_X$ with $A$.
Its zeroth projection is onto
$B_0=\mathcal O_{E_0}/\mathfrak m_{E_0}^2$. Compose it with ordinary
reduction $B_0\to k_{E_0}$. Both maps are onto, so their composite is
onto and sends an original family $\alpha$ to the residue of
$\alpha_{E_0}$. The actual DVR normalized valuation agrees with the stable field valuation
by the uniformizer theorem. The kernel therefore consists exactly of the
families for which $v_{E_0}(\alpha_{E_0})>0$. The stable-coordinate valuation identity
identifies this condition with $v_X(\alpha)>0$, including zero;
it is exactly membership in the original maximal ideal
$\mathfrak m_X$. The first isomorphism theorem therefore gives the
displayed isomorphism. Take the specified starting field of the tower
to contain a residue-stability witness. Its first field then lies
above that witness, and the actual residue inclusion
$k_{E_0}\to k_L$ is bijective by the intrinsic residue-field lemma.
Composing the two isomorphisms proves the last assertion.
\end{proof}

\begin{theorem}[Completeness for the existing canonical valuation uniformity]
\label{fw:norm-limit-valuation-complete}
\source{wintenberger:page-0012,wintenberger:page-0013}
\dcref{Proposition 2.2.4 transported through Proposition 2.3.1}
\uses{fw:norm-limit-truncation-equivalence,fw:truncated-norm-inverse-limit,fw:norm-limit-field-valuation,fw:norm-limit-uniformizer}
\lean{FontaineWintenberger.normLimitTruncationHom_isUniformInducing, FontaineWintenberger.normLimitIntegral_completeSpace, FontaineWintenberger.normLimit_completeSpace, FontaineWintenberger.normLimit_t2Space, FontaineWintenberger.normLimitIntegral_t2Space}
\leanok
For the uniformity of the existing stable field valuation,
$\mathcal O_X$ and $X_K(L)$ are complete and Hausdorff.
The actual map $\xi$ is a uniform isomorphism between the induced
valuation uniformity on $\mathcal O_X$ and the closed-product
uniformity on $A$.
\end{theorem}
\begin{proof}
For an integer $N\geq0$ put
$V_N=\{\alpha\in\mathcal O_X:v(\alpha)\geq N\}$.
The actual uniformizer has multiplicative value $\exp(-1)$, so its
powers represent the radii $\exp(-N)$ in the existing value group.
Every nonzero radius in that group exceeds $\exp(-N)$ for sufficiently
large $N$. Thus these subgroups form a neighborhood basis at zero for
the induced valuation uniformity, whose entourages are
$v(\alpha-\beta)\geq N$. The proved ring isomorphism and exact kernel
identity give
\[
 \xi(\alpha)_n=\xi(\beta)_n
 \quad\Longleftrightarrow\quad
 v(\alpha-\beta)\geq h_n.
\]
Since $h_n=n+2$ is cofinal among the nonnegative integers, and single
projection kernels are cofinal among the neighborhoods on $A$, this
identity proves uniform continuity of $\xi$ and of its inverse.
It therefore transfers the already proved Hausdorff completeness of
$A$ to the original valuation ring with its existing uniformity.

To pass to the whole field, let $\mathcal F$ be a nontrivial Cauchy
filter for that same valuation uniformity. The entourage
$v(x-y)\geq0$ supplies a set $S\in\mathcal F$ whose pairwise
differences are integral. Choose $a\in S$, which is possible because
the filter is nontrivial. Then $S\subseteq a+\mathcal O_X$.
After translating by $-a$, the filter is eventually in
$\mathcal O_X$ and induces there a nontrivial Cauchy filter.
Its completeness gives a limit $b\in\mathcal O_X$.
Translating back gives convergence of $\mathcal F$ to $a+b$ in the
original field. The trivial filter converges automatically.
Hausdorffness follows from $v(x)=+\infty$ only for $x=0$.
Thus the completeness just proved belongs to the original valuation
uniformity, rather than to a separately chosen topology on a model.
\end{proof}

\begin{lemma}[The intrinsic residue field of the abstract algebraic extension]
\label{fw:abstract-extension-residue-field}
\source{wintenberger:page-0008}
\dcref{Original-field interpretation of the residue field in paragraph 2.1.2}
\uses{fw:spectral-integers,fw:spectral-integer-residue-maps,fw:apf-residue-stabilization,fw:norm-limit-multiplicative}
\lean{FontaineWintenberger.integralClosure_valuationRing, FontaineWintenberger.mem_integralClosure_inclusion_iff, FontaineWintenberger.integralClosureInclusion_maximalIdeal_comap, FontaineWintenberger.integralClosureSpectralResidueEquiv, FontaineWintenberger.IsAPF.exists_finiteLayer_integralClosureResidue_bijective, FontaineWintenberger.integralClosureResidueEquiv, FontaineWintenberger.IsAPF.perfectField_integralClosureResidue}
\leanok
Let $\mathcal O_L$ be the integral closure of $\mathcal O_K$ in the
original abstract separable algebraic field $L$.
It is a valuation ring, with maximal ideal $\mathfrak m_L$, and put
$k_L=\mathcal O_L/\mathfrak m_L$.
For every actual finite $E\subseteq L$,
\[
 \mathcal O_L\cap E=\mathcal O_E,\qquad
 \mathfrak m_L\cap E=\mathfrak m_E.
\]
On a residue-stable finite tail the actual inclusions induce
isomorphisms $k_E\simeq k_L$. In particular $k_L$ is perfect, with no
finiteness assumption on its cardinality.
\end{lemma}
\begin{proof}
Every $x\in L$ lies in an actual finite layer $K(x)$.
The spectral integer-ring theorem identifies integrality over
$\mathcal O_K$ on that layer with membership in its original valuation
ring. Thus $\mathcal O_L$ is the directed union of these original
integer rings. Applied to nonzero $x$, the valuation-ring property in
$K(x)$ says that $x$ or $x^{-1}$ is integral, so this union is a
valuation ring. Restriction of the spectral valuation to each finite
layer proves the first intersection identity. An element of
$\mathcal O_E$ becomes a unit in $\mathcal O_L$ only if its inverse,
lying in the same field $E$, is integral; this proves the second identity.
Consequently each finite residue inclusion $k_E\to k_L$ is injective.

Take a finite residue-stability witness $T$. For $E\supseteq T$,
any residue class of an integral $x\in L$ is represented in a finite
$F\supseteq E$ containing $x$. The proved finite residue inclusion
$k_E\to k_F$ is surjective, so some $y\in\mathcal O_E$ has
$x-y\in\mathfrak m_F\subseteq\mathfrak m_L$.
This proves surjectivity onto $k_L$ and hence the asserted actual
isomorphism. The field $k_E$ is a finite extension of the perfect field
$k_K$, so it is perfect. Its isomorphism with $k_L$ proves perfectness
without asserting finiteness as a set.
\end{proof}

\begin{lemma}[Finite maximal unramified and tame parts]
\label{fw:apf-finite-unramified-tame-parts}
\source{wintenberger:page-0008}
\dcref{The finiteness and normalization of the fields in paragraph 2.1.2}
\uses{fw:apf-abstract,fw:apf-products,fw:absolute-upper-projections,fw:profinite-galois-foundations,fw:inertia-relative-index,fw:wild-ramification-p-group,fw:tame-inertia-quotient,fw:apf-residue-stabilization,fw:abstract-extension-residue-field}
\lean{FontaineWintenberger.apfMaximalUnramified, FontaineWintenberger.apfMaximalTame, FontaineWintenberger.le_apfMaximalUnramified_iff, FontaineWintenberger.isGreatest_apfMaximalUnramified, FontaineWintenberger.apfMaximalUnramified_le_maximalTame, FontaineWintenberger.apfMaximalUnramifiedResidueEquiv, FontaineWintenberger.apfMaximalUnramifiedResidueEquiv_apply, FontaineWintenberger.le_apfMaximalTame_iff, FontaineWintenberger.isGreatest_apfMaximalTame, FontaineWintenberger.finrank_eq_prime_pow_of_apfMaximalTame_le, FontaineWintenberger.apfMaximalTame_wild_tail}
\leanok
There are unique maximal unramified and maximal tamely ramified
subextensions $K_0/K$ and $K_1/K$ inside the original $L/K$.
They are finite, $K_0\subseteq K_1$, and $k_{K_0}=k_L$ under the
actual residue inclusions. Every finite $E\supseteq K_1$ in $L$
is totally ramified over $K_1$ of degree a power of $p$.
\end{lemma}
\begin{proof}
Choose an actual separable-closure embedding of $L$ and write $H=G_L$.
Put $I=G_K^0$, and let $P$ be the closed subgroup generated by
$G_K^u$ for all $u>0$. It is normal and contained in $I$.
In any finite Galois quotient with group $G$, the image of $I$ is
$G_0$. The image of $P$ is exactly $G_1$: every positive upper group
is contained in $G_1$, and for a sufficiently small positive parameter
it is $G_1$, by continuity and strict increase of the finite Herbrand
parameter. Thus the closed subgroup just defined has precisely the
usual finite wild-inertia images; this proves the property of $P$
needed below.

The products $HI$ and $HP$ are closed because they are products of
compact subgroups with a normal factor. APF makes $HI$ open.
Since $G_K^1\subseteq P$, the subgroup $HP$ contains the open APF
product $HG_K^1$ and is open as well.
Their fixed fields $K_0$ and $K_1$ are finite and contained in $L$,
and $P\subseteq I$ gives $K_0\subseteq K_1$.
The zero-stage residue-stabilization theorem identifies the actual
residue field of $K_0$ with the intrinsic $k_L$ just constructed.

For a finite subfield $M/K$ inside $L$, choose a finite Galois
normal closure $\Omega/K$ containing it and write
$J=\operatorname{Gal}(\Omega/M)$, $D=G_0$, and $P_\Omega=G_1$.
The actual ramification index is $[D:J\cap D]$.
It is one exactly when $D\subseteq J$, so the unramified subfields
are exactly those fixed by $I$. It is prime to $p$ exactly when
$P_\Omega\subseteq J$. Indeed, if this containment holds the index
divides $[D:P_\Omega]$, which is prime to $p$. Conversely
$[(J\cap D)P_\Omega:J\cap D]$ is both a power of $p$ and a divisor
of $[D:J\cap D]$; if that last index is prime to $p$, the former
index is one and $P_\Omega\subseteq J$.
Since the residue fields are perfect, prime-to-$p$ ramification index
is exactly the finite tame condition. This identifies the tame
subfields with those fixed by $P$.
The fields fixed by $HI$ and $HP$ are therefore the unique maximal
ones, which also proves that their original subfields are independent
of the initial separable embedding.

Finally let finite $E\supseteq K_1$ lie in $L$, and take a finite
Galois $\Omega/K$ containing $E$. Let $H_\Omega$ be the image of $H$.
The actual subgroup fixing $K_1$ in this quotient is
$H_\Omega P_\Omega$, and the subgroup $J$ fixing $E$ contains
$H_\Omega$ and is contained in that product.
The index $[H_\Omega P_\Omega:H_\Omega]$ is a power of $p$, since
its cosets are the $P_\Omega/(H_\Omega\cap P_\Omega)$ cosets.
Its divisor $[H_\Omega P_\Omega:J]=[E:K_1]$ is therefore a power
of $p$. Both residue fields are already $k_L$ because both fields
contain $K_0$, so this actual finite extension is totally ramified.
\end{proof}

\begin{lemma}[Perfect-residue multiplicative representatives in the original fields]
\label{fw:perfect-multiplicative-residue-lifts}
\source{wintenberger:page-0008,wintenberger:page-0011}
\dcref{Construction of the multiplicative representatives used in paragraphs 2.1.2 and 2.2.2.6}
\uses{fw:ideal-powers-topology,fw:ideal-power-valuation,fw:adic-complete-unit-ball,fw:spectral-integer-residue-maps}
\lean{FontaineWintenberger.residueTeichmuller, FontaineWintenberger.residue_residueTeichmuller, FontaineWintenberger.residueTeichmuller_unique, FontaineWintenberger.residueTeichmuller_frobenius, FontaineWintenberger.residueTeichmuller_sModEq, FontaineWintenberger.residueTeichmuller_naturality, FontaineWintenberger.spectralTeichmuller, FontaineWintenberger.residue_spectralTeichmuller, FontaineWintenberger.spectralTeichmuller_unique, FontaineWintenberger.spectralTeichmuller_naturality, FontaineWintenberger.coe_spectralTeichmuller_naturality}
\leanok
Let $E$ be a complete discrete valuation field with perfect residue
field $k$ of characteristic $p$. There is a unique multiplicative
section $x\mapsto[x]$ of $\mathcal O_E\to k$ preserving zero and one.
It commutes with Frobenius and with the actual inclusions of complete
finite valued extensions with perfect residue fields.
\end{lemma}
\begin{proof}
If $a,b\in\mathcal O_E$ and $v_E(a-b)\geq q\geq1$, the binomial
formula and divisibility of $\binom pj$ by $p$ for $0<j<p$ give
\[
 v_E(a^p-b^p)\geq q+1.
\]
Indeed the term $(a-b)^p$ has valuation at least $pq\geq q+1$,
and every other term has valuation at least $v_E(p)+q\geq q+1$;
if $p=0$ in $E$ those other terms vanish. Iteration gives
$v_E(a^{p^m}-b^{p^m})\geq q+m$.
For $x\in k$, choose arbitrary lifts $b_m$ of $x^{p^{-m}}$.
The residues of $b_{m+1}^p$ and $b_m$ agree, so
\[
 v_E(b_{m+1}^{p^{m+1}}-b_m^{p^m})\geq m+1.
\]
Thus $b_m^{p^m}$ is Cauchy and has an integral limit $[x]$.
Two choices of $b_m$ give terms differing to valuation at least $m+1$,
so the limit is independent of all choices. Its residue is $x$.
For products one may choose $b_m(x)b_m(y)$ as the lift of
$(xy)^{p^{-m}}$; continuity of multiplication gives
$[xy]=[x][y]$. Choosing constant zero and one lifts gives $[0]=0$
and $[1]=1$, and multiplicativity gives $[x^p]=[x]^p$.
Any multiplicative section $s$ satisfies
$s(x)=s(x^{p^{-m}})^{p^m}$. Comparing these lifts with $b_m$ by the
same estimate forces $s(x)=[x]$, proving uniqueness.
Actual finite field inclusions are continuous and send these lifts
to valid lifts of the included residue roots. Passing to the same
limits proves compatibility with those original inclusions.
\end{proof}

\begin{lemma}[Canonical coefficients in a characteristic-$p$ truncated ring]
\label{fw:perfect-truncated-coefficients}
\source{wintenberger:page-0012}
\dcref{Finite nilpotent form of the coefficient construction preceding Proposition 2.2.4}
\uses{fw:perfect-multiplicative-residue-lifts}
\lean{FontaineWintenberger.truncatedResidueCoefficients, FontaineWintenberger.truncatedResidueCoefficients_eq_residueTeichmuller, FontaineWintenberger.residue_truncatedResidueCoefficients, FontaineWintenberger.truncatedResidueCoefficients_unique, FontaineWintenberger.truncatedResidueCoefficients_eq_pow, FontaineWintenberger.truncated_residue_pow_independent, FontaineWintenberger.truncatedResidueCoefficients_naturality}
\leanok
Let $B$ be a commutative local ring of characteristic $p$, with perfect residue
field $k$ and maximal ideal $\mathfrak n$ satisfying
$\mathfrak n^h=0$ for some $h\geq1$.
For $x\in k$, choose $m$ with $p^m\geq h$ and any lift $b\in B$
of $x^{p^{-m}}$. Then
\[
 c_B(x)=b^{p^m}
\]
is independent of both choices and defines a unital ring section
$k\to B$. For a local unital ring homomorphism $B\to B'$ between
such rings, these sections commute with the induced residue map.
\end{lemma}
\begin{proof}
If $b,b'$ have the same residue, their difference is in $\mathfrak n$.
In characteristic $p$ their $p^m$th powers differ by
$(b-b')^{p^m}=0$ whenever $p^m\geq h$.
If $b'$ lifts $x^{p^{-m-1}}$, then $(b')^p$ and a lift of
$x^{p^{-m}}$ have the same residue; the preceding independence shows
that the constructions at $m$ and $m+1$ agree. Comparing any two
admissible exponents at their maximum proves complete independence.
The inverse Frobenius on the perfect residue field is a field
automorphism. Consequently sums and products of chosen lifts lift
the corresponding inverse-Frobenius roots. Frobenius additivity and
multiplicativity in $B$ then give
$c_B(x+y)=c_B(x)+c_B(y)$ and $c_B(xy)=c_B(x)c_B(y)$.
Constant zero and one lifts give a unital ring map, and its reduction
is $x$. For a local homomorphism to $B'$, choose one $m$ large enough
for both nilpotence exponents. The image of a chosen lift is a lift
of the image residue root, since field homomorphisms commute with
Frobenius and its inverse. Taking its $p^m$th power proves naturality.
\end{proof}

\begin{theorem}[The source residue constants and the actual residue-field isomorphism]
\label{fw:source-residue-constants-field-embedding}
\source{wintenberger:page-0008,wintenberger:page-0009,wintenberger:page-0012}
\dcref{The coefficient embedding in Theorem 2.1.3 and Proposition 2.2.4}
\uses{fw:residue-constants,fw:abstract-extension-residue-field,fw:apf-finite-unramified-tame-parts,fw:perfect-multiplicative-residue-lifts,fw:perfect-truncated-coefficients,fw:deep-truncated-finite-tower,fw:norm-limit-truncation-equivalence,fw:truncated-norm-inverse-limit,fw:norm-limit-uniformizer,fw:norm-limit-residue-field}
\lean{FontaineWintenberger.sourceResidueConstantsFieldEmbedding, FontaineWintenberger.sourceResidueConstantsFieldEmbedding_apply, FontaineWintenberger.sourceResidueConstantsFieldEmbedding_injective, FontaineWintenberger.sourceResidueConstantsResidueEquiv, FontaineWintenberger.sourceResidueConstantsResidueEquiv_apply, FontaineWintenberger.normLimitEval_sourceResidueConstants_any, FontaineWintenberger.normLimitTruncation_sourceResidueConstantsInteger, FontaineWintenberger.sourceResidueConstantsInteger_add}
\leanok
Identify $k_{K_0}$ with the intrinsic $k_L$ by the actual residue map.
For finite $E\supseteq K_1$, write $[E:K_1]=p^{d_E}$ and put
\[
 f_{L/K}(x)_E=[x^{p^{-d_E}}]\in K_0\subseteq E.
\]
These coordinates extend uniquely to an original norm-compatible family.
They define a field embedding $f_{L/K}:k_L\to X_K(L)$, and its
reduction induces an isomorphism from $k_L$ onto the actual residue
field of the original valuation ring $\mathcal O_X$.
\end{theorem}
\begin{proof}
The preceding two original-field lemmas supply the actual finite
$K_0,K_1$, their residue identification, the $p$-power degrees, and
the multiplicative representatives. For $K_1\subseteq E\subseteq F$,
the displayed element at $F$ is already in $K_0\subseteq E$.
The scalar norm identity therefore gives
\[
 N_{F/E}([x^{p^{-d_F}}])
 =[x^{p^{-d_F}}]^{p^{d_F-d_E}}=[x^{p^{-d_E}}].
\]
The tail above $K_1$ is cofinal. For arbitrary finite $D$, take a
finite $E\supseteq DK_1$ and norm its prescribed coordinate down to
$D$. Passing to a common upper field proves independence of $E$
and every original norm compatibility. Thus these formulas construct
an element of the same original norm-limit carrier. Evaluating at a
stable upper layer containing $K_1$ gives nonnegative stable valuation,
so this family belongs to $\mathcal O_X$. Multiplication, zero and one
are preserved by construction.

Choose the deep truncated tower with its specified starting field
containing $K_1$, and write $[E_n:K_1]=p^{d_n}$.
The residue field of each $B_n$ is canonically the actual $k_L$.
It has characteristic $p$ and nilpotent maximal ideal, so the preceding
coefficient lemma gives a ring section $c_n:k_L\to B_n$.
For any $m$ with $p^m\geq h_n$, the original representative satisfies
\[
 [x^{p^{-d_n}}]=[x^{p^{-d_n-m}}]^{p^m}.
\]
Thus its class in $B_n$ is exactly $c_n(x^{p^{-d_n}})$.
Inverse Frobenius is additive and multiplicative on $k_L$.
Each coordinate of $\xi\circ f_{L/K}$ is therefore a unital ring map.
The proved injectivity of $\xi$ gives additivity of $f_{L/K}$ itself,
so this is the source coefficient field embedding into the original field.

Explicitly, the composition
$\mathcal O_X\xrightarrow{\xi}A\xrightarrow{\operatorname{pr}_0}B_0\to k_L$
is a surjective ring homomorphism $\rho$; it is original evaluation at
$E_0$ followed by residue reduction. Surjectivity follows from the actual
surjectivity of $\operatorname{pr}_0:A\to B_0$ and of $B_0\to k_L$.
Its kernel consists exactly of the families with stable valuation at
least one, by the stable-coordinate identity at $E_0$.
This is the actual maximal ideal of $\mathcal O_X$, so $\rho$ induces
an isomorphism of its residue field with $k_L$.
For the particular source normalization above,
\[
 \rho(f_{L/K}(x))=x^{p^{-d_0}},
\]
an automorphism of the perfect field $k_L$.
Hence $f_{L/K}$ is injective and induces an isomorphism onto the full
actual residue field. This also records the Frobenius twist caused by
using a finite starting field above $K_1$, rather than changing the
source normalization of $f_{L/K}$.
\end{proof}

\begin{theorem}[The field of norms]
\label{fw:field-norms}
\dcref{Wintenberger1983, Theorem 2.1.3}
\source{wintenberger:page-0008,wintenberger:page-0009}
\uses{fw:norm-limit-field,fw:norm-limit-uniformizer,fw:norm-limit-valuation-complete,fw:source-residue-constants-field-embedding,fw:norm-limit-residue-field}
\lean{FontaineWintenberger.normLimitField, FontaineWintenberger.normLimitField_charP, FontaineWintenberger.normLimitFieldValuation, FontaineWintenberger.normLimitValued, FontaineWintenberger.normLimit_completeSpace, FontaineWintenberger.normLimitFieldValuation_isRankOneDiscrete, FontaineWintenberger.normLimitFieldValuation_isDiscreteValuationRing, FontaineWintenberger.normLimitValuation_surjective, FontaineWintenberger.sourceResidueConstantsFieldEmbedding, FontaineWintenberger.sourceResidueConstantsFieldEmbedding_injective, FontaineWintenberger.sourceResidueConstantsResidueEquiv, FontaineWintenberger.sourceResidueConstantsResidueEquiv_apply}
\leanok
For every infinite APF extension $L/K$, the preceding addition,
multiplication, and valuation make $X_K(L)$ a complete discretely valued
field of characteristic $p$, with $v(X_K(L)^\times)=\mathbb Z$.
The map $f_{L/K}:k_L\longrightarrow X_K(L)$ is a field embedding and
induces an isomorphism of $k_L$ with the residue field of $X_K(L)$.
\end{theorem}
\begin{proof}
The field structure and characteristic $p$ are the already established
ones on the original multiplicative norm limit. The actual family
uniformizer makes its value group all of $\mathbb Z$ and its valuation
ring a DVR, whose normalized valuation is the same stable valuation.
The completeness theorem applies to precisely that existing valuation
uniformity on the same field. Finally the constructed source constants
$f_{L/K}$ give a field embedding of the intrinsic original $k_L$ and
induce an isomorphism onto the actual residue field of this DVR.
These assertions give every part of the theorem for arbitrary perfect
residue field and for every residue prime, including $p=2$.
\end{proof}

\begin{theorem}[Laurent coordinates attached to every uniformizer]
\label{fw:uniformizer-laurent-coordinates}
\source{wintenberger:page-0008,wintenberger:page-0009}
\dcref{The complete equicharacteristic discrete valuation field in Theorem 2.1.3, with its coefficient embedding}
\uses{fw:field-norms}
\lean{FontaineWintenberger.powerSeriesUniformizerEquiv, FontaineWintenberger.powerSeriesUniformizerEquiv_C, FontaineWintenberger.powerSeriesUniformizerEquiv_X, FontaineWintenberger.powerSeriesUniformizerEquiv_unique, FontaineWintenberger.continuous_powerSeriesUniformizerEquiv, FontaineWintenberger.valuedPowerSeriesEquiv, FontaineWintenberger.valuation_valuedPowerSeriesEquiv, FontaineWintenberger.laurentUniformizerEquiv, FontaineWintenberger.valuation_laurentUniformizerEquiv, FontaineWintenberger.continuous_laurentUniformizerEquiv, FontaineWintenberger.continuous_laurentUniformizerEquiv_symm, FontaineWintenberger.laurentUniformizerEquiv_unique, FontaineWintenberger.normLimitPowerSeriesEquiv, FontaineWintenberger.normLimitLaurentEquiv, FontaineWintenberger.normLimitLaurentEquiv_C, FontaineWintenberger.normLimitLaurentEquiv_X, FontaineWintenberger.normLimitLaurentEquiv_valuation, FontaineWintenberger.continuous_normLimitLaurentEquiv, FontaineWintenberger.continuous_normLimitLaurentEquiv_symm, FontaineWintenberger.normLimitLaurentEquiv_unique}
\leanok
Let $F$ be a complete field with an integer-valued multiplicative valuation
$v$, let $\mathcal O=\{x:v(x)\leq1\}$ be its actual valuation ring, and
let $k$ be a field. Suppose $c:k\to\mathcal O$ is a unital ring map whose
composition with the original residue quotient is an isomorphism
$\bar c:k\overset\sim\longrightarrow\mathcal O/\mathfrak m$.
For every $\pi\in\mathcal O$ with $v(\pi)=\exp(-1)$, evaluation is an
isomorphism
\[
 E_{c,\pi}:k[[T]]\overset\sim\longrightarrow\mathcal O,
 \qquad \sum_{n\geq0}a_nT^n\longmapsto\sum_{n\geq0}c(a_n)\pi^n.
\]
It is the unique continuous unital ring map sending each constant $a$
to $c(a)$ and $T$ to $\pi$. Its canonical extension to fraction fields is
a continuous valuation-preserving field isomorphism
$k((T))\overset\sim\longrightarrow F$, sending the Laurent parameter to
$\pi$ and constants to the original coefficient embedding. Its inverse
is continuous. In particular, for every uniformizer $u$ of the original
field of norms, the original source coefficient embedding $f_{L/K}$
and its induced residue isomorphism give these isomorphisms with
$F=X_K(L)$ and $\pi=u$; the field operations, valuation and topology are
the already constructed ones on this same norm-limit carrier.
\end{theorem}
\begin{proof}
The actual valuation ring is a closed subset of the complete valued field,
hence complete and Hausdorff for its inherited uniformity. Valuation balls
in it are ideals: addition satisfies the ultrametric inequality and
multiplication by an integral element cannot increase the valuation.
These balls give a basis of neighborhoods of zero. Since
$v(\pi^n)=\exp(-n)\to0$, evaluation of power series converges in this
original integer ring and is a continuous unital ring homomorphism.

To prove surjectivity, start with $q_0=x\in\mathcal O$. Given $q_n$, put
$a_n=\bar c^{-1}(q_n\bmod\mathfrak m)$. The difference $q_n-c(a_n)$
has zero residue, so belongs to $\mathfrak m=(\pi)$. There is therefore
a unique $q_{n+1}\in\mathcal O$ such that
$q_n-c(a_n)=\pi q_{n+1}$; uniqueness uses $\pi\ne0$.
Induction gives the exact identity
\[
 x-\sum_{j<n}c(a_j)\pi^j=\pi^nq_n.
\]
Since $q_n$ is integral, the valuation of the right side is at most
$\exp(-n)$, hence the partial sums tend to $x$. The power series
$\sum a_nT^n$ therefore evaluates to $x$ by uniqueness of limits.

If $A\ne0$ has first nonzero coefficient at $r$, factor it as
$A=T^rB$, where $B$ has nonzero constant coefficient. Thus $B$ is a
unit of $k[[T]]$, and its image is a unit of the actual integer ring.
Its value is consequently $1$, so
$v(E_{c,\pi}(A))=\exp(-r)$. In particular this image is nonzero.
The ring map is injective, and the surjectivity already proved makes it
an isomorphism. Polynomials are dense for the coefficientwise topology
on $k[[T]]$, with $k$ discrete; continuous maps agreeing on constants
and the parameter agree on every polynomial and hence on all series.
This proves the asserted uniqueness.

Every element of $k((T))$ is a quotient of two power series with nonzero
denominator, and $F$ is the fraction field of its actual valuation ring.
The ring isomorphism therefore extends uniquely by
$A/B\mapsto E_{c,\pi}(A)/E_{c,\pi}(B)$ to a field isomorphism.
The just-proved value formula gives preservation of the valuation on
these quotients, including zero. The preimages, and the images under
the inverse map, of valuation balls of every radius are the corresponding
valuation balls. Both maps are therefore continuous for precisely the
existing valuation topologies. Finally the field-of-norms theorem supplies
the complete original field, and the source coefficient embedding induces
its actual residue isomorphism. Applying the general construction to each
of its uniformizers proves the specialization without changing any field
or topological structure.
\end{proof}

\begin{theorem}[Continuous Laurent maps and power-series substitution]
\label{fw:laurent-power-series-substitution}
\source{wintenberger:page-0008,wintenberger:page-0009}
\dcref{Continuity and coefficient expansion in the complete equicharacteristic field of Theorem 2.1.3}
\uses{fw:uniformizer-laurent-coordinates}
\lean{FontaineWintenberger.continuous_ofPowerSeries_valued, FontaineWintenberger.continuous_powerSeries_subst, FontaineWintenberger.map_ofPowerSeries_eq_subst, FontaineWintenberger.comp_apply_parameter, FontaineWintenberger.iterate_ofPowerSeries_eq_subst_iterate, FontaineWintenberger.iterate_apply_parameter_eq_subst_iterate}
\leanok
Let $k$ be any field, with the discrete topology on its coefficients,
and give $k((T))$ its original normalized valuation topology.
Write $j:k[[T]]\hookrightarrow k((T))$ for the usual inclusion.
Let $F:k((T))\to k((T))$ be a continuous unital ring homomorphism
fixing $k$. If $g\in k[[T]]$ satisfies $g(0)=0$ and $F(T)=j(g)$, then
\[
 F(j(A))=j(A\circ g)\qquad(A\in k[[T]]).
\]
No condition on the linear coefficient of $g$ is required.
For any unital ring homomorphism $G$ with $G(T)=j(h)$, this implies
\[
 (F\circ G)(T)=j(h\circ g).
\]
Moreover, writing $S_g(A)=A\circ g$, for every $n\geq0$ and $A\in k[[T]]$
one has
\[
 F^{\circ n}(j(A))=j(S_g^{\circ n}(A)).
\]
In particular this describes the image of $T$ under every iterate.
\end{theorem}
\begin{proof}
These are consequences of continuity and the usual coefficient expansion;
the cited field-of-norms theorem supplies the complete field when this
assertion is applied to its Laurent coordinates.

First the usual inclusion $j$ is continuous from the coefficientwise
topology on $k[[T]]$ to the original valuation topology.
Indeed, if the first $N$ coefficients of $A$ vanish, then
$v(j(A))\leq\exp(-N)$. The conditions that these finitely many
coefficients vanish form an open neighborhood of zero because $k$ is
discrete. Since $\exp(-N)$ tends to zero, these neighborhoods map into
every valuation neighborhood of zero. Additivity then gives continuity
at every point.

Substitution by $g$ is a continuous unital ring endomorphism of
$k[[T]]$. To see that it is defined, note that $g^r$ has order at least
$r$, so each coefficient in $\sum_{r\geq0}a_rg^r$ is a finite sum.
In particular the coefficient of degree $m$ depends only on
$a_0,\ldots,a_m$; this proves continuity in the coefficientwise topology.
The ring laws follow by finite coefficient sums, and constants and $T$
map respectively to the same constants and to $g$.

The maps $F\circ j$ and $j\circ S_g$ are therefore continuous into
the Hausdorff valued field $k((T))$. They agree on constants and $T$,
hence on polynomials by the ring laws. Every power series is the
coefficientwise limit of its polynomial truncations. Continuity and
uniqueness of limits give their equality on every power series.

Applying this equality to $A=h$ and using $G(T)=j(h)$ proves the
composition formula, including its stated order of substitution.
Finally induction on $n$ proves the iterate formula: the case $n=0$
is the identity, and applying $F$ to the inductive equality and then the
already established formula for $S_g^{\circ n}(A)$ gives the successor
case. Taking $A=T$ gives the parameter formula.
\end{proof}

\section{Functoriality and separable extensions}

The norm-family constructions in this section concern actual separable
algebraic extensions of $K$. Whenever $X_K(L)$ is used as a field,
$L/K$ is infinite APF and $K$ retains the standing complete discretely
valued hypotheses with arbitrary perfect residue field of characteristic $p$.
The first two algebraic descent lemmas retain their explicitly stated
more general field hypotheses.

\begin{lemma}[Cofinal finite tensor descent]
\label{fw:cofinal-finite-tensor-descent}
\source{wintenberger:page-0014}
\dcref{Expansion of the cofinality assertion in paragraph 3.1.1}
\lean{FontaineWintenberger.exists_finiteLayer_tensor_descent, FontaineWintenberger.adjoinTensorEquiv, FontaineWintenberger.adjoinTensorEquiv_tmul, FontaineWintenberger.adjoin_inf_fieldRange_eq_bot, FontaineWintenberger.finiteDescentIntersection, FontaineWintenberger.finiteDescentIntersection_map, FontaineWintenberger.finiteDescentUpperTensorEquiv, FontaineWintenberger.finiteDescentUpperTensorEquiv_tmul, FontaineWintenberger.exists_finiteLayer_tensor_descent_tower, FontaineWintenberger.exists_common_finiteLayer_tensor_descent_tower}
\leanok
Let $M/K$ be algebraic and let $N/M$ be finite and separable.
The finite $K$-subfields $D\subseteq N$ for which, putting $C=M\cap D$,
\[
 M\otimes_C D\longrightarrow N
\]
is an isomorphism are cofinal among all finite $K$-subfields of $N$.
For a finite separable tower $M\subseteq N\subseteq P$, one can choose
cofinally many finite $D\subseteq P$ good simultaneously for $P/M$
and $P/N$, with $N\cap D$ good for $N/M$.
\end{lemma}
\begin{proof}
Choose a primitive element $\theta$ for $N/M$, and put $d=[N:M]$.
Given a finite $K$-subfield $D_0\subseteq N$, choose finitely many
$K$-generators of $D_0$ and express each in the basis
$1,\theta,\ldots,\theta^{d-1}$ over $M$.
Let $C_0$ be the $K$-field generated by all these coefficients and all
coefficients of the monic minimal polynomial of $\theta$ over $M$.
It is finite over $K$ because $M/K$ is algebraic. Put $D=C_0(\theta)$.
The chosen expressions show $D_0\subseteq D$.
The displayed monic polynomial has degree $d$ over $C_0$, and a smaller
relation over $C_0$ would also be a smaller relation over $M$.
Thus $[D:C_0]=d$. The canonical tensor map onto $M(\theta)=N$
is surjective between $M$-vector spaces of dimension $d$, hence is an
isomorphism. Its injectivity implies $M\cap D=C_0$: for an element
of the intersection, the two corresponding scalars in the tensor
product are equal, which, after extending $1$ to a $C_0$-basis of $D$,
forces every nonconstant basis coefficient to vanish.
This proves cofinality; applying it to finite composita also gives
common good upper fields.

For the tower assertion choose a primitive element $a$ for $N/M$.
Apply the preceding construction to $P/M$, requiring $D$ to contain
$a$ and the prescribed finite subfield of $P$.
Put $C=M\cap D$ and $B=N\cap D$. Then $MB=N$, since $a\in B$,
and $ND=P$, since $MD=P$. Tensor injectivity for $M,D$ over $C$
restricts to tensor injectivity for $M,B$; as $MB=N$, it gives
$[B:C]=[N:M]$. The degree identities
$[D:C]=[P:M]$ and $[P:M]=[P:N][N:M]$ now give
$[D:B]=[P:N]$. Consequently the surjective canonical maps
$M\otimes_C B\to N$ and $N\otimes_B D\to P$ are isomorphisms by
equality of dimensions. This is the required simultaneous choice.
\end{proof}

\begin{lemma}[The original algebraic norm under tensor base change]
\label{fw:algebraic-norm-tensor-basechange}
\source{wintenberger:page-0014}
\dcref{The determinant calculation underlying the norm identity in paragraph 3.1.1}
\uses{fw:cofinal-finite-tensor-descent}
\lean{FontaineWintenberger.norm_tensor_includeRight, FontaineWintenberger.norm_eq_of_tensorEquiv, FontaineWintenberger.norm_adjoin_of_minpoly_lifts}
\leanok
Let $R,A,B$ be commutative rings, let $A,B$ be $R$-algebras, and
suppose $B$ is finite free over $R$. For $b\in B$ one has
\[
 N_{A\otimes_R B/A}(1\otimes b)
   =\iota_{R,A}(N_{B/R}(b)).
\]
If $e:A\otimes_R B\simeq_A S$ is an algebra isomorphism, the same
identity holds for $N_{S/A}(e(1\otimes b))$. In particular it applies
to the actual canonical multiplication isomorphism of any descended
primitive field in Lemma~\ref{fw:cofinal-finite-tensor-descent}.
\end{lemma}
\begin{proof}
Choose a finite $R$-basis $(b_i)_i$ of $B$.
The elements $(1\otimes b_i)_i$ form an $A$-basis of $A\otimes_R B$.
If multiplication by $b$ has matrix $(r_{ij})$ in the first basis,
multiplication by $1\otimes b$ has matrix
$(\iota_{R,A}(r_{ij}))$ in the second. The determinant commutes
with the ring map $\iota_{R,A}$, which proves the first identity
by the definition of the original algebraic norm.
Transporting the second basis through $e$ gives exactly the same
multiplication matrix in $S$, proving the second identity.
The canonical isomorphism constructed by finite descent sends
$1\otimes d$ to the original element $d$, so this conclusion is
an identity of the original relative field norms.
\end{proof}

\begin{lemma}[Original norm base change on good finite pairs]
\label{fw:good-finite-pair-norm-basechange}
\source{wintenberger:page-0014}
\dcref{The norm identity in paragraph 3.1.1}
\uses{fw:cofinal-finite-tensor-descent,fw:algebraic-norm-tensor-basechange}
\lean{FontaineWintenberger.finiteLayerPullback, FontaineWintenberger.IsGoodFiniteLayer, FontaineWintenberger.exists_goodFiniteLayer_ge, FontaineWintenberger.exists_goodFiniteLayer_ge_pair, FontaineWintenberger.goodFinitePairTensorEquiv, FontaineWintenberger.goodFinitePairTensorEquiv_tmul, FontaineWintenberger.linearDisjoint_goodFinitePair, FontaineWintenberger.goodFinitePair_sup_eq, FontaineWintenberger.goodFinitePair_norm}
\leanok
Let $N/M$ be as above and let $D\subseteq D'$ be good finite fields.
Put $C=M\cap D$ and $C'=M\cap D'$. Then $D'=DC'$, the fields
$D,C'$ are linearly disjoint over $C$, and for $a\in C'$,
\[
 N_{D'/D}(a)=N_{C'/C}(a)
\]
under the actual inclusions in $N$.
\end{lemma}
\begin{proof}
Injectivity of $M\otimes_C D\to N$ restricts to
$C'\otimes_C D\to N$, proving linear disjointness and
$[DC':D]=[C':C]$. Write $d=[N:M]$.
Goodness gives $[D:C]=[D':C']=d$. Multiplying degrees in the
finite tower over $C$ therefore gives $[D':D]=[C':C]$.
Since $DC'\subseteq D'$, equality of degrees gives $D'=DC'$.
Choose a $C$-basis of $C'$. Linear disjointness makes its image a
$D$-basis of $D'$. The matrix of multiplication by $a$ in the second
basis is the image of its matrix in the first basis. Its determinant,
which defines the original field norm, is therefore the image of
$N_{C'/C}(a)$. This proves the stated norm identity.
\end{proof}

\begin{lemma}[Transport of original good finite fields through isomorphisms]
\label{fw:good-finite-field-isomorphism}
\source{wintenberger:page-0014}
\dcref{The coordinate identification in paragraph 3.1.1, with its isomorphism transport made explicit}
\uses{fw:good-finite-pair-norm-basechange}
\lean{FontaineWintenberger.finiteLayerPullback_mem_map_iff, FontaineWintenberger.finiteLayerPullback_map, FontaineWintenberger.finiteLayerPullbackEquiv, FontaineWintenberger.coe_finiteLayerPullbackEquiv, FontaineWintenberger.finiteLayerTensorAddEquiv, FontaineWintenberger.finiteLayerTensorAddEquiv_tmul, FontaineWintenberger.finiteLayerTensorMap_equiv, FontaineWintenberger.isGoodFiniteLayer_map_iff}
\leanok
Let $M\to N$ and $M'\to N'$ be field embeddings over $K$, and let
$e:M\simeq_K M'$ and $f:N\simeq_K N'$ commute with these embeddings.
For a finite original $K$-field $D\subseteq N$, put $D'=f(D)$,
$C=M\cap D$ and $C'=M'\cap D'$, with intersections interpreted as
original pullbacks. Then $e$ restricts to an isomorphism $C\simeq_K C'$.
The canonical multiplication map $M\otimes_C D\to N$ is bijective
if and only if the canonical multiplication map
$M'\otimes_{C'}D'\to N'$ is bijective.
\end{lemma}
\begin{proof}
If $c\in C$, its image in $N$ belongs to $D$. Applying the commuting
square shows that the image of $e(c)$ in $N'$ belongs to $f(D)$,
so $e(c)\in C'$. Applying the inverse square proves the converse,
hence gives the asserted restriction $e_C:C\simeq C'$.
The restriction $f_D:D\simeq D'$ is the actual original map on elements.
Both $e$ and $f_D$ are semilinear over the same scalar isomorphism
$e_C$: for $c\in C$ and an element $z$, multiplicativity and the
commuting square give $e(cz)=e_C(c)e(z)$ and
$f_D(cz)=e_C(c)f_D(z)$ in the corresponding original fields.
Their tensor product therefore gives a semilinear additive equivalence
\[
 U:M\otimes_C D\longrightarrow M'\otimes_{C'}D',
 \qquad m\otimes d\longmapsto e(m)\otimes f_D(d).
\]
Indeed, the displayed scalar identities respect the balancing relation;
additivity respects the additive tensor relations. The maps obtained
from the two inverse isomorphisms give an inverse, as their composites
fix every pure tensor, and pure tensors additively generate both
products. If $\mu,\mu'$ denote the original multiplication maps,
then on pure tensors the commuting square gives
\[
 \mu'(U(m\otimes d))=e(m)f_D(d)=f(md)=f(\mu(m\otimes d)).
\]
Both sides are additive, so $\mu'\circ U=f\circ\mu$ on every tensor.
As $U$ and $f$ are bijections, this identity makes bijectivity of
$\mu$ and $\mu'$ equivalent. No normality, valuation or APF hypothesis
is used.
\end{proof}

\begin{lemma}[Extension from any cofinal set of original norm coordinates]
\label{fw:cofinal-norm-family-extension}
\uses{fw:norm-limit-multiplicative}
\lean{FontaineWintenberger.NormLimit.ext_cofinal, FontaineWintenberger.existsUnique_normLimit_of_compatible_cofinal, FontaineWintenberger.normLimitOfCofinal, FontaineWintenberger.normLimitEval_normLimitOfCofinal, FontaineWintenberger.normLimitEval_normLimitOfCofinal_any}
\leanok
Let $L/K$ be any field extension, and let $\mathcal S$ be a cofinal
set of finite $K$-subfields of $L$. Suppose $a_E\in E$ is prescribed
for each $E\in\mathcal S$ and
\[
 N_{F/E}(a_F)=a_E
 \qquad(E,F\in\mathcal S,\ E\subseteq F).
\]
There is a unique element $a$ of the original multiplicative norm
limit with $a_E$ as its coordinate at every $E\in\mathcal S$.
For every original finite field $D/K$, every $F\in\mathcal S$ with
$D\subseteq F$ computes its coordinate by
\[
 a_D=N_{F/D}(a_F).
\]
Thus cofinal original coordinates also distinguish elements of this
same norm limit.
\end{lemma}
\begin{proof}
For any finite $D/K$, choose $F\in\mathcal S$ containing $D$ and set
$b_D=N_{F/D}(a_F)$. If $F'$ is another choice, cofinality supplies
$G\in\mathcal S$ containing the finite compositum $FF'$.
Compatibility on $\mathcal S$ and actual norm transitivity give
\[
 N_{F/D}(a_F)=N_{G/D}(a_G)=N_{F'/D}(a_{F'}).
\]
The choice is immaterial. For $D\subseteq D'$ choose such a $G$
containing both chosen upper fields. The same transitivity proves
$N_{D'/D}(b_{D'})=b_D$. The original compatible-family
description of the norm limit therefore gives a unique element
with all coordinates $b_D$: the zero test for finite norms makes
these coordinates either all zero or all nonzero, so this is
precisely an element of the original zero-adjoined unit limit.
For $D\in\mathcal S$, compatibility gives $b_D=a_D$.
Finally any original norm family with the prescribed coordinates
has $D$-coordinate $N_{F/D}(a_F)$; joint injectivity of the original
evaluations proves uniqueness and the last assertion.
\end{proof}

\begin{lemma}[Naturality of the original finite coordinate maps]
\label{fw:finite-functor-naturality}
\source{wintenberger:page-0014}
\dcref{The coordinate construction in paragraph 3.1.1, with its isomorphism naturality made explicit}
\uses{fw:cofinal-finite-tensor-descent,fw:good-finite-pair-norm-basechange,fw:cofinal-norm-family-extension,fw:good-finite-field-isomorphism,fw:norm-field-isomorphism}
\lean{FontaineWintenberger.finiteNormLimitEmbeddingMap_algEquiv, FontaineWintenberger.finiteNormLimitMap_naturality, FontaineWintenberger.finiteNormLimitEmbeddingMap_naturality}
\leanok
For a finite separable embedding $i:M\to N$ over $K$, where $M/K$
is algebraic, denote by $\Phi_i$ the map on the original norm
families prescribed by
\[
 (\Phi_i\alpha)_D=i(\alpha_{i^{-1}D})
 \qquad\text{at every original good finite field }D\subseteq N.
\]
Let $i':M'\to N'$ be another such embedding and let
$e:M\simeq_K M'$ and $f:N\simeq_K N'$ satisfy $f\circ i=i'\circ e$.
Writing $e_*,f_*$ for the original isomorphism maps on norm families,
one has
\[
 f_*\circ\Phi_i=\Phi_{i'}\circ e_*.
\]
For an isomorphism $e:M\simeq_K M'$, the finite-embedding map
$\Phi_e$ is exactly the original isomorphism map $e_*$.
\end{lemma}
\begin{proof}
The good-pair norm identity makes the displayed coordinates compatible
on all good finite pairs. Good fields are cofinal by finite descent,
so the cofinal coordinate extension lemma constructs a unique original
norm family with these coordinates. This establishes the meaning of
$\Phi_i$ without using any functor law.

First suppose $i,i'$ are the original inclusions, and let $D$ be good
for $N/M$. Put $D'=f(D)$ and $C=i^{-1}D$.
The preceding good-field transport lemma says that $D'$ is good for
$N'/M'$ and its original pullback is exactly $e(C)$.
The original isomorphism maps therefore give, at $D'$, respectively
\[
 (f_*\Phi_i\alpha)_{D'}=f(i(\alpha_C)),\qquad
 (\Phi_{i'}e_*\alpha)_{D'}=i'(e(\alpha_C)).
\]
The commuting square equates these elements of the same original
field $D'$. The fields $f(D)$ are cofinal: pull any prescribed finite
field back through $f$, take a good upper field, and apply $f$ again.
Thus equality on these coordinates gives the stated equality of norm
families. For general $i,i'$, apply this calculation to their actual
image inclusions. The restrictions of $f$ to the actual images give
the coefficient-field isomorphism, and the commuting square identifies
their preimages in $M,M'$ through $e$. Hence the same two displayed
formulas hold, proving the assertion for the original embeddings.

Finally, if $i=e$ is an isomorphism, its coordinate at every good
$D\subseteq M'$ is $e(\alpha_{e^{-1}D})$. This is the defining
coordinate of $e_*\alpha$. Cofinality of the good fields equates
$\Phi_e$ with $e_*$ on the same original norm families.
\end{proof}

\begin{lemma}[The actual covariant map on norm families]
\label{fw:covariant-norm-family-map}
\source{wintenberger:page-0014}
\dcref{Construction in paragraph 3.1.1, including functor laws}
\uses{fw:cofinal-finite-tensor-descent,fw:good-finite-pair-norm-basechange,fw:norm-limit-multiplicative,fw:cofinal-norm-family-extension,fw:norm-field-isomorphism,fw:finite-functor-naturality}
\lean{FontaineWintenberger.finiteNormLimitEmbeddingMap, FontaineWintenberger.coe_normLimitEval_finiteNormLimitEmbeddingMap, FontaineWintenberger.finiteNormLimitEmbeddingMap_refl, FontaineWintenberger.finiteNormLimitEmbeddingMap_trans, FontaineWintenberger.finiteNormLimitEmbeddingMap_unique}
\leanok
For a finite separable $K$-embedding $\tau:L\to L'$ there is a unique
zero-, one- and multiplication-preserving map on the original norm
families with coordinate $\tau(\alpha_{\tau^{-1}D})$ at every good
finite $D\subseteq L'$. These maps preserve identities and composition.
\end{lemma}
\begin{proof}
Identify $L$ with its actual image $M=\tau(L)$ for the finite-field
calculation, retaining $\tau$ in the coordinate formulas.
For good $D$, prescribe $b_D=\tau(\alpha_{\tau^{-1}D})$.
The preceding original norm identity proves compatibility on all good
pairs. For arbitrary finite $A\subseteq L'$, choose good $D\supseteq A$
and set $b_A=N_{D/A}(b_D)$.
Any two choices admit a common good upper field; norm transitivity and
the proved good-pair compatibility give the same $b_A$.
The same argument proves every original norm compatibility, hence gives
an element of the original norm limit. It also proves uniqueness, by
norming from a good upper field. Norms preserve zero, one and products,
so the resulting map preserves these operations.
For the identity, the preceding isomorphism comparison identifies the
map with the original identity reindexing.
For composable finite embeddings $\tau:L\to M$ and $\sigma:M\to N$,
consider the actual tower $\sigma(\tau L)\subseteq\sigma M\subseteq N$.
Both finite extensions remain separable under the actual image
isomorphisms. The simultaneous clause of finite descent gives cofinal
finite fields $D\subseteq N$ that are good for both inclusions and
whose pullback to $\sigma M$ is good over $\sigma(\tau L)$.
The preceding naturality lemma transports these same fields and their
original pullbacks through the image isomorphisms.
At each such $D$, the coordinate of the successive maps and that of
the composite are therefore both
\[
 \sigma\bigl(\tau(\alpha_{\tau^{-1}(\sigma^{-1}D)})\bigr).
\]
Here $\tau^{-1}(\sigma^{-1}D)=(\sigma\tau)^{-1}D$ as literal
original subfields of $L$. Cofinality and uniqueness prove composition.
\end{proof}

\begin{lemma}[Additivity and the original valuation topology]
\label{fw:finite-norm-field-embedding}
\source{wintenberger:page-0008,wintenberger:page-0013,wintenberger:page-0014}
\dcref{Expanded embedding and continuity assertion of paragraph 3.1.3.3}
\uses{fw:covariant-norm-family-map,fw:good-finite-pair-norm-basechange,fw:field-norms,fw:addition-limit}
\lean{FontaineWintenberger.finiteNormLimitEmbeddingRingHom, FontaineWintenberger.finiteNormLimitEmbeddingRingHom_apply, FontaineWintenberger.finiteNormLimitEmbeddingMap_injective, FontaineWintenberger.coe_normLimitEval_finiteNormLimitEmbeddingRingHom, FontaineWintenberger.exists_finiteNormLimitEmbeddingMap_valuation_scaling, FontaineWintenberger.uniformContinuous_finiteNormLimitEmbeddingMap}
\leanok
The actual covariant map for a finite separable $K$-embedding between
infinite APF extensions is an injective field homomorphism, continuous
for the existing stable valuation topologies.
\end{lemma}
\begin{proof}
Fix a good finite target field $D$, with source preimage $C$.
Index by the actual ordered set
$S_D=\{D'\supseteq D:D'\text{ is good}\}$.
It is nonempty, since $D$ belongs to it, and directed: take a good
upper field of the finite compositum of two prescribed fields.
Its inclusion into the finite target fields above $D$ is monotone
and cofinal. The source preimage map into the finite source fields
above $C$ is also monotone and cofinal: for any prescribed source
field $F$, choose a good target field containing both $D$ and
$\tau(F)$. Its preimage contains $F$.
Restrict both already proved norm-sum limits to this same directed
set. The preceding base-change identity identifies each target term
with the image of the corresponding source term at $C$.
The actual finite field embedding $C\to D$ is continuous for their
spectral valuations. The two limits in the Hausdorff field $D$
therefore agree, proving additivity at $D$; uniqueness on good
coordinates gives additivity of the original map.
It is unital and its source and target are the original fields, hence
it is injective.

Choose good $D$ containing both the target stable-valuation witness and
the image of the source stable-valuation witness, and again put
$C=\tau^{-1}D$. The coordinate formula and ordinary finite valuation
scaling give, for all original norm families $\alpha$,
\[
 v_{X(L')}(X(\tau)\alpha)=e(D/\tau C)\,v_{X(L)}(\alpha),
\]
including zero, where $e(D/\tau C)>0$ is the actual finite
ramification index of the integer-ring inclusion.
For the multiplicative normalized valuations $w$, the exact same
identity is
$w_{X(L')}(X(\tau)\alpha)=w_{X(L)}(\alpha)^{e(D/\tau C)}$.
The positive finite exponent is fixed. Given a target valuation
neighborhood, choose an integer exponent sufficiently small in the
multiplicative value group that its $e(D/\tau C)$-th power lies in
that neighborhood. The corresponding source valuation neighborhood
maps inside it, proving continuity at zero and hence uniform
continuity of the additive homomorphism for the existing topologies.
\end{proof}

\begin{lemma}[Fixed fields of the finite relative Galois action]
\label{fw:finite-norm-field-fixed-field}
\source{wintenberger:page-0014}
\dcref{Fixed-field assertion in paragraph 3.1.3.3}
\uses{fw:finite-norm-field-embedding,fw:covariant-norm-family-map,fw:good-finite-pair-norm-basechange,fw:profinite-galois-foundations,fw:cofinal-finite-aut-stable-layers,fw:norm-field-isomorphism}
\lean{FontaineWintenberger.galoisStableGoodFiniteLayers_cofinal, FontaineWintenberger.galoisNormDescent, FontaineWintenberger.finiteNormLimitMap_galoisNormDescent, FontaineWintenberger.mem_range_finiteNormLimitMap_iff_fixed, FontaineWintenberger.finiteNormLimitGaloisAction, FontaineWintenberger.finiteNormLimitEmbeddingMap_algEquiv}
\leanok
For finite Galois $L'/L$ with group $G$, the actual covariant maps give
an action of $G$ on $X_K(L')$, whose fixed field is precisely the image
of the actual embedding $X_K(L)\to X_K(L')$.
\end{lemma}
\begin{proof}
Restricting scalars gives a homomorphism
$G\to\operatorname{Aut}_K(L')$, and the proved original norm-field
automorphisms give the actual action. Choose a primitive element for
$L'/L$ and a finite $K$-field $D_0\subseteq L'$ containing it and
any prescribed finite set. Replace $D_0$ by the actual compositum of
all its $K$-embedding images in $L'$. The proved finite stable-field
lemma makes this field finite and stable under every original
$K$-automorphism, without assuming $L'/K$ normal.
Its $G$-action is faithful: an element fixing this field fixes the
primitive element and $L$, hence fixes $L'=L(\alpha)$.
Its fixed field is $C_0=L\cap D_0$, and finite Artin theory gives
$[D_0:C_0]=|G|=[L':L]$. Since $LD_0=L'$, the canonical tensor map
is an isomorphism by equality of dimensions, so $D_0$ is good.
This constructs a cofinal $G$-stable good subsystem.

For an original source family $\alpha$, the actual image at each
stable good $D$ is the original inclusion of $\alpha_C$, with
$C=L\cap D$. This lies in $L$ and is fixed by $G$. The original
coordinate covariance and cofinality therefore show that the entire
image family is fixed.

Conversely, if an original family $\beta$ is fixed by $G$, its
coordinate at each such $D$ is fixed by every $L$-automorphism of
$L'$. The ordinary finite Galois fixed-field theorem puts that actual
element in $L$, hence in the original pullback $C=L\cap D$.
Write $a_D\in C$ for its unique preimage. These $C$ are cofinal in
$L$: require $D$ to contain the image of any prescribed finite source
field. For stable good fields $D,D'$, choose a common stable good
upper field $F$ and put $B=L\cap F$. Actual good-pair norm base
change and injectivity of the original inclusions give
$a_D=N_{B/C}(a_F)$ and $a_{D'}=N_{B/C'}(a_F)$.
If $C=C'$, these identities give the same element of $C$, independently
of $D$. If $C\subseteq C'$, norm transitivity gives
$N_{C'/C}(a_{D'})=N_{B/C}(a_F)=a_D$.
Thus the original coordinates define a well-defined compatible family
on the actual cofinal source fields.
Extending by actual norms to arbitrary finite source fields gives an
original family $\alpha\in X_K(L)$. Its actual image agrees with
$\beta$ on the cofinal subsystem, hence equals $\beta$.
This proves the reverse fixed-field inclusion.
\end{proof}

\begin{lemma}[Faithfulness for a finite relative Galois extension]
\label{fw:finite-norm-field-faithful}
\source{wintenberger:page-0014}
\dcref{Lemma 3.1.3.1 and its proof in paragraph 3.1.3.2}
\uses{fw:faithful-action,fw:profinite-galois-foundations}
\lean{FontaineWintenberger.finiteNormLimitGaloisAction_injective}
\leanok
For finite Galois $L'/L$ the actual action of its finite Galois group
on $X_K(L')$ is faithful.
\end{lemma}
\begin{proof}
Restriction of scalars gives an injective homomorphism
$\operatorname{Aut}_L(L')\to\operatorname{Aut}_K(L')$:
two relative automorphisms with the same underlying field map agree.
The actual relative action on $X_K(L')$ is the composition of this
homomorphism with the proved faithful original action of
$\operatorname{Aut}_K(L')$. Its action on each original finite
coordinate is the same coordinate covariance formula; no new action
or field is introduced. The composite of two injective homomorphisms
is injective. This reuses the proved separation by actual source
constants and an original uniformizer in the general faithfulness
lemma, so supplies the same source argument without repeating it.
\end{proof}

\begin{lemma}[Finite degree and separability by normal closure]
\label{fw:finite-norm-field-degree-separable}
\source{wintenberger:page-0014,wintenberger:page-0015}
\dcref{Completion of the proof of Theorem 3.1.2}
\uses{fw:finite-norm-field-fixed-field,fw:finite-norm-field-faithful,fw:finite-norm-field-embedding,fw:apf-finite-extension}
\lean{FontaineWintenberger.finiteNormLimit_isGaloisGroup, FontaineWintenberger.finiteDimensional_finiteNormLimit_of_isGalois, FontaineWintenberger.isGalois_finiteNormLimit, FontaineWintenberger.finrank_finiteNormLimit_of_isGalois, FontaineWintenberger.finiteNormLimitGaloisEquiv, FontaineWintenberger.finiteNormLimitGaloisEquiv_apply, FontaineWintenberger.finiteNormLimitRingHom_isScalarTower, FontaineWintenberger.finiteNormLimit_finite_separable_degree, FontaineWintenberger.finiteNormLimitEmbedding_finite_separable_degree}
\leanok
The actual embedding attached to a finite separable $L'/L$ has degree
$[L':L]$ and is separable. If $L'/L$ is Galois, it is Galois and the
actual covariant action identifies the two Galois groups.
\end{lemma}
\begin{proof}
In the Galois case, the action is faithful and its fixed field is the
actual image of $X_K(L)$. Finite Artin theory gives a finite Galois
extension of that fixed field, degree equal to $|\operatorname{Gal}(L'/L)|$,
and identifies the actual acting group with its Galois group.
This is $[L':L]$.
For general finite separable $L'/L$, form the normal closure $N$ of
$L'/L$ inside an algebraic closure of $L'$. The original maps give
an actual tower $L\subseteq L'\subseteq N$; the normal closure is
finite Galois over $L$, and hence also finite Galois over $L'$.
It is separable over $K$, infinite over $K$ because it contains $L$,
and APF by finite inheritance. Apply the Galois result to $N/L$ and
$N/L'$.

Composition of the actual norm-field inclusion maps gives the scalar
tower $X_K(L)\subseteq X_K(L')\subseteq X_K(N)$ on the original
norm fields. The injection of $X_K(L')$ into the finite-dimensional
$X_K(L)$-vector space $X_K(N)$ is linear for these scalar structures,
so $X_K(L')/X_K(L)$ is finite. The two degree tower formulas give
\[
 [X_K(L'):X_K(L)] [N:L']
 = [X_K(N):X_K(L)]
 = [N:L]
 = [L':L][N:L'].
\]
Since $[N:L']$ is a positive finite integer, cancellation proves the
claimed degree. Separability descends along this same scalar tower
from the Galois extension $X_K(N)/X_K(L)$.
For an arbitrary finite separable embedding, factor it through its
actual image and the original source-to-image isomorphism. The
proved norm-field isomorphism and actual range-factorization of the
embedding identify its degree and separability with this inclusion
case.
\end{proof}

\begin{lemma}[Isomorphisms of the original norm fields]
\label{fw:norm-field-isomorphism}
\dcref{Wintenberger1983, 3.1.1, the isomorphism case; multiplicative reindexing generalized to arbitrary field extensions}
\source{wintenberger:page-0014}
\uses{fw:norm-limit-multiplicative,fw:addition-limit,fw:norm-limit-stable-valuation,fw:norm-limit-field,fw:norm-limit-field-valuation,fw:ramification-field-equiv}
\lean{FontaineWintenberger.FiniteLayer.equivMap_norm, FontaineWintenberger.normLimitMap, FontaineWintenberger.normLimitMulEquiv, FontaineWintenberger.normLimitEval_normLimitMap, FontaineWintenberger.normLimitMap_refl, FontaineWintenberger.normLimitMap_trans, FontaineWintenberger.normLimitMap_add, FontaineWintenberger.normLimitValuation_normLimitMap, FontaineWintenberger.normLimitRingEquiv, FontaineWintenberger.normLimitUniformEquiv, FontaineWintenberger.normLimitAut, FontaineWintenberger.normLimitAut_valuation}
\leanok
For arbitrary field extensions $L/K$ and $M/K$ and a $K$-isomorphism
$e:L\simeq_K M$, the original multiplicative norm limits admit a unique
multiplicative equivalence $X_K(e)$ such that
\[
 (X_K(e)(x))_F=e(x_{e^{-1}(F)})
\]
for every actual finite $F/K$ in $M$. These multiplicative equivalences
preserve identity and composition without APF or valuation hypotheses.
Under the standing complete discrete-valuation hypotheses, if both
extensions are separable algebraic APF, they also preserve the original
norm-limit additions and normalized valuations. For infinite extensions
they are therefore isomorphisms of the original valued norm fields.
In particular they define
an action of $\operatorname{Aut}_K(L)$ by valuation-preserving
automorphisms of the original norm field.
\end{lemma}
\begin{proof}
The assignments $E\mapsto e(E)$ and $F\mapsto e^{-1}(F)$ are inverse
order isomorphisms of the actual finite intermediate fields. Their
restrictions $e_E:E\simeq_K e(E)$ commute with every actual inclusion.
For $E\subseteq E'$ and $a\in E'$, transport a basis of $E'/E$
through $e$. The multiplication matrix of $e(a)$ in the transported
basis is obtained by applying $e_E$ to every entry of the multiplication
matrix of $a$. Taking determinants gives
\[
 e_E(N_{E'/E}(a))=N_{e(E')/e(E)}(e_{E'}(a)).
\]
Consequently the displayed coordinate formula takes each original
norm-compatible family to an original norm-compatible family. It
preserves multiplication, one and zero; the image of a family of units
is again a family of units. The evaluations jointly distinguish norm
families, which proves uniqueness.

At a finite target field $F$, applying first $e$ and then another
isomorphism $f:M\simeq_K N$ gives the coordinate
$f(e(x_{e^{-1}(f^{-1}(F))}))$. This is exactly the coordinate for
$f\circ e$. The identity gives the original coordinate at every layer.
Joint injectivity therefore gives the identity and composition laws,
and in particular $X_K(e^{-1})=X_K(e)^{-1}$.

Fix a finite $F/K$ in $M$ and put $E=e^{-1}(F)$.
For each finite $F'\supseteq F$, set $E'=e^{-1}(F')$.
The two original norm-sum nets satisfy the exact identity
\[
 N_{F'/F}\bigl((X_K(e)x)_{F'}+(X_K(e)y)_{F'}\bigr)
   =e_E\bigl(N_{E'/E}(x_{E'}+y_{E'})\bigr).
\]
The finite-layer order isomorphism takes the upper-tail filter to the
upper-tail filter: a tail above $E_0\supseteq E$ is taken to the tail
above $e(E_0)\supseteq F$, and the inverse has the same property.
The map $e_E$ is continuous for the canonical spectral norms, since
it is a $K$-linear map between finite-dimensional normed spaces over
the complete field $K$. Thus it takes the defining source sum-net
limit to the target sum-net limit. Uniqueness of these limits in
each original finite field, and then joint injectivity of evaluations,
prove that $X_K(e)$ preserves the original norm-limit addition.

The uniqueness of the extension of the base norm identifies the
spectral integer rings under $e_E$. The induced ring isomorphism
identifies their units and maximal ideals, and sends a uniformizer
$\pi$ to a uniformizer. For a nonzero $a\in E$, write
$a=\pi^m u$ with $m\in\mathbf Z$ and $u$ an integer-ring unit.
Its image is $e_E(\pi)^m e_E(u)$, with the same integer $m$ and a
unit $e_E(u)$. Hence their normalized valuations agree; at zero
both valuations are zero in multiplicative notation.

Choose the proved source and target residue-stable finite tails.
Take a finite target layer $F$ containing the target starting field
and the image of the source starting field. Then $e^{-1}(F)$ is on
the source tail. Both original norm-limit valuations are computed
at these respective coordinates, so the finite valuation equality
and the displayed coordinate formula give
\[
 w_M(X_K(e)x)=w_L(x).
\]
For infinite extensions the additive and multiplicative structures
are the original norm-field structures already constructed.
The isomorphism and its inverse preserve the valuation of every
difference, so they take each original valuation neighborhood and
uniform entourage to the corresponding one. This proves continuity
and uniform continuity for the existing valuation topologies.
Finally the identity and composition laws, with $M=L$, give the
asserted automorphism action.
\end{proof}

\begin{lemma}[Continuity of the original automorphism action]
\label{fw:norm-field-action-continuity}
\uses{fw:norm-field-isomorphism,fw:norm-limit-truncation-equivalence,fw:norm-limit-valuation-complete,fw:norm-limit-uniformizer}
\lean{FontaineWintenberger.normLimitEval_normLimitMap_of_mem_fixingSubgroup, FontaineWintenberger.normLimitValuation_sub_le_of_mem_deep_fixingSubgroup, FontaineWintenberger.continuous_normLimitAut_apply_of_integral, FontaineWintenberger.continuous_normLimitAut_apply}
\leanok
\source{false-tate-lift:page-0012,wintenberger:page-0014}
Let $L/K$ be an infinite separable APF extension under the standing
complete discrete-valuation hypotheses with perfect residue field.
Give $\operatorname{Aut}_K(L)$ its original Krull topology and
$X=X_K(L)$ its original valuation topology. The original action has
continuous orbit maps $g\mapsto X(g)(x)$ for every $x\in X$.
More precisely, fix the original deep tower $(E_n)$ starting at $K$
with quotient precision $n+2$. Its finite-field fixer
$H_n=\operatorname{Aut}_{E_n}(L)$ is open, and for every $g\in H_n$
and every original integral family $x$ one has
\[
 w\bigl(X(g)(x)-x\bigr)\leq \exp(-(n+2)).
\]
Here $w$ is the original multiplicative normalized valuation.
\end{lemma}
\begin{proof}
A finite field $E/K$ has finitely many algebraic generators. The condition
that an automorphism fix them is one of the defining open conditions
of the Krull topology, so its pointwise fixer is open. This does not
require $E/K$ or $L/K$ to be normal.
If $g$ fixes $E$ pointwise, then $g^{-1}(E)=E$, and the original
coordinate formula gives
\[
 (X(g)x)_E=g(x_{g^{-1}(E)})=g(x_E)=x_E.
\]
Thus $g\in H_n$ fixes the $E_n$ coordinate of every original family.
Let $q_n$ be the already constructed ring homomorphism from the original
norm-field valuation ring to
$\mathcal O_{E_n}/\mathfrak m_{E_n}^{n+2}$.
Valuation preservation makes $X(g)x$ integral whenever $x$ is integral.
Their equal original coordinates give $q_n(X(g)x)=q_n(x)$.
Additivity of this same ring homomorphism gives
$q_n(X(g)x-x)=0$, and its exact valuation-kernel identity gives the
asserted bound, uniformly over all integral $x$.

Fix an integral $x$ and an automorphism $g_0$. For $g_0^{-1}g\in H_n$,
the action law and valuation preservation give
\[
 w\bigl(X(g)x-X(g_0)x\bigr)
 =w\bigl(X(g_0)(X(g_0^{-1}g)x-x)\bigr)
 \leq\exp(-(n+2)).
\]
The set $g_0H_n$ is an open neighborhood of $g_0$. The displayed
integer-radius closed balls form a neighborhood basis at zero in the
original norm-field topology, so this proves continuity at every $g_0$
for every integral $x$. Finally every element of $X$ is $a/b$ with
$a,b$ integral and $b\ne0$: multiply numerator and denominator by a
sufficiently large nonnegative power of an original uniformizer.
The action consists of field automorphisms, so its orbit at $a/b$
is the quotient of its continuous orbits at $a$ and $b$.
The denominator never vanishes, and division is continuous in the same
valued-field topology. This proves all the asserted orbit continuities.
\end{proof}

\begin{theorem}[Finite covariant functoriality]
\label{fw:finite-functor}
\dcref{Wintenberger1983, 3.1.1 and Theorem 3.1.2}
\source{wintenberger:page-0014,wintenberger:page-0015}
\uses{fw:covariant-norm-family-map,fw:finite-norm-field-embedding,fw:finite-norm-field-degree-separable,fw:norm-field-isomorphism}
\lean{FontaineWintenberger.finiteNormLimitEmbeddingRingHom, FontaineWintenberger.coe_normLimitEval_finiteNormLimitEmbeddingRingHom, FontaineWintenberger.continuous_finiteNormLimitEmbeddingMap, FontaineWintenberger.finiteNormLimitEmbeddingMap_refl, FontaineWintenberger.finiteNormLimitEmbeddingMap_trans, FontaineWintenberger.finiteNormLimitEmbedding_finite_separable_degree, FontaineWintenberger.isGalois_finiteNormLimitEmbedding, FontaineWintenberger.finiteNormLimitEmbeddingGaloisEquiv, FontaineWintenberger.finiteNormLimitEmbeddingGaloisEquiv_apply, FontaineWintenberger.finiteNormLimitEmbeddingMap_algEquiv}
\leanok
Let $L/K$ and $L'/K$ be infinite APF and let
$\tau:L\longrightarrow L'$ be a $K$-embedding with $L'/\tau(L)$ finite
and separable. There is a continuous finite separable embedding
\[
 X_K(\tau):X_K(L)\longrightarrow X_K(L')
\]
characterized by
\[
 (X_K(\tau)(\alpha))_{E'}=\tau(\alpha_{\tau^{-1}(E')})
\]
on the cofinal set of finite $E'/K$ in $L'$ for which
$\tau(L)\otimes_{\tau(L)\cap E'}E'\longrightarrow L'$ is an isomorphism.
These maps preserve identities and composition, and
\[
 [X_K(L'):X_K(\tau)(X_K(L))]=[L':\tau(L)].
\]
If $L'/\tau(L)$ is Galois, so is
$X_K(L')/X_K(\tau)(X_K(L))$, and $\sigma\mapsto X_K(\sigma)$ identifies
their Galois groups.
\end{theorem}

\begin{proof}
Factor the original embedding as the actual $K$-isomorphism
$L\simeq\tau(L)$ followed by the actual inclusion $\tau(L)\subseteq L'$.
The norm-field isomorphism theorem handles the first factor and the
inclusion results above handle the second. Their composite agrees with
the actual covariant norm-family construction on every good coordinate,
so its unique characterization gives the displayed formula as well as
the identity and composition laws.
The embedding lemma proves additivity, injectivity and continuity
for the original fields and their existing valuation topologies.
The normal-closure lemma gives the degree and separability assertions
and, in the Galois case, exactly the stated actual Galois-group
identification. All constructions use the same original norm families.
\end{proof}

\begin{lemma}[Discriminants, integral roots, and a resultant identity]
\label{fw:separable-polynomial-discriminant}
\dcref{Algebraic details used in Wintenberger1983, paragraphs 3.2.5.1--3.2.5.5}
\source{wintenberger:page-0015,wintenberger:page-0016}
\lean{FontaineWintenberger.monic_discr_resultant, FontaineWintenberger.monic_discr_map, FontaineWintenberger.exists_mul_add_mul_eq_discr, FontaineWintenberger.monic_derivative_resultant, FontaineWintenberger.monic_discr_ne_zero_iff, FontaineWintenberger.monic_root_valuation_le_one, FontaineWintenberger.valuation_discr_le_max, FontaineWintenberger.valuation_discr_le_derivative, FontaineWintenberger.exists_monic_two_root_factor, FontaineWintenberger.sq_sub_dvd_discr, FontaineWintenberger.valuation_discr_le_sub_sq, FontaineWintenberger.additiveValuation_two_nsmul_sub_le_discr, FontaineWintenberger.valuation_mapped_discr_le_sub_sq, FontaineWintenberger.additiveValuation_discr_nonneg, FontaineWintenberger.additiveValuation_discr_ne_top, FontaineWintenberger.valuation_discr_le_sub_sq_of_extension, FontaineWintenberger.additiveValuation_two_nsmul_sub_le_discr_of_extension, FontaineWintenberger.monicDiscriminantPolynomial, FontaineWintenberger.monicDiscriminantPolynomial_eval}
\leanok
Let $C$ be a discretely valued field and let
$f(T)=T^k+\sum_{i<k}a_iT^i\in\mathcal O_C[T]$ be separable, with
$k\geq1$. Write $\Delta(f)$ for its discriminant and
$D=v_C(\Delta(f))$. Then $D$ is finite and nonnegative.
Every root is integral, and, for any two distinct roots $a,a'$ in a
valued splitting field, with the valuation extended from $C$,
\[
 2v(a-a')\leq D.
\]
There are $A,B\in\mathcal O_C[T]$ such that
\[
 Af+Bf'=(-1)^{k(k-1)/2}\Delta(f).
\]
The discriminant of a monic degree-$k$ polynomial is a polynomial with
integer coefficients in its $k$ coefficients. In particular it
commutes with every unital ring homomorphism on the coefficient ring.
For any monic polynomial of positive degree, its discriminant is
nonzero if and only if the polynomial is separable.
\end{lemma}
\begin{proof}
If a root $a$ had negative valuation, $a^k$ would have strictly smaller
valuation than each $a_i a^i$, contradicting $f(a)=0$ and the
ultrametric inequality. Thus all roots are integral. In a splitting
field write $f=\prod_{j=1}^k(T-a_j)$. Separability makes the $a_j$
distinct, and
\[
 \Delta(f)=\prod_{i<j}(a_i-a_j)^2.
\]
Every factor has nonnegative valuation, and none is zero. Taking
valuations proves both assertions concerning $D$ and each pair of roots.
The same product formula, with roots counted with multiplicity, makes
the discriminant zero exactly when two roots coincide. This proves
the stated criterion for an arbitrary monic polynomial.

Equivalently, the bound for a particular pair can be obtained directly
from integral resultants, without choosing the other roots. In the
valuation ring of the valued splitting field, monic division gives
$f=(T-a)(T-a')g$ with $g$ monic and integral: after dividing by $T-a$,
evaluation at $a'$ and $a'\ne a$ show that the quotient vanishes at $a'$.
Multiplicativity of the resultant in its first argument gives
\[
 \operatorname{Res}(f,f')
 =f'(a)f'(a')\operatorname{Res}(g,f')
 =-(a-a')^2g(a)g(a')\operatorname{Res}(g,f').
\]
All three remaining factors are integral, and the signed
discriminant is this resultant. Taking valuations gives exactly the
same pair bound. This argument also proves the multiplicative
valuation inequality when the discriminant is zero.

For the integral identity use the Sylvester matrix of $f$ and $f'$,
allowing the leading coefficients of $f'$ to vanish. It is the matrix
of
\[
 (A,B)\longmapsto Af+Bf'
\]
from polynomials of respective degree less than $k-1$ and $k$ to
polynomials of degree less than $2k-1$, in their monomial bases.
Its determinant is the resultant of the monic $f$ and $f'$, namely
\[
 \prod_j f'(a_j)=(-1)^{k(k-1)/2}\Delta(f).
\]
The product formula follows by evaluating multiplication by $f'$ on
$C[T]/(f)$: after passing to the splitting field, evaluation at the
$k$ distinct roots diagonalizes that map, with diagonal entries
$f'(a_j)$. Polynomial division by the monic $f$ identifies its
determinant with the Sylvester determinant. Apply the adjugate matrix
to the vector of the constant polynomial $1$. Multiplication by the
Sylvester matrix gives its determinant times that vector, which is
exactly the asserted identity with integral coefficients. For $k=1$
it reads $0f+1f'=1$. Finally the Sylvester determinant is an integer
polynomial in the coefficients, so the same holds for the discriminant
and proves its compatibility with coefficient homomorphisms.
\end{proof}

\begin{lemma}[Newton lifting with a nonunit derivative]
\label{fw:separable-newton-lifting}
\uses{fw:adic-complete-unit-ball}
\lean{FontaineWintenberger.unitBall_polynomial_norm_sub_le, FontaineWintenberger.unitBall_polynomial_norm_sub_eq_of_norm_lt_derivative, FontaineWintenberger.unitBall_hensel_nonunit, FontaineWintenberger.discreteValued_hensel_nonunit, FontaineWintenberger.discreteValued_hensel_nonunit_additive, FontaineWintenberger.discreteValued_hensel_nonunit_ambient, FontaineWintenberger.discreteValued_hensel_nonunit_ambient_additive}
\leanok
Let $C$ be a complete discretely valued field,
$g\in\mathcal O_C[T]$, and $z\in\mathcal O_C$. Suppose
$d=v(g'(z))<\infty$ and $v(g(z))>2d$. There is a unique root $a$ of
$g$ satisfying $v(a-z)>d$. If $g(z)\ne0$, then
\[
 v(a-z)=v(g(z))-d.
\]
\end{lemma}
\begin{proof}
For integral $x,h$, expansion of each monomial gives
\[
 g(x+h)=g(x)+g'(x)h+h^2R,
 \qquad g'(x+h)-g'(x)=hS
\]
with $R,S$ integral. If $g(z)=0$, take $a=z$; uniqueness follows
from the last argument below. Otherwise start at $z_0=z$ and, until
an exact root is reached, set $z_{n+1}=z_n-g(z_n)/g'(z_n)$.
Write $m_n=v(g(z_n))$. Inductively the Taylor identities give
\[
 v(g'(z_n))=d,\qquad
 v(z_{n+1}-z_n)=m_n-d>d,\qquad
 m_{n+1}\geq2(m_n-d).
\]
Indeed the derivative changes by an element of valuation greater than
$d$, and the linear term cancels in the first identity. In particular
$m_n-2d\geq2^n(m_0-2d)$. The increments tend to zero with strictly
increasing valuations, so the sequence is Cauchy and integral.
Completeness gives a limit $a$, polynomial continuity gives $g(a)=0$,
and the first increment has strictly smaller valuation than every
later increment. Thus $v(a-z)=m_0-d$. If the iteration reaches a root,
the identical argument applies to the finite sum of its increments.

For uniqueness, if $x,y$ both lie in $v(T-z)>d$, the identity
\[
 g(x)-g(y)=(x-y)\bigl(g'(z)+R_{x,y}\bigr)
\]
has $v(R_{x,y})>d$: expand each divided difference about $z$;
every term beyond the constant derivative contains $x-z$ or $y-z$
and otherwise only integral factors. The parenthesized factor is
nonzero of valuation $d$. Consequently two roots in this ball are
equal. This also proves uniqueness when $z$ itself is a root.
\end{proof}

\begin{lemma}[Extension of separable embeddings]
\label{fw:separable-embedding-extension}
\lean{FontaineWintenberger.exists_aut_extends_intermediate_embedding, FontaineWintenberger.exists_aut_extends_separableClosure_embedding, FontaineWintenberger.nonempty_separablyClosed_embeddings}
\leanok
Let $P$ be a field. Any $P$-embedding of a subextension of a separable
closure $\Omega/P$ into $\Omega$ extends to a $P$-automorphism of
$\Omega$. More generally, every separable algebraic extension of $P$
embeds over $P$ into any separably closed separable algebraic extension
of $P$.
\end{lemma}
\begin{proof}
Order extensions of the given embedding by inclusion of their domains.
The union along a chain is a field embedding, so a maximal extension
exists. If its domain omits an element of the separable algebraic
source, the minimal polynomial of that element over the domain is
separable. Its transported polynomial has a root in the separably
closed target. Sending the omitted element to that root extends the
embedding to the simple extension, contradicting maximality. Hence
the domain is the whole source. This proves the second assertion and
the existence part of the first. In the first case the resulting
endomorphism of $\Omega$ is onto: for each separable polynomial over
$P$, it maps its finite set of roots in $\Omega$ injectively to itself,
hence permutes that set. Every element of $\Omega$ belongs to such a
root set, so every element is in its image.
\end{proof}

\begin{lemma}[Krasner descent and the unique nearby embedding]
\label{fw:separable-krasner-descent}
\uses{fw:separable-embedding-extension,fw:spectral-norm-extension}
\dcref{Krasner argument used in Wintenberger1983, paragraph 3.2.5.8}
\source{wintenberger:page-0017}
\lean{FontaineWintenberger.spectralIntegerRing_eq_separable_extension_valuationSubring, FontaineWintenberger.discreteValued_krasner_mem_adjoin_and_unique, FontaineWintenberger.discreteValued_krasner_mem_adjoin_and_unique_additive}
\leanok
Let $C$ be complete and discretely valued and let $a,b$ belong to a
separable closure of $C$, with its extended valuation. If
\[
 v(a-b)>\max_{a'\ne a}v(a-a'),
\]
where $a'$ ranges over the other $C$-conjugates of $a$ and the
maximum over the empty set is $-\infty$, then $a\in C(b)$.
Moreover the inclusion of $C(a)$ is the unique $C$-embedding into the
separable closure whose image of $a$ has distance from $b$ satisfying
this same strict valuation bound.
\end{lemma}
\begin{proof}
Completeness makes the extension of the valuation to every finite
extension unique, so every $C$-automorphism of the separable closure
preserves it. If $\sigma$ fixes $b$, then
$v(\sigma a-b)=v(a-b)$. The ultrametric inequality therefore gives
$v(\sigma a-a)\geq v(a-b)$, which contradicts the displayed bound
unless $\sigma a=a$. If $a\notin C(b)$, its separable minimal
polynomial over $C(b)$ has a different root. The embedding sending
$a$ to that root extends, by the preceding lemma, to an automorphism
fixing $C(b)$, a contradiction. Thus $a\in C(b)$. If another
$C$-conjugate $a'$ had the asserted closeness to $b$, the ultrametric
inequality applied through $b$ would give
$v(a-a')>\max_{a''\ne a}v(a-a'')$, again a contradiction. An embedding
of $C(a)$ is determined by its value on $a$, proving uniqueness.
\end{proof}

\begin{lemma}[Extraction of an exact separable root from an approximate root]
\label{fw:separable-exact-root-extraction}
\uses{fw:separable-polynomial-discriminant,fw:separable-newton-lifting,fw:separable-krasner-descent,fw:spectral-norm-extension}
\lean{FontaineWintenberger.unitBall_root_eq_of_norm_sub_lt_discr, FontaineWintenberger.unitBall_hensel_of_discr, FontaineWintenberger.exists_root_near_of_discr, FontaineWintenberger.exists_root_mem_adjoin_of_discr, FontaineWintenberger.discreteValued_exists_root_mem_adjoin_of_discr, FontaineWintenberger.discreteValued_exists_root_mem_adjoin_of_discr_additive}
\leanok
Let $C$ be complete and discretely valued, let
$f\in\mathcal O_C[T]$ be monic separable of positive degree, and put
$D=v_C(\Delta(f))$. Give a finite separable extension $C'/C$ the
valuation extending $v_C$. For an integral $z\in C'$, if
\[
 v(f(z))>2D,
\]
then $f$ has a root $a\in C(z)\subseteq C'$. It is the unique root
with $v(a-z)>D$.
\end{lemma}
\begin{proof}
If $f(z)=0$, take $a=z$ and use root separation for uniqueness.
Otherwise evaluate the resultant identity at $z$. Since its two
polynomial coefficients have integral values and $v(f(z))>D$,
$v(f'(z))\leq D$: if both summands had valuation greater than $D$,
their sum could not have valuation $D$. In particular the derivative
is nonzero. The finite extension $C'$ is complete. Apply Newton
lifting there, multiplying all displayed valuations by its
ramification index if a normalized integer valuation is desired.
It gives a root $a\in C'$ with
\[
 v(a-z)=v(f(z))-v(f'(z))>D.
\]
Every other $C$-conjugate of $a$ is a different root of $f$, and root
separation bounds its distance from $a$ by $D/2$. Krasner descent
therefore gives $a\in C(z)$. Two roots at distance of valuation
greater than $D$ from $z$ would have mutual distance of valuation
greater than $D$, contrary to the same separation bound. This proves
both the field containment and uniqueness.
\end{proof}

\begin{lemma}[Actual coordinate projections at a prescribed high break]
\label{fw:separable-high-break-projections}
\dcref{Wintenberger1983, 2.2.3.3 and 3.2.5.6, at a prescribed finite coordinate}
\source{wintenberger:page-0011,wintenberger:page-0012,wintenberger:page-0017}
\uses{fw:absolute-upper-projections,fw:wild-ramification-p-group,fw:local-inertia-size,fw:apf-extended-bound,fw:finite-layer-norm-precision,fw:truncated-wild-norm-map,fw:deep-truncated-finite-tower,fw:truncated-norm-inverse-limit,fw:norm-limit-truncation-equivalence,fw:norm-limit-stable-valuation,fw:norm-valuation,fw:addition-limit,fw:ideal-powers-topology}
\lean{FontaineWintenberger.FiniteLayer.HasFullUpper, FontaineWintenberger.FiniteLayer.hasFullUpper_iff_embedding, FontaineWintenberger.FiniteLayer.hasFullUpper_apfPsi_iff, FontaineWintenberger.FiniteLayer.hasFullUpper_iff_stage_le, FontaineWintenberger.finrank_eq_prime_pow_of_apfFiniteStage_pos_le, FontaineWintenberger.FiniteLayer.HasFullUpper.of_le, FontaineWintenberger.FiniteLayer.HasFullUpper.residue_surjective, FontaineWintenberger.FiniteLayer.HasFullUpper.finrank_eq_prime_pow, FontaineWintenberger.FiniteLayer.HasFullUpper.le_relativeFirstBound, FontaineWintenberger.normLimitValuation_eq_coordinate_of_fullUpper, FontaineWintenberger.FiniteLayer.HasFullUpper.truncatedNorm, FontaineWintenberger.FiniteLayer.HasFullUpper.truncatedNorm_mk, FontaineWintenberger.FiniteLayer.HasFullUpper.truncatedNorm_surjective, FontaineWintenberger.normLimitIntegralCoordinateAt, FontaineWintenberger.normLimitHighBreakProjection, FontaineWintenberger.normLimitHighBreakProjection_apply, FontaineWintenberger.normLimitHighBreakProjection_surjective, FontaineWintenberger.normLimitHighBreakProjection_eq_zero_iff, FontaineWintenberger.normLimitHighBreakProjection_eq_iff, FontaineWintenberger.ker_normLimitHighBreakProjection}
\leanok
Let $V/K$ be infinite separable APF, let $E/K$ be a finite subfield,
and work in a common separable closure. Put
$G=\operatorname{Gal}(K^{\mathrm{sep}}/E)$ and
$H=\operatorname{Gal}(K^{\mathrm{sep}}/V)$. Suppose $b>0$ and
$HG^b=G$. If
\[
 1\leq r\leq\left\lceil\frac{p-1}{p}b\right\rceil,
\]
then the original coordinate map induces a surjective unital ring map
\[
 q_{E,r}:\mathcal O_{X_K(V)}\longrightarrow
                \mathcal O_E/\mathfrak m_E^r,
 \qquad a\longmapsto[a_E],
\]
whose kernel is exactly $\{a:v_{X_K(V)}(a)\geq r\}$.
For every original family $a$, including nonintegral families,
$v_{X_K(V)}(a)=v_E(a_E)$.
\end{lemma}
\begin{proof}
For every finite $F/E$ in $V$, the fixing subgroup $H_F$ contains $H$,
so $H_FG^b=G$. Projection to a finite Galois overfield shows that the
wild ramification group acts transitively on the embeddings of $F$:
the positive upper group is contained in the wild group. Thus
$F/E$ is totally ramified, by the inertia-index formula, and its
degree is a power of $p$, since an orbit of a finite $p$-group has
$p$-power size. The same product identity gives $i(F/E)\geq b$.
Consequently every such finite norm is additive and surjective modulo
$\mathfrak m_E^r$. The identity extension has these properties as well.

Coordinatewise multiplication preserves products, zero and one.
For integral $a,c$, every term
$N_{F/E}(a_F+c_F)-a_E-c_E$ in the defining norm-sum net lies in the
closed ideal $\mathfrak m_E^r$. Taking the original limit proves
$q_{E,r}(a+c)=q_{E,r}(a)+q_{E,r}(c)$. Thus $q_{E,r}$ is a unital
ring map on the original norm families.

To prove surjectivity, take a deep cofinal tower $D_n$ containing $E$,
with precisions $h_n=n+2$, and choose $m$ with $h_m\geq r$.
The actual finite norm induces a surjection
\[
 \mathcal O_{D_m}/\mathfrak m_{D_m}^{h_m}
    \longrightarrow\mathcal O_E/\mathfrak m_E^r
\]
by the truncated-norm lemma and the first-bound inequality just proved.
Choose a preimage of the prescribed target class. The inverse-limit
projection at $m$ is surjective; the actual truncation isomorphism
therefore supplies an original integral norm family $a$ having this
preimage as its $D_m$-class. Norm compatibility makes $[a_E]$ the
prescribed target class. Finally all finite extensions above $E$ have
residue degree one. For any original family $a$, the normalized
norm-valuation identity therefore makes $v_F(a_F)=v_E(a_E)$ on this
whole tail, including zero and nonintegral families. This
constant tail value is precisely the stable norm-field valuation.
Hence $q_{E,r}(a)=0$ is equivalent to $v(a)\geq r$, as asserted.
The upper-group product replaces a specified tame starting field: it
provides the wild finite-tail hypotheses needed for the cited
truncated-norm construction.
\end{proof}

\begin{lemma}[Deep finite coordinates with an actual upper-group product]
\label{fw:separable-deep-coordinate-stages}
\uses{fw:apf-finite-stages,fw:absolute-finite-base-upper,fw:apf-herbrand-homeomorphism,fw:countable-apf-finite-layers}
\lean{FontaineWintenberger.IsAPF.exists_finiteLayer_fullUpper_tail, FontaineWintenberger.IsAPF.exists_monotone_cofinal_fullUpper_stages}
\leanok
For infinite APF $L/K$ and any sequence of positive integers
$b_n\to\infty$, there is an increasing cofinal sequence of actual
finite $K$-subfields $E_n\subset L$ such that
\[
 \operatorname{Gal}(K^{\mathrm{sep}}/L)
 \operatorname{Gal}(K^{\mathrm{sep}}/E_n)^{b_n}
       =\operatorname{Gal}(K^{\mathrm{sep}}/E_n).
\]
In particular every finite step above $E_n$ in $L$ is totally
ramified of $p$-power degree.
\end{lemma}
\begin{proof}
Put $H=\operatorname{Gal}(K^{\mathrm{sep}}/L)$ and choose a cofinal
sequence $C_n$ of finite subfields. Recursively choose $u_n\geq b_n$
so large that the finite APF stage $E_n=(K^{\mathrm{sep}})^{H G_K^{u_n}}$
contains $C_n$ and $E_{n-1}$ when $n>0$. This is possible by
cofinality and monotonicity of the stages. Its actual fixing group is
$H G_K^{u_n}$. Finite-base upper comparison identifies
$G_K^{u_n}$ with the relative upper group at
$\psi_{E_n/K}(u_n)$. Multiplying by $H$ gives the whole fixing
group of $E_n$. Since $\psi_{E_n/K}(u_n)\geq u_n\geq b_n$,
decreasingness of upper groups proves the displayed product at $b_n$.
Containment of $C_n$ proves cofinality. For a finite step $F/E_n$,
replace $H$ by its larger fixing subgroup and use the wild-group
orbit and inertia argument in the preceding lemma.
\end{proof}

\begin{lemma}[Stability of the actual coordinate discriminant]
\label{fw:separable-discriminant-stability}
\dcref{Wintenberger1983, Lemma 3.2.5.2 and its proof in 3.2.5.3}
\source{wintenberger:page-0015,wintenberger:page-0016}
\uses{fw:separable-high-break-projections,fw:separable-deep-coordinate-stages,fw:separable-polynomial-discriminant}
\lean{FontaineWintenberger.normLimitIntegralCoordinateAt_zero, FontaineWintenberger.normLimitIntegralCoordinateAt_one, FontaineWintenberger.normLimitCoordinatePolynomial, FontaineWintenberger.coeff_normLimitCoordinatePolynomial, FontaineWintenberger.natDegree_normLimitCoordinatePolynomial, FontaineWintenberger.monic_normLimitCoordinatePolynomial, FontaineWintenberger.monic_map_normLimitCoordinatePolynomial, FontaineWintenberger.natDegree_map_normLimitCoordinatePolynomial, FontaineWintenberger.map_normLimitCoordinatePolynomial, FontaineWintenberger.normLimitHighBreakProjection_discr, FontaineWintenberger.normLimitCoordinatePolynomial_discr_sub_mem, FontaineWintenberger.normLimitCoordinatePolynomial_discr_valuation, FontaineWintenberger.normLimitCoordinatePolynomial_discr_additiveValuation, FontaineWintenberger.separable_normLimitCoordinatePolynomial, FontaineWintenberger.exists_normLimit_discr_precision, FontaineWintenberger.eventually_normLimitCoordinatePolynomial_discr_stable, FontaineWintenberger.exists_cofinal_normLimitCoordinatePolynomial_discr_stable}
\leanok
Let $X=X_K(L)$ and let
$f(T)=T^k+\sum_{i<k}a_iT^i\in\mathcal O_X[T]$ be separable,
with $k\geq1$. Put $D=v_X(\Delta(f))$. For a deep coordinate
sequence as above, set
\[
 f_n(T)=T^k+\sum_{i<k}(a_i)_{E_n}T^i,
 \qquad r_n=\left\lceil\frac{p-1}{p}b_n\right\rceil.
\]
Then, for all sufficiently large $n$,
\[
 v_{E_n}(\Delta(f_n))=D.
\]
In particular $f_n$ is separable.
\end{lemma}
\begin{proof}
The actual coordinate quotient is a ring map. Its compatibility with
the integer polynomial defining the discriminant gives
\[
 (\Delta(f))_{E_n}-\Delta(f_n)\in\mathfrak m_{E_n}^{r_n}.
\]
The stable-coordinate valuation identity in that same projection
lemma gives $v_{E_n}((\Delta(f))_{E_n})=D$. Since $p\geq2$ and
$b_n\to\infty$, also $r_n\to\infty$. For $r_n>D$ the displayed
error has strictly greater valuation than $(\Delta(f))_{E_n}$.
The ultrametric equality for unequal valuations therefore gives
$v_{E_n}(\Delta(f_n))=D<\infty$, which also proves separability.
\end{proof}

\begin{lemma}[A discriminant bound annihilates a splitting-field upper group]
\label{fw:separable-splitting-upper-bound}
\dcref{Root-separation version of Wintenberger1983, paragraph 3.2.5.5}
\source{wintenberger:page-0016}
\uses{fw:separable-polynomial-discriminant,fw:absolute-upper-projections,fw:finite-upper-groups,fw:herbrand-functions,fw:normalized-valuation-scaling}
\lean{FontaineWintenberger.finrank_splittingField_le_factorial, FontaineWintenberger.finrank_le_factorial_of_isSplittingField, FontaineWintenberger.herbrandFunction_le_self, FontaineWintenberger.le_inverseHerbrandFunction, FontaineWintenberger.spectral_ramificationIdx_le_finrank, FontaineWintenberger.addVal_discr_ne_top, FontaineWintenberger.upperRamificationGroup_fixes_integral_polynomial_root, FontaineWintenberger.upperRamificationGroup_splittingField_eq_bot, FontaineWintenberger.absoluteUpperRamificationGroup_fixes_integral_polynomial_root}
\leanok
Let $f\in\mathcal O_E[T]$ be monic separable of degree $k\geq1$,
where $E$ is complete and discretely valued. Put
$D=v_E(\Delta(f))$. If $b\geq k!D$, then the absolute upper group
$G_E^b$ fixes every root of $f$.
\end{lemma}
\begin{proof}
Let $P/E$ be its finite splitting field. Adjoining the roots in order
generates $P$. Since $f$ is separable, $P/E$ is Galois and its actual
root set has $k$ elements. Its Galois group acts faithfully on this
root set: an automorphism fixing every root fixes the field they
generate. The resulting injection into the permutation group of the
root set gives
\[
 [P:E]=|\operatorname{Gal}(P/E)|\leq |S_k|=k!.
\]
The fundamental degree formula $[P:E]=e(P/E)f(P/E)$ and $f(P/E)\geq1$
then give $e(P/E)\leq[P:E]\leq k!$.
For $\sigma\in G_E^b$, its restriction to $P$ belongs to
$\operatorname{Gal}(P/E)^b$. In lower numbering this is the group
at $\psi_{P/E}(b)\geq b$. Hence for every integral root $a$,
\[
 v_P(\sigma a-a)\geq\psi_{P/E}(b)+1>b.
\]
If $\sigma a\ne a$, root separation instead gives
\[
 v_P(\sigma a-a)\leq e(P/E)D/2\leq k!D/2\leq b,
\]
a contradiction. Every root is therefore fixed. The factorial bound
is a deliberate larger bound than the nonnormal embedding bound in
the cited paragraph; its elementary proof suffices for the required
cofinal approximation.
\end{proof}

\begin{lemma}[Disjoint root adjunction and its base-changed first bound]
\label{fw:separable-disjoint-root-base-change}
\dcref{A sufficient lower-bound form of Wintenberger1983, 3.2.5.4--3.2.5.5}
\source{wintenberger:page-0016,wintenberger:page-0017}
\uses{fw:separable-splitting-upper-bound,fw:separable-embedding-extension,fw:apf-base-separable-closure,fw:absolute-finite-base-upper,fw:apf-inheritance,fw:apf-extended-bound}
\lean{FontaineWintenberger.exists_fixing_aut_root_of_fullUpper, FontaineWintenberger.minpoly_natDegree_eq_of_fullUpper, FontaineWintenberger.minpoly_map_eq_of_fullUpper, FontaineWintenberger.linearDisjoint_root_of_fullUpper, FontaineWintenberger.finrank_root_adjunction_eq_of_fullUpper, FontaineWintenberger.finrank_root_adjunction_le_polynomial_degree, FontaineWintenberger.rootFieldToCompositum, FontaineWintenberger.rootAdjunctionTensorMap, FontaineWintenberger.rootAdjunctionTensorMap_bijective_of_fullUpper, FontaineWintenberger.rootAdjunctionTensorEquivOfFullUpper, FontaineWintenberger.rootAdjunctionTensorEquivOfFullUpper_tmul, FontaineWintenberger.absoluteUpper_le_root_fixingSubgroup, FontaineWintenberger.apfSubgroup_root_compositum_eq_fixingSubgroup, FontaineWintenberger.isAPF_rootAdjunction, FontaineWintenberger.rootAdjunction_not_finite, FontaineWintenberger.rootAdjunction_fullUpper, FontaineWintenberger.apfPsi_le_rootAdjunction_firstBound, FontaineWintenberger.le_rootAdjunction_firstBound, FontaineWintenberger.apfSubgroup_comap_closureEquiv_eq_top, FontaineWintenberger.linearDisjoint_root_of_fullUpper_commonClosure, FontaineWintenberger.rootAdjunctionTensorMap_surjective_commonClosure, FontaineWintenberger.rootAdjunctionTensorMap_bijective_commonClosure, FontaineWintenberger.rootAdjunctionTensorEquivCommonClosure, FontaineWintenberger.rootAdjunctionTensorEquivCommonClosure_tmul, FontaineWintenberger.finrank_root_adjunction_eq_commonClosure, FontaineWintenberger.upper_le_root_fixingSubgroup_commonClosure, FontaineWintenberger.upperProduct_root_compositum_eq_commonClosure, FontaineWintenberger.isAPF_rootAdjunction_commonClosure, FontaineWintenberger.rootAdjunction_not_finite_commonClosure, FontaineWintenberger.rootAdjunction_fullUpper_commonClosure, FontaineWintenberger.apfPsi_le_rootAdjunction_firstBound_commonClosure, FontaineWintenberger.le_rootAdjunction_firstBound_commonClosure}
\leanok
Let $E\subset L\subset K^{\mathrm{sep}}$, with $E/K$ finite and
$L/K$ infinite APF. Suppose $b>0$ and $G_LG_E^b=G_E$.
Let $f\in\mathcal O_E[T]$ be monic separable of degree $k$ with
$b\geq k!v_E(\Delta(f))$, and let $x$ be a root. Put
$E'=E(x)$ and $M=L(x)$. Then $L$ and $E'$ are linearly disjoint
over $E$, $[M:L]=[E':E]\leq k$, and
\[
 G_MG_{E'}^{\psi_{E'/E}(b)}=G_{E'},
 \qquad i(M/E')\geq\psi_{E'/E}(b)\geq b.
\]
\end{lemma}
\begin{proof}
Write $G=G_E$, $H=G_L$, $J=G_{E'}$ and $U=G_E^b$.
The preceding lemma gives $U\subset J$, and the hypothesis is $HU=G$.
Thus $H$ acts transitively on $G/J$: every $g=hu$ has $gJ=hJ$.
These cosets are the $E$-embeddings of $E'$ into the separable closure.
Their single $H$-orbit has size $[E':E]$, so the minimal polynomial
of $x$ over $L$ has this same degree. Equivalently
$[LE':L]=[E':E]$, which is linear disjointness of the finite separable
$E'/E$ and $L/E$. The original multiplication map
$L\otimes_E E'\to M$ is therefore an isomorphism.

Since $U\subset J$, intersecting $HU=G$ with $J$ gives
$J=(H\cap J)U$: if $hu\in J$, then $h\in J$ as well.
Finite-base upper comparison identifies
$U=J\cap G_E^b$ with $G_{E'}^{\psi_{E'/E}(b)}$, while
$H\cap J=G_M$. This proves the displayed product and hence the
first-bound lower estimate. The second estimate is the pointwise
lower bound for the inverse Herbrand function.
This proof only needs this lower bound. In particular it does not
use the equality of upper groups for $u\leq i(L/E)$ printed in
3.2.5.5: upper groups decrease, so containment in $G_{E'}$ follows
from containment at $i(L/E)$ for $u\geq i(L/E)$.
\end{proof}

\begin{lemma}[Approximate roots in actual finite-extension norm fields]
\label{fw:separable-approximate-roots}
\dcref{Wintenberger1983, paragraphs 3.2.5.6--3.2.5.7}
\source{wintenberger:page-0017}
\uses{fw:separable-discriminant-stability,fw:separable-deep-coordinate-stages,fw:separable-disjoint-root-base-change,fw:separable-high-break-projections,fw:finite-functor,fw:normalized-valuation-scaling,fw:norm-limit-uniformizer}
\lean{FontaineWintenberger.exists_normLimit_approximate_root_normalized, FontaineWintenberger.exists_normLimit_approximate_root, FontaineWintenberger.exists_normLimit_approximate_root_sequence, FontaineWintenberger.normLimit_approximation_precision_tendsto, FontaineWintenberger.separableNormUnionStageExponent_eq_ramificationIdx, FontaineWintenberger.separableNormUnionStageExponent_le_finrank}
\leanok
For every monic separable $f\in\mathcal O_{X_K(L)}[T]$ of positive
degree $k$ and every real $B$, there are an actual finite separable
extension $M/L$, with $[M:L]\leq k$, and an integral element
$z\in X_K(M)$ such that
\[
 v_M(f(z))>eB,\qquad
 1\leq e=e(X_K(M)/X_K(L))\leq[M:L]\leq k.
\]
Here $v_M$ is the original normalized valuation of $X_K(M)$ and the
valuation extending the normalized valuation of $X_K(L)$ is
$\widetilde v_M=v_M/e$. Thus the displayed strict bound is exactly
$\widetilde v_M(f(z))>B$.
The witnesses $M$ may be chosen inside any fixed separable closure of
$K$ containing $L$.

Equivalently, one can choose finite separable $M_n/L$, with
$[M_n:L]\leq k$, and integral $z_n\in X_K(M_n)$ such that
\[
 \widetilde v_{M_n}(f(z_n))\geq\frac1k
       \left\lceil\frac{p-1}{p}b_n\right\rceil\longrightarrow\infty,
\]
where $b_n$ is any sequence of positive integers tending to infinity
and eventually satisfying $b_n\geq k!v(\Delta(f))$.
The coordinate construction also gives
$v_{M_n}(f(z_n))\geq\lceil(p-1)b_n/p\rceil$ in the original
normalized valuations.
\end{lemma}
\begin{proof}
Use the deep fields $E_n$ and coordinate polynomials $f_n$ in the
discriminant-stability lemma, discarding finitely many indices so that
$v_{E_n}(\Delta(f_n))=D=v(\Delta(f))$. Choose a root $x_n$ of $f_n$
in $K^{\mathrm{sep}}$, and put $E'_n=E_n(x_n)$ and $M_n=L(x_n)$.
The disjoint-root lemma gives the actual tensor isomorphism
$L\otimes_{E_n}E'_n\simeq M_n$ and the relative upper-group product
at $\psi_{E'_n/E_n}(b_n)\geq b_n$. Thus the prescribed-coordinate
projection for $M_n$ at $E'_n$ is a surjective ring map at precision
$r_n=\lceil(p-1)b_n/p\rceil$. Lift $[x_n]$ by this map to an original
integral norm family $z_n\in X_K(M_n)$.

The tensor isomorphism makes $E'_n$ a good coordinate for the
original norm-field inclusion. Hence each coefficient $a_i$ of $f$
has $E'_n$-coordinate exactly its original $E_n$-coordinate, included
in $E'_n$. Applying the same quotient ring map to $f(z_n)$ gives
\[
 q_{E'_n,r_n}(f(z_n))=f_n(x_n)\bmod\mathfrak m_{E'_n}^{r_n}=0.
\]
Its exact kernel gives $v_{X_K(M_n)}(f(z_n))\geq r_n$.
The original inclusion has ramification index
$e_n=e(E'_n/E_n)\leq[E'_n:E_n]=[M_n:L]\leq k$:
this follows by evaluating a source norm-family uniformizer at these
totally ramified tail coordinates. Dividing the last inequality by
$e_n$ gives $v(f(z_n))\geq r_n/e_n\geq r_n/k$.
Since $p\geq2$, this tends to infinity. In particular the weaker
bound $(p-1)b_n/(pk)$ follows; it is the factor $1/k$ dictated by
valuation scaling, in place of the terminal factor $k$ printed in
3.2.5.7.

For the finite-precision assertion, choose one index with $r_n/k>B$.
Then $\widetilde v_{M_n}(f(z_n))>B$, which is exactly
$v_{M_n}(f(z_n))>e_nB$ after multiplication by the positive $e_n$.
Conversely, the finite-precision assertion allows a choice of a witness
for each bound $B=r_n$. Those witnesses have
$\widetilde v_{M_n}(f(z_n))>r_n\geq r_n/k$ and
$v_{M_n}(f(z_n))>e_n r_n\geq r_n$, giving both sequence estimates.
Thus selecting a sequence adds no existence assumption.
\end{proof}

\begin{lemma}[The actual directed union and its maps]
\label{fw:separable-norm-union}
\dcref{Wintenberger1983, paragraph 3.2.1}
\source{wintenberger:page-0015}
\uses{fw:finite-functor,fw:apf-inheritance,fw:finite-norm-field-embedding,fw:norm-limit-uniformizer}
\lean{FontaineWintenberger.SeparableNormUnion,
  FontaineWintenberger.separableNormUnionField,
  FontaineWintenberger.separableNormUnionOf_injective,
  FontaineWintenberger.separableNormUnion_exists_finiteStage,
  FontaineWintenberger.separableNormUnion_exists_polynomial,
  FontaineWintenberger.separableNormUnionAlgebra,
  FontaineWintenberger.separableNormUnion_isSeparable,
  FontaineWintenberger.separableNormUnionStageExponent_pos,
  FontaineWintenberger.separableNormUnionStageExponent_spec,
  FontaineWintenberger.separableNormUnionStageExponent_unique,
  FontaineWintenberger.separableNormUnionTransition_valuation_scaling,
  FontaineWintenberger.separableNormUnionValuation,
  FontaineWintenberger.separableNormUnionValuation_of,
  FontaineWintenberger.separableNormUnionValuation_base,
  FontaineWintenberger.separableNormUnionValued,
  FontaineWintenberger.separableNormUnionAlgMap,
  FontaineWintenberger.separableNormUnionMap_of,
  FontaineWintenberger.separableNormUnionMap_injective,
  FontaineWintenberger.separableNormUnionMap_base,
  FontaineWintenberger.separableNormUnionMap_valuation,
  FontaineWintenberger.separableNormUnionMap_id,
  FontaineWintenberger.separableNormUnionMap_comp,
  FontaineWintenberger.separableNormUnionFiniteAlgEquiv,
  FontaineWintenberger.separableNormUnionFiniteEquiv_of,
  FontaineWintenberger.separableNormUnionFiniteEquiv_base,
  FontaineWintenberger.separableNormUnionFiniteEquiv_valuation}
\leanok
Let $K$ be a complete discretely valued field with perfect residue field
of characteristic $p>0$, and let $L/K$ be an infinite separable APF
extension. Put $X=X_K(L)$ and normalize its additive valuation by
$v_X(\pi_X)=1$ for a uniformizer $\pi_X$.
For every separable algebraic $M/L$, the directed system of original
norm fields of finite $L$-subextensions defines a field
\[
 X_{L/K}(M)=\varinjlim_{L\subset N\subset M,\ [N:L]<\infty}X_K(N).
\]
Its structure maps are injective, every finite subset lies in one
such image, and it is separable algebraic over $X_K(L)$.
These fields and the original maps on finite subextensions define a
covariant functor on all separable algebraic extensions and their
$L$-embeddings. If $M/L$ is finite, its canonical identification with
$X_K(M)$ respects all these structure maps.
Moreover, for each finite $N/L$ in $M$, let $v_N$ be the normalized
integer valuation of the original $X_K(N)$. There is a unique positive
integer $e_N$ such that $v_N(i_N(x))=e_Nv_X(x)$, where
$i_N:X\longrightarrow X_K(N)$ is the actual finite norm-field map.
The rational-valued valuations $v_N/e_N$ define a valuation on
$X_{L/K}(M)$ extending $v_X$. Every structure map and every induced
$L$-embedding preserves these compatible valuations. For finite
$M/L$, the canonical identification gives the valuation $v_M/e_M$
on $X_K(M)$.
\end{lemma}
\begin{proof}
Use pairs $(N,x)$, with $x\in X_K(N)$, identifying two pairs when
their images agree in a common finite upper subextension. The finite
composition laws make this an equivalence relation: for transitivity
pass to a common upper field for the two witnessing equalities.
Compute addition and multiplication in any common upper field.
The same composition laws make these operations independent of the
choices and of representatives. The inverse of a nonzero pair is
represented by $(N,x^{-1})$. Injectivity of the finite maps ensures
that a pair is zero precisely when its entry is zero. Thus the
result is a field and every structure map is injective.
Finitely many representatives have a common finite compositum,
proving the finite-subset assertion. Every element lies in a finite
separable extension $X_K(N)/X_K(L)$, so the union is separable algebraic.

An embedding $t:M\to M'$ sends each finite $N$ isomorphically onto
its actual finite image $t(N)$. Apply the original norm-field map of
this isomorphism and then the structure map for $t(N)\subset M'$.
Finite naturality makes these maps agree on overlaps, so they define
a field embedding on the displayed union. The finite identity and
composition laws prove the corresponding laws on every representative.
When $M/L$ is finite, $M$ is a greatest index of the system; the map
represented by inclusion into this index gives the asserted canonical
identification and respects the structure maps by construction.

For the valuation, the actual finite-embedding scaling theorem supplies
each positive integer $e_N$. Evaluating the scaling identity at the
original uniformizer $\pi_X$ gives
$e_N=v_N(i_N(\pi_X))$, which proves uniqueness. If $N\subset P$,
the same theorem supplies a positive integer $c_{NP}$ such that
$v_P(j_{NP}(x))=c_{NP}v_N(x)$. Apply the proved finite composition
identity $j_{NP}\circ i_N=i_P$ to $\pi_X$ to obtain
\[
 e_P=c_{NP}e_N.
\]
Consequently $v_P(j_{NP}(x))/e_P=v_N(x)/e_N$ for nonzero $x$;
zero has value $+\infty$ on both sides. Hence
\[
 v([N,x])=\begin{cases}v_N(x)/e_N&x\ne0,\\+\infty&x=0\end{cases}
\]
is independent of the representative. Compute products and sums in
a common finite stage to obtain the valuation axioms. The scaling
identity for $i_N$ shows that its restriction to the base field is
exactly $v_X$.

For an $L$-embedding $t:M\to M'$, the original isomorphism
$X_K(N)\simeq X_K(t(N))$ preserves their normalized integer
valuations and intertwines their maps from $X$. Evaluation on
$\pi_X$ therefore gives $e_{t(N)}=e_N$. This proves valuation
preservation for the induced embedding on every representative.
The greatest-index identification in the finite case carries the
same formula, giving exactly $v_M/e_M$. This normalization over
the base field is the compatible valuation on the directed union;
the individual normalized integer valuations need not agree under
finite inclusions. This equips the directed-limit construction of
paragraph 3.2.1 with its valuation normalized over the base field.
\end{proof}

\begin{lemma}[Full faithfulness on finite separable extensions]
\label{fw:separable-finite-full-faithful}
\dcref{Finite Galois form of Wintenberger1983, paragraph 3.2.6}
\source{wintenberger:page-0018}
\uses{fw:finite-functor,fw:separable-embedding-extension}
\lean{FontaineWintenberger.finiteNormLimitAlgHom, FontaineWintenberger.exists_algHom_factor_of_fixed, FontaineWintenberger.exists_common_finiteGalois_embeddings, FontaineWintenberger.finiteNormLimitAlgHom_comp, FontaineWintenberger.finiteNormLimitAlgHom_algebraMap, FontaineWintenberger.finiteNormLimitAlgHom_injective_in_galoisOverfield, FontaineWintenberger.finiteNormLimitAlgHom_surjective_in_galoisOverfield, FontaineWintenberger.finiteNormLimitAlgHom_bijective, FontaineWintenberger.finiteNormLimitHomEquiv}
\leanok
For finite separable $M_1,M_2/L$, the original functor induces a
bijection
\[
 \operatorname{Hom}_L(M_1,M_2)\simeq
 \operatorname{Hom}_{X_K(L)}(X_K(M_1),X_K(M_2)).
\]
\end{lemma}
\begin{proof}
Embed the fields into a common finite Galois extension $N/L$ by
choosing a separable closure and taking the compositum of their
normal closures. The original isomorphism maps identify both Hom
sets with the corresponding sets for their images, so it suffices
to treat subfields of $N$. Put $G=\operatorname{Gal}(N/L)$ and
$H_i=\operatorname{Gal}(N/M_i)$. Finite functoriality identifies
$G$ with $\operatorname{Gal}(X_K(N)/X_K(L))$, and its application
to the Galois extensions $N/M_i$ gives
\[
 X_K(M_i)=X_K(N)^{H_i}
\]
as actual image subfields.

An $L$-embedding of $M_1$ into $M_2$ extends to an automorphism of
the separable closure. Normality preserves $N$, so it is restriction
of some $\sigma\in G$. The containment $\sigma(M_1)\subset M_2$
is equivalent to $\sigma^{-1}H_2\sigma\subset H_1$. Two such
restrictions agree precisely when their automorphisms belong to the
same coset $\sigma H_1$. Thus this Hom set is
\[
 \{\sigma\in G:\sigma^{-1}H_2\sigma\subset H_1\}/H_1.
\]
The identical argument in the finite Galois extension of norm fields
gives the identical coset set for the second Hom set. The original
functor sends a restriction of $\sigma$ to the restriction of
$X_K(\sigma)$ by its composition law. Its map on Hom sets is therefore
exactly the identity under these two identifications, proving bijectivity.
\end{proof}

\begin{lemma}[The norm-field union of a separable closure is separably closed]
\label{fw:separable-norm-union-closed}
\dcref{Wintenberger1983, paragraphs 3.2.5.1--3.2.5.8}
\source{wintenberger:page-0015,wintenberger:page-0016,wintenberger:page-0017}
\uses{fw:separable-norm-union,fw:separable-approximate-roots,fw:separable-exact-root-extraction,fw:norm-limit-valuation-complete,fw:norm-limit-uniformizer,fw:apf-inheritance}
\lean{FontaineWintenberger.separableNormUnion_isSepClosed,FontaineWintenberger.separableNormUnion_isSepClosure}
\leanok
Let $\Omega=K^{\mathrm{sep}}$ contain $L$ and put
$Y=X_{L/K}(\Omega)$. Then $Y$ is a separable closure of $X_K(L)$.
\end{lemma}
\begin{proof}
It is already separable algebraic over $X_K(L)$. Let a nonconstant
separable polynomial over $Y$ be given. Its finitely many coefficients
lie in $X_K(N)$ for some finite $N/L$ in $\Omega$. Divide by its
leading coefficient to make it monic of degree $k$. Its discriminant
is nonzero in $Y$, hence already nonzero in $X_K(N)$, so it is
separable over $X_K(N)$ as well. If necessary
choose $s\in X_K(N)^\times$ of sufficiently large positive valuation
and replace $f(T)$ by $s^k f(T/s)$. Each lower coefficient then has
valuation $v(a_i)+(k-i)v(s)\geq0$. This replacement preserves
separability, and a root of the new polynomial divided by $s$ is a
root of the old one. It therefore suffices to treat a monic integral
separable $f$ over $C=X_K(N)$.

The field $N/K$ is infinite APF by finite inheritance. Approximate
roots supply a finite separable $M/N$ in $\Omega$ and integral
$z\in X_K(M)$ with $v(f(z))>2v_C(\Delta(f))$. Both actual norm
fields are complete for their original valuations. Here $v$ extends
$v_C$; in the original normalized valuation $v_M$ the inequality is
\[
 v_M(f(z))>2e\,v_C(\Delta(f))=2v_M(\Delta(f)),
 \qquad e=e(X_K(M)/C)\leq[M:N]\leq k.
\]
Exact root
extraction gives a root in $C(z)\subseteq X_K(M)\subset Y$.
As $M/L$ is finite, this is an element of the same directed union.
Every nonconstant separable polynomial over $Y$ has thus acquired
a root there. Repeated division by that root's linear factor shows
that every such polynomial splits in $Y$, since the quotient of a
separable polynomial is again separable. Hence $Y$ is separably
closed and is a separable closure of $X_K(L)$.
This uses Newton lifting inside the actual complete finite-extension
norm field and Krasner descent; it requires no assertion that an
arbitrary approximate root generates that whole norm field.
\end{proof}

\begin{lemma}[Reconstruction of actual finite subfields]
\label{fw:separable-finite-subfield-reconstruction}
\dcref{Finite Galois reconstruction in Wintenberger1983, Theorem 3.2.2(i)}
\source{wintenberger:page-0015,wintenberger:page-0017,wintenberger:page-0018}
\uses{fw:separable-norm-union,fw:finite-functor,fw:separable-finite-full-faithful,fw:separable-automorphism-krull}
\lean{FontaineWintenberger.separableNormUnion_existsUnique_finite_subfield_preimage,FontaineWintenberger.separableNormUnionFiniteImage_le_iff}
\leanok
Inside $Y=X_{L/K}(\Omega)$, every finite $X_K(L)$-subfield $F$ is
$X_K(M)$ for a unique actual finite $L$-subfield $M\subset\Omega$,
where $X_K(M)$ denotes its canonical image in $Y$.
This correspondence preserves and reflects inclusion.
Here $\Omega$ is an original separable closure containing $L$;
no separable-closedness of $Y$ is needed.
\end{lemma}
\begin{proof}
Put $X=X_K(L)$. The finitely many images of an $X$-basis of $F$
lie in one actual finite stage $X_K(Q)$ of $Y$. The canonical
stage map is $X$-linear, so its image contains their span and hence
all of $F$. Take the normal closure of $Q/L$ inside the original
separable closure $\Omega$, obtaining an actual finite Galois
$P/L$ containing $Q$. Write $j_P:X_K(P)\hookrightarrow Y$ for
the canonical map and put
\[
 F'=j_P^{-1}(F)\subseteq X_K(P).
\]
The containment just proved makes $j_P:F'\to F$ an $X$-isomorphism.
Finite functoriality gives the actual group isomorphism
\[
 \gamma:\operatorname{Gal}(P/L)\simeq
 \operatorname{Gal}(X_K(P)/X).
\]
Let $H$ be the inverse image under $\gamma$ of the subgroup fixing
$F'$ pointwise. Let $M\subset\Omega$ be the image of the actual
fixed field $P^H$ under the original inclusion $P\subset\Omega$.
This is a finite $L$-subfield. Finite Galois correspondence gives
\[
 F'=X_K(P)^{\gamma(H)}.
\]
The original extension $P/P^H$ is finite Galois with group $H$.
Its actual finite norm-field fixed-field theorem identifies the
image of $X_K(P^H)$ in $X_K(P)$ with this same fixed field. The
canonical isomorphism $P^H\simeq M$ and finite composition laws
identify its composite with $j_P$ with the canonical image of
$X_K(M)$ in $Y$. Thus that image is exactly $F$.

For uniqueness and inclusion reflection, write $B_N$ for the canonical
image of $X_K(N)$ in $Y$, and let
\[
 \Theta:\operatorname{Aut}_L(\Omega)\simeq
 \operatorname{Aut}_X(Y)
\]
be the original automorphism correspondence. Its finite-fixer property,
proved by finite full faithfulness, says that $\Theta(\sigma)$ fixes
$B_N$ pointwise if and only if $\sigma$ fixes $N$ pointwise.
Suppose $B_N\subseteq B_M$. Every automorphism fixing $B_M$
then fixes $B_N$, so this property gives
\[
 \operatorname{Gal}(\Omega/M)\subseteq
 \operatorname{Gal}(\Omega/N).
\]
The original $\Omega/L$ is Galois. Its Galois correspondence gives
$N=\Omega^{\operatorname{Gal}(\Omega/N)}$ and
$M=\Omega^{\operatorname{Gal}(\Omega/M)}$, and hence $N\subseteq M$.
Conversely, $N\subseteq M$ makes the canonical map $X_K(N)\to Y$
factor through $X_K(M)$ by the finite composition law, giving
$B_N\subseteq B_M$. Thus the correspondence preserves and reflects
inclusion. Applying reflection to two equal images in both directions
proves equality of the original fields and the required uniqueness.
\end{proof}

\begin{lemma}[Every finite separable target is an original finite norm field]
\label{fw:separable-finite-essential-image}
\dcref{Finite-extension consequence of Wintenberger1983, Theorem 3.2.2(i)}
\source{wintenberger:page-0015,wintenberger:page-0017}
\uses{fw:separable-norm-union-closed,fw:separable-finite-subfield-reconstruction,fw:separable-embedding-extension}
\lean{FontaineWintenberger.exists_finite_normLimit_algEquiv_of_separableClosure,FontaineWintenberger.exists_finite_normLimit_algEquiv}
\leanok
Every finite separable extension $F/X_K(L)$ is isomorphic over
$X_K(L)$ to $X_K(M)$ for some finite separable $M/L$.
\end{lemma}
\begin{proof}
Choose an original separable closure $\Omega$ containing $L$.
The norm-field union $Y=X_{L/K}(\Omega)$ is a separable closure
of $X_K(L)$ by the preceding separable-closedness theorem.
Consequently the given finite separable $F/X_K(L)$ has an
$X_K(L)$-embedding into $Y$. Its actual image is a finite
$X_K(L)$-subfield of $Y$. Apply actual finite-subfield
reconstruction to that image to obtain a finite original
$M\subset\Omega$ whose canonical norm-field image is exactly
this subfield. The given embedding identifies $F$ with its image,
and the injective canonical map identifies $X_K(M)$ with the same
image. Their composite gives the required $X_K(L)$-isomorphism.
\end{proof}

\begin{lemma}[Gluing the unique embeddings on finite subextensions]
\label{fw:separable-algebraic-full-faithful}
\dcref{Directed-union form of Wintenberger1983, Theorem 3.2.2(ii)}
\source{wintenberger:page-0015,wintenberger:page-0018}
\uses{fw:separable-norm-union,fw:separable-finite-full-faithful}
\lean{FontaineWintenberger.separableNormUnionAlgMap_bijective,
  FontaineWintenberger.separableNormUnionHomEquiv}
\leanok
For arbitrary separable algebraic $M_1,M_2/L$, the original directed
union functor induces a bijection
\[
 \operatorname{Hom}_L(M_1,M_2)\simeq
 \operatorname{Hom}_{X_K(L)}(X_{L/K}(M_1),X_{L/K}(M_2)).
\]
\end{lemma}
\begin{proof}
Given an embedding $t$ on the right, restrict it to $X_K(Q)$ for
any finite $L$-subfield $Q\subset M_1$. The images of a finite
$X_K(L)$-basis lie in some $X_K(R)$, with $R/L$ finite in $M_2$.
Their span contains the whole image of $X_K(Q)$. Finite full
faithfulness therefore gives a unique $L$-embedding $t_Q:Q\to R$
inducing that restriction. The resulting map into $M_2$ does not
depend on $R$: for two choices pass to their finite compositum in
$M_2$ and use the injectivity part of finite full faithfulness there.
If $Q\subset Q'$, place the images of both finite maps in a common
finite target field. Their induced norm-field maps agree on $X_K(Q)$,
so finite faithfulness gives $t_{Q'}|_Q=t_Q$.

The maps $t_Q$ thus define a function $\widetilde t:M_1\to M_2$.
Every finite collection of elements of $M_1$ lies in one finite
subextension, where $t_Q$ preserves the field operations and fixes
$L$. Hence $\widetilde t$ is an $L$-embedding. Its induced map on
each $X_K(Q)$ is the original $t$, so its directed union map is $t$.
Finally, if two $L$-embeddings induce the same map on directed unions,
fix $Q$ and place both finite images in a common finite $R/L$ in
$M_2$. The structure map $X_K(R)\to X_{L/K}(M_2)$ is injective,
so their finite norm-field maps agree. Finite faithfulness makes
their restrictions to $Q$ equal. Varying $Q$ proves the embeddings
equal, which proves injectivity as well as surjectivity.
\end{proof}

\begin{lemma}[Automorphisms of the original separable norm-field union]
\label{fw:separable-automorphism-correspondence}
\dcref{Automorphism-group consequence of Wintenberger1983, Theorem 3.2.2(ii)}
\source{wintenberger:page-0015,wintenberger:page-0018}
\uses{fw:separable-algebraic-full-faithful,fw:separable-norm-union}
\lean{FontaineWintenberger.separableNormUnionAutEquiv, FontaineWintenberger.separableNormUnionAutEquiv_of, FontaineWintenberger.separableNormUnionAutEquiv_valuation, FontaineWintenberger.separableNormUnion_algEquiv_valuation}
\leanok
Let $K$ be complete discretely valued with perfect residue field of
characteristic $p>0$, let $L/K$ be infinite separable APF, and let
$M/L$ be separable algebraic. Put $X=X_K(L)$ and $U=X_{L/K}(M)$.
The original union functor induces a group isomorphism
\[
 \operatorname{Aut}_L(M)\simeq\operatorname{Aut}_X(U).
\]
For $\sigma\in\operatorname{Aut}_L(M)$, a finite $L$-subfield
$Q\subset M$, and $x\in X_K(Q)$, its image acts by
\[
 i_Q(x)\longmapsto
 i_{\sigma(Q)}\bigl(X_K(\sigma|_Q)(x)\bigr).
\]
Every $X$-automorphism of $U$ preserves the compatible rational valuation.
\end{lemma}
\begin{proof}
Full faithfulness identifies the endomorphism sets of $M$ and $U$.
The functor's identity and composition laws make this bijection an
isomorphism of endomorphism monoids. The units of a field's algebra
endomorphism monoid are precisely its algebra automorphisms: an
invertible endomorphism has its inverse as a two-sided inverse function,
and an automorphism and its inverse give such a unit. Restricting the
monoid isomorphism to units therefore gives the displayed group
isomorphism. Its underlying map on every finite-stage representative
is the original functor map, which gives the stated formula.
The original induced embeddings preserve the compatible rational
valuations. Surjectivity of the group isomorphism now shows the same
for every $X$-automorphism of $U$.
\end{proof}


\begin{lemma}[Krull topology and finite norm-field fixers]
\label{fw:separable-automorphism-krull}
\dcref{The finite-fixer topology argument in Wintenberger1983, Corollary 3.2.3}
\source{wintenberger:page-0015}
\uses{fw:separable-automorphism-correspondence,fw:separable-algebraic-full-faithful,fw:separable-norm-union,fw:finite-functor}
\lean{FontaineWintenberger.separableNormUnionFiniteImage_finite, FontaineWintenberger.separableNormUnionAutEquiv_mem_fixingSubgroup_iff, FontaineWintenberger.separableNormUnion_exists_finiteImage_ge, FontaineWintenberger.separableNormUnionKrullEquiv}
\leanok
Let $K$ be complete discretely valued with perfect residue field of
characteristic $p>0$, let $L/K$ be infinite separable APF, and let
$M/L$ be separable algebraic. Put $X=X_K(L)$ and $U=X_{L/K}(M)$,
and let $\Theta:\operatorname{Aut}_L(M)\simeq\operatorname{Aut}_X(U)$
be the original automorphism correspondence. For a finite
$L$-subfield $Q\subset M$, let $B_Q$ be the actual image of
$X_K(Q)$ in $U$. Then
\[
 \Theta^{-1}\bigl(\operatorname{Aut}_{B_Q}(U)\bigr)
       =\operatorname{Aut}_Q(M),
\]
where both sides denote pointwise fixing subgroups in the respective
ambient automorphism groups. The fields $B_Q$ are finite over $X$
and are cofinal among the finite $X$-subfields of $U$.
Consequently $\Theta$ is an isomorphism of topological groups for
the two Krull topologies.
\end{lemma}
\begin{proof}
For the original inclusion $\iota:Q\hookrightarrow M$, identify
$X_{L/K}(Q)$ with $X_K(Q)$ by the canonical finite terminal
comparison. Functoriality identifies the two maps induced by
$\sigma\circ\iota$ and $\iota$ with
$\Theta(\sigma)\circ X(\iota)$ and $X(\iota)$.
The latter maps are equal exactly when $\Theta(\sigma)$ fixes
$B_Q$ pointwise. Full faithfulness says that they are equal exactly
when $\sigma\circ\iota=\iota$, which says that $\sigma$ fixes $Q$
pointwise. This proves the displayed equality.

Finite functoriality makes $X_K(Q)$ finite over $X$, so its actual
image $B_Q$ is finite over $X$. Let $D\subset U$ be any finite
$X$-subfield and choose a finite $X$-basis $d_1,\ldots,d_n$ of $D$.
Each basis element has an actual finite-stage representative in the
original directed union. A finite compositum of those original
$L$-subfields is a finite $Q/L$ whose image $B_Q$ contains every
$d_j$. It also contains $X$, so closure under $X$-linear combinations
shows $D\subset B_Q$. This proves cofinality.

The pointwise fixing subgroups of finite subfields form a
neighborhood basis at the identity for each Krull topology.
The displayed equality identifies the source basis with the
cofinal target basis just proved. Thus the inverse image under
$\Theta$ of every target identity neighborhood is a source
identity neighborhood, and its image of every source identity
neighborhood contains a target identity neighborhood.
It follows that $\Theta$ and $\Theta^{-1}$ are continuous at the
identity. Since both are group homomorphisms and translation is a
homeomorphism, they are continuous everywhere.
\end{proof}

\begin{lemma}[Isomorphisms of original norm-field unions]
\label{fw:separable-norm-union-congruence}
\lean{FontaineWintenberger.separableNormUnionCongr,FontaineWintenberger.separableNormUnionCongr_toAlgHom,FontaineWintenberger.separableNormUnionCongr_apply,FontaineWintenberger.separableNormUnionCongr_symm_apply,FontaineWintenberger.separableNormUnionCongr_valuation,FontaineWintenberger.separableNormUnionCongr_refl_apply,FontaineWintenberger.separableNormUnionCongr_trans_apply,FontaineWintenberger.separableNormUnionULiftEquiv}
\dcref{The isomorphic-copy identifications in Wintenberger1983, 3.2.1 and 3.2.6}
\source{wintenberger:page-0015,wintenberger:page-0018}
\uses{fw:separable-norm-union}
\leanok
Let $L/K$ be infinite separable APF, where $K$ is complete discretely
valued with perfect residue field of characteristic $p>0$.
For separable algebraic extensions $M/L$ and $N/L$, every
$L$-isomorphism $e:M\xrightarrow{\sim}N$ induces an
$X_K(L)$-isomorphism
\[
 X(e):X_{L/K}(M)\xrightarrow{\sim}X_{L/K}(N).
\]
Its underlying map is the original induced union map, its inverse is
the original map induced by $e^{-1}$, and it preserves the rational
valuation. The induced isomorphism of an identity is the identity,
and induced isomorphisms commute with composition. In particular,
replacing an original extension by any isomorphic copy gives this
same canonical comparison of its original norm-field unions.
\end{lemma}
\begin{proof}
The original union maps are $X_K(L)$-algebra homomorphisms.
Their proved composition and identity formulas give
\[
 X(e^{-1})X(e)=X(e^{-1}e)=1,
 \qquad X(e)X(e^{-1})=X(ee^{-1})=1.
\]
Thus they are inverse algebra isomorphisms on the same original
union fields. The valuation assertion is the proved valuation formula
for the original map on each finite-stage representative.
The identity and composition assertions follow from those same
map formulas, since the underlying maps determine algebra
isomorphisms. The final assertion applies this construction to the
given isomorphism with the chosen copy.
\end{proof}

\begin{lemma}[The functor on separable algebraic field extensions]
\label{fw:separable-algebraic-functor}
\lean{FontaineWintenberger.SeparableExtensionCat,FontaineWintenberger.separableNormUnionObjectMapOf_apply,FontaineWintenberger.separableNormUnionObjectMapOf_valuation,FontaineWintenberger.separableNormUnionObjectFiniteAlgEquiv,FontaineWintenberger.separableNormUnionObjMap_bijective,FontaineWintenberger.separableNormUnionFunctor,FontaineWintenberger.separableNormUnionFunctorFullyFaithful}
\dcref{The functor and full-faithfulness assertions of Wintenberger1983, Theorem 3.2.2}
\source{wintenberger:page-0015,wintenberger:page-0018}
\uses{fw:separable-norm-union,fw:separable-algebraic-full-faithful,fw:finite-functor}
\leanok
For a field $R$, let $\mathcal S(R)$ be the category of separable
algebraic field extensions of $R$, with $R$-algebra embeddings as
morphisms. Let $K$ be complete discretely valued with perfect residue
field of characteristic $p>0$, and let $L/K$ be infinite separable APF.
The original union construction defines a fully faithful functor
\[
 \mathcal X_{L/K}:\mathcal S(L)\longrightarrow\mathcal S(X_K(L)),
 \qquad M\longmapsto X_{L/K}(M).
\]
On every embedding it is exactly the original induced map on the
same directed norm-field union. These maps preserve the compatible
rational valuations. For finite $M/L$, the canonical terminal comparison
identifies its value with the original field $X_K(M)$ over $X_K(L)$.
\end{lemma}
\begin{proof}
Identity maps and composition of algebra embeddings satisfy the category
laws on the original fields. Every object $M$ of $\mathcal S(L)$ gives
the original norm-field union $X_{L/K}(M)$. Its proved field structure
and separable algebraicity over $X_K(L)$ make it an object of the target
category. For an $L$-embedding $f:M\hookrightarrow N$, the original
finite-image transports induce an $X_K(L)$-algebra embedding between
these same unions. The proved identity and composition formulas for
these induced maps are precisely the two functor laws. Their rational
valuation preservation is the already established finite-image
valuation formula on every stage representative.
For every pair of objects, the map on morphisms is the original union
map on embeddings. The established full Hom correspondence says that
this map is bijective. Injectivity is faithfulness and surjectivity is
fullness, so the functor is fully faithful. Finally, when $M/L$ is finite,
the directed system has terminal field $M$, and its canonical terminal
comparison is an $X_K(L)$-isomorphism with $X_K(M)$, as asserted.
\end{proof}

\begin{theorem}[Equivalence of separable extensions]
\label{fw:separable-equivalence}
\lean{FontaineWintenberger.exists_separableNormUnion_algEquiv_of_separableClosure,FontaineWintenberger.exists_separableNormUnion_algEquiv,FontaineWintenberger.separableNormUnionFunctorEssSurj,FontaineWintenberger.separableNormUnionFunctorIsEquivalence,FontaineWintenberger.separableNormUnionEquivalence,FontaineWintenberger.separableNormUnionEquivalence_functor}
\dcref{Wintenberger1983, 3.2.1 and Theorem 3.2.2}
\source{wintenberger:page-0015,wintenberger:page-0016,wintenberger:page-0017,wintenberger:page-0018}
\uses{fw:separable-algebraic-functor,fw:separable-norm-union-congruence,fw:separable-norm-union,fw:separable-norm-union-closed,fw:separable-embedding-extension,fw:separable-finite-subfield-reconstruction,fw:separable-algebraic-full-faithful}
\leanok
Let $L/K$ be infinite APF. For every separable algebraic extension $M/L$,
set
\[
 X_{L/K}(M)=\varinjlim_{L'\in\mathcal E_{M/L}}X_K(L').
\]
This construction, with the induced maps on embeddings, is a covariant
equivalence from the category of separable algebraic extensions of $L$
and $L$-embeddings to that of separable algebraic extensions of $X_K(L)$
and $X_K(L)$-embeddings. Explicitly, every separable algebraic extension
of $X_K(L)$ is isomorphic to some $X_{L/K}(M)$, and for all $M_1,M_2/L$
the induced map
\[
 \operatorname{Hom}_L(M_1,M_2)
   \longrightarrow
 \operatorname{Hom}_{X_K(L)}(X_{L/K}(M_1),X_{L/K}(M_2))
\]
is bijective. For finite $M/L$, one has $X_{L/K}(M)=X_K(M)$ under the
canonical identification.
\end{theorem}
\begin{proof}
The preceding fully faithful functor has the canonical finite
identifications. It remains to prove essential surjectivity. Let $F/X_K(L)$ be any separable
algebraic extension and embed it in the separable closure
$Y=X_{L/K}(\Omega)$. Identify $F$ with its image for this argument.
For every finite subfield $F_0\subset F$, let $M_{F_0}\subset\Omega$
be its unique reconstructed finite $L$-subfield. If $F_0\subset F_1$,
then $M_{F_0}\subset M_{F_1}$. Since finite subfields of $F$ are
directed by compositum, the fields $M_{F_0}$ form a directed family.
Put $M=\bigcup_{F_0\subset F}M_{F_0}$. This is a separable algebraic
extension of $L$ inside $\Omega$.

Every finite $L$-subfield $Q\subset M$ is contained in some
$M_{F_0}$: finitely many algebraic generators of $Q$ lie in finitely
many members of this directed family, and a common upper member
contains all of them. Its canonical image $X_K(Q)$ is therefore
contained in $X_K(M_{F_0})=F_0$. Conversely every element of $F$
belongs to a finite $F_0$, which equals the image of
$X_K(M_{F_0})$ in $X_{L/K}(M)$. The canonical embedding
$X_{L/K}(M)\to Y$ consequently has image exactly $F$, giving the
required isomorphism. Undoing the initial embedding proves the
assertion for the original abstract $F$.
Replacing $M$ by an $L$-isomorphic field preserves this comparison:
the induced maps on the two original unions are inverse algebra
isomorphisms.

This is the equivalence of \cite{Wintenberger1983}. The proof passes
through the same original finite norm fields, but makes the finite
fixed-field reconstruction and the compatibility of all finite
embeddings explicit before taking directed unions.
\end{proof}

\begin{corollary}[The absolute Galois group]
\label{fw:absolute-galois}
\dcref{Wintenberger1983, Corollary 3.2.3}
\source{wintenberger:page-0015}
\uses{fw:separable-norm-union-closed,fw:separable-automorphism-krull}
\lean{FontaineWintenberger.separableNormUnion_isSepClosure,FontaineWintenberger.separableNormUnionKrullEquiv,FontaineWintenberger.normUnionAbsoluteGaloisEquiv,FontaineWintenberger.normUnionAbsoluteGaloisEquiv_apply,FontaineWintenberger.normLimitAbsoluteGaloisEquiv,FontaineWintenberger.normLimitAbsoluteGaloisEquiv_apply,FontaineWintenberger.fieldOfNormsAbsoluteGaloisEquiv}
\leanok
For infinite separable APF $L/K$, let $\Omega$ be a separable closure
of $K$ containing $L$, and put $X=X_K(L)$ and
$U=X_{L/K}(\Omega)$. Then $U$ is a separable closure of $X$.
The original union map induces an isomorphism of profinite groups
\[
 \Theta:\Gal(\Omega/L)\xrightarrow{\ \sim\ }\Gal(U/X).
\]
Consequently, for chosen separable closures $L^{\mathrm{sep}}$ and
$X^{\mathrm{sep}}$, there is an isomorphism of topological groups
\[
 G_L=\Gal(L^{\mathrm{sep}}/L)
       \xrightarrow{\ \sim\ }
 G_X=\Gal(X^{\mathrm{sep}}/X).
\]
More precisely, choose $L$- and $X$-isomorphisms
$a:L^{\mathrm{sep}}\simeq\Omega$ and
$b:U\simeq X^{\mathrm{sep}}$. The last map sends $\sigma$ to
$b\Theta(a\sigma a^{-1})b^{-1}$, and hence is the same original
union correspondence after these separable-closure comparisons.
\end{corollary}
\begin{proof}
The original field $\Omega$ is separably closed. It is separable
algebraic over $L$, because it is separable algebraic over $K$ and
$K\subset L\subset\Omega$. Thus it is also a separable closure
of $L$. The norm-field union closure theorem proves that the same
$U$, with its original $X$-algebra structure, is a separable closure
of $X$. The original automorphism correspondence and its proved
finite-fixer comparison therefore give exactly the displayed
Krull-topological group isomorphism $\Theta$.

Any two separable closures of a field are isomorphic over that
field: extend an embedding of the base to one closure, and use
normality and algebraicity to see that its image is all of the
other closure. Choose the resulting $a$ and $b$ as in the statement.
For any field isomorphism $e:A\simeq_F B$, conjugation sends an
$F$-automorphism $g$ of $A$ to $ege^{-1}$. This is a group
isomorphism with inverse conjugation by $e^{-1}$. If $D/F$ is a
finite intermediate field of $B/F$, its inverse image $e^{-1}(D)$
is finite over $F$, and conjugation pulls back the pointwise fixer
of $D$ to the pointwise fixer of $e^{-1}(D)$. Conversely, the image
under $e$ of every finite intermediate field of $A/F$ is finite.
Since these fixers form the identity-neighborhood bases of the
Krull topologies, conjugation and its inverse are continuous at
the identity and therefore everywhere. Composing conjugation by
$a$, $\Theta$, and conjugation by $b$ gives the required
isomorphism and the stated formula.

This is Corollary 3.2.3 of \cite{Wintenberger1983}, using directly
the separable-closure theorem and the finite-fixer proof of its
automorphism correspondence; the categorical equivalence is not
needed for this consequence.
\end{proof}

\section{Ramification transfer and the ambient action}

\begin{lemma}[A cofinal test for the valuation of a norm-field difference]
\label{fw:cofinal-displacement-valuation}
\uses{fw:norm-limit-coordinate-approximation,fw:norm-limit-field,fw:norm-limit-field-valuation,fw:norm-valuation}
\lean{FontaineWintenberger.normLimitEval_sub_sub_mem, FontaineWintenberger.normLimitValuation_sub_eq_of_cofinal_coordinate}
\leanok
Let $L/K$ be infinite APF, and write $w$ for the normalized multiplicative
valuation of the original norm field $X=X_K(L)$, with
$w(\pi_X)=\exp(-1)$. For $x,y\in X$ and every $s\in\mathbb N$, there is
an actual finite $T/K$ in $L$ such that, for every finite $E\supseteq T$,
\[
 (x-y)_E-x_E+y_E\in\mathfrak m_E^s.
\]
Here the ideal is viewed in the original finite field $E$, and the
normalized multiplicative valuation $w_E$ satisfies
$w_E(\mathfrak m_E^s)\leq\exp(-s)$.

In particular, let $n\in\mathbb Z$. If for every actual finite $T/K$ in
$L$ there is an actual finite $E\supseteq T$ such that
\[
 w_E(x_E-y_E)=\exp(-n),
\]
then $w(x-y)=\exp(-n)$.
\end{lemma}
\begin{proof}
Apply the coordinate approximation for addition to the two original
families $x-y$ and $y$. Their sum in the constructed field is exactly
$x$. Thus, on an actual finite tail,
\[
 x_E-(x-y)_E-y_E\in\mathfrak m_E^s.
\]
Negating this element proves the first assertion.

For the second, choose a natural number $s>n$, and an actual finite
stage $T_0$ above which the stable norm-field valuation is computed by
each coordinate. Choose also a stage $T_1$ for the preceding difference
estimate at precision $s$. The cofinal hypothesis supplies a finite
$E\supseteq T_0T_1$ with $w_E(x_E-y_E)=\exp(-n)$. Therefore
\[
 w_E\big((x-y)_E-(x_E-y_E)\big)
 \leq\exp(-s)<\exp(-n)=w_E(x_E-y_E).
\]
The strict ultrametric inequality gives
$w_E((x-y)_E)=w_E(x_E-y_E)=\exp(-n)$.
Since $E\supseteq T_0$, its left side is exactly $w(x-y)$.
This proves the claim for every signed $n$, with no stability assumption
on a distinguished cofinal family of finite fields.
\end{proof}

\begin{lemma}[Tate averages for arbitrary finite embeddings]
\label{fw:finite-embedding-tate-average}
\dcref{Wintenberger1983, 1.1.1; Tate's integral polynomial argument without normality of the intermediate field}
\source{wintenberger:page-0004}
\uses{fw:embedding-pair-generator,fw:embedding-pair-ambient-invariance,fw:scaled-tate-identity,fw:local-inertia-size,fw:herbrand-coset-max}
\lean{FontaineWintenberger.embeddingOrder_intermediate_eq_sum_div, FontaineWintenberger.exists_embeddingNumber_intermediate_eq_herbrand, FontaineWintenberger.closureEmbeddingOrder_restriction_eq_sum, FontaineWintenberger.closureEmbeddingNumber_restriction_mul_index_eq_sum}
\leanok
Let $\Omega/K$ be finite Galois, $E$ any actual intermediate field,
$H=\operatorname{Gal}(\Omega/E)$, and $s\in\operatorname{Gal}(\Omega/K)$.
With the actual two embeddings $s|_E$ and the inclusion, one has
\[
 J_{E,\Omega}(s|_E,\mathrm{inc})
 =\frac{1}{e(\Omega/E)}\sum_{h\in H}j_\Omega(sh).
\]
No stability of $E$ under $s$ is required. If $s|_E$ is not the
inclusion, choose $s_0\in sH$ at which $j_\Omega$ is maximal and put
$t=j_\Omega(s_0)-1$. Then $t$ is finite and
\[
 i_E(s)=\varphi_H(t).
\]
For a finite separable tower $E\subset F$ and an embedding
$\sigma:F\to\overline K$, grouping the actual restriction fibers gives
\[
 i_E(\sigma)+1
 =\frac1{e(F/E)}\sum_{\tau|_E=\sigma|_E}(i_F(\tau)+1).
\]
\end{lemma}
\begin{proof}
Choose actual integral generators $x$ of $\mathcal O_\Omega$ and
$y$ of $\mathcal O_E$ over $\mathcal O_K$. Put
$f=\min_{\mathcal O_E}(x)$, $a=s(y)-y$, and
$b=s(f)(x)$, where $s(f)$ applies the embedding into
$\mathcal O_\Omega$ to each coefficient. This coefficient operation
is defined even when $s(E)\ne E$.
For every $c\in\mathcal O_E=\mathcal O_K[y]$, the polynomial identity
$P(s(y))-P(y)$ shows that $a$ divides $s(c)-c$ in
$\mathcal O_\Omega$. Since $f(x)=0$, evaluating $s(f)-f$ at $x$
therefore gives $a\mid b$.
Conversely, choose $g\in\mathcal O_K[T]$ with $g(x)=y$.
Monic minimal-polynomial division in $\mathcal O_E[T]$ gives
$g(T)-y=f(T)q(T)$ with integral $q$. Applying the original embedding
$s$ to coefficients and evaluating at $x$ gives
$y-s(y)=s(f)(x)s(q)(x)$, hence $b\mid a$.
Thus $a$ and $b$ have equal normalized valuations, including zero.
The genuine Galois product is
\[
 f(T)=\prod_{h\in H}(T-h(x)),\qquad
 b=\prod_{h\in H}(x-sh(x)).
\]
The integral-generator formulas for the actual embedding pair and
the automorphisms of $\Omega$ now give
\[
 e(\Omega/E)J_{E,\Omega}(s|_E,\mathrm{inc})
 =v_\Omega(a)=v_\Omega(b)=\sum_{h\in H}j_\Omega(sh).
\]
The denominator is the positive integer $|H_0|$ by the actual
inertia-cardinality theorem. Division proves the first formula.
If the restriction is the inclusion, the coset contains the identity
and both sides are infinite.

Otherwise no element of $sH$ is the identity, so its finitely many
orders are finite and a maximum exists. For $h\in H$, the strict
ultrametric product rule and maximality imply
\[
 j_\Omega(s_0h)=\min(j_\Omega(h),j_\Omega(s_0)).
\]
Indeed, unequal values give this equality by the strict rule; for
equal values the ordinary lower bound and maximality give equality.
Substitute into the proved average. The finite truncated-sum formula
for $\varphi_H$, including $t=-1$, gives the second assertion.
Finally choose one finite Galois overfield containing $F$ and the
image of the given embedding. Its fixing subgroups satisfy
$H_F\subset H_E$. Partition $sH_E$ into its right $H_F$-cosets.
The first formula on each coset computes the corresponding actual
embedding of $F$; each embedding occurs once. The ramification-index
tower $e(\Omega/E)=e(\Omega/F)e(F/E)$ gives the last formula.
\end{proof}

\begin{lemma}[Restriction bounds and the relative first bound for embeddings]
\label{fw:finite-embedding-restriction-bounds}
\source{wintenberger:page-0004,wintenberger:page-0018,wintenberger:page-0019}
\dcref{The finite comparisons in 1.1.1 and 3.3.3.1, expressed through the actual subgroup Herbrand functions}
\uses{fw:abstract-normed-finite-base,fw:finite-embedding-tate-average,fw:finite-galois-overfield-apf,fw:herbrand-subgroup-integral,fw:apf-herbrand-homeomorphism,fw:apf-extended-bound,fw:finite-totally-ramified-first-bound,fw:normalized-valuation-scaling,fw:ramification-field-equiv}
\lean{FontaineWintenberger.closureEmbeddingNumber_restriction_le_apfPsi, FontaineWintenberger.apfRelativeBound_le_embeddingPairNumber_of_normedAlgebra, FontaineWintenberger.embeddingNumber_eq_restriction_of_lt_originalBound}
\leanok
For finite separable $K\subset E\subset F$ and a $K$-embedding
$\sigma:F\to\overline K$ whose restriction to $E$ is not the
inclusion,
\[
 i_F(\sigma)\leq\psi_{F/E}(i_E(\sigma)).
\]
Every $E$-embedding $\tau:F\to\overline K$ satisfies
$i_F(\tau)\geq i(F/E)$, with the usual infinite value for the
inclusion. In particular, if $F/E$ is totally ramified,
$\sigma$ acts trivially on the residue field of $E$, and
\[
 i(F/E)>i_E(\sigma),
\]
then $i_F(\sigma)=i_E(\sigma)$; here $\sigma(E)$ and $\sigma(F)$
need not equal $E$ and $F$.
\end{lemma}
\begin{proof}
Choose a genuine finite Galois overfield $\Omega/K$ containing the
two images of $F$, and an actual lift $s$ of $\sigma$.
Put $D=\operatorname{Gal}(\Omega/E)$ and
$H=\operatorname{Gal}(\Omega/F)\subset D$.
Let $a$ and $b$ be the largest raw orders minus one on $sD$ and
$sH$, respectively. They are finite, since the restriction on $E$
is nonidentity, and $b\leq a$.
The preceding lemma gives
$i_E(\sigma)=\varphi_D(a)$ and
$i_F(\sigma)=\varphi_H(b)$. The actual finite nonnormal change
of parameter, applied over $E$, is
\[
 \psi_{F/E}=\varphi_H\circ\psi_D.
\]
All rings here carry the original extended valuation: finite-base
valuation compatibility identifies its normalizations before this
change of parameter. Monotonicity of $\varphi_H$ therefore gives
\[
 i_F(\sigma)\leq\varphi_H(a)
 =\varphi_H(\psi_D(i_E(\sigma)))
 =\psi_{F/E}(i_E(\sigma)).
\]

For the second assertion use a finite Galois overfield over $E$ and
write $G$ for its actual Galois group and $H$ for the fixing subgroup
of $F$. The inclusion has infinite value, so consider a nonidentity
embedding represented by $s\notin H$. For any $r\geq0$ with
$HG^r=G$, every element of $sH$ can be represented modulo $H$
by an element of $G^r$. Hence the maximal raw order minus one on
$sH$ is at least $\psi_G(r)$. The preceding maximum formula gives
\[
 i_F(s)\geq\varphi_H(\psi_G(r))=\psi_{F/E}(r)=r.
\]
For the last equality, $HG^v=G$ for $0\leq v\leq r$, so the
actual inertia-relative index is one on that interval.
If no nonnegative $r$ is admissible the first bound is $-1$ and
$i_F(s)\geq-1$ already proves the claim. Otherwise take the supremum
of these $r$ in the defining extended first bound. This includes
its infinite endpoint; a nontrivial finite extension with positive
bound has the finite attained endpoint proved earlier.

For the final assertion set $c=i_E(\sigma)$ and
$d=[F:E]=e(F/E)$. The finite average is an average of $d$ finite
numbers, so one extension $\alpha$ of $\sigma|_E$ satisfies
$i_F(\alpha)\leq c$. All these extensions have the same residue
embedding, because $F/E$ is totally ramified and $\sigma|_E$ is
inertia; in particular their numbers are nonnegative.
For any two extensions, extend their maps to automorphisms of a
common finite Galois overfield. Their relative embedding fixes the
image of $E$ and therefore has order at least
$i(\sigma F/\sigma E)=i(F/E)>c$ by the second assertion and actual
isomorphism invariance. The strict ultrametric rule applied to
$\alpha(x)-x$ and the difference of these two images, first on an
integral generator and then using the pair-order formula, shows
that all the numbers equal $i_F(\alpha)$.
The average now equals this common number. Thus every extension,
including the specified $\sigma$, has number $c$.
\end{proof}

\begin{lemma}[Ramification numbers of valuation-preserving automorphisms]
\label{fw:valued-automorphism-ramification-number}
\source{wintenberger:page-0018}
\dcref{Paragraph 3.3.1, with uniformizer independence proved by the perfect-residue criterion}
\uses{fw:perfect-residue-uniformizer-criterion,fw:ideal-power-valuation,fw:normalized-field-valuation}
\lean{FontaineWintenberger.valuedRamificationNumber, FontaineWintenberger.valuedRamificationNumber_eq_iInf, FontaineWintenberger.valuedRamificationNumber_eq_top_iff, FontaineWintenberger.valuedRamificationNumber_eq_nat_sub_one, FontaineWintenberger.neg_one_le_valuedRamificationNumber, FontaineWintenberger.valuedRamificationNumber_eq_neg_one_of_residue_ne, FontaineWintenberger.integerRamificationNumber_eq_uniformizer, FontaineWintenberger.normalizedAdditiveValue_coe, FontaineWintenberger.valuedRamificationNumber_eq_uniformizer_ratio, FontaineWintenberger.uniformizer_ratio_normalizedAdditiveValue_eq}
\leanok
Let $F$ be a discretely valued field with perfect residue field of
characteristic $p>0$, and let $g$ preserve its valuation. On the actual
valuation ring define
\[
 j_F(g)=\inf_{a\in\mathcal O_F}v_F(g(a)-a),\qquad
 i_F(g)=j_F(g)-1\in\mathbb R\cup\{+\infty\}.
\]
Then $i_F(g)=-1$ when its actual residue action is nonidentity.
Otherwise every uniformizer $\pi$ computes
\[
 j_F(g)=v_F(g(\pi)-\pi),\qquad
 i_F(g)=v_F(g(\pi)/\pi-1).
\]
Thus the latter expression is independent of the uniformizer.
The value is infinite exactly when $g$ is the identity, and otherwise
is a finite integer at least $-1$.
\end{lemma}
\begin{proof}
The induced map on the actual DVR is an automorphism, since valuation
preservation applies to $g$ and its inverse. The infimum is a minimum:
it is either infinite, or the least of a nonempty set of natural
numbers. It is positive exactly when all integral differences belong
to the maximal ideal, equivalently when the residue action is the
identity. If some residue moves, a lift has difference a unit and
hence the minimum is zero.
In the inertia case, for every integer $m\geq0$, the proved
perfect-residue uniformizer criterion says that every integral
difference belongs to $\mathfrak m_F^{m+1}$ if and only if
$g(\pi)-\pi$ does. These are precisely the same lower bounds on
$j_F(g)$ and $v_F(g(\pi)-\pi)$; they are equal, also at infinity.
Since $v_F(\pi)=1$, division by $\pi$ subtracts one from the
additive value, proving the formula. This argument applies separately
to every actual uniformizer. Infinite minimum means every integral
element is fixed; passage to the fraction field makes $g$ the identity.
Conversely the identity has only zero differences. For a nonidentity,
the finite minimum is a natural number, so subtracting one has the
asserted finite signed value. The inclusion of the nonnegative
extended values into the extended reals and subtraction by the finite
number one preserve the infinite endpoint and never introduce
negative infinity.
\end{proof}

\begin{lemma}[The actual residue action and source constants]
\label{fw:norm-action-residue-identification}
\source{wintenberger:page-0019}
\dcref{The residue-field case in 3.3.3.2, using the canonical multiplicative lifts of 2.1.2}
\uses{fw:source-residue-constants-field-embedding,fw:perfect-multiplicative-residue-lifts,fw:norm-limit-residue-field,fw:norm-field-isomorphism,fw:apf-finite-unramified-tame-parts,fw:cofinal-finite-aut-stable-layers}
\lean{FontaineWintenberger.normLimitMap_sourceResidueConstants, FontaineWintenberger.normLimitResidueAut_sourceResidueConstantsResidueEquiv, FontaineWintenberger.normLimitResidueAut_eq_refl_iff_algebraicResidueAut, FontaineWintenberger.finiteLayer_residue_eq_iff_algebraicResidueAut}
\leanok
Let $L/K$ be infinite APF, $X=X_K(L)$ and
$\sigma\in\operatorname{Aut}_K(L)$. The proved source-constant
embedding $f:k_L\to\mathcal O_X$ satisfies
\[
 X_K(\sigma)(f(a))=f(\overline\sigma(a))\qquad(a\in k_L).
\]
Consequently its actual residue isomorphism $k_L\simeq k_X$
intertwines the two actual residue actions. For every finite
$E\subset L$ containing the maximal tame subfield $K_1$, the two
residue embeddings of $E$ attached to $\sigma|_E$ and the inclusion
agree exactly when $\overline\sigma$ is the identity on $k_L$.
The field $E$ need not be stable under $\sigma$.
\end{lemma}
\begin{proof}
Take the cofinal family of actual finite $\sigma$-stable fields
$F\supset K_1$ supplied by the finite all-image compositum construction.
Write $[F:K_1]=p^{d_F}$ and let $\sigma_F$ be the actual restriction
of $\sigma$. The original residue inclusion $r_F:k_F\to k_L$ is an
isomorphism and commutes with the residue actions: this follows by
reducing the commuting maps from $\mathcal O_F$ into the actual
integral closure of $\mathcal O_K$ in $L$. Consequently
$r_F^{-1}\overline\sigma=\overline\sigma_F r_F^{-1}$.
The automorphism $\overline\sigma_F$ commutes with inverse Frobenius,
because $k_F$ is perfect, and $\sigma_F$ commutes with the canonical
multiplicative residue lifts by their proved naturality. Applying
these identities to the actual defining coordinate
\[
 f(a)_F=[r_F^{-1}(a)^{p^{-d_F}}]_F
\]
gives exactly
$\sigma_F(f(a)_F)=f(\overline\sigma(a))_F$; the same exponent
$d_F$ occurs on both sides because the original field $F$ is stable.
Original coordinate covariance identifies the left side with
$(X_K(\sigma)(f(a)))_F$. Uniqueness of norm families on this cofinal
set proves the displayed identity on the original field of norms.
Reducing it and using the proved exact reduction formula for this
same $f$ proves residue conjugacy. Finally, put the two finite images
of $E$ in their actual finite compositum inside $L$. Its residue field
embeds into the intrinsic $k_L$. Under this injection the two residue
maps from $k_E$ become the inclusion and $\overline\sigma$ composed
with the same surjective map $k_E\to k_L$. They agree if and only if
$\overline\sigma$ is the identity. This proves both directions,
including the case of an arbitrary perfect infinite residue field.
\end{proof}

\begin{lemma}[Eventual exact valuations of coordinate differences]
\label{fw:eventual-coordinate-displacement}
\source{wintenberger:page-0019}
\dcref{Quantitative form of the coordinate comparison in 3.3.3.2, using the already constructed norm-field operations}
\uses{fw:cofinal-displacement-valuation,fw:norm-limit-stable-valuation,fw:norm-limit-field-valuation,fw:ideal-power-valuation}
\lean{FontaineWintenberger.normLimitValuation_sub_eventually_coordinate}
\leanok
Let $L/K$ be infinite APF, $X=X_K(L)$, and $x,y\in X$ with
$x\ne y$. Put $m=v_X(x-y)\in\mathbb Z$, for the original normalized
additive valuation. There is an actual finite $T\subset L$ such that
\[
 v_E(x_E-y_E)=m\qquad(T\subseteq E\subseteq L,\ [E:K]<\infty).
\]
\end{lemma}
\begin{proof}
Choose an integer $N\geq0$ with $N>m$. The proved finite-coordinate
approximation supplies an actual finite $T_1$ such that
\[
 r_E=(x-y)_E-x_E+y_E\in\mathfrak m_E^N
 \qquad(E\supseteq T_1).
\]
The original stable norm valuation supplies an actual finite $T_0$
such that $v_E((x-y)_E)=v_X(x-y)=m$ for every $E\supseteq T_0$.
Take $T=T_0T_1$. For every $E\supseteq T$, the original normalized
valuation satisfies $v_E(r_E)\geq N>m$. Since
$x_E-y_E=(x-y)_E-r_E$, the strict ultrametric equality gives
$v_E(x_E-y_E)=m$. This argument includes negative $m$ and zero
coordinates individually; the assumption $x\ne y$ is used only to
make $m$ finite.
\end{proof}

\begin{lemma}[An actual finite tail of constant embedding order]
\label{fw:apf-embedding-order-stabilization}
\source{wintenberger:page-0018,wintenberger:page-0019}
\dcref{Paragraphs 3.3.3.1--3.3.3.2, applying the independently constructed norm-field approximation first and the finite restriction comparison to every upper field}
\uses{fw:finite-embedding-restriction-bounds,fw:apf-deep-finite-tail,fw:cofinal-finite-aut-stable-layers,fw:norm-action-residue-identification,fw:eventual-coordinate-displacement,fw:valued-automorphism-ramification-number,fw:faithful-action,fw:norm-limit-uniformizer,fw:norm-families,fw:embedding-pair-uniformizer,fw:norm-field-isomorphism}
\lean{FontaineWintenberger.exists_finiteLayer_stableRamificationNumber_eq, FontaineWintenberger.exists_finiteLayer_inertial_ramificationNumber_eq}
\leanok
If $L/K$ is infinite APF and $\sigma\ne1$ is a $K$-automorphism of
$L$ acting trivially on $k_L$, there are a finite actual $E\subset L$
stable under $\sigma$ and a finite real $c\geq0$ such that
\[
 i_{E'}(\sigma|_{E'})=c\qquad(E\subseteq E'\subseteq L, [E':K]<\infty).
\]
One may take $c=i_{X_K(L)}(X_K(\sigma))$.
\end{lemma}
\begin{proof}
Put $g=X_K(\sigma)$. Its proved faithfulness gives $g\ne1$;
residue conjugacy makes $g$ inertia. Let $\Pi$ be the actual original
norm-family uniformizer. The every-uniformizer criterion on this
same field gives
\[
 c=i_X(g)=v_X(g(\Pi)-\Pi)-1\in\mathbb Z_{\geq0}.
\]
In particular $n=c+1$ is a positive finite integer and
$g(\Pi)-\Pi\ne0$. Apply the preceding coordinate-difference lemma
to $g(\Pi)$ and $\Pi$, obtaining an actual finite $T$ with
\[
 v_F((g(\Pi))_F-\Pi_F)=n\qquad(F\supseteq T).
\]
For every actual $\sigma$-stable $F\supseteq T,K_1$, original
coordinate covariance identifies $(g(\Pi))_F$ with
$\sigma_F(\Pi_F)$. Since $F\supset K_0$, the same normalization
makes $\Pi_F$ a genuine uniformizer. Residue conjugacy makes
$\sigma_F$ inertia, so the finite every-uniformizer criterion gives
$i_F(\sigma)=n-1=c$. These stable fields are cofinal by the finite
all-image compositum construction, without normality of $L/K$.

Take the actual APF deep finite tail at the real number $c+1$ and
choose one actual stable $E$ containing its starting field, $T$ and
$K_1$. Then $i_E(\sigma)=c$ is finite, so the restriction to $E$ is
nonidentity. For every finite $E'\supseteq E$, the deep-tail theorem
on the original fields and their original compatible norms gives
\[
 i(E'/E)\geq c+1>i_E(\sigma).
\]
All these extensions are totally ramified because $E\supseteq K_1$.
The actual finite embedding comparison therefore gives
$i_{E'}(\sigma)=i_E(\sigma)=c$. Its proof uses an actual finite
Galois ambient field for the two images, where the original relative
first bound separates every restriction fiber; it does not require
$E'$ to be stable. This proves the full finite tail with the asserted
value of $c$.
\end{proof}


\begin{proposition}[Stabilization of finite-layer ramification numbers]
\label{fw:ramification-transfer}
\dcref{Wintenberger1983, 3.3.1 and Proposition 3.3.2; the coordinate comparison and finite restriction argument use arbitrary cofinal stable fields in place of the elementary tower}
\source{wintenberger:page-0018,wintenberger:page-0019}
\uses{fw:apf-embedding-order-stabilization,fw:valued-automorphism-ramification-number,fw:norm-action-residue-identification,fw:embedding-pair-ambient-invariance}
\lean{FontaineWintenberger.exists_finiteLayer_ramificationNumber_eq, FontaineWintenberger.finiteLayerRamificationNumber_eq_closure, FontaineWintenberger.normLimitRamificationNumber_eq_uniformizer_ratio}
\leanok
Let $L/K$ be infinite APF, and let $\sigma$ be a $K$-automorphism of $L$.
There is a finite $E/K$ in $L$ such that for every finite $E'/E$ in $L$,
\[
 i_{E'}(\sigma|_{E'})=i_{X_K(L)}(X_K(\sigma)).
\]
On the left the two actual maps from $E'$ are compared in the finite
compositum $E'\sigma(E')\subset L$. Ambient independence identifies
this intrinsic value with the value in any separable closure.
The field $E'$ need not be stable under $\sigma$.
For a continuous automorphism $g$ of a local field $F$, the notation on
the right means $i_F(g)=-1$ if $g$ acts nontrivially on $k_F$, and otherwise
\[
 i_F(g)=v_F(g(\pi_F)/\pi_F-1),
\]
where $\pi_F$ is any uniformizer, with $i_F(\mathrm{id})=+\infty$.
In particular, in the inertia case the unshifted number
$v_F(g(\pi_F)-\pi_F)$ is $i_F(g)+1$.
\end{proposition}
\begin{proof}
For $\sigma=1$, both actual maps on every original layer are the
inclusion and the norm-field action is the identity, so both numbers
are $+\infty$.
If the actual residue action on $k_L$ is nonidentity, take $E=K_1$.
The proved residue identification makes every finite embedding above
$E$ noninertia and makes the actual norm-field residue action
nonidentity. Both numbers are therefore $-1$.

For $\sigma\ne1$ acting trivially on $k_L$, the preceding lemma gives
an actual finite stable $E$ and the value
$c=i_{X_K(L)}(X_K(\sigma))$ such that every finite $E'\supseteq E$
has intrinsic number $c$, including nonstable $E'$. This is the
asserted equality on the original norm field and original finite
image composita. Finally the valuation-preserving automorphism
criterion applies to every uniformizer separately, so the stated
right-hand expression is independent of its choice, with the same
infinite endpoint. No extension of the valuation to an arbitrary
nonnormed infinite ambient field is needed: finite ambient
independence supplies the displayed closure formulation.
\end{proof}

\begin{lemma}[Cofinal finite fields stable under the original automorphisms]
\label{fw:cofinal-finite-aut-stable-layers}
\lean{IntermediateField.exists_finite_aut_stable, IntermediateField.map_normalClosure_eq, IntermediateField.normalClosureRestrict, IntermediateField.coe_normalClosureRestrict}
\leanok
Let $L/K$ be an algebraic field extension. Every finite intermediate
field $D$ is contained in a finite intermediate field $C(D)$ which
is stable under every $K$-automorphism of $L$. One may take the actual
compositum
\[
 C(D)=\bigvee_{f:D\hookrightarrow_K L}f(D)\subseteq L.
\]
No normality assumption on $L/K$ is needed, and no normality of
$C(D)/K$ is asserted. Restriction gives an actual $K$-automorphism
of $C(D)$ for every original $K$-automorphism of $L$.
\end{lemma}
\begin{proof}
There are finitely many $K$-embeddings $D\to L$. Indeed, fix a
finite $K$-basis $b_1,\ldots,b_r$ of $D$. Each image of $b_i$
is a root in $L$ of its nonzero minimal polynomial over $K$,
so has only finitely many possibilities. The tuple of these
images determines the embedding, since every element is a
$K$-linear combination of the basis. The set of embeddings
therefore injects into a finite product of finite root sets.
Each $f(D)$ has finite degree, equal to $[D:K]$; the compositum
of this finite family is finite, by repeated multiplication of
finite field-degree bounds. The original inclusion of $D$ into
$L$ is one of the embeddings, so $D\subseteq C(D)$.

For any $K$-automorphism $\sigma$ of $L$, composition
$f\mapsto\sigma\circ f$ permutes the embeddings $D\to L$,
with inverse composition by $\sigma^{-1}$. Hence $\sigma$
permutes the fields generating $C(D)$ and maps $C(D)$ onto
itself. The restrictions of $\sigma$ and $\sigma^{-1}$ to
that actual subfield are inverse $K$-algebra maps. Finally the
construction can be applied to the compositum of any two
prescribed finite fields, proving the asserted cofinality.
\end{proof}

\begin{corollary}[Faithfulness on the base norm field]
\label{fw:faithful-action}
\dcref{Wintenberger1983, Corollary 3.3.4, faithfulness assertion; Remark 3.2.4}
\source{wintenberger:page-0015}
\uses{fw:cofinal-finite-aut-stable-layers,fw:norm-field-isomorphism,fw:source-residue-constants-field-embedding,fw:norm-families,fw:norm-limit-uniformizer,fw:embedding-pair-uniformizer,fw:finite-embedding-pair-order}
\lean{FontaineWintenberger.finiteLayer_aut_eq_of_normLimitMap_fixed, FontaineWintenberger.normLimitMap_eq_self_iff, FontaineWintenberger.normLimitAut_injective}
\leanok
If $L/K$ is infinite APF and Galois, then
\[
 \Gal(L/K)\longrightarrow\operatorname{Aut}(X_K(L)),
 \qquad \sigma\longmapsto X_K(\sigma),
\]
is injective. This is the faithfulness assertion of Corollary 3.3.4;
the closed-image and filtration assertions of that corollary are stronger.
More generally, for every infinite separable APF extension $L/K$,
the original group $\operatorname{Aut}_K(L)$ acts faithfully on
$X_K(L)$, without a normality hypothesis.
\end{corollary}
\begin{proof}
Suppose $X_K(\sigma)$ is the identity. Fix any $a\in L$ and
apply the preceding lemma to the actual finite compositum
$K(a)K_1$ in $L$, where $K_1$ is the proved maximal tame
subfield. The resulting actual finite field $E$ contains $a$
and $K_1$, and is stable under $\sigma$ and $\sigma^{-1}$.
Write $s=\sigma|_E$, an actual $K$-automorphism of $E$.

The original coordinate covariance formula for norm families,
specialized to this stable field, is
\[
 (X_K(\sigma)x)_E=s(x_E).
\]
The assumed identity therefore makes $s$ fix every actual
$E$-coordinate of every norm family. In particular it fixes
the $E$-coordinates of all source constants. Their residue map
is the actual isomorphism $k_L\simeq k_E$: its inverse-Frobenius
factor is bijective because $k_L$ is perfect, and its original
inclusion from $K_0$ is surjective because $E\supseteq K_1$.
Thus these coordinate residues exhaust $k_E$, and $s$ acts
trivially on the actual residue field of $E$.

Take the proved original norm-family uniformizer $\Pi$.
Since $E\supseteq K_1\supseteq K_0$, the every-layer
normalization identity gives $v_E(\Pi_E)=v(\Pi)=1$.
The coordinate $\Pi_E$ is integral and is an actual uniformizer:
in a DVR an integral element of value $1$ has factorization
$u\pi^n$ with $n=1$, hence generates the maximal ideal.
It too is fixed by $s$. The perfect-residue embedding-pair
criterion, applied to $s$ and the identity of $E$, now gives
\[
 r_E(s,\mathrm{id})=v_E(s(\Pi_E)-\Pi_E)=+\infty.
\]
The actual embedding-order infinity criterion gives $s=1$.
Consequently $\sigma(a)=a$. As $a$ was arbitrary, $\sigma=1$.
The original norm-field maps form a group homomorphism, so a
trivial kernel proves injectivity. This direct argument adapts
the source faithfulness assertion by separation through actual
source constants and a uniformizer; it does not use the general
ramification-transfer or finite-extension functor theorems.
\end{proof}

\begin{theorem}[The absolute base action on the original norm field]
\label{fw:base-absolute-galois-action}
\dcref{Wintenberger1983, restriction to the base norm field in Remark 3.2.4}
\source{wintenberger:page-0015}
\uses{fw:norm-field-isomorphism,fw:norm-field-action-continuity,fw:faithful-action}
\lean{FontaineWintenberger.normLimitAbsoluteAut, FontaineWintenberger.normLimitAbsoluteAut_apply, FontaineWintenberger.coe_normLimitEval_normLimitAbsoluteAut, FontaineWintenberger.normLimitAbsoluteAut_valuation, FontaineWintenberger.continuous_normLimitAbsoluteAut_apply, FontaineWintenberger.normLimitAbsoluteAut_ker, FontaineWintenberger.normLimitAbsoluteAut_range}
\leanok
Let $K$ be a complete discretely valued field with perfect residue field
of characteristic $p>0$, fix a separable closure $\Omega$ of $K$, and let
$K\subset L\subset\Omega$ be an infinite normal APF extension.
Put $X=X_K(L)$ and $G_K=\operatorname{Aut}_K(\Omega)$.
There is a canonical homomorphism
\[
 \rho_{K,L}:G_K\longrightarrow\operatorname{Aut}(X),\qquad
 \rho_{K,L}(\sigma)=X_K(\sigma|_L).
\]
For every finite intermediate field $E/K$ in $L$ and original norm
family $x$, its original coordinates satisfy
\[
 \bigl(\rho_{K,L}(\sigma)x\bigr)_E
   =\sigma\bigl(x_{\sigma^{-1}(E)}\bigr),
\]
where equality is taken in $\Omega$. Each $\rho_{K,L}(\sigma)$
preserves the normalized valuation, and every orbit map
$G_K\longrightarrow X$, $\sigma\longmapsto\rho_{K,L}(\sigma)x$,
is continuous for the Krull and valuation topologies.
The kernel is exactly $\operatorname{Aut}_L(\Omega)$, and the image
equals the image of the faithful original action of $\operatorname{Gal}(L/K)$.
This assertion concerns the stable base field $X$ in Remark 3.2.4.
\end{theorem}
\begin{proof}
Normality makes every $\sigma\in G_K$ preserve $L$, and restriction is
a group homomorphism $r:G_K\longrightarrow\operatorname{Gal}(L/K)$.
Compose $r$ with the proved homomorphism $g\mapsto X_K(g)$ on the
original norm field. This defines $\rho_{K,L}$ and proves the identity
and composition laws. Applying the original coordinate covariance
formula to $g=\sigma|_L$ gives the displayed identity, after including
its two sides from $L$ into $\Omega$. The already proved valuation
preservation applies to this same $g$ and the same original family.

To see continuity of $r$, take a finite intermediate field $D/K$ in
$L$. The inverse image under $r$ of its pointwise fixer consists exactly
of the automorphisms of $\Omega$ fixing the same finitely generated
field $D$. That fixer is open in the Krull topology on $G_K$.
These fixers form a neighborhood basis of the identity; translation
therefore proves continuity everywhere. Composing $r$ with the proved
continuous original orbit map gives continuity of each orbit of
$\rho_{K,L}$.

Faithfulness of the original norm-field action gives
$\rho_{K,L}(\sigma)=1$ if and only if $r(\sigma)=1$, which means
that $\sigma$ fixes every element of $L$. This proves the kernel
identity. Finally each $K$-automorphism of $L$ extends to a
$K$-automorphism of $\Omega$: extend the corresponding embedding
into the separable closure, whose $K$-endomorphisms are automorphisms
because $\Omega/K$ is normal and algebraic. Hence $r$ is surjective,
and composition with $r$ does not change the image of the original
norm-field action. This proves the image assertion.
\end{proof}

\begin{theorem}[The ambient action on the directed norm-field union]
\label{fw:separable-union-absolute-action}
\dcref{Wintenberger1983, finite-subextension action in Remark 3.2.4}
\source{wintenberger:page-0015}
\uses{fw:separable-norm-union,fw:base-absolute-galois-action,fw:finite-functor,fw:norm-field-isomorphism}
\lean{FontaineWintenberger.separableNormUnionAbsoluteAut,FontaineWintenberger.separableNormUnionAbsoluteAut_of,FontaineWintenberger.separableNormUnionAbsoluteAut_base,FontaineWintenberger.separableNormUnionAbsoluteAut_restrictScalars,FontaineWintenberger.separableNormUnionAbsoluteAut_valuation}
\leanok
Let $K$ be a complete discretely valued field with perfect residue field
of characteristic $p>0$, fix a separable closure $\Omega$ of $K$, and let
$K\subset L\subset\Omega$ be an infinite normal APF extension.
Put $X=X_K(L)$, $U=X_{L/K}(\Omega)$, and
$G_K=\operatorname{Aut}_K(\Omega)$.
Write $i:X\longrightarrow U$ and $j_N:X_K(N)\longrightarrow U$
for the actual structure maps, where $N/L$ is finite in $\Omega$.
There is a canonical homomorphism
\[
 \widetilde\rho_{K,L}:G_K\longrightarrow\operatorname{Aut}(U)
\]
characterized on every original finite-stage family by
\[
 \widetilde\rho_{K,L}(\sigma)(j_N(x))
   =j_{\sigma(N)}\bigl(X_K(\sigma|_N)(x)\bigr).
\]
It preserves the rational valuation of $U$ normalized over $X$ and
extends the original absolute base action:
\[
 \widetilde\rho_{K,L}(\sigma)\circ i
    =i\circ\rho_{K,L}(\sigma).
\]
For $\tau\in\operatorname{Aut}_L(\Omega)$, its value at the underlying
$K$-automorphism is the union map $X_{L/K}(\tau)$ constructed from the
actual $L$-embedding $\tau$.
\end{theorem}
\begin{proof}
Normality makes $\sigma$ restrict to a $K$-automorphism $\sigma_L$ of
$L$. Its image $\sigma(N)$ contains $L$, since for each $a\in L$
the element $\sigma_L^{-1}(a)$ belongs to $N$ and maps to $a$.
If $b_1,\ldots,b_d$ is an $L$-basis of $N$, then
$\sigma(b_1),\ldots,\sigma(b_d)$ spans $\sigma(N)$ over $L$:
apply $\sigma$ to the basis expansion of an element of $N$ and send
its coefficients through $\sigma_L$. Conversely an $L$-linear
relation between these images, after applying $\sigma^{-1}$, is a
relation between the $b_i$ with coefficients sent through
$\sigma_L^{-1}$. All these coefficients vanish, and applying
$\sigma_L$ makes the original coefficients vanish. Thus
$\sigma(N)/L$ is finite. These are the actual image subfields of
$\Omega$, and the restriction $\sigma|_N:N\simeq\sigma(N)$ is a
$K$-isomorphism. Both fields are infinite APF over $K$ by finite
extension of $L$.

Let $\theta_{\sigma,N}=X_K(\sigma|_N)$. If $N\subset P$, the
commuting original field inclusions give, by finite naturality,
\[
 \theta_{\sigma,P}\circ j_{NP}
   =j_{\sigma(N),\sigma(P)}\circ\theta_{\sigma,N}.
\]
Hence the maps $j_{\sigma(N)}\circ\theta_{\sigma,N}$ agree on
representatives identified in a common finite upper field. They
define a ring map on the same directed union $U$; sums and products
are computed in a common finite stage, where the assertion follows
from the original norm-field isomorphisms.
The equalities $1(N)=N$ and
$\sigma(\tau(N))=(\sigma\tau)(N)$ hold as actual subfields of
$\Omega$. The identity and composition laws for the original
isomorphism maps therefore give the same laws on every representative
of $U$. The map associated with $\sigma^{-1}$ is an inverse, so the
maps are automorphisms and define the displayed homomorphism.

The original square between $L$, $N$, $\sigma_L$ and $\sigma|_N$
commutes. Finite naturality gives
$\theta_{\sigma,N}\circ i_N=i_{\sigma(N)}\circ X_K(\sigma_L)$.
After the structure maps into $U$, this proves the base restriction
formula, since $X_K(\sigma_L)=\rho_{K,L}(\sigma)$.
Let $v_N$ be the normalized integer valuation and $e_N>0$ its
base scaling factor. The original isomorphism $\theta_{\sigma,N}$
preserves normalized valuations, as does $X_K(\sigma_L)$ on $X$.
Evaluate the preceding square on a normalized uniformizer of $X$.
The two valuation formulas give $e_{\sigma(N)}=e_N$. Consequently
$v_{\sigma(N)}(\theta_{\sigma,N}(x))/e_{\sigma(N)}=v_N(x)/e_N$
for every nonzero $x$; both sides take value $+\infty$ at zero.
The defining formula of the union valuation proves its preservation.

Finally, if $\tau$ fixes $L$, the actual image fields $\tau(N)$ and
the restrictions $\tau|_N$ are exactly those used in the union
functor on $L$-embeddings. Its map and the automorphism just
constructed therefore agree on every $j_N(x)$. Every element of
$U$ has such a representative, proving the last assertion.
\end{proof}

\begin{theorem}[The ambient absolute Galois action]
\label{fw:galois-action}
\dcref{Wintenberger1983, Remark 3.2.4}
\source{wintenberger:page-0015}
\uses{fw:separable-union-absolute-action,fw:separable-automorphism-correspondence,fw:faithful-action,fw:absolute-galois}
\lean{FontaineWintenberger.normUnionGaloisSubgroup, FontaineWintenberger.mem_normUnionGaloisSubgroup_iff_stable_image, FontaineWintenberger.normUnionGaloisSubgroup_base_unique, FontaineWintenberger.normUnionGaloisRestriction, FontaineWintenberger.normUnionGaloisRestriction_base, FontaineWintenberger.mem_normUnionGaloisRestriction_ker_iff, FontaineWintenberger.normUnionGaloisRelativeKernelEquiv, FontaineWintenberger.normUnionGaloisKernelEquiv, FontaineWintenberger.normUnionAmbientGaloisEquiv, FontaineWintenberger.normUnionAmbientGaloisEquiv_apply, FontaineWintenberger.normUnionAmbientGaloisHom_restriction, FontaineWintenberger.normUnionAmbientGaloisHom_relative}
\leanok
Let $L/K$ be infinite APF and Galois, choose a separable closure
$\Omega$ of $K$ containing $L$, and put $X=X_K(L)$ and
$U=X_{L/K}(\Omega)$. Write $i:X\longrightarrow U$ for the
original base embedding and $\Gamma$ for the image of the faithful
original action $\operatorname{Gal}(L/K)\longrightarrow\operatorname{Aut}(X)$.
The finite-level maps induce an action of $G_K=\operatorname{Aut}_K(\Omega)$
on $U$ and an isomorphism of groups
\[
 G_K\xrightarrow{\ \sim\ }S,
 \qquad
 S=\{\beta\in\operatorname{Aut}(U):
       \beta(i(X))=i(X),\ i^{-1}\beta i\in\Gamma\}.
\]
Equivalently, $\beta\in S$ if and only if there is an original
$\sigma\in\operatorname{Gal}(L/K)$ such that
\[
 \beta(i(x))=i(X_K(\sigma)(x))\qquad(x\in X).
\]
This $\sigma$ is unique. It defines a restriction homomorphism
$r:S\longrightarrow\operatorname{Gal}(L/K)$, whose kernel is
canonically $\operatorname{Aut}_X(U)$. The restriction of the
displayed isomorphism to $G_L=\operatorname{Aut}_L(\Omega)$ is
exactly the preceding original isomorphism onto
$\operatorname{Aut}_X(U)$, and its action on the stable subfield $X$
factors through ordinary Galois restriction to $\operatorname{Gal}(L/K)$.
\end{theorem}
\begin{proof}
The same-union ambient action theorem supplies an actual group
homomorphism $A:G_K\longrightarrow\operatorname{Aut}(U)$ and proves
\[
 A(g)(i(x))=i(X_K(g|_L)(x)).
\]
It also identifies $A(h)$ for $h\in G_L$ with the original
relative automorphism correspondence
$\Theta:G_L\simeq\operatorname{Aut}_X(U)$.

First verify the asserted description of $S$. If $\beta$ satisfies
the displayed identity for $\sigma$, then it maps $i(X)$ into
$i(X)$, and surjectivity of $X_K(\sigma)$ gives equality of these
subfields. Its transported restriction is exactly $X_K(\sigma)$.
Conversely, membership of the transported restriction in $\Gamma$
provides such a $\sigma$. Injectivity of $i$ and faithfulness of the
original action make $\sigma$ unique. The identity automorphism has
witness $1$. If $\beta,\gamma$ have witnesses $\sigma,\tau$, their
product has witness $\sigma\tau$, by the original action law.
Applying $\beta^{-1}$ to the same identity, with
$x=X_K(\sigma^{-1})(y)$, proves that $\beta^{-1}$ has witness
$\sigma^{-1}$. Thus $S$ is an actual subgroup of
$\operatorname{Aut}(U)$, and uniqueness proves that its witness map
$r$ is a homomorphism. Its kernel consists exactly of the
$\beta$ fixing every $i(x)$: faithfulness gives the converse from
pointwise fixation to witness $1$. These are precisely the
$X$-algebra automorphisms of the original $X$-algebra $U$;
forgetting the algebra structure is injective and identifies their
group with this kernel.

The proved base square gives $A(g)\in S$ with
$r(A(g))=g|_L$. Suppose $A(g)=1$. Faithfulness on $X$ gives
$g|_L=1$, so $g$ is the underlying $K$-automorphism of an actual
$h\in G_L$. The relative comparison gives $\Theta(h)=1$.
Injectivity of $\Theta$ yields $h=1$ and hence $g=1$.
Thus $A$ is injective.

For surjectivity take $\beta\in S$ with witness $q=r(\beta)$.
Normality of $L/K$ and the separable-closure extension theorem
supply $g_0\in G_K$ with $g_0|_L=q$. The automorphism
$\delta=\beta A(g_0)^{-1}$ has witness $qq^{-1}=1$, so it fixes
$i(X)$ pointwise. Regard it, by the preceding kernel
identification, as an actual $X$-algebra automorphism of $U$.
Surjectivity of the original $\Theta$ supplies $h\in G_L$ with
$\Theta(h)=\delta$. The actual relative comparison and the group
law for $A$ give $A(hg_0)=\delta A(g_0)=\beta$.
This proves surjectivity and the required isomorphism.
The same proved base square gives the base factorization, and the
relative comparison gives its exact restriction to $G_L$.
All maps use the original union and its original base embedding;
no choice of a second separable-closure carrier is needed.
\end{proof}
